\documentclass[b5paper,11pt]{book}
% ...geometry тохиргоо...
\hoffset=0cm
\voffset=-1.4cm
\textwidth=13cm
\textheight=19cm

\usepackage{polyglossia}
\setdefaultlanguage{mongolian}
\setotherlanguage{english}
\usepackage{fontspec}
\usepackage{unicode-math}
\defaultfontfeatures{Ligatures=TeX}
\setmainfont{Libertinus Serif}
\setsansfont{Libertinus Sans}
\setmonofont{Libertinus Mono}
\setmathfont{Libertinus Math}
\usepackage{amsthm, latexsym}
\makeatletter
\makeatother
\usepackage{relsize,exscale}
\usepackage{titlesec}
\usepackage{titletoc}
\usepackage{calrsfs}
\usepackage{imakeidx}
\usepackage{graphicx}
\usepackage{caption}
\counterwithout{figure}{chapter}
\usepackage{tocloft}
\usepackage{pstricks-add,tikz}
\usepackage{tkz-euclide}
\usepackage{pgfplots}
\usepackage{multicol}
\usepackage{enumitem}
\pgfplotsset{compat=1.15}
\usepackage{mathrsfs}
\usetikzlibrary{arrows,math,tikzmark,calc}
\usetikzlibrary{patterns.meta}
\pagestyle{empty}
\theoremstyle{plain}
\newtheorem{thm}{Теорем}[chapter]
\newtheorem{lem}{Лемм}[chapter]
\newtheorem{cor}{Мөрдлөгөө}[chapter]
\newtheorem{axm}{Аксиом}[chapter]
\newtheorem{dfn}{Тодорхойлолт}[chapter]
\newtheorem{sen}{Чанар}[chapter]
\newtheorem{exm}{Жишээ}[chapter]
\newtheorem{prop}{Өгүүлбэр}[chapter]
\theoremstyle{remark}
\newtheorem*{sol}{\normalfont{\textbf{Бодолт}}}
\newtheorem*{hnt}{\normalfont{\textbf{Санамж}}}
\newtheorem{prb}{\normalfont{\textbf{Бодлого}}}
\newtheorem{note}{\normalfont{\textbf{Тэмдэглэл}}}
\newtheorem*{prf}{\normalfont{\textbf{Баталгаа}}}
\DeclareMathOperator{\Ln}{\textrm{Ln }}
\DeclareMathOperator{\Arg}{\textrm{Arg }}
\DeclareMathOperator{\Arctan}{\textrm{Arctan }}
\DeclareMathOperator{\Arccot}{\textrm{Arccot }}
\DeclareMathOperator{\Arccos}{\textrm{Arccos }}
\DeclareMathOperator{\Arcsin}{\textrm{Arcsin }}
\DeclareMathOperator{\re}{\textrm{\normalfont{Re }}}
\DeclareMathOperator{\im}{\textrm{\normalfont{Im }}}
%\DeclareMathOperator{\ch}{ch}
%\DeclareMathOperator{\sh}{sh}
\usepackage{fancyhdr}
\numberwithin{equation}{chapter}
\newcommand{\ml}[1]{\mathcal{L}\{#1\}}
\newcommand{\ind}[1]{{\textcolor{blue}{\it{#1}}} \index{#1}}
\newcommand{\z}{$z=x+iy$}
\newcommand{\hz}{$\overline{z}=x-iy$}
\newcommand{\bint}{\displaystyle\int}
\newcommand{\undertilde}[1]{\underset{\widetilde{}}{#1}}
\newcommand*\circled[1]{\tikz[baseline=(char.base)]{
            \node[shape=circle,draw,inner sep=2pt] (char) {#1};}}
\newcommand{\colgur}[1]
{\begin{multicols}{3}#1\end{multicols}}
\newcommand{\colhoy}[1]
{\begin{multicols}{2}#1\end{multicols}}
\newcommand{\morshilj}{\textcolor{white}{mor shiljne  oor um bhgui\newline}}
%\usepackage{xltxtra}
%\setmainfont[Mapping=tex-text]{Times New Roman}
\renewcommand\qedsymbol{$\blacksquare$}
\setcounter{chapter}{0}
\titleformat{\chapter}
  {\bf\Large\filcenter}{\Large{Лекц }\thechapter.}{1em}{}
\titleformat{\section}{\bf\large\filcenter}{\thesection.}{1em}{}
\setcounter{tocdepth}{3}
\pagestyle{fancy}
\fancyhf{}
\fancyhead[LE,RO]{Комплекс анализ}
\fancyhead[RE,LO]{}
\fancyfoot[LE,RO]{\thepage}
%\titleformat{\part}{\normalfont\HUGE\bfseries}{ ХЭСЭГ \thepart. \hfill}{}{$\star$}
\title{\LARGE \bf  Комплекс анализ}
\author{Д.Даянцолмон}
\makeindex
\begin{document}
\frontmatter
\begin{titlepage}
\begin{center}

%\includegraphics[width=0.15\textwidth]{muislogo.png}

\textsc{\bf МОНГОЛ УЛСЫН ИХ СУРГУУЛЬ}\\[5cm]



%\textsc{\Large Final year project}\\[0.5cm]

% Title

{ \LARGE \bfseries Комплекс анализын лекц} \\[3cm]

{ \LARGE \bfseries Д.Даянцолмон}\\%[0.4cm]%
\end{center}
\end{titlepage}
\newpage
%
\section*{MATH311 Комплекс анализ хичээлийн ерөнхий мэдээлэл}

\begin{tabular}{ll}
Хичээлийн түвшин: & Бакалавр (Математикийн анги) \\
Кредит: & 3 кредит \\
Лекц: & 2 цаг / долоо хоног \\
Семинар: & 2 цаг / долоо хоног \\
Өмнөх мэдлэг: & Анализ I, II, Олон хувьсагчийн функцийн онол
\end{tabular}

%-------------------------------------------------

\section*{MATH311 Комплекс анализ хичээлийн тодорхойлолт}

Энэ хичээлээр комплекс тоо болон комплекс хувьсагчийн функцийн үндсэн онолыг судална. Комплекс функцийн хязгаар, тасралтгүй чанар, уламжлал, интеграл, цуваа, онцгой цэгүүд, үлдэгдэцийн теорем болон түүний хэрэглээг онол–нотолгоо, бодлого бодолттой хослуулан үзнэ. Хичээл нь цаашдын математик анализ, дифференциал тэгшитгэл, математик физикийн хичээлүүдийн суурь болно.

%-------------------------------------------------

\section*{Сургалтын үр дүн}

Хичээлийг амжилттай дүүргэсний дараа оюутан:
\begin{itemize}
\item Комплекс тооны алгебр ба геометр шинжийг тайлбарлаж чадна.
\item Комплекс функцийн хязгаар, тасралтгүй чанарыг шалгаж чадна.
\item Коши--Риманы тэгшитгэлийг ашиглан уламжлал орших эсэхйиг тогтооно.
\item Голоморф ба гармоник функцийн уялдааг ойлгоно.
\item Кошигийн интеграл теорем, томьёог хэрэглэнэ.
\item Тэйлор, Лораны цуваагаар функц задлан шинжилнэ.
\item Онцгой цэгүүдийг ангилж, үлдэц ашиглан интеграл бодно.
\end{itemize}

%-------------------------------------------------
\newpage
\section*{Лекцийн сэдвийн хөтөлбөр (16 долоо хоног)}

\begin{enumerate}[label=\textbf{Хичээл \arabic*.}]

\item Комплекс тоо ба комплекс хавтгай

\item $\mathbb{C}$-ийн топологийн суурь ойлголтууд  
(дотоод, гадаад, хязгаарын цэг, тусгаарлагдсан цэг, цэгийн орчин, задгай ба битүү олонлог,  холбоост ба тусгаарлагдсан олонлог, муж)

\item Комплекс функцийн жишээнүүд, олон утгатай функцүүд  
(олон гишүүнт, рационал, экспоненциал, тригонометр функцүүд)

\item Комплекс функцийн хязгаар ба тасралтгүй чанар

\item Комплекс уламжлал ба Коши--Риманы тэгшитгэл  
(зайлшгүй ба хүрэлцээтэй нөхцөл)

\item Голоморф ба гармоник функцүүд, дээд эрэмбийн уламжлал

\item Конформ дүрслэл ба уламжлалын геометр утга

\item Комплекс интеграл

\item Кошигийн интеграл теорем  

\item Кошигийн интеграл томьёо

\item Тэйлорын цуваа, аналитик ба гармоник функц, тэдгээрийн уялдаа

\item Лораны цуваа ба онцгой цэгүүд

\item Үлдэц (Residue)

\item Үлдэцийн теоремын хэрэглээ

\item Конформ дүрслэл ба Мёбиусын хувиргалт

\item Геометрийн зарчмууд (аргументын зарчим, хамгийн их модулийн зарчим, Шварцын лемма, Алгебрийн үндсэн теоремын баталгаа)
\end{enumerate}

%-------------------------------------------------
\newpage
\section*{Үнэлгээний журам (108 оноо)}

\subsection*{Ирц — 21 оноо}
\begin{itemize}
\item Лекцийн ирц: 15 оноо (1 тасалбал 1 оноо хасна)
\item Семинарын ирц: 6 оноо (1 тасалбал 1 оноо хасна)
\item Лекц эсвэл семинарыг \textbf{6 удаа тасалбал} \textbf{F*} үнэлгээ авна
\end{itemize}

\subsection*{Семинарын оноо  — 25 оноо}
Хичээлийн явцад семинарын хичээл дээр 5 удаа самбарт аман асуултад хариулна.\\
Нэг удаагийн асуулт: 5 оноо\\ Асуултын агуулга:
\begin{itemize}
\item Комплекс тооны тооцоолол, зэрэг, логарифм
\item Уламжлал, Коши--Риманы тэгшитгэл
\item Голоморф, гармоник функц, конформ дүрслэл
\item Кошигийн теорем ба томьёо
\item Цуваа, онцгой цэг, үлдэц
\end{itemize}

\subsection*{Гэрийн даалгавар — 7 оноо}

\subsection*{Явцын шалгалт — 25 оноо}
\begin{itemize}
\item 9-р долоо хоногт авна
\item Эхний 7 лекцийн агуулга (интегралаас өмнөх)
\end{itemize}

\subsection*{Улиралын шалгалт — 30 оноо}
\begin{itemize}
\item 17-р долоо хоногт авна
\item Сүүлийн 9 лекцийн агуулга (интегралаас хойших)
\end{itemize}

\subsection*{Нийт оноо}
108 онооноос авсан оноогоор дүн гаргана.  
100-аас дээш оноо авсан тохиолдолд 100 гэж дүгнэнэ.

%-------------------------------------------------

\section*{Сурах бичиг, материал}

\begin{itemize}
\item С.Чимэдцэрэн, С.Батболд — \emph{Нэг хувьсагчийн функцийн комплекс анализ}
\item С.Чимэдцэрэн — \emph{Комплекс хувьсагчийн функцийн онол}
\item L. Ahlfors — \emph{Complex Analysis}
\item J. B. Conway — \emph{Functions of One Complex Variable}
\item E. Stein, R. Shakarchi — \emph{Complex Analysis}
\end{itemize}

%-------------------------------------------------

\tableofcontents
\mainmatter
%\chapter{Комплекс тоо, хавтгай}
\section{Комплекс тоо, үндсэн үйлдлүүд}
\begin{dfn}
Бодит $x$, $y$ тоонуудын эрэмбэлэгдсэн $(x,y)$ хосыг комплекс тоо гэнэ. Комплекс тоонуудын олонлогийг $$\mathbb{C}=\{z=(x,y)\mid x, y \in\mathbb{R}\}$$-ээр тэмдэглэе. 
$x$-ийг $z$-ийн бодит, $y$-ийг хуурмаг хэсэг \index{Бодит (ба хуурмаг) хэсэг} гэх ба
$$x=\re z, \qquad y=\im z$$
тэмдэглэл хэрэглэдэг. 
\end{dfn}
Комплекс тоонуудын адилтгал, тэнцэл, нийлбэр, үржвэр дараах байдлаар тодорхойлогдоно:
\begin{enumerate}
    \item [1)] $y=0$ үед $(x,0)$ тоо бодит $x$-тэй адилтгагдана: 
    $$(x,0)=x, \qquad \forall x\in\mathbb{R}$$
Энд $(0,0)=0$, $(1,0)=1$, $(0,1)=i$ тоонуудыг харгалзан тэг, нэгж ба хуурмаг нэгж \index{Хуурмаг нэгж} хэмээн нэрлэнэ.
\item [2)] $z_1=(x_1,y_1)$, $z_2=(x_2,y_2)$ байхад $x_1=x_2$ ба $y_1=y_2$ байвал $z_1=z_2$ гэдэг.
\item [3)] Комплекс $(x_1+x_2, y_1+y_2)$ тоог тэдгээрийн нийлбэр \index{Комплекс тоонуудын нийлбэр} гэнэ:
\begin{equation}\label{1.1}
    z_1+z_2=(x_1+x_2,y_1+y_2)
\end{equation} 
\item[4)] Комплекс $(x_1x_2-y_1y_2$, $x_1y_2+x_2y_1)$ тоог тэдгээрийн үржвэр\index{Комплекс тоонуудын үржвэр} гэнэ:
\begin{equation}\label{1.2}
    z_1\cdot z_2=(x_1x_2-y_1y_2 , x_1y_2+x_2y_1)
\end{equation} 
\end{enumerate}
(\ref{1.1}), (\ref{1.2}) ёсоор дараах арифметик хуулиуд шууд биелэгдэх нь илэрхий: 
\begin{align*}
&z_1+z_2=z_2+z_1  - \text {нэмэхийн байр солих хууль};\\
&z_1\cdot z_2=z_2\cdot z_1 - \text {үржүүлэхийн байр солих хууль}; \\
&(z_1+z_2)+z_3=z_1+(z_2+z_3) - \text {нэмэхийн хэсэглэн нэгтгэх хууль}; \\
&(z_1\cdot z_2)z_3=z_1(z_2\cdot z_3) - \text {үржүүлэхийн хэсэглэн нэгтгэх хууль}; \\ 
&(z_1+z_2) z_3=z_1z_3+z_2z_3 - \text {хаалт задлах хууль}.
\end{align*}
Мөн
$$z+0=z, \quad z\cdot 0 =0,\quad z\cdot 1=z, \quad i\cdot i=i^2=-1$$
байх нь ч мэдээж. $z=(x,y)\in \mathbb{C}$ тоо $z=x+iy$ гэсэн алгебр хэлбэртэй тавигдана.

Учир нь
$$z=(x,y)=(x,0)+(0,y)=(x,0)+(0,1)(y,0)=x+iy$$
\begin{dfn} $z=(x,y)$ тооны хувьд $(-x,-y)$ тоог $z=(x,y)$ тооны эсрэг тоо гэнэ.

Хасах үйлдэл нь эсрэг тоог нэмэх үйлдэл гэж тодорхойлогдоно. Өөрөөр хэлбэл $z_1-z_2$ ялгавар \index{Комплекс тоонуудын ялгавар} нь $z_1+(-z_2)$  тоо байна. Хэрэв $z_1=(x_1,y_1)$, $z_2=(x_2,y_2)$ гэвэл $z_3=z_1-z_2=(x_1-x_2, y_1-y_2).$   
\end{dfn}
\begin{dfn}
    $(x,-y)=x-iy$ тоог $z=(x,y)$ тооны хосмог гэх бөгөөд $\overline{z}=x-iy$ гэж тэмдэглэнэ. Өгөгдсөн комплекс тооноос түүний хосмогт шилжихийг хосмогийн үйлдэл гэнэ. 
\begin{align*}
&\overline{\overline{z}}=z, \quad \overline{z_1+z_2}=\overline{z_1}+\overline{z_2},  \quad \overline{z_1\cdot z_2}=\overline{z_1}\cdot \overline{z_2}, \\ &z+\overline{z}=2x, \quad z-\overline{z}=2iy,  \quad z\cdot \overline{z}=|z|^2=x^2+y^2\geq 0 
\end{align*} байх нь мэдээж.
\end{dfn}
\begin{dfn}
    Тэгээс ялгаатай комплекс тоонд хуваах үйлдэл нь үржүүлэхийн урвуу үйлдэл гэж тодорхойлогдоно. Өөрөөр хэлбэл: $\dfrac{z_1}{z_2}$, $z_2\neq 0$ ноогдвор\index{Комплекс тооны ноогдвор} нь  $z_2\cdot z_4=z_1$ чанарт нийцэх
$z_4=(x_4,y_4)$ тоо байна. Иймд
\begin{align*}
z_4=\dfrac{z_1}{z_2}&=\dfrac{x_1+iy_1}{x_2+iy_2}=\dfrac{(x_1+iy_1)(x_2-iy_2)}{x_2^2+y_2^2} =\\
&=\left(\frac{x_1x_2+y_1y_2}{x_2^2+y^2_2}, \dfrac{x_2y_1-x_1y_2}{x_2^2+y_2^2}\right).
\end{align*}
\end{dfn}
\section{Комплекс тооны геометр дүрслэл. \\
Модуль ба аргумент. Туйлын хэлбэр}
Хавтгай дээр тэгш өнцөгт координатын $Oxy$ систем өгөгдсөн байг. Тэгвэл $z=x+iy$ гэсэн комплекс тоо нь уул хавтгайн $x$—абсцисстай, $y$—ординаттай  $(x,y)$ цэгээр дүрслэгдэнэ. 
\begin{figure}[h]
    \centering
\begin{tikzpicture}[line cap=round,line join=round,line width=0.8pt,>=triangle 45,x=1cm,y=1cm]
\clip(0,0) rectangle (6,4.5);
\coordinate (x1) at (0,1);
\coordinate (y1) at (1,0);
\coordinate (x2) at (5.5,1);
\coordinate (y2) at (1,4);
\coordinate (o) at (1,1);
\coordinate (z) at (4,3);
\draw[-stealth,line width=0.5pt] (x1)  -- (x2) node[right] {$x$};
\draw[-stealth,line width=0.5pt] (y1)  -- (y2) node[right] {$y$};
\draw[stealth-,line width=0.5pt] (z) node[above right ] {$z=(x,y)$} --(o);
\draw (o) node[anchor=north east] {$O$};
\draw (o)+(1.4,0.8) node[anchor=north west] {$\varphi$};
\draw ($(o)!0.6!(z)$) node[anchor=south] {$r$};
\draw[-stealth] (3,1) arc(0:atan(2/3):2);
\draw[dashed] (4,1) node[below] {$x$} -- (z) -- (1,3) node[left] {$y$};
\end{tikzpicture}
  \captionof{figure}{}
    \label{z1}
\end{figure}
Комплекс тоо ба хавтгайн цэгийн олонлогуудын энэ харгалзаа харилцан нэгэн утгатай нь мэдээж.

Комплекс $z=x+iy$ тоо авах бүрд, түүнд $x,y$ координат бүхий вектор (чөлөөт) харгалзлаа ч гэж ярьж болно. Учир нь 
$$z=x+iy=x\cdot 1+y\cdot i =x\cdot (1,0)+y\cdot (0,1)$$
болж $z$ нь тэнхлэгүүд дээрх нэгж $1=(1,0)$ ба $i=(0,1)$ векторуудаар задлагдлаа. Энэ харгалзаа мөн л харилцан нэгэн утгатай.

Иймд комплекс тоонуудын, хавтгайн цэгүүдийн, чөлөөт векторуудын олонлог өөр хоорондоо эквивалент (чацуу чанартай) бөлгөө.

Зураг 2 дээр комплекс тоонуудын нийлбэр, ялгавар
\begin{figure}[h]
    \centering
\begin{tikzpicture}[line cap=round,line join=round,line width=0.8pt,>=triangle 45,x=1cm,y=1cm]
\clip(0,0) rectangle (6,4.5);
\coordinate (x1) at (0,1);
\coordinate (y1) at (1,0);
\coordinate (x2) at (5.5,1);
\coordinate (y2) at (1,4);
\coordinate (o) at (1,1);
\coordinate (z2) at (1.7,2.5);
\coordinate (z1) at (2.8,1.5);
\draw[-stealth,line width=0.5pt] (x1)  -- (x2) node[right] {$x$};
\draw[-stealth,line width=0.5pt] (y1)  -- (y2) node[right] {$y$};
\draw (o) node[anchor=north east] {$O$};
\draw[-stealth,line width=0.5pt] (o)--(z1) node[below right ] {$z_1$};
\draw[-stealth,line width=0.5pt] (o)--(z2) node[above right ] {$z_2$};
\draw[-stealth,line width=0.5pt] (o) -- ($(z1)+(z2)-(o)$) node[above right ] {$z_1+z_2$};
\draw[-stealth,line width=0.5pt] (o)-- ($(z1)-(z2)+(o)$) node[above right] {$z_1-z_2$};
\end{tikzpicture}
  \captionof{figure}{}
    \label{z2}
\end{figure}
олохыг векторуудын нэмэх, хасах дүрэм ашиглан геометр байдлаар харуулав. Эндээс

\begin{equation}
\tag{a}
    |z_1+z_2|\leq |z_1|+|z_2|
\end{equation} бөгөөд энэ нь
\begin{equation}
\tag{b}
    |z_1-z_2|\geq |z_1|-|z_2|
\end{equation} тэнцэтгэл биштэй эквивалент. Учир нь (а) үнэн гэвэл $$|z_1|=|z_1-z_2+z_2|\leq |z_2|+|z_1-z_2|$$
 бөгөөд эндээс (b) гарав. Иймд $(a)\sim (b).$
\begin{dfn}
$r=\sqrt{x^2+y^2}=\sqrt{z\cdot \overline{z}}$ тоог комплекс $z$ тооны модуль гэж нэрлэх бөгөөд $|z|$-аар тэмдэглэнэ (Зураг 1). Энэ тоо нь координатын эхээс $z$ хүртлэх зай юм. Иймд $r=|z|\geq 0$ байх нь ойлгомжтой. Зөвхөн $z=0$ үед тэнцэл нь биелэгдэнэ.    
\end{dfn}
\begin{dfn}
$Ox$ тэнхлэгийн эерэг чиглэлтэй $z$ векторын үүсгэх $\varphi$ өнцгийг комплекс $z=x+iy$ тооны аргумент гэж нэрлэх бөгөөд $\Arg  z$-ээр тэмдэглэнэ (Зураг 1). Ингэхдээ $Ox$ тэнхлэгээс нар буруу (цагийн зүүний эсрэг) тоологдож байвал уул өнцгийн хэмжээг эерэг, харин нар зөв бол сөрөг гэж тооцдог. $z=0$ тооны хувьд аргумент тодорхойлогдохгүй.
\end{dfn}
Комплекс тооны аргументийн
$$-\pi <\varphi_0\leq \pi$$
нөхцөлд тохирох цор ганц $\varphi_0$ утгыг аргументийн гол утга гэх ба $\varphi_0=\arg  z$ гэж тэмдэглэнэ.
\begin{equation} \label{1.3}
\begin{split}
    &\arg  z=\arctan \frac{y}{x}+\varepsilon \pi; \quad x>0 \quad \text{үед} \quad \varepsilon=0;\\
    &x<0\quad y\geq 0 \quad \text{үед} \quad \varepsilon=1; \\
    &x<0\quad y< 0 \quad \text{үед} \quad \varepsilon=-1; \\
    &\arg  z=\frac{\pi}{2}, \quad x=0.
\end{split}
\end{equation}
Комплекс тооны аргумент олон утгат бөгөөд тэр нь гол утгаасаа $2\pi$ өнцөгт бүхлээр хуваагдахуйц нэмэгдэхүүнээр ялгагдах нь илэрхий. 
\begin{exm}
$z=-\sin\dfrac{\pi}{3}-i\cos \dfrac{\pi}{3}$ бол түүний модуль ба аргументийг ол.
\end{exm}
\begin{sol}
$x=-\sin \dfrac{\pi}{3}<0$, $y=-\cos \dfrac{\pi}{3}<0$ тул (\ref{1.3}) ёсоор 
$$\arg  z =-\pi +\arctan \left(\cot\frac{\pi}{3}\right) =-\pi +\arctan \left(\tan \left(\frac{\pi}{2}-\frac{\pi}{3}\right)\right)$$ тул
\begin{align*}
    &\Arg z=-\frac{5\pi}{6}+2k\pi , \quad k=0,\pm 1,... \\
    & |z|=\sqrt{\sin^2\frac{\pi}{3}+\cos^2\frac{\pi}{3}}=1
\end{align*}
\end{sol}

\begin{exm}
$z=-1+i\sqrt{3}$ бол $|z|$, $\arg  z$-ыг ол. 
\end{exm}
\begin{sol}
\begin{align*}
    & |z|=\sqrt{(-1)^2+(\sqrt{3})^2}=2\\
    & \arg  z=\arctan \frac{\sqrt{3}}{-1}+\pi \quad (x<0 \quad \text{тул}) =\pi -\frac{\pi}{3}=\frac{2\pi}{3}.
\end{align*}
\end{sol}
\begin{center}
    \textbf{Дасгал}
\end{center}
1. Дараах комплекс тоонуудын модуль ба аргументын гол утгыг ол.
\begin{enumerate}
\colgur{
    \item [a.] $z=3+4i$;
    \item [б.] $z=-2+2\sqrt{3}i$;
    \item [в.] $-5-i$;}
\colhoy{    \item [г.] $z=-\cos \frac{\pi}{5}+i\sin \frac{\pi}{5}$;
    \item[д.] $z=-3i$;}
    \colgur{
    \item [е.] $z=-4$;
    \item [ж.] $z=3$;
    \item [з.] $z=\frac{3}{2}i.$}
\end{enumerate}
\section{Комплекс тооны тригонометр дүрслэл}
Зураг 1-ээс харахад $z$ цэгийн декарт ба туйлын координатуудын холбоо
\begin{equation} \label{1.5}
\begin{cases}
x=r\cos \varphi \\
y=r\sin\varphi
\end{cases}
\end{equation}
байх нь ойлгомжтой. (\ref{1.5}) ёсоор комплекс $z$ тоо
\begin{equation} \label{1.6}
z=x+iy=r(\cos \varphi+i\sin \varphi)
\end{equation}
дүрстэй тавигдана. (\ref{1.6}) бичиглэлийг комплекс тооны тригонометр хэлбэр гэнэ. Бичиглэлийн товч байдлыг харгалзан
\begin{equation}\tag{$*$}
\varepsilon({i\varphi})=\cos\varphi +i\sin\varphi 
\end{equation}
тэмдэглэл орлуулбал
\begin{equation} \label{1.7}
z=r\varepsilon({i\varphi})
\end{equation}
бөгөөд үүнийг комплекс тооны илтгэгч хэлбэр гэнэ. (\ref{1.6}) ба (\ref{1.7}) хэлбэрийг заримдаа туйлын хэлбэр гэж нэрлэнэ. Тэмдэглэл байдлаар оруулсан (${*}$) тэнцлийг Эйлер (1707-1783)- ийн томьёо гэж нэрлэх бөгөөд түүнийг $$\varepsilon({i\varphi})=\cos\varphi +i\sin\varphi =e^{i\varphi}$$ үнэн болохыг хожуу батална. Энд $e^{2\pi i}=1$, $e^{\pi i}=-1$, $e^{\frac{\pi}{2}i}=i$, $e^{-\frac{\pi}{2}i}=-i$ байна. 

$(*)$ тэнцэлд $\varphi$-ийн оронд $-\varphi$ тавихад
\begin{equation}
\tag{$\overline{*}$}
e^{-i\varphi}=\cos\varphi -i\sin\varphi 
\end{equation}
болж $(*)$ ба $(\overline{*})$-ийг хооронд нь нэмэх, хасах үед 
\begin{equation} 
\cos\varphi =\dfrac{e^{i\varphi}+e^{-i\varphi}}{2}
\end{equation}
\begin{equation}
\sin\varphi =\dfrac{e^{i\varphi}-e^{-i\varphi}}{2i}
\end{equation}
болно.

Комплекс тооны илтгэгч хэлбэр ашиглаж, зэрэг \index{Комплекс тоонуудын зэрэг} дэвшүүлэх болон язгуур \index{Комплекс тоонуудын язгуур} гаргах үйлдлүүд гүйцэтгэхэд тохиромжтой.

$z_1=r_1e^{i\varphi_1}$, $z_2=r_2e^{i\varphi_2}$ авахад
\begin{equation} \label{1.10}
z_1\cdot z_2=r_1\cdot r_2 e^{i(\varphi_1+\varphi_2)}
\end{equation}
\begin{equation} \label{1.11}
\dfrac{z_1}{z_2}=\dfrac{r_1}{r_2}\cdot e^{i(\varphi_1-\varphi_2)}, \quad r_2\neq 0
\end{equation} $z=re^{i\varphi}$, $z\neq 0$ гэвэл 
\begin{equation} \label{1.12}
z^n=r^ne^{in\varphi}, \quad n\in\mathbb{N}
\end{equation}
\begin{equation} \label{1.13}
\sqrt[n]{z}=\sqrt[n]{r}e^{i\frac{\varphi+2k\pi}{n}}, \quad (k=0,1,2,\ldots, n-1)
\end{equation}
байхыг хялбархан баталж болно. Жишээ нь 
\begin{align*}
    z_1\cdot z_2&=r_1\cdot r_2[(\cos \varphi_1\cdot \cos \varphi_2-\sin \varphi_1\cdot \sin\varphi_2)+ \\ &+i(\cos \varphi_1\cdot\sin \varphi_2+\cos \varphi_2\cdot \sin\varphi_1)] \\
    &= r_1\cdot r_2 [\cos (\varphi_1+\varphi_2)+i(\sin \varphi_1+\varphi_2)]\\
    &=r_1r_2\cdot e^{i(\varphi_1+\varphi_2)}
\end{align*}
болж (\ref{1.10}) батлагдлаа. Үүнтэй төстэйгөөр (\ref{1.11}) батлагдана.

Үржүүлэх дүрэм дээр математик индукцийн арга хэрэглэхэд (\ref{1.12}) батлагдана.

Одоо комплекс тооноос язгуур гаргах (\ref{1.13}) томьёог баталъя. $n\in N$ байг. Өгөгдсөн комплекс $z$ тооны хувьд $z=W^n$ тэгшитгэл бодъё. Энэ тэгшитгэлийн шийдийн олонлогийг $\sqrt[n]{z}$ гэж тэмдэглэе. $z$ ба $W$ тоонуудыг тригонометр хэлбэрт бичвэл 
$$z=r(\cos\varphi +i\sin \varphi); \qquad W=\rho (\cos v+ i\sin v)$$ 
Эндээс (\ref{1.12}) ёсоор
$$r(\cos\varphi +i\sin \varphi)=\rho^n(\cos nv+ i\sin nv)$$
болж хоёр комплекс тооны тэнцэх нөхцөлөөр
$$r=\rho^n, \qquad nv=\varphi +2k\pi $$
бөгөөд улмаар
$$W=\sqrt[n]{z}=\sqrt[n]{r} \left(\cos \frac{\varphi +2k\pi}{n}+i\sin \frac{\varphi+2k\pi}{n}\right), \quad k=0, \pm 1, \pm 2 ,...$$ болно. \\
Үүний бие биеэсээ ялгаатай язгуурын тоо нь $k=0,1,2,\ldots ,n-1$ утгуудад харгалзах $n$ ширхэг гэж хялбархан тогтоогдоно. Эдгээр язгуур нь координатын эх дээр төвтэй $\sqrt[n]{|z|}$ радиустай дугуйд багтсан зөв $n$-өнцөгтийн оройн цэгүүд болно. Эндээс 
\begin{align*}
    \eta_k&=\sqrt[n]{z}=\sqrt[n]{|z|}\left(\cos \frac{\arg  z+2k\pi}{n}+i\sin \frac{\arg  z+2k\pi}{n}\right) \\
    &=\sqrt[n]{r}e^{\frac{\varphi+2k\pi}{n}},
\end{align*}
$$k=0,1,2,\ldots ,n-1$$
(\ref{1.12}) ба (\ref{1.13})-ыг Муавр (1667-1754)-ын томьёо \index{Муаврын томьёо} гэнэ. 

(\ref{1.10})-аас 
$$|z_1\cdot z_2|=|z_1|\cdot |z_2|,$$
$$\Arg  (z_1\cdot z_2)=\Arg  z_1+\Arg  z_2$$
(\ref{1.11})-ээс 
$$\left|\frac{z_1}{z_2}\right|=\frac{|z_1|}{|z_2|}, \quad (z_2\neq 0)$$
(\ref{1.12})-оос
$$|z^n|=|z|^n,$$
$$\Arg  (z^n)=n\cdot \Arg  z$$
мөрдөлгөөнүүд гаргаж болно.
\begin{exm}
$z=-1-i\sqrt{3}$ тоог тригонометр хэлбэрт бич.
\end{exm}
\begin{sol}
$|z|=2$, $\arg(-1-i\sqrt{3})=-\pi +\arctan \frac{-\sqrt{3}}{-1}=-\frac{2\pi}{3}$ тул 
$$-1-i\sqrt{3}=2\left[\cos\left(-\frac{2\pi}{3}+2k\pi\right) +i\sin \left(-\frac{2\pi}{3}+2k\pi\right)\right],  $$
$$k=0,\pm 1, \ldots \quad .$$
\end{sol}
\begin{exm}
$z^{n-1}=\overline{z}$ нөхцөлд тохирох бүх $z\quad (z\neq 0)$ тоонуудыг ол.
\end{exm}
\begin{sol}
$z=re^{i\varphi}$ гэвэл $r^{n-1}\cdot e^{i(n-1)\varphi}=re^{-i\varphi}$ буюу \\ $r^{n-2}e^{in \varphi}=1$, эндээс $r^{n-2}=1$, $in\varphi =2k\pi i$, өөрөөр хэлбэл $$r=1, \qquad \varphi=\dfrac{2k\pi}{n}, \qquad k=0,1,\ldots, n-1.$$ Иймд $z_k=e^{i\frac{2k\pi}{n}}$, $k=0,1,\ldots n-1.$
\end{sol}
\begin{exm}
$(\sqrt{3}-3i)^6$ илэрхийллийн утгыг ол.
\end{exm}
\begin{sol}
$|\sqrt{3}-3i|=\sqrt{12}$ 
$$\arg  (\sqrt{3}-3i)=\arctan \frac{-3}{\sqrt{3}}=-\arctan \frac{\sqrt{3}}{1}=-\frac{\pi}{3}.$$
Иймд $(\sqrt{3}-3i)^6=(\sqrt{12})^6\cdot e^{-6\cdot \frac{\pi}{3}i}=1728. $
\end{sol}
\begin{exm}
$\sqrt{-1}$ утгыг ол.
\end{exm}
\begin{sol}
$$\sqrt{-1}=1\cdot \left(\cos\frac{\pi+2k\pi}{2}+i\sin\frac{\pi +2k\pi}{2}\right), \qquad k=0,1$$ тул
$$\eta_0=\cos\frac{\pi}{2}+i\sin\frac{\pi}{2}=i, \qquad \eta_1=\cos\frac{2\pi}{2}+i\frac{3\pi}{2}=-i$$
\end{sol}
\newpage
\section*{Семинар I} 
\begin{enumerate}
\item Дараах комплекс тоонуудыг тригонометр, илтгэгч хэлбэрт бич.
\begin{enumerate}
\begin{multicols}{2}
        \item $-2$
        \item $2i$
        \item  $2$
        \item $-i$
        \item $-\sqrt{2}+i\sqrt{2}$
        \item $-1+i\sqrt{3}$
        \item $(1-\sqrt{3}i)^{30}$
        \item $(2-2i)^7$
        \item $\left(\dfrac{1+3i}{1-i}\right)^4$
        \item $\sqrt{-2-2\sqrt{3}i}$
\end{multicols}
    \end{enumerate}
    \item Үйлдлийг гүйцэтгэн бодит ба хуурмаг хэсэг, модуль, аргументийн гол утгыг олж тригонометр хэлбэрт бич.
\begin{enumerate}
\begin{multicols}{2}
\item $\displaystyle \frac{2+i}{3-2i}-\frac{2-i}{3+2i}$
\item $\displaystyle \frac{12+i}{3-i}-\frac{2+3i}{3+i}$
\item $\displaystyle \frac{2+3i}{3-2i}+\frac{3-i}{3+2i}$
\item $\displaystyle \frac{4+3i}{1-2i}+\frac{5-3i}{2-i}$
\item $\displaystyle \frac{7-i}{3i+4}-\frac{4-3i}{1+2i}$
\item $\displaystyle \frac{7+i}{3-4i}+\frac{3-5i}{3+2i}$
\item $\displaystyle \frac{-7+i}{3-4i}+\frac{2+3i}{3+2i}$
\item $\displaystyle \frac{12-5i}{3+4i}-\frac{3-i}{3-2i}$
\item $\displaystyle \frac{-12-5i}{4-3i}-\frac{5-4i}{2+3i}$
\item $\displaystyle \frac{-12+5i}{3-4i}+\frac{2-3i}{3+2i}$
\end{multicols}
\end{enumerate}
\item Үйлдлийг гүйцэтгэж илтгэгч хэлбэрт бич. (Комплекс язгуур ол.)
\begin{enumerate}
\begin{multicols}{2}
\item $\sqrt[4]{-1}$
\item $\sqrt[3]{i}$
\item $\sqrt[4]{16i}$
\item $\sqrt[3]{\dfrac{-1+i\sqrt{3}}{2}}$
\item $\sqrt[4]{16}$
\item $\sqrt[4]{-\dfrac{1}{16}}$
\item $\sqrt[3]{\dfrac{i}{8}}$
\item $\sqrt[3]{-\dfrac{1}{8}}$
\item $\sqrt[3]{27i}$
\item $\sqrt[3]{\dfrac{i}{27}}$
\end{multicols}
\end{enumerate}
\item Өгөгдсөн тэнцэтгэл бишүүдээр тодорхойлогдох олонлогийг дүрсэл.
\begin{enumerate}
\item $|z-1|\le 1,\quad |z+1|>2$
\item $|z+1|\ge 1,\quad |z+i|<1$
\item $|z-1-i|\le 1,\quad \mathrm{Im}\,z>1,\ \mathrm{Re}\,z\ge 1$
\item $|z-1-i|\ge 1,\quad 0\le \mathrm{Re}\,z<2,\ 0<\mathrm{Im}\,z\le 2$
\item $|z-1-i|<1,\quad |\arg z|\le \dfrac{\pi}{4}$
\item $|z|\ge 2,\quad |\arg z|\le \dfrac{\pi}{3}$
\item $|z+1-i|<1,\quad |\arg z|<\dfrac{\pi}{4}$
\item $\mathrm{Re}\,z<\mathrm{Im}\,z,\quad \dfrac{\pi}{4}<\arg z<\pi$
\item $|z-1|-|z-1|<3$
\item $|z-a|<r,\quad |z-b|<r;\quad a<b<a+r$
\end{enumerate}
\item Муруйн дүрс тодорхойлж, зураг зур.
\begin{enumerate}
\begin{multicols}{2}
\item $\mathrm{Im}\,z^2=3$
\item $\mathrm{Im}(z^2-z)=2-\mathrm{Im}\,z$
\item $|z-i|+|z+i|=3$
\item $|z|-3\,\mathrm{Im}\,z=3$
\item $|z-2i|=|z+i|$
\item $z\bar z+i(z-\bar z)-2=0$
\item $|z-1|=|\mathrm{Re}\,z|$
\item $|z-1|-|z-2+i|=6$
\item $\arg\left(\dfrac{z-1}{z-2i}\right)=\dfrac{\pi}{3}$
\item $\mathrm{Re}\,z+\mathrm{Im}\,z=3$
\end{multicols}
\end{enumerate}
\end{enumerate}
\newpage
\section*{Гэрийн даалгавар I} 
\begin{enumerate}
\item Дараах комплекс тоонуудыг тригонометр, илтгэгч хэлбэрт бич.
\begin{enumerate}
\begin{multicols}{2}
        \item $-0.5$
        \item $2i+2\sqrt{3}$
        \item  $e$
        \item $1-i$
        \item $-\sqrt{2}+i\sqrt{2}$
        \item $-1+i\sqrt{3}$
        \item $(1+\sqrt{3}i)^{30}$
        \item $(2+2i)^7$
        \item $\left(\dfrac{1-3i}{1-i}\right)^4$
        \item $\sqrt{2-2\sqrt{3}i}$
\end{multicols}
    \end{enumerate}
    \item Үйлдлийг гүйцэтгэн бодит ба хуурмаг хэсэг, модуль, аргументийн гол утгыг олж тригонометр хэлбэрт бич.
\begin{enumerate}
\begin{multicols}{2}
\item $\displaystyle \frac{2+3i}{3-2i}-\frac{3-i}{3+2i}$
\item $\displaystyle \frac{4+3i}{1-2i}-\frac{5-3i}{2-i}$
\item $\displaystyle \frac{2+2i}{3+4i}+\frac{2-i}{1+2i}$
\item $\displaystyle \frac{7+i}{3-4i}-\frac{3-5i}{3+2i}$
\item $\displaystyle \frac{-7+i}{3-4i}-\frac{2+3i}{3+2i}$
\item $\displaystyle \frac{1+7i}{3+4i}+\frac{3-i}{3+2i}$
\item $\displaystyle \frac{-7-i}{4-3i}+\frac{5-3i}{2+3i}$
\item $\displaystyle \frac{5-12i}{3i+4}-\frac{4+3i}{1+2i}$
\item $\displaystyle \frac{12-5i}{3+4i}+\frac{3-i}{3-2i}$
\item $\displaystyle \frac{-12-5i}{4-3i}+\frac{5-4i}{2+3i}$
\end{multicols}
\end{enumerate}
\item Үйлдлийг гүйцэтгэж илтгэгч хэлбэрт бич. (Комплекс язгуур ол.)
\begin{enumerate}
\begin{multicols}{2}
\item $\sqrt[4]{\dfrac{-1+i\sqrt{8}}{2}}$
\item $\sqrt[4]{1}$
\item $\sqrt[3]{-i}$
\item $\sqrt[3]{8}$
\item $\sqrt[3]{-8}$
\item $\sqrt[4]{-8+8\sqrt{3}\,i}$
\item $\sqrt[4]{\dfrac{1}{16}}$
\item $\sqrt[3]{-\dfrac{i}{8}}$
\item $\sqrt[3]{-\dfrac{1}{27}}$
\item $\sqrt[3]{-\dfrac{i}{27}}$
\end{multicols}
\end{enumerate}
\item Өгөгдсөн тэнцэтгэл бишүүдээр тодорхойлогдох олонлогийг дүрсэл.
\begin{enumerate}
\item $|z+i|\ge 1,\quad |z|>2$
\item $|z+1|<1,\quad |z-i|\le 1$
\item $|z-1+i|\ge 1,\quad \mathrm{Re}\,z<1,\ \mathrm{Im}\,z\ge -1$
\item $|z+i|<2,\quad 0<\mathrm{Re}\,z\le 1$
\item $|z+i|>1,\quad -\dfrac{\pi}{4}\le \arg z<0$
\item $|z-2|<2,\quad -\pi<\arg z\le -\dfrac{\pi}{3}$
\item $|z|<1,\quad \mathrm{Re}\,z<0,\quad \mathrm{Im}\,z<0$
\item $\mathrm{Re}\,z+\mathrm{Im}\,z>1,\quad \arg z<\dfrac{\pi}{4}$
\item $|z+i|+|z-i|>3,\quad \mathrm{Re}\,z>0$
\item $z\bar z<1,\quad |z-1|<1,\quad \mathrm{Im}\,z<0$
\end{enumerate}
\item Муруйн дүрс тодорхойлж, зураг зур.
\begin{enumerate}
\begin{multicols}{2}
\item $\mathrm{Im}\,\dfrac{1}{z}=\dfrac{1}{2}$
\item $z^2+\bar z^{\,2}=1$
\item $|z-i|-|z+i|=2$
\item $|z-1+i|=|z-1-i|$
\item $\mathrm{Re}(z^2-\bar z)=0$
\item $\mathrm{Re}\,\dfrac{z-1}{z+1}=0$
\item $|z-1|=|z-2|$
\item $6=|z-3|$
\item $\mathrm{Re}(z^2-4)=0$
\item $|z|=\mathrm{Re}\,z+1$
\end{multicols}
\end{enumerate}
\end{enumerate}
%\begin{center}
    \Large{\textbf{Лекц 2: Комплекс тооны зэрэг, Комплекс хавтгай дээрх топологи}}
\end{center}
\begin{center}
    \textbf{Комплекс тооны комплекс зэрэг, логарифм}
\end{center}
\begin{thm}
\begin{equation}\label{eulereq1}
    \cos\varphi+i\sin\varphi=e^{i\varphi}
\end{equation}
    байна.
\end{thm}
\begin{prf}
   $ \varepsilon(\varphi)=\cos\varphi+i\sin\varphi$ функцийн уламжлалыг авч үзье. 
   $$\dfrac{d \varepsilon(\varphi)}{d\varphi}=\dfrac{d (\cos\varphi+i\sin\varphi)}{d\varphi}=-\sin\varphi+i\cos\varphi=i(\cos\varphi+i\sin\varphi)$$ болно. Иймд $\varepsilon(\varphi)$ функц 
   \begin{equation}\label{difeq1}
      y^{\prime}=iy, \quad y(0)=1    
   \end{equation}
  дифферинциал тэгшитгэлийн шийд болж байна. (\ref{difeq1}) тэгшитгэл $y=e^{i\varphi}$ цор ганц шийдтэй.
\end{prf}
(\ref{eulereq1})-г Эйлерийн томъёо гэдэг. Бид энэ томъёог ашиглан комплекс тооноос логарифм авах, зэрэгт дэвшүүлэх үйлдэл хийж болно.

$z=(x,y)=re^{i\varphi} , -\pi<\varphi\leq \pi$ байг.  Тэгвэл $z=e^{\ln r+i\varphi}$ гэж бичиж болох тул $$z^{a+ib}=(e^{\ln r+i\varphi})^{a+ib}=e^{(a\ln r-b\varphi)+i(a\varphi+b\ln r)}$$
 болно. Мөн $$\ln z=\ln r+i\varphi$$ гэж тодорхойлоё.
 \begin{exm}
     $(1+i)^{i}$-г тооцоолоё. 
$(1+i)^{i} = e^{i \ln(1+i)}$

Одоо, $\ln(1+i)$-г тооцъё. Үүний тулд $1+i$-г экспоненциал хэлбэрт шилжүүлье:

$1+i = \sqrt{2} (\cos \frac{\pi}{4} + i \sin \frac{\pi}{4}) = \sqrt{2} e^{i \frac{\pi}{4}}$

Ингээд натурал логарифм авбал:

$\ln(1+i) = \ln(\sqrt{2}) + i \frac{\pi}{4} = \frac{1}{2} \ln 2 + i \frac{\pi}{4}$

Энд $\ln(\sqrt{2}) = \frac{1}{2} \ln 2$ болохыг ашигласан.

Ингээд үр дүнг эхний томъёонд орлуулбал:

$(1+i)^{i} = e^{i (\frac{1}{2} \ln 2 + i \frac{\pi}{4})} = e^{-\frac{\pi}{4}} e^{i \frac{1}{2} \ln 2}$

Эйлерийн томъёог ашиглан задалбал:

$e^{i \frac{1}{2} \ln 2} = \cos (\frac{1}{2} \ln 2) + i \sin (\frac{1}{2} \ln 2)$

Тэгвэл,

$(1+i)^{i} = e^{-\frac{\pi}{4}} [\cos (\frac{1}{2} \ln 2) + i \sin (\frac{1}{2} \ln 2)]$

$= e^{-\frac{\pi}{4}} \cos (\frac{1}{2} \ln 2) + i e^{-\frac{\pi}{4}} \sin (\frac{1}{2} \ln 2)$

Ийнхүү $(1+i)^{i} \approx 0.7688 + 0.3592i$ болно.
 \end{exm}
\section{Битүү комплекс хавтгай \index{Битүү комплекс хавтгай} \\ Риманий бөмбөлөг}

Ойлгоцыг гүнзгийрүүлэхийн үүднээс комплекс тоон $\mathbb{C}$ хавтгайг компактчилах шаардлага гардаг. Энэ нь төгсгөлгүй алсын цэг \index{Төгсгөлгүй алсын цэг} гэж нэрлэгдэх $z=\infty$ гэсэн тэмдэгэн элементээр түүнийг өргөтгөх явдал юм. Компактчилагдсан\index{ Компактчлал} хавтгайг битүү комплекс хавтгай гэх бөгөөд үүнийг $\overline{\mathbb{C}}$-ээр тэмдэглэнэ. Харин $\mathbb{C}$-ийг цаашид задгай хавтгай гэе. Иймд $\mathbb{C}\cup \{\infty\}=\overline{\mathbb{C}}.$ $z=\infty$ дээр алгебр үйлдэл гүйцэтгэх дүрэм өвөрмөц.

$a$ ба $b (b\neq 0)$ төгсгөлөг комплекс тоонууд бол 
$$\infty \pm a=a\pm \infty =\infty , $$ 
$$a \cdot \infty =\infty \cdot a = \infty \cdot \infty=\infty, $$ 
$$\dfrac{a}{\infty}=0 , \dfrac{\infty}{a}=\infty , \dfrac{b}{0}=\infty$$ бөгөөд $\infty \pm \infty$, $0\cdot \infty$, $\dfrac{0}{0}$, $\dfrac{\infty}{\infty}$ нь утгагүйд тооцогдоно.

$\infty$ тооны геометр дүрслэлийг гарган авахын тулд комплекс тоог бөмбөлгийн цэг болгон тавих хэрэгтэй юм.

Энэ зорилгоор $\mathbb{C}$-хавтгайн $O$ цэгт төвтэй, нэгж радиусгай $S$ бөмбөлөг авч улмаар огторгуйн тэгш өнцөгт декарт  $O\xi \eta \zeta$ системийн $O\xi$, $O\eta$ тэнхлэгүүд харгалзан $Ox$, $Oy$-тэй давхцаж, $O\zeta$ нь хойд туйл $N$ цэгийг дайрахаар байгуулалт хийе (Зураг 3). 
\begin{figure}[h]
    \centering
   \begin{tikzpicture}[line cap=round,line join=round,>=triangle 45,x=1cm,y=1cm]
\clip(0,0) rectangle (8,6.5);
\coordinate (o) at (3,3);
\coordinate (x1) at (1,1);
\coordinate (y1) at (6.5,3);
\coordinate (z1) at (3,5.9);
\coordinate (z) at (6,1);
\coordinate (n) at (3,4.6);
\coordinate (Z) at (5,3);
\coordinate (m) at (5,3);
\draw[-stealth,line width=1pt] (o)  -- (x1) node[left] {$x,\xi$};
\draw[-stealth,line width=1pt] (o)  -- (y1) node[right] {$y,\eta$};
\draw[-stealth,line width=1pt] (o)  -- (z1) node[right] {$\zeta$};
\draw (o) node[anchor=south east] {$O$};
\fill(n) circle (0.05);
\fill(z) circle (0.05);
\draw (n) node[anchor=south west] {$N$};
\draw[black](3,3) ellipse (2cm and 1.6cm);
\draw[dashed] (5,3) arc (0:180:2cm and 0.6cm);
\draw (1,3) arc (-180:0:2cm and 0.6cm);
\draw ($(z)!-0.1!(n)$) -- ($(z)!0.75!(n)$);
\draw[dashed] ($(z)!0.75!(n)$) -- (n);
\draw ($(n)!-0.1!(z)$) -- (n);
\draw ($(z)!0.75!(n)$) node[right] {$Z(\xi,\eta,\zeta)$};
\draw (z) node[right] {$z(x,y,0)$};
\fill ($(z)!0.75!(n)$) circle (0.05);
\end{tikzpicture}
     \captionof{figure}{}
    \label{fig:my_label}
\end{figure}
Иймд 
\begin{equation} \label{1.14}
S:\xi ^2+\eta^2+\zeta^2=1
\end{equation}
бөгөөд $\forall z\in\mathbb{C}$ цэгийг $N$-тэй холбоход $Z$ цэгээр $S$-ийг цоолно. $N(0,0,1)$,$Z(\xi, \eta, \zeta)$, $z(x,y,0)$ цэгүүд  шулуун дээр орших учир $$\dfrac{\xi-0}{x-0}=\dfrac{\eta -0}{y-0}=\dfrac{\zeta-1}{0-1}=t$$ болж, $Nz$ шулууны параметрт тэгшитгэл:
\begin{equation} \label{1.15}
\xi=xt, \qquad \eta =yt,\qquad \zeta =1-t
\end{equation}
(\ref{1.15})-ыг (\ref{1.14})-т тавибал $t=\dfrac{2}{1+|z|^2}$ бөгөөд улмаар
\begin{equation} \label{1.16}
\xi=\dfrac{2x}{1+|z|^2}, \quad \eta= \dfrac{2y}{1+|z|^2}, \quad \zeta=\dfrac{|z|^2-1}{1+|z|^2}
\end{equation}
болно. (\ref{1.16})-ийн сүүлчийн тэнцлээс $1-\zeta =\dfrac{2}{1+|z|^2}$ тул 
\begin{equation} \label{1.17}
x=\dfrac{\xi}{1-\zeta}, \qquad y=\dfrac{\eta}{1-\zeta}
\end{equation}
(\ref{1.16}), (\ref{1.17}) томьёонуудаас харахад $z\rightarrow Z$ харгалзаа $\mathbb{C}$ олонлогийг $S\setminus N$ олонлогт харилцан нэгэн утгатай буулгах нь илэрхий. $N$ цэгт $z=\infty$ цэг харгалзсан гэж үзвэл  $\overline{\mathbb{C}}$ олонлог $S$-тэй харилцан нэгэн утгатай. $S$-ийг Риманий бөмбөлөг \index{Риманий бөмбөлөг} гэнэ.

Комплекс тоог геометрт дүрслэх хоёр аргатай уялдаж, $\mathbb{C}$ олонлогт хоёр янзын зай оруулна. Нэг нь ердийн евклид зай:\index{Евклид зай} $z_1=(x_1,y_1)$ ба $z_2=(x_2,y_2)$ цэгүүдийн хоорондох зайг
\begin{eqnarray} \label{1.18}
 |z_2-z_1|=\sqrt{(x_2-x_1)^2+(y_2-y_1)^2}
\end{eqnarray}
гэж, нөгөө нь бөмбөлгийн зай \index{Бөмбөлгийн зай}: тэдгээр цэгүүдийн $S$ бөмбөлөг дээрх $Z_1$ ба $Z_2$ дүрүүдийн хоорондын евклид зайг
\begin{equation} \label{1.19}
\rho(z_1,z_2)=\dfrac{2|z_2-z_1|}{\sqrt{1+|z_1|^2}\cdot \sqrt{1+|z_2|^2}}
\end{equation}
гэж авна. (\ref{1.19}) нь (\ref{1.16}) томьёонууд ашиглан жижигхэн гаргалгаа хийсний дараах бичиглэл.

$\mathbb{C}$ олонлог дээр бичигдсэн (\ref{1.19}) томьёог $z$ ба $\infty$-ийн хувьд
\begin{equation}
\rho(z,\infty)=\dfrac{2}{\sqrt{1+|z|^2}} \label{1.20} \footnote{(\ref{1.20}) томьёо нь (\ref{1.19}) дээр $z_1=z$ гэж аваад хүртвэр хуваарийг $z_2$-т хувааж $z_2\rightarrow \infty$-ээр тооцож гаргасан.}
\end{equation}
гэж аваад $\overline{\mathbb{C}}$ олонлогт дэлгэрүүлж болно.

$\forall z_1,z_2\in \overline{\mathbb{C}}$ авахад $\rho (z_1,z_2)\leq 2$ байх нь мэдээж. (\ref{1.18}), (\ref{1.19}) зайн алийг нь ч авсан $\mathbb{C}$ олонлог метрик огторгуй үүсгэнэ. Өөрөөр хэлбэл зайн ердийн аксиомууд биелэгдэнэ. Тухайлбал (\ref{1.18}) зайн хувьд гурвалжны аксиом нь
$$|z_1+z_2|\leq |z_1|+|z_2|. \qquad (2(a)- \text{г \ үз})$$

Бэхлэн авсан $\{|z|<R\}$, $R<\infty$ дугуйд орших зааглагдсан $M\subset \mathbb{C}$ олонлогт Евклид ба бөмбөлөг зайнууд чацуу чанартай байна. Учир нь хэрэв $M\subset \{|z|\leq R\}$ гэвэл $\forall z_1$, $z_2\in M$-ийн хувьд (\ref{1.19}) томьёо ёсоор
\begin{equation}
\dfrac{2|z_2-z_1|}{1+R^2}\leq \rho (z_1,z_2) \leq 2|z_2-z_1|
\end{equation}
байх хоёрлосон тэнцэл биш биелэгдэнэ. Иймд бөмбөлгийн зай нь зааглагдаагүй олонлогт ердийн байдлаар хэрэглэгдэнэ.

$\mathbb{C}$ дахь олонлогт евклид, харин $\overline{\mathbb{C}}$ дахид бөмбөлөг зай авдаг.

Комплекс тоон $\mathbb{C}$ ба $\overline{\mathbb{C}}$ олонлогт зай оруулах замаар метрик огторгуй үүсгэхтэй танилцлаа. Одоо эдгээр олонлогт топологи тодорхойлоё.

$\varepsilon >0$ дурын тоо байг. $z_0\in\mathbb{C}$ цэгийн $\varepsilon$-орчин \index{$\varepsilon$-орчин} $U(z_0,\varepsilon)$ гэж евклид зайн тохиолдолд энэ $z_0$ цэгт төвтэй $\varepsilon$ радиустай дугуйг, өөрөөр хэлбэл
\begin{equation}
|z-z_0|<\varepsilon
\end{equation}
тэнцэтгэл биш хангах  $z\in \mathbb{C}$ цэгүүдийн олонлогийг ойлгоно. Харин $z_0\in\overline{\mathbb{C}}$ цэгийн $\varepsilon$-орчин гэж
\begin{equation}
\rho (z,z_0)<\varepsilon
\end{equation}
нөхцөлд тохирох $z\in \overline{\mathbb{C}}$ цэгүүдийн олонлогийг хэлнэ. 

(\ref{1.20}) ёсоор $\rho (z,\infty)<\varepsilon$ нь $|z|>\sqrt{\dfrac{4}{\varepsilon^2}-1}$-тэй адил тул хавтгайн төгсгөлгүй алсын цэгийн $\varepsilon$-орчин гэдэг нь координатын эхэд төвтэй дугуйн гаднах цэгүүдийн олонлог ($z=\infty$ цэгийг агуулсан) юм.

Хэрэв $\forall z_0\in E\subset \mathbb{C}$ (буюу $\overline{\mathbb{C}}$) цэг авахад $U(z_0,\varepsilon)\subset E$ байх $\varepsilon$-орчин олдож байвал $E$-г задгай олонлог гэнэ. Ингэж задгайн тухай ойлголт оруулснаар $\mathbb{C}$ ба $\overline{\mathbb{C}}$ олонлог топологи огторгуй үүсгэхийг хялбархан харуулж болно. $\mathbb{C}$ ба $\overline{\mathbb{C}}$ олонлогт 
\begin{equation}
0<|z-z_0|<\varepsilon , \quad 0<\rho (z,z_0)<\varepsilon
\end{equation}
байх $z$-ийн олонлогийг $z_0$ цэгийн цоорхойт орчин \index{Цоорхойт орчин} гэнэ.
\begin{dfn}
$z_0\in\mathbb{C}$  (тэрчлэн $\overline{\mathbb{C}}$) цэгийн цоорхойт аливаа орчинд $\mathbb{C}$  (тэрчлэн $\overline{\mathbb{C}}$)-ийн топологи утгаар $M\subset \mathbb{C}$ (тэрчлэн $\overline{\mathbb{C}}$) олонлогоос ядаж нэг цэг олдож байвал $z_0$ цэгийг $M$ олонлогийн хязгаарын цэг гэнэ. Өөрийнхөө хязгаарын бүх цэгээ агуулсан $M$ олонлогийг битүү олонлог \index{Битүү олонлог} гэнэ. $M$ олонлогт өөрийн нь хязгаарын бүх цэгийг нэгтгэн гаргасан олонлогийг $M$-ийн битүүрэлт\index{Олонлог битүүрэл} гэх ба түүнийг $\overline{M}$ гэж тэмдэглэнэ. 
\end{dfn}
\begin{exm}
$M=\{0,\pm 1,\pm 2,\ldots \}$ нь $\mathbb{C}$-д хязгаарын цэг байхгүй (тул битүү олонлог), харин $\overline{\mathbb{C}}$-д өөртөө үл агуулагдах хязгаарын ганц $z=\infty$ цэгтэй (иймд битүү биш олонлог).

$\overline{\mathbb{C}}$ дэх төгсгөлгүй аливаа олонлог ядаж нэг хязгаарын цэгтэй байна (Компактын зарчим) \index{Компактын зарчим}. $\mathbb{C}$-д компактын зарчим биелэгдэхгүй (Энэ нь өмнөх жишээнээс илэрхий). Харин зааглагдсан, төгсгөлгүй, өөрөөр хэлбэл, ямар нэг $\{|z|\leq R\}$, $R<\infty$ орших олонлогийн хувьд энэ зарчим биелэгдэнэ. Ийм олонлогийг компакт олонлог \index{Компакт олонлог}гэнэ.
\end{exm}
\begin{exm}
Риманы $S$ бөмбөлөг дээр $z_0=1-\sqrt{3}i$ цэгийн дүрийг ол.
\end{exm}
\begin{sol}
(\ref{1.16}) томьёог ашиглая. $|z_0|=|1-\sqrt{3}i|=2$, \\ $\xi =\dfrac{2\cdot 1}{1+4}=\dfrac{2}{5}$, $\eta =\dfrac{-2\sqrt{3}}{5}$, $\zeta =\dfrac{3}{5}$. Иймд $z_0$ цэг $\left(\dfrac{2}{5}, \dfrac{-2\sqrt{3}}{5}, \dfrac{3}{5}\right)$ цэгт буулгана.
\end{sol}
\begin{exm}
$z$ цэг $S$ бөмбөрцөг дээр $(\xi, \eta,\zeta)$ цэгт буух бол $\overline{z}$, $-z$ цэгүүд юунд буух вэ?
\end{exm}
\begin{sol}
$\overline{z}=x-iy$ үед түүний $S$ дээрх дүр нь:
$$\left(\frac{2x}{1+|z|^2}, \frac{-2y}{1+|z|^2}, \frac{|z|^2-1}{1+|z|^2}\right)=(\xi, -\eta, \zeta)$$
$$-z=-(x+iy) \ \text{үед} \ \left(-\frac{2x}{1+|z|^2}, \frac{-2y}{1+|z|^2}, \frac{|z|^2-1}{1+|z|^2}\right)=(-\xi, -\eta, \zeta).$$
\end{sol}
\begin{center}
   \textbf{ Дасгал}
\end{center}
\begin{enumerate}
    \item [7.] $z_1=2$, $z_2=-i$, $z_3=-1+2i$ цэгүүдийн $S$ дээрх дүрийг ол.
    \item [8.] $S$ бөмбөрцөг дээрх $(-2,2,-1)$ ба $(1,1,1)$ цэгүүдийн эх дүр болох цэгүүдийг ол.
    \item [9.] $|z|=r$ ба $y=x$ шугамуудын $S$-дээрх дүрийг ол.
\end{enumerate}
\section{Семинар 2}
\begin{enumerate}
    \item Тооцоол.
    \begin{enumerate}
        \item $\ln(3+4i)$
        \item $(3+4i)^{3+2i}$
        \item $(5+12i)^i$
    \end{enumerate}
     \item $z=x+iy , -\pi<y\leq \pi$ бол $e^z:=e^x\cos y+ie^x\sin y$ гэж тодорхойлоё. $$\ln(e^z)=z$$ батал.
   \item  $w^z=e^{z\cdot\ln w}$ гэе. Дараах тэнцэтгэлүүдийг батал.
   \begin{enumerate}
        \item $w^{z_1}\cdot w^{z_2}=w^{z_1+z_2}$
        \item $w^{z_1}: w^{z_2}=w^{z_1-z_2}$
        \item $(w^{z_1})^{z_2}=w^{z_1\cdot z_2}$
    \end{enumerate}
     \item $z\in\mathbb{C}$ үед $$\cos z:=\dfrac{e^{iz}+e^{-iz}}{2}$$
     $$\sin z:=\dfrac{e^{iz}-e^{-iz}}{2i}$$ гэе.
    \begin{enumerate}
        \item $\cos^z+\sin^2 z=1$ батал.
        \item $\cos 2z=\cos^2 z-\sin^2 z$ батал.
        \item $\sin 2z=2\sin z\cos z$ батал.
        \item $|\cos z|=10$  тэгшитгэл шийдтэй юу?
    \end{enumerate}
     \item 
    \begin{enumerate}
        \item $\mathbb{N}\subset\mathbb{C}$ хувьд $\mathbb{N}$ олонлог битүү, задгай, компакт эсэхийг тодорхойл. $\mathbb{N}^{o}, \overline{\mathbb{N}} , Lim(\mathbb{N})$ ол.
        \item $U(z,\varepsilon)$ задгай олонлог байхыг батал.
        \item $A$ задгай байх $\longleftrightarrow$ түүний бүх цэг дотоод цэг байх .
         \item $A$ битүү байх $\longleftrightarrow$ түүний бүх хязгаарын цэг нь $A$ д агуулагдах.
         \item $\partial A:=\overline{A}\cap\overline{A^c}$ гэе. $\partial U(0,1)$ ол.
    \end{enumerate}
\end{enumerate}

Гэрийн даалгавар:


\begin{enumerate}
    \item Тооцоол.

a) $\ln(2-3i)$\\
b) $(1+2i)^{2-i}$\\
c) $(4-3i)^{3i}$
\item
$z=x+iy$ бол $\cos(-z)=\cos z$  болохыг батал. $Re(\sin z) , Im(\sin z)$ ол.
\item
$w^z=e^{z\cdot\ln w}$ гэе. Дараах тэнцэтгэлүүдийг батал.
a) $(w_1\cdot w_2)^z=w_1^z\cdot w_2^z$\\
b) $(w_1:w_2)^z=w_1^z:w_2^z$\\
c) $(\sqrt[n]{w})^z=\sqrt[n]{w^z}$
\item
$z\in\mathbb{C}$ үед
$$\tan z:=\frac{\sin z}{\cos z}$$
$$\cot z:=\frac{\cos z}{\sin z}$$ гэе.
a) $\tan^2 z+1=\sec^2 z$ батал.\\
b) $\cot^2 z+1=\csc^2 z$ батал.\\
c) $\tan(z_1+z_2)=\frac{\tan z_1+\tan z_2}{1-\tan z_1\tan z_2}$ батал.\\
d) $|\sin z|=2$ тэгшитгэл шийдтэй юу?
\item
a) $\mathbb{Z}\subset\mathbb{C}$ хувьд $\mathbb{Z}$ олонлог битүү, задгай, компакт эсэхийг тодорхойл. $\mathbb{Z}^{o}, \overline{\mathbb{Z}} , Lim(\mathbb{Z})$ ол.\\
b) $B(z_0,r)=\{z\in\mathbb{C}| |z-z_0|\leq r\}$ битүү олонлог байхыг батал.\\
c) $A$ компакт байх $\longleftrightarrow$ $A$ битүү ба хязгаарлагдсан байх.\\
d) $A$ компакт байх $\longleftrightarrow$ $A$-ын бүх дэд дараалал нь $A$ дотроо хуримтлагдах дэд дараалалтай байх.\\
e) $\partial B(0,1)$ ол.
\end{enumerate}
\section{Дараалал, түүний хязгаар. \index{Дарааллын хязгаар} \\
Муруй. Муж}

Комплекс тоон $\{a_n\}$ дараалал нь сөрөг биш бүхэл тоон олонлогийг $\mathbb{C}$ (буюу $\overline{\mathbb{C}}$) хавтгай руу буулгадаг өвөрмөц функц юм. 
\begin{dfn}
$a\in\mathbb{C}$ (буюу $\overline{\mathbb{C}}$) цэгийн аливаа орчинд $\mathbb{C}$ (буюу  $\overline{\mathbb{C}}$)-ийн топологи утгаар $\{a_n\}$ дарааллын төгсгөлгүй олон элемент оршиж байвал уул дарааллыг $a$-хязгаартай гэж нэрлэнэ. $\{a_n\}$ дараалал $\mathbb{C}$-д цор ганц хязгаарын $a$ цэгтэй байвал $a$-руу нийллээ гэх бөгөөд
\begin{equation} \label{1.25}
\lim_{n\rightarrow \infty} a_n=a 
\end{equation}
гэж бичдэг.

Комплекс (\ref{1.25}) тэнцэл хоёр бодит тоон тэнцэлтэй чацуу чанартай. $a\neq \infty$ гэвэл $a_n\neq \infty$ байх тул
$$\lim_{n\rightarrow \infty} \re a_n=\re a, \qquad \lim_{n\rightarrow \infty} \im a_n=\im a$$
\end{dfn}
гэж хялбархан батлагдана.

Ямар нэгэн тасралтгүй $z=\gamma (t)$ буулгалтын $[\alpha, \beta]$ хэрчмийн дүр болох цэгүүдийн $\gamma\subset \overline{\mathbb{C}}$ олонлогийг муруй\index{Муруй} гэж ойлгоно. Өөрөөр хэлбэл муруй нь бодит $t\in [\alpha, \beta]$ аргументийн комплекс утгат тасралтгүй $z=\gamma (t)=x(t)+iy (t)$ функц.

$\gamma (\alpha)$ муруйн эхэн, $\gamma (\beta)$-төгсгөл гэнэ. Эхэн ба төгсгөл давхацвал уул муруйг битүү муруй\index{Битүү муруй} гэнэ. Параметрийн өсөлтөөр муруйн эерэг чиглэл тогтоогдоно. $[\alpha,\beta]$ хэрчмийн дотор орших ялгаатай хоёр буюу түүнээс дээш утгуудад муруйн ганц цэг харгалзаж байвал энэ цэгийг давхардлын цэг \index{Давхардлын цэг} гэнэ. Давхардлын цэггүй муруйг энгийн буюу Жорданы муруй \index{Жорданы муруй}гэнэ.

$\gamma: [\alpha, \beta]\rightarrow \mathbb{C}$ муруй $t\in [\alpha,\beta]$ бүр дээр тасралтгүй $\gamma '(t)= x' (t)+iy '(t)$ уламжлал оршин байвал түүнийг тасралтгүй дифференциалчлагдах муруй гэнэ. Харин $t\in [\alpha,\beta]$ бүхэн дээр
$\gamma ' (t)\neq 0$ нөхцөл (өөрөөр хэлбэл, онцгой цэггүй) биелэгдэж байвал гөлгөр муруй \index{Гөлгөр муруй} гэнэ. Бараг бүх $[\alpha, \beta]$ дээр $\gamma '(t)$ оршин байхын зэрэгцээ Лебегийн утгаар абсолют интегралчлагдаж \index{Лебегийн утгаар абсолют интегралчлагдах муруй} (өөрөөр хэлбэл 
$$\int\limits_{\alpha}^\beta |\gamma ' (t)|dt= \int\limits_\alpha ^\beta \sqrt{ [x'(t)]^2+[y'(t)]^2}dt$$
муруйн урт оршин байх) байвал уул муруйг шулуутгагдах муруй \index{Шулуутгагдах муруй} гэнэ.

Функцийн (муруйн) гөлгөр байхыг тодорхойлоход нийтлэг хэллэгийн нэр томьёо ашиглаж тасралтгүй функцийг $C^0$ анги, тасралтгүй дифференциалчлагдах функцийг \index{Тасралтгүй дифференциалчлагдах муруй} $C^1$ ба $n$-эрэмбийн уламжлалын хамтаар тасралтгүй функцийг $C^n$ ангийн функцүүд \index{$C^n$ ангийн функцүүд} гэх мэт цаашид нэрлэе.
\begin{dfn}
$\mathbb{C}$ (буюу $\overline{\mathbb{C}}$) цэгүүдийн дараах хоёр чанартай олонлог $D$-г муж \index{Муж} гэнэ.
\begin{enumerate}
    \item [a)] Аливаа $a\in D$ цэгийн хувьд энэ цэгийн $D$-д орших орчин олдох (задгай байх чанар)
    \item [б)] Аливаа $a,b\in D$ цэгүүдийн хувьд тэдгээрт төгсгөлүүдтэй зам $D$-д олдох (холбоост байх чанар).
\end{enumerate}

$D$-д үл орших атлаа $D$-ийн цэгүүдийн хязгаарын цэг (өөрөөр хэлбэл тэдгээрийн аливаа орчинд $D$-д орших цэгүүдээс гадна $D$-д үл орших ядаж нэг цэг олдож байх) болох $\overline{\mathbb{C}}$-ийн цэгийг $D$-ийн хилийн цэг \index{Хилийн цэг} гэнэ. $D$-ийн хилийн бүх цэгүүдийн цогцыг $D$-олонлогийн хил гэх бөгөөд үүнийг $\partial D$ гэж тэмдэглэнэ. \qquad $D\cup \partial D=\overline{D}$ байна.
\end{dfn}
\begin{thm}
Аливаа $D$ мужийн $\partial D$ хил битүү олонлог байна. 
\end{thm}
\begin{prf}
$\zeta _0$ нь $\partial D$ олонлогийн хязгаарын аль нэг цэг байг. $\zeta _0\in \partial D$ гэж батлах ёстой. $U(\zeta, \zeta_0)>0$ байх орчин авъя. Энэ орчинд $\zeta \in\partial D$ цэг оршино. Иймд $\zeta$ цэгт $V\subset U(\zeta, \zeta_0)$ байх орчин олдоно. $V$-д тэгэхлээр $U(\zeta, \zeta_0)$-д $D$-ийн ч, бас $D$-ийн биш ч цэгүүд олдоно. Иймд, $\zeta _0$ нь $D$-ийн хилийн цэг болох нь нотлогдов.
\hfill $\blacksquare$\end{prf}
Холбоост чанарын тухай өмнө оруулсан ойлголтыг өргөтгөе.
\begin{dfn}
Хэрэв $M$ олонлогийг хоосон биш  $M_1$, $M_2$ гэж $\overline{M}\cap M_2$ ба $M_1\cap \overline{M_2}$ огтлол бүр хоёулаа хоосон байхаар задлан тавьж болохгүй байвал $M$-ийг холбоост олонлог\index{Холбоост олонлог} гэнэ. Тухайлбал, битүү олонлогийг үл огтлолцсон битүү (бас хоосон биш) хоёр дэд олонлог болгон задлан тавьж болохгүй байвал түүнийг холбоост гэнэ. Битүү холбоост олонлогийг континум \index{Континум} гэж нэрлэнэ.
Хэрэв $\partial D$ хил холбоост олонлог байвал $D\subset \overline{\mathbb{C}}$ мужийг нэг холбоост \index{Олонлог нэг холбоост}гэнэ. Эсрэг тохиолдолд олон холбоост \index{Олонлог олон холбоост} гэнэ. Холбоосын тоо төгсгөлөг байвал энэ тоог $D$ мужийн холбоосын эрэмбэ \index{Холбоосын эрэмбэ} гэнэ.

$D$ мужийг агуулж байх $|z|<R$, $R<\infty$ дугуй оршин байвал түүнийг компакт муж \index{Компакт муж} гэнэ. $\overline{M}$ олонлог $D$-д агуулагдаж байвал $M$ нь $D$-д компактаар оршиж байна гэж ярьдаг. Компакт оршихыг $\subset$ гэж тэмдэглэдэг. Тэгэхлээр $\overline{M}\subset D$ гэвэл $M\subset D$ юм. $D$ нь компакт олонлог гэдгээс $D\subset \mathbb{C}$ байна.

Цаашид нэг биш удаа хэрэглэх дараах теорем авч үзье. 
\end{dfn}
\begin{thm}
$M\subset \overline{\mathbb{C}}$ нь холбоост олонлог, $N$ нь түүний хоосон биш дэд олонлог байг. Хэрэв $M$-ийн харьцангуй топологит \footnote{$\overline{\mathbb{C}}$-ийн топологи дахь цэгүүдийн орчны $M$-тай үүсгэх огтлол нь мөнхүү цэгүүдийн орчин байх топологийг $M\subset \overline{\mathbb{C}}$ олонлогийн харьцангуй топологи гэнэ.}  $N$ нь нэгэн зэрэг битүү, бас задгай байвал $M=N$ байна.
\end{thm}
\begin{prf}
Эсрэгээс нь, тухайлбал, $N'=M\setminus N$ олонлог хоосон биш гэе. $\overline{\mathbb{C}}$-ийн топологи дахь $\overline{N}$ битүүрэл нь $M$-ийн топологи дахь түүний $(\overline{N})_M$ битүүрлийн цэгүүдээс тогтох нь мэдээж юм.

Иймд $\overline{N}\cap N'=(\overline{N})_M\cap N'$, харин $N$ нь $M$-ийн топологит битүү тул $(\overline{N})_M=N$ бөгөөд эндээс $\overline{N}\cap N'=N\cap N'=\emptyset$.

$N$ нь $M$-ийн тэрхүү топологит битүү байна. ($N'$-ийн хязгаарын цэгүүд $N$-ийн задгай чанараар $N$-д агуулагдахгүй. Эндээс тэдгээр нь $N'$-т оршино). Иймд $\overline{N}\cap N'$ огтлолд хэрэглэж байсны адилаар сэтгэж $\overline{N}'\cap N$ огтлолд хандвал $\overline{N}'\cap N$ нь хоосон болов. Энэ нь $M$-ийн холбоостын тодорхойлолтод харшилна.
\hfill $\blacksquare$\end{prf}
\begin{exm}
Дараах параметрт тэгшитгэлүүдээр ямар шугам дүрслэгдэх вэ? 
\begin{enumerate}
    \item $z=a+(b-a)t$, \qquad $t\in [0,1];$
    \item $z=t+it^2$, \qquad $t\in (-\infty, 0);$
    \item $z=ae^{it}+\dfrac{1}{a}e^{-it}$, \quad $t\in (0, 2\pi]$, $a>1$;
\end{enumerate}
\end{exm}
\begin{sol}
\begin{enumerate}
    \item $\left.\begin{array} {c}
x(t)=a+(b-a)t, \\
y(t)=0
\end{array}\right\}$, $t\in [0,1]$ тул \\
$z=a(t=0),$ $z=b(t=1)$ байх учир $z(t):[a,b].$
\item 
Өгөгдсөнөөс харахад $ \left.
\begin{array}{c}
x(t)=t \\
y(t)=t^2
\end{array} \right\},$ $t\in (-\infty, 0)$ бөгөөд \\
$y=x^2$, $x<0$, өөрөөр хэлбэл $y=x^2$ параболын зүүн хагас; зурагдах чиглэл нь $z=\infty$ цэгээс $z=0$ рүү. 
\item 
\begin{align*}
    z&=a(\cos t +i\sin t)+\dfrac{1}{a}(\cos t-i\sin t) \\
    &=\left( a+\dfrac{1}{a}\right)\cos t+i\left(a-\dfrac{1}{a}\right)\sin t
\end{align*}
 тул 
 $$\left.\begin{array}{c}
x=x(t)=\left(a+\dfrac{1}{a}\right)\cos t, \\
y=y(t)=\left(a-\dfrac{1}{a}\right) \sin t
\end{array} \right\}, \quad t\in [0,2\pi], \quad a>1.$$

Эндээс нар буруу тойролттой 
$$\frac{x^2}{\left(a+\frac{1}{a}\right)^2}+\frac{y^2}{\left(a-\frac{1}{a}\right)^2}=1, \quad a>1$$  эллипс болой. 
\end{enumerate} 
\end{sol}
\begin{exm}
Дараах параметрт тэгшитгэлүүдээр $\partial D$ хил нь тодорхойлогдох $D$ муж дүрсэл.
\begin{enumerate}
    \item $z=a+Re^{i\varphi}, \quad 0\leq \varphi \leq 2\pi.$
    \item $z=-it$, \quad $-\infty <t<\infty$
    \item $z=t^2$, \quad $-\infty <t<\infty.$
\end{enumerate}
\end{exm}
\begin{sol}
\begin{enumerate}
    \item $|z-a|=R$ тул $D$ нь $a$ цэгт төвтэй, $R$ радиустай дугуй: $|z-a|<R.$
    \item $\left.\begin{array}{c}
x=0, \\
y=-t
\end{array}\right\}$, $t\in (-\infty, \infty)$ болж $OY$ тэнхлэгийн сөрөг чиглэл дүрслэгдэх тул $D$ нь баруун хагас хавтгай.\\
Учир нь: Сөрөг чиглэлтэй $OY$ шулуун нь $\infty$-цэгийг дайрсан $\partial D$ тойрог. $D$ муж $\partial D$-ийн тойрох чиглэлийн зүүн гарт үлдэх ёстойг тооцов. 
\item $z=t^2$, $t\in (-\infty, \infty)$ гэдгээс 
$\left.\begin{array}{c}
x=x(t)=t^2\geq 0, \\
y= y(t)=0
\end{array}\right\}$, $t\in (-\infty, \infty)$ бөгөөд эндээс $z=x\geq 0$ болж $t\in (-\infty, 0)$ ба $t\in (0,\infty)$ тохиолдолд бодит тэнхлэгийн эерэг чиглэлд давхацсан хоёр цацраг үүснэ. Эдгээр цацраг параметрийн өсөх дарааллаар $z=\infty$ цэгээс $z=0$ цэгийг дайран $z=\infty$ цэгт дахин хүрэх аж. Иймд $D$ муж: $0<\arg  z<2\pi $ байна. Өөрөөр хэлбэл $OX$ тэнхлэгийн эерэг чиглэлд хагалбар бүхий бүх хавтгай болой
\end{enumerate}
\end{sol}
\begin{center}
   \textbf{ Дасгал}
\end{center}
\begin{enumerate}
    \item [10.]Дараах тэгшитгэлүүдээр өгөгдсөн шугамыг тодорхойл.
    \begin{enumerate}
        \item [a)] $z=a(\cos t+i\sin t)$, \quad $\dfrac{\pi}{2}\leq t\leq \dfrac{3\pi}{2}$, $a>0.$
        \item [б)] $z=t+i\sqrt{ 1-t^2}$, \quad $-1\leq t\leq 1.$
    \end{enumerate}
    \item [11. ] Дараах тэгшитгэлүүдээр $\partial D$ хил нь тодорхойлогдох $D$ муж дүрсэл.
    \begin{enumerate}
        \item [а)] $z=a+\rho e^{-it}$
        \item [b)] $z=t+it^2$, \quad $-\infty<t<\infty$
        \item [c)] $z=\dfrac{3t}{1+t^3}+i\dfrac{3t^2}{1+t^3}$, \quad $-1\leq t\leq 0.$
    \end{enumerate}
\end{enumerate}
\setcounter{chapter}{5}
%\chapter{Комплекс функцийн хязгаар ба тасралтгүй чанар}

\section{Комплекс хувьсагчийн функц}

\begin{dfn}
$D\subset\mathbb{C}$ олонлог дээр $f$ дүрэм өгөгдөж, $z\in D$ бүрт комплекс тоо $w$-г
$w=f(z)$ гэж харгалзуулж байвал $f$-г \emph{комплекс хувьсагчийн функц} гэнэ:
\[
f: D\to \mathbb{C},\qquad z\mapsto w=f(z).
\]
$D$-г \emph{тодорхойлогдох муж} гэнэ.
\end{dfn}

Хэрэв $z=x+iy$, $w=u+iv$ ($x,y,u,v\in\mathbb{R}$) гэвэл
\[
f(z)=u(x,y)+iv(x,y),
\]
өөрөөр хэлбэл комплекс функц нь хоёр хувьсагчийн бодит утгат $u(x,y)$, $v(x,y)$ функцүүдээр тодорхойлогдоно:
\[
u(x,y)=\Re f(z),\qquad v(x,y)=\Im f(z).
\]

\begin{exm}
$w=z^2$ функцийн бодит ба хуурмаг хэсгийг ол.
\end{exm}

\begin{sol}
$z=x+iy$ гэвэл
\[
z^2=(x+iy)^2=x^2-y^2+2xy\,i.
\]
Иймээс
\[
u(x,y)=x^2-y^2,\qquad v(x,y)=2xy.
\]
\end{sol}

%---------------------------------------------------

\section{Функцийн хязгаар}

\begin{dfn}
$z_0\in D$ байг. Хэрэв $A\in\mathbb{C}$ тоо олдож
\[
\forall \varepsilon>0\ \exists \delta>0:\ 
0<|z-z_0|<\delta \Rightarrow |f(z)-A|<\varepsilon
\]
нөхцөл биелбэл $A$-г $f(z)$-ийн $z\to z_0$ үед авах \emph{хязгаар} гэнэ.
Тэмдэглэгээ:
\[
\lim_{z\to z_0} f(z)=A.
\]
\end{dfn}

\begin{thm}[Хязгаарын бодит ба хуурмаг хэсэг]
$f(z)=u(x,y)+iv(x,y)$, $z_0=x_0+iy_0$, $A=\alpha+i\beta$ гэе.
Тэгвэл
\[
\lim_{z\to z_0} f(z)=A
\]
байх нь
\[
\lim_{(x,y)\to(x_0,y_0)}u(x,y)=\alpha,\qquad
\lim_{(x,y)\to(x_0,y_0)}v(x,y)=\beta
\]
байхтай эквивалент.
\end{thm}

\begin{proof}
$(\Rightarrow)$ гэж үзье. 
\[
\lim_{z\to z_0} f(z)=A
\]
гэж авъя. Тэгвэл $\forall \varepsilon>0$ тоо бүрт $\exists \delta>0$ олдож
\[
0<|z-z_0|<\delta \ \Rightarrow\ |f(z)-A|<\varepsilon
\]
биелнэ.

Энд
\[
f(z)-A=(u(x,y)-\alpha)+i(v(x,y)-\beta)
\]
учир
\[
|u(x,y)-\alpha|\le |f(z)-A|,\qquad |v(x,y)-\beta|\le |f(z)-A|
\]
гэж үнэлж болно. Иймээс дээрх нөхцлөөс
\[
0<|z-z_0|<\delta \ \Rightarrow\ |u(x,y)-\alpha|<\varepsilon
\]
мөн
\[
0<|z-z_0|<\delta \ \Rightarrow\ |v(x,y)-\beta|<\varepsilon
\]
гарна. Харин $|z-z_0|=\sqrt{(x-x_0)^2+(y-y_0)^2}$ тул энэ нь яг
\[
\lim_{(x,y)\to(x_0,y_0)}u(x,y)=\alpha,\qquad
\lim_{(x,y)\to(x_0,y_0)}v(x,y)=\beta
\]
гэдгийг илэрхийлнэ.

$(\Leftarrow)$ Одоо эсрэгээр нь
\[
\lim_{(x,y)\to(x_0,y_0)}u(x,y)=\alpha,\qquad
\lim_{(x,y)\to(x_0,y_0)}v(x,y)=\beta
\]
гэж үзье. Тэгвэл $\forall \varepsilon>0$ тоо бүрт $\exists \delta_1>0$, $\exists \delta_2>0$ олдож
\[
0<\sqrt{(x-x_0)^2+(y-y_0)^2}<\delta_1 \Rightarrow |u(x,y)-\alpha|<\frac{\varepsilon}{\sqrt2},
\]
\[
0<\sqrt{(x-x_0)^2+(y-y_0)^2}<\delta_2 \Rightarrow |v(x,y)-\beta|<\frac{\varepsilon}{\sqrt2}
\]
биелнэ. $\delta=\min\{\delta_1,\delta_2\}$ гэж авбал
\[
0<|z-z_0|<\delta \Rightarrow 
|u(x,y)-\alpha|<\frac{\varepsilon}{\sqrt2},\quad
|v(x,y)-\beta|<\frac{\varepsilon}{\sqrt2}.
\]

Тэгвэл
\[
|f(z)-A|^2
=
|(u-\alpha)+i(v-\beta)|^2
=
(u-\alpha)^2+(v-\beta)^2
<
\frac{\varepsilon^2}{2}+\frac{\varepsilon^2}{2}
=\varepsilon^2.
\]
Иймээс
\[
0<|z-z_0|<\delta \Rightarrow |f(z)-A|<\varepsilon.
\]
Энэ нь
\[
\lim_{z\to z_0} f(z)=A
\]
гэдгийг баталж байна.
\end{proof}


\begin{sen}
Хэрэв $f(z)$-ийн хязгаар орших бол $z_0$ руу ямар ч замаар ойртсон
утга ижил байна.
\end{sen}

\begin{exm}
\[
\lim_{z\to 0} \frac{z}{\overline{z}}
\]
орших уу?
\end{exm}

\begin{sol}
$z=re^{i\theta}$ гэж бичвэл
\[
\frac{z}{\overline{z}}=e^{2i\theta}.
\]
$z\to 0$ гэдэг нь $r\to 0$ гэсэн үг. Гэтэл $\dfrac{z}{\overline{z}}=e^{2i\theta}$
нь $\theta$-оос хамаарна. Иймээс хязгаар $r\to 0$ үед $\theta$ хамаарсан тул хязгаар оршихгүй.
\end{sol}

\begin{exm}
\[
\lim_{(x,y)\to(0,0)} \frac{x^2}{x^2+y^2}
\]
орших уу?
\end{exm}

\begin{sol}
$y=0$ үед утга $1$,
$x=0$ үед утга $0$.
Хоёр өөр утга гарч байна.
Иймээс хязгаар оршихгүй.
\end{sol}

\begin{sen}[Шахах зарчим]
Хэрэв
\[
|f(z)|\le g(z)
\]
мөн
\[
\lim_{z\to z_0} g(z)=0
\]
бол
\[
\lim_{z\to z_0} f(z)=0.
\]
\end{sen}

\begin{exm}
\[
\lim_{z\to 0} z^2 \sin \frac{1}{z}
\]
утгыг ол.
\end{exm}

\begin{sol}
\[
|z^2\sin \tfrac{1}{z}| \le |z|^2.
\]
\[
\lim_{z\to 0} |z|^2=0.
\]
Иймээс
\[
\lim_{z\to 0} z^2 \sin \frac{1}{z}=0.
\]
\end{sol}
\subsection*{Хязгаар $z\to\infty$ үед}

\begin{dfn}
Хэрэв
\[
\forall \varepsilon>0\ \exists R>0:\ |z|>R \Rightarrow |f(z)-A|<\varepsilon
\]
бол
\[
\lim_{z\to\infty} f(z)=A.
\]
\end{dfn}

\begin{exm}
\[
\lim_{z\to\infty} \frac{1}{z}
\]
утгыг ол.
\end{exm}

\begin{sol}
\[
\left|\frac{1}{z}\right|=\frac{1}{|z|}\to 0.
\]
Иймээс
\[
\lim_{z\to\infty} \frac{1}{z}=0.
\]
\end{sol}
\begin{thm}[Хязгаарын арифметик үйлдлүүд]
Хэрэв
\[
\lim_{z\to z_0} f(z)=A,\qquad
\lim_{z\to z_0} g(z)=B
\]
бол дараах тэнцэтгэлүүд биелнэ:
\begin{enumerate}
    \item ямар ч тогтмол $c\in\mathbb{C}$-ийн хувьд.
$\lim\limits_{z\to z_0} c\,f(z)=cA$

\item $\lim\limits_{z\to z_0} (f(z)+g(z))=A+B$
\item $\lim\limits_{z\to z_0} (f(z)-g(z))=A-B$
\item$\lim\limits_{z\to z_0} (f(z)g(z))=AB$
\item $\lim\limits_{z\to z_0} \dfrac{f(z)}{g(z)}=\dfrac{A}{B} \qquad B\neq 0.$
\end{enumerate}
\end{thm}
\begin{exm}
Дараах хязгаарыг батал:
\[
\lim_{z\to z_0}\sin z=\sin z_0.
\]
\end{exm}

\begin{sol}
$z=x+iy$, $z_0=x_0+iy_0$ гэж бичье.
Дараах томьёог ашиглая:
\[
\sin(x+iy)=\sin x\,\ch y+i\cos x\,\sh y.
\]

Тэгвэл
\[
\sin z-\sin z_0
=
(\sin x\,\ch y-\sin x_0\,\ch y_0)
+i(\cos x\,\sh y-\cos x_0\,\sh y_0).
\]

Бодит хэсгийг задлавал
\[
\sin x\,\ch y-\sin x_0\,\ch y_0
=
(\sin x-\sin x_0)\ch y
+
\sin x_0(\ch y-\ch y_0).
\]

$x\to x_0$, $y\to y_0$ үед
\[
\sin x-\sin x_0\to 0,\qquad
\ch y-\ch y_0\to 0.
\]
Иймээс бодит хэсэг 0 рүү тэмүүлнэ.

Хуурмаг хэсгийг мөн адил задлахад
\[
\cos x\,\sh y-\cos x_0\,\sh y_0
=
(\cos x-\cos x_0)\sh y
+
\cos x_0(\sh y-\sh y_0),
\]
эндээс мөн 0 гарна.

Иймээс
\[
\lim_{z\to z_0}\sin z=\sin z_0.
\]
\end{sol}
\begin{thm}[Давхар функцийн хязгаар]
Хэрэв
\[
\lim_{z\to z_0} f(z)=w_0,
\qquad
\lim_{w\to w_0} g(w)=L,
\]
бол
\[
\lim_{z\to z_0} g(f(z)) = L.
\]
\end{thm}
\begin{exm}
\[
\lim_{z\to 1} \sin(z^2)
\]
утгыг ол.
\end{exm}

\begin{sol}
\[
\lim_{z\to 1} z^2 = 1.
\]
Синус функц бүх цэгт тасралтгүй тул
\[
\lim_{z\to 1} \sin(z^2)=\sin(1).
\]
\end{sol}

%---------------------------------------------------

\section{Тасралтгүй чанар}

\begin{dfn}
Хэрэв
\[
\lim_{z\to z_0} f(z)=f(z_0)
\]
бол $f$-г $z_0$ дээр тасралтгүй гэнэ.
\end{dfn}

\begin{thm}
$f(z)=u(x,y)+iv(x,y)$ бол $f$ нь $z_0$ дээр тасралтгүй байх нь
\[
u(x,y)\ \text{ба}\ v(x,y)
\]
нь $(x_0,y_0)$ дээр тасралтгүй байхтай эквивалент.
\end{thm}

\begin{thm}
Хэрэв $f$ ба $g$ нь $z_0$ дээр тасралтгүй бол
\[
f+g,\quad fg,\quad \frac{f}{g}\ (g(z_0)\neq 0)
\]
нь мөн $z_0$ дээр тасралтгүй.
\end{thm}

\begin{thm}

$P(z)$ олон гишүүнт нь $\mathbb{C}$ дээр бүх цэгт тасралтгүй.
Рациональ функц
\[
R(z)=\frac{P(z)}{Q(z)}
\]
нь
\[
Q(z)\neq 0
\]
байх бүх цэгт тасралтгүй.
\end{thm}
\begin{thm}
Комплекс экспонент функц
\[
f(z)=e^z
\]
нь $\mathbb{C}$ хавтгай дээр бүх цэгт тасралтгүй.
\end{thm}

\begin{proof}
$z=x+iy$ гэж бичвэл
\[
e^z=e^{x+iy}=e^x(\cos y+i\sin y).
\]

Иймээс
\[
u(x,y)=e^x\cos y,\qquad v(x,y)=e^x\sin y.
\]

Бодит функцууд $e^x,\cos y,\sin y$ нь $\mathbb{R}$ дээр тасралтгүй.
Үржвэр нь мөн тасралтгүй тул $u(x,y)$ ба $v(x,y)$ нь
$(x_0,y_0)$ дээр тасралтгүй.

Тэгвэл “$u,v$ тасралтгүй $\Rightarrow$ $f$ тасралтгүй” шалгуураар
$e^z$ нь бүх $z_0\in\mathbb{C}$ дээр тасралтгүй.
\end{proof}

\begin{thm}
Логарифм
\[
\ln z=\ln|z|+i\arg z
\]
нь
\[
\Omega=\mathbb{C}\setminus\{x\le 0,\ y=0\}
\]
муж дээр тасралтгүй.
\end{thm}

\begin{proof}
$z=x+iy$ гэж бичье.

Тэгвэл
\[
\ln z=\frac12\ln(x^2+y^2)
+i\,\arg(z).
\]

1) $\ln(x^2+y^2)$ нь $(0,0)$-оос бусад бүх цэг дээр тасралтгүй.

2) $\Arg(z)$ нь сөрөг бодит тэнхлэгийг хассан муж дээр тасралтгүй.

Иймээс бодит ба хуурмаг хэсгүүд нь $\Omega$ муж дээр тасралтгүй.

“$u,v$ тасралтгүй $\Rightarrow$ $f$ тасралтгүй” шалгуураар
$\Ln z$ нь $\Omega$ муж дээр тасралтгүй.
\end{proof}
\section{Олон утгат функц}

\begin{dfn}
Хэрэв комплекс тоо $z$ бүрт нэгээс илүү $w$ утга харгалзаж байвал
түүнийг \emph{олон утгат комплекс функц} гэнэ.
\end{dfn}
\begin{exm}
Комплекс тооны аргумент.
\end{exm}

Комплекс тоо $z\neq 0$-г
\[
z=re^{i\varphi}
\]
гэж бичвэл
\[
\arg z=\varphi+2k\pi,\qquad k\in\mathbb{Z}.
\]

Иймээс $\arg z$ нь төгсгөлгүй олон утгатай функц.
\begin{exm}
Комплекс логарифм.
\end{exm}

$z\neq 0$ үед
\[
\Ln z=\ln|z|+i(\arg z+2k\pi),\qquad k\in\mathbb{Z}.
\]

Иймээс $\Ln z$ нь олон утгат функц.
\begin{exm}
Комплекс язгуур.
\end{exm}

$z=re^{i\varphi}$ гэж бичвэл
\[
\sqrt[n]{z}
=
r^{1/n}
e^{\,i\frac{\varphi+2k\pi}{n}},
\qquad k=0,1,\dots,n-1.
\]

Иймээс $n$-зэрэгт язгуур нь $n$ утгатай функц.
\begin{dfn}
Олон утгат функцийн утгуудын дундаас нэгийг нь
сонгон авснаар үүсэх нэг утгат функцийг
\emph{салбар} гэнэ.
\end{dfn}
\begin{dfn}
Хэрэв
\[
-\pi<\Arg z\le \pi
\]
гэж сонговол үүсэх логарифмыг
\emph{логарифмын гол утга} гэнэ:
\[
\ln z=\ln|z|+i\Arg z.
\]
\end{dfn}
\section{Семинар}
\begin{enumerate}

\item $f(z)=z^3$ функцийн бодит ба хуурмаг хэсгийг ол.

\item $f(z)=e^z$ функцийн бодит ба хуурмаг хэсгийг ол.

\item $\displaystyle \lim_{z\to 2} (3z^2-1)$ утгыг ол.

\item $\displaystyle \lim_{z\to 1+i} z^2$ утгыг ол.

\item $\displaystyle \lim_{z\to 0} \frac{z}{z+1}$ утгыг ол.

\item $\displaystyle \lim_{z\to 0} \frac{z^2+2z}{z}$ утгыг ол.

\item $\displaystyle \lim_{z\to z_0} \cos z$ утгыг ол.

\item $\displaystyle \lim_{z\to z_0} e^z$ утгыг ол.

\item $\displaystyle \lim_{z\to \infty} \frac{1}{z}$ утгыг ол.

\item $\displaystyle \lim_{z\to \infty} \frac{2z+1}{z}$ утгыг ол.

\item Дараах хязгаар орших эсэхийг шалга:
\[
\lim_{z\to 0} \frac{|z|}{z}.
\]

\item Дараах хязгаар орших эсэхийг шалга:
\[
\lim_{(x,y)\to(0,0)} \frac{x^2-y^2}{x^2+y^2}.
\]

\item $\displaystyle \lim_{z\to 0}z^{-2}(1-\cos z))$ утгыг ол.

\item $f(z)=\frac{z^2}{|z|}$, $z\neq 0$, $f(0)=0$ гэж тодорхойлов.
$f$ функц 0 дээр тасралтгүй юу?

\item $\Ln z$ функц сөрөг бодит тэнхлэг дээр яагаад тасралтгүй биш вэ?
Тайлбарла.

\item Дараах хязгаар орших эсэхийг шалга:
\[
\lim_{z\to 0} \frac{xy}{x^2+y^2}.
\]

\item $f(z)=|z|$ функц $z_0=0$ дээр тасралтгүй эсэхийг $\varepsilon -\delta$ тодорхойлолтоор батал.

\item $\displaystyle \lim_{z\to 0} \frac{\sin z}{z}$ утгыг ол.
(Синусын тодорхойлолтыг ашиглаж бод.)

\item $\sqrt{z}$ функцийн бүх утгыг ол:
\[
z=-4.
\]
Мөн гол утгыг тодорхойл.

\end{enumerate}
\section*{Даалгаврын бодлого}

\begin{enumerate}

\item $\displaystyle \lim_{z\to 0} f(z)=L$ байх $\Longleftrightarrow $ нь $\forall t\in [0;2\pi[$ хувьд $\displaystyle \lim_{r\to 0} f(re^{it})=L$ байхыг батал. Үүнийг ерөнхий тохиолдолд  
$$\exists\displaystyle \lim_{z\to a} f(z) \Longleftrightarrow \displaystyle \lim_{r\to 0} f(re^{it})$$ хязгаар $t$-с үл хамааран оршин байх гэж бичиж болно. Энд $z-a=re^{it}.$ $\displaystyle\lim_{r\to 0} f(re^{it})=\lim_{r\to 0} u(r\cos t, r\sin t)+i\lim_{r\to 0} v(r\cos t, r\sin t)$ гэсэн нэг хувьсагчийн бодит функцийн хязгаар юм. 

\item 
\begin{equation}
\Arccos z= -i \Ln (z+\sqrt{z^2-1}) 
\end{equation} функцийн бодит ба хуурмаг хэсгийг ол. 

\item $\forall z\in\mathbb{C}$ хувьд $\cos(\Arccos(z))=z$ байхыг батал. Өөрөөр хэлбэл $\cos z$ урвуу функц олон утгат функц байна.

\item $-1\leq x\leq 1$ үед $\Arccos(x)$ нь бодит хувьсагчийн $\arccos(x)$ функц болохыг харуул.

\item $\sin z$ урвуу функц $\Arcsin(z)$ ол.

\item $\cos z, \sin z$ функцүүд $\mathbb{C}$ дээр тасралтгүй байхыг батал.

\item  $\tan z, \cot z$ функцүүд тасралтгүй байх олонлогийг ол.

\item $|z|\geq |Re(z)|$ батал.

\item $|z_1|+|z_2|\geq |z_1+z_2|\geq ||z_1|-|z_2||$ батал.

\item $|z|^2\geq 2Re(z)Im(z)$ батал.

\end{enumerate}

%\chapter{Комплекс функц, түүний хязгаар ба тасралтгүй чанар}
\section{Комплекс хувьсагчийн функц}
\begin{dfn}
$z\in D$ цэгт комплекс $w\in E$ тоо харгалзаж байх $f$ хууль өгөгдөж байвал $D\subset \overline{\mathbb{C}}$ олонлог дээр $f(z)$ функц өгөгдлөө гэх бөгөөд
\begin{equation}
f:D\rightarrow \overline{\mathbb{C}} \text{ буюу } w=f(z)
\end{equation}
тэмдэглэл хийдэг.

$D$ мужийг функцийн тодорхойлогдох,  $E$ мужийг өөрчлөгдөх муж гэж ярих ба $D$-г $\circled{z}$ хавтгайд $E$-г $\circled{w}$ хавтгайд авч заншжээ. Иймд, $w=f(z)$ функц $\circled{z}$ хавтгайн $D$ мужийг $\circled{w}$ хавтгайн $E$ мужид нэгэн утгатай буулгана. Энэ утгаар $w$-ийг $z$-ийн дүр, $z$-ийг эх дүр гэдэг.
\end{dfn}
\begin{exm}
$w=z^2$ функц I мөчийг хаана буулгах вэ?
\end{exm}
\begin{sol}
Эхлээд, $Oz$ цацрагийн дүр энэ буулгалтаар хаана гарахыг олъё.
\begin{center}
    \includegraphics[width=9cm]{Zurag/4.jpg}
     \captionof{figure}{}
\end{center}
$\Arg  \varphi$ гэвэл $\Arg  w=\Arg(z^2)=2\cdot \Arg  x=2\varphi$ байх тул \\ $0\leq \varphi \leq \dfrac{\pi}{2}$ учир $0\leq \Arg  w \leq \pi$. Иймд $w=z^2$ функц I мөчийг дээд хагас хавтгайд буулгана. 
\end{sol}
$z=x+iy$, $w=u+iv$ гэвэл $f(z)=u(x,y)+iv(x,y)$ бичиглэл үйлдэж болно. Өөрөөр хэлбэл $z$ аргументийн комплекс\\ $w=f(z)$ функц нь хоёр хувьсагчийн $u(x,y)$, $v(x,y)$ хоёр функцээр тодорхойлогдох болж байна:
\begin{equation}
u(x,y)=\re f(z), v(x,y)=\im f(z)
\end{equation}
Дээрх $w=z^2$ функцийн бодит ба хуурмаг хэсгийг ол.
\begin{align*}
    &w=z^2=(x+iy)^2=x^2-y^2+2xyi \ \text{тул}\\
    &u(x,y)=\re (z^2)=x^2-y^2 \\
    &v(x,y) =\im (z^2)=2xy.
\end{align*}
\begin{hnt}
Хэрэв хоёр хувьсагчийн $u(x,y)$, $v(x,y)$ гэсэн хоёр функц $D(x,y)$ муж дээр өгөгдсөн гэвэл $w=u(x,y)+iv(x,y)$ нь комплекс $z=x+iy$ аргументийн комплекс функц байна.
\end{hnt}
\section{Функцийн хязгаар} \index{Функцийн хязгаар}
\begin{dfn}
    $z_0=(x_0,y_0)\in D$ байг. $w=f(z)$ нь $D\setminus\{z_0\}$ дээр тодорхойлогдсон функц байг. Хэрэв дараах нөхцөл хангах $A\in \mathbb{C}$ тоо олдож байвал $A$ тоог $w=f(z)$ функцийн $z\to z_0$ хязгаар гэнэ:
\begin{equation} \label{хязгаар тэмдэг}
    \forall \varepsilon >0 \ \  \exists \delta (\varepsilon, z_0)>0: 0<|z-z_0|<\delta \Rightarrow |f(x)-A|<\varepsilon 
    \end{equation}
$w=f(z)$ функцийн $z\to z_0$ хязгаарыг 
\begin{equation} \label{хязгаар}
\lim_{z\rightarrow z_0} f(z)=A
\end{equation}
гэж тэмдэглэнэ.
\end{dfn}
\begin{equation} \label{хязгаар}
\lim_{z\rightarrow z_0} f(z)=A=\alpha +i\beta
\end{equation}
 (\ref{хязгаар}) нь
\begin{equation}
\lim_{(x,y)\rightarrow (x_0,y_0)} u(x,y)=\alpha \qquad \lim_{(x,y)\rightarrow (x_0,y_0)} v(x,y)=\beta
\end{equation}
гэж хоёр хувьсагчийн хоёр функцийн хязгаар бодоход хүргэнэ. 

$z_0=\infty$ тохиолдолд функцийн хязгаар
$$\forall \varepsilon >0 \quad \exists R_0(\varepsilon): \quad |z|>R_0\Rightarrow |f(z)-A|<\varepsilon$$
гэж тодорхойлогдох бөгөөд
\begin{equation}
\lim_{z\rightarrow \infty} f(z)=A.
\end{equation}
\section{Функцийн хязгаар. Тасралтгүй чанар} \index{Тасралтгүй чанар}
Комплекс хувьсагчийн $w=f(z)$ функцийн $z_0\in D$ цэг дээр тасралтгүй байх нь
\begin{equation} \label{2.6}
\lim_{z\rightarrow z_0} f(z)=f(z_0)
\end{equation}
гэсэн хязгаарын тэнцлээр тодорхойлогдоно. Комплекс $z$ хувьсагчид $\Delta z=z-z_0$ өөрчлөлт өгвөл $w=f(z)$ функц харгалзан
$$\Delta w= \Delta f(z_0)=f(z)-f(z_0)$$
өөрчлөлт авч (\ref{2.6}) ёсоор
\begin{equation}
\lim_{\Delta z\rightarrow 0} \Delta w =0
\end{equation}
болж, өөрөөр хэлбэл аргументийн төгсгөлгүй бага өөрчлөлтөд функцийн төгсгөлгүй бага өөрчлөлт харгалзаж байвал\\ $w=f(z)$ функц $z_0$ цэг дээр тасралтгүй байна гэж ярьдаг.

Бид зөвхөн нэг утгат функцийн талаар авч үзлээ. Гэвч онолын болон олон тооны хавсарга ухаануудад олон хувьсагчийн комплекс функц авч үзэх шаардлага тулгардаг. Наад захын жишээ татахад л $\Arg  z$ функц байна. Энэ функц $z\neq 0$ утга бүхэнд нэг нь нөгөөгөөс $2\pi $ тоонд бүхлээр хуваагдах нэмэгдэхүүнээр ялгаатай төгсгөлгүй олон утга авна. Комплекс тооны язгуурын $w=\sqrt[n]{z}$ функц $n$ утгат функц болохтой өмнө танилцсан билээ. 

Нэг утгат тасралтгүй $w=f(z)$ функц $D\subset \mathbb{C}$ мужид тодорхойлогдсон байг. $\forall z_1\in D$, $\forall z_2\in D$ авахад $f(z_1)\neq f(z_2)$ байх нь зөвхөн $z_1\neq z_2$ байхад биелэгдэж байвал $f(z)$-ийг $D$ дээр харилцан нэг утгат буюу нэг хуудаст функц гэнэ.

Иймд, $w=f(z)$ нь $D$ мужид нэг хуудаст бөгөөд $D$-г $D^*$-д буулгах бол $D^*$ мужид урвуу $z=f^{-1}(w)$ функц оршин байх болов. Өөрөөр хэлбэл  $\forall z\in D$ хувьд $z=f^{-1}(f(z))$ ба $\forall w\in D^*$ хувьд $w=f(f^{-1}(w))$.
\section{Үндсэн элементар функцүүд}
Нэг утгат хялбар функцүүдийн нэг нь $w=z^n$, $n\in \mathbb{N}$ гэсэн зэрэгт функц. Зэрэгт функцийн тусламжтайгаар комплекс $z$ хувьсагчийн олон гишүүнт функц үүсгэнэ:
\begin{equation}
P_n(z)=a_0z^n+a_1z^{n-1}+\cdots +a_n, \qquad a_0\neq 0
\end{equation}
$a_0, \ldots , a_n$ коэффициентууд нь тогтмол комплекс тоонууд. Олон гишүүнт нь $\overline{\mathbb{C}}$ дээр тодорхойлогдоно. Энд $P_n(\infty)=\infty$ гэж үзнэ. Хялбар олон гишүүнт нь шугаман функц:
$$w=az+b, \quad a\neq 0.$$

Олон гишүүнтийг өөрөөр бүхэл рациональ функц гэнэ. Хоёр олон гишүүнтийн ноогдвороор бутархай рациональ $R(z) $ функц дүрслэгдэнэ:
\begin{equation}
R(z)=\dfrac{P_m(z)}{P_n(z)}=\dfrac{b_0z^m+b_1z^{m-1}+\cdots +b_m}{a_0z^n+a_1z^{n-1}+\cdots +a_n}, \quad a_0\neq 0
\end{equation}
Бутархай рационал хялбар
\begin{equation}
L(z)=\dfrac{az+b}{cz+d}, \quad
\begin{vmatrix}
a&b\\
c&d
\end{vmatrix}\neq 0
\end{equation}
функцийг бутархай шугаман буулгалт гэнэ. \\ $w=e^z$ \quad $(z=x+iy)$-ийг комплекс экспонент буюу илтгэгч функц \index{Илтгэгч функц}гэнэ. Энэ функц бодит илтгэгчийн бүхий л дүрэмд захирагдана. Харин комплекс хувьсагчийн илтгэгч функц үет функц байна. Учир нь
\begin{align*}
    e^{z+2n\pi i}&= e^{x+i(y+2n\pi)}=e^x[\cos (y+2n\pi)+i\sin (y+2n\pi)] \\
    &=e^x(\cos y+i\sin y)=e^{x+iy}=e^z
\end{align*}
болж $e^z$ нь $T=2n\pi i$ үет функц болох нь нотлогдов.

Эйлерийн томьёо ашиглаж тригонометр
\begin{equation}
\cos z=\dfrac{e^{iz}+e^{-iz}}{2}, \quad \sin z= \dfrac{e^{iz}-e^{-iz}}{2i}
\end{equation}
функцүүд $z\in\mathbb{C}$ дээр тодорхойлогдох бөгөөд
\begin{equation}
\tan z=\dfrac{\sin z}{\cos z}, \quad \cot z= \dfrac{\cos z}{\sin z}
\end{equation}
гэж тодорхойлогдох нь $\tan z$ нь $z_n=(2n+1)\dfrac{\pi}{2}$, $n=0, \pm 1, ...$ ба харин $\cot z$ нь $z_n=n\pi $ $(n=0, \pm 1,...)$ цэг дээр тодорхойлогдохгүй.

Гиперболлог синус, косинус функцүүд \index{Гиперболлог функцүүд } $\mathbb{C}$ хавтгайд
\begin{equation}
\sh z=\dfrac{e^z-e^{-z}}{2}, \quad \ch z= \dfrac{e^z+e^{-z}}{2}
\end{equation}
гэж тодорхойлогдох бөгөөд $\ch z =\cos iz$, $\sh z=-i\sin z$ байна. Комплекс аргументийн тригонометр мөн гиперболлог функцүүдийн хувьд бодит аргументийн эдгээр функцийн бүхий л чанар хадгалагдана.
$$\cos (z_1\pm z_2)=\cos z_1\cos z_2 \mp \sin z_1 \sin z_2$$
$$\sin (z_1\pm z_2)=\sin z_1\cos z_2 \pm \sin z_2 \cos z_1$$
томьёонууд ашиглаж тригонометр функцийн бодит ба хуурмаг хэсгийг олдог.
\begin{exm}
\begin{align*}
    \cos z&=\cos(x+iy)=\cos x\cos (iy)-\sin x \sin (iy) \\
    &= \cos x ch y -i\sin x \sh y =\cos x ch y -i\sin x \sh y
\end{align*}
байх тул  $$ \re (\cos z)= \cos x \ch y, \quad \im (\cos z)=-\sin x \sh y$$ болно. 
\end{exm}
Комплекс мужид илтгэгч функцийн урвуу функцийг логарифм функц \index{Логарифм функц} гэдэг. Өөрөөр хэлбэл хэрэв $w= ln z$ бол $z=e^w$ тэгшитгэлийг хангах $w$ олдоно: $w=u+iv$, $z=re^{i\varphi}=re^{i \Arg  z}$
гэж авбал 
$$z=re^{i \Arg  z}= e^{w}=e^u\cdot e^{iv}$$
бөгөөд эндээс $r=e^u$, $e^{i \Arg  z}=e^{iv}$ эхний тэгшитгэлээс $z\neq 0$ авбал $u=\ln r$ гэж нэг утгатай олдох ба харин дараах тэгшитгэлээс $u=\Arg  z$ гарна.

Иймд $z\neq 0$, $z\neq \infty$ байхад $w=\Ln z$ гэсэн олон утгат логарифм функц тодорхойлогдоно.
\begin{equation} \label{2.13}
w=\Ln z= \ln |z|+i (\arg  z+2n \pi), \quad n=0, \pm 1, ...
\end{equation}
(\ref{2.13})-аас $n=0$ тохиолдолд логарифмын гол утга гэж нэрлэгдэх
\begin{equation}
\ln z= \ln |z|+i \arg  z
\end{equation}
функц $\mathbb{C} \setminus \{ -\infty < x \leq 0, \quad y=0 \}$ өөрөөр хэлбэл $Ox$-ийн сөрөг чиглэлд хагалбартай бүх хавтгай муж дээр тасралтгүй, нэг утгатай байхаар оршин байна.

Тригонометр $\cos  z$, $\sin z$, $\tan z$-г функцүүдийн \index{Тригонометрийн урвуу функцүүд} хувьд тэдгээрийн урвуу болох $\Arccos z$, $\Arcsin z$, $\Arctan z$ зэрэг олон утгат функцүүд тодорхойлогдоно.
\begin{equation}
\Arccos z= -i \Ln (z+\sqrt{z^2-1}) 
\end{equation}
\begin{equation}
\Arcsin z= -i \Ln (iz-\sqrt{1-z^2}) 
\end{equation}
\begin{equation}
\Arctan z= -\dfrac{i}{2}\Ln \dfrac{1+iz}{1-iz}
\end{equation}
Мөн ерөнхий зэрэгт $z^a$, $a=\alpha +i\beta$  функц экспонентоор
\begin{equation}
z^a=e^{a\Ln z}
\end{equation}
ба гол утга нь
\begin{equation}
z^a=e^{a\ln z}
\end{equation}
гэж тодорхойлогдоно. Мөн ерөнхий илтгэгч $a^z$, $a\neq 0$ экспонентоор
\begin{equation}
a^z=e^{z\Ln a}
\end{equation}
ба гол утга нь
\begin{equation}
a^z=e^{z\ln a}
\end{equation}
\begin{exm}
$\cos (2+i)$ тооны бодит ба хуурмаг хэсгийг олъё.
\end{exm}
\begin{sol}
$\cos (2+i)=\cos 2\cdot \ch 1-i \sin 2\cdot \sh1$ тул $$\re \cos (2+i)=\cos 2 \sh1, \quad \im \cos (2+i)=-\sin 2\cdot \sh1$$
\end{sol}
\begin{exm}
Дараах утгуудыг бодож гарга.

\begin{enumerate}
\colgur{
    \item $\Ln (-1)$
    \item $ \Arccos \dfrac{1}{2}$
    \item $2^i$}
    \end{enumerate}
\end{exm}
\begin{sol}
\begin{enumerate}
    \item \begin{align*}
    \Ln(-1)&=\ln|-1|+i \arg  (-1)+2 k\pi i \\
    &= 0+\pi i +2 k \pi i= (2k+1)\pi i, \qquad k=0,\pm 1, ...
\end{align*}
\item 
\begin{align*}
    \Arccos\dfrac{1}{2}&= -i \Ln \left(\dfrac{1}{2}+\dfrac{\sqrt{3}}{2}i\right) \\
    &=-i \left(\ln\left|\dfrac{1}{2}+i\dfrac{\sqrt{3}}{2}\right|=\pm i\cdot \dfrac{\pi}{3}+2k\pi i\right)  \\
    &=2 k\pi \pm \dfrac{\pi}{3}
\end{align*}
\item 
\begin{align*}
    2^i&=e^{i \Ln 2} =e^{i(\ln|2|+0+2k\pi i)}=e^{i(\ln2+2k\pi i)}\\
    &=e^{-2k \pi} (\cos \ln 2+ i\sin \ln 2).
\end{align*}
\end{enumerate}
\end{sol}
\newpage
\begin{center}
    \textbf{Дасгал}
\end{center}
\begin{enumerate}
    \item [11.] Дараах илэрхийллийн бодит ба хуурмаг хэсгийг ол. \\
   \textbf{ 1.} $e^{z^2}$; \qquad \textbf{2. }$\ch (z-i)$; \qquad \textbf{3.} $\tan z$;
   \item [12.] Бодож гарга. \\
    \textbf{ 1.} $\Ln (-i)$; \qquad \textbf{2. }$\ln(-3-2i)$; \qquad \textbf{3.} $i^{\frac{1}{i}}$; \qquad \textbf{4.} $\left ( \dfrac{1+i}{\sqrt{2}}\right)2i$
    \item [13.] $\sin z=3$ тэгшитгэл бод.
    \item [14.] Алгебр хэлбэрт бич:\\
    \textbf{ 1.} $\sin \pi i$; \qquad \textbf{2. }$ \tan \dfrac{\pi}{2}i$; \qquad \textbf{3.} $\Arccos i$; \qquad \textbf{4.} $\Arctan \dfrac{i}{3}$;
\end{enumerate}
%\chapter{Комплекс функцийн дифференциалчлал}
\section{Уламжлал ба дифференциал. \\ \index{Уламжлал}
Эйлер-Даламберийн нөхцөл, \index{Эйлер-Даламберийн нөхцөл}
Голоморф функц}
Комплекс хувьсагчийн функцийн уламжлалын тодорхойлолт бодит хувьсагчийн функцийнх шиг формаль байдлаар өгөгдөнө.\index{Формаль уламжлал} $z_0$-цэгийн ямар нэг $U=U(z_0,\delta)=\{|z-z_0|<\delta\}$ орчинд нэг утгат $w=f(z)=u+iv$ функц тодорхойлогдсон байг. Удирдан хувьсах $z$ хэмжигдэхүүн өөрийн $U$ орчноос үл гарах бүх боломжит
$$\Delta z=z-z_0=(x-x_0)+i(y-y_0)=\Delta z+i\Delta y\neq 0$$
өөрчлөлт авч байг. $\Delta z$ өөрчлөлт бүрд харгалзах функцийн өөрчлөлт
$$\Delta w=\Delta f(z_0)=f(z)-f(z_0) =\Delta u+i\Delta v$$
байна гэж үзье. Функцийн $\Delta w$ өөрчлөлтийг удирдан хувьсах хэмжигдэхүүний $\Delta z$ өөрчлөлтөд хувааж $\Delta z\rightarrow 0$ байхаар хязгаарт шилжье. Хэрэв энэ хязгаар оршин байвал түүнийг\\ $w=f(z)$ функцийн $z_0$ цэг дээрх уламжлал гэж нэрлэх бөгөөд 
\begin{eqnarray} \label{3.1}
 \dfrac{dw}{dz}\Big |_{z=z_0}=f'(z_0)=\lim\limits_{\Delta z\rightarrow 0}\dfrac{\Delta w}{\Delta z}
\end{eqnarray}

тэмдэглэл хэрэглэнэ. 

$w=f(z)$ функц $z_0$ цэгийн ямар нэг $U=U(z_0,\delta)$ орчинд тодорхойлогдохын хамт түүний $\Delta w$ өөрчлөлтийг $\Delta z$-ийн хувьд шугаман $A\cdot \Delta z$ хэсэг ба $\Delta z$-тэй жишихэд дээд эрэмбийн төгсгөлгүй багасах хэмжигдэхүүний нийлбэрийн
\begin{eqnarray} \label{3.2}
 \Delta w=\Delta f(z_0)=A\cdot \Delta z+\circ (\Delta z)=A\cdot \Delta z+\alpha (\Delta z)\cdot \Delta z
\end{eqnarray}

дүрстэй бичигдэж байвал $w=f(z)$-ийг $z_0$ цэг дээр дифференциалчлагдах функц гэнэ. Энд $A$ тоо $\Delta z$-ээс үл хамаарах бөгөөд харин $\alpha (\Delta z)$ нь $\Delta z\rightarrow 0$ үед төгсгөлгүй бага функц, өөрөөр хэлбэл $\lim\limits_{\Delta z\rightarrow 0}\alpha (\Delta z)=0$, иймд $\circ (\Delta z)=\alpha (\Delta z)\cdot \Delta z$ нь $\Delta z$-тэй жишихэд дээд эрэмбийн төгсгөлгүй багасах хэмжигдэхүүн.
\begin{thm}
$w=f(z)$ функц $z_0$ цэг дээр дифференциалчлагдах зайлшгүй бөгөөд хүрэлцээтэй нөхцөл\index{Дифференциалчлагдах нөхцөл} нь уул функц $z_0$ дээр уламжлалтай, тухайлбал $A=f'(z_0)$ байх явдал мөн.
\end{thm}
\begin{prf}
\textbf{Зайлшгүй:} (\ref{3.2}) биелэгдэж байна гэдгээс
$$\frac{\Delta w}{\Delta z}=A+\frac{\circ (\Delta z)}{\Delta z}$$ бөгөөд $\Delta z \rightarrow 0$ үед 
$$A=\lim\limits_{\Delta z \rightarrow 0} \frac{\Delta w}{\Delta z}=w'(z_0)$$ болов. \\
\textbf{Хүрэлцээтэй нь:} $w'(z_0)$ оршин байна гэдгээс $\Delta z\rightarrow 0$ үед
$$\Delta w=A\cdot \Delta z+\alpha (\Delta z)\cdot \Delta z$$ болж батлагдлаа.
\hfill $\blacksquare$\end{prf}
$w=f(z)$ функцийн $z_0$ цэг дээрх дифференциал \index{Дифференциал} гэж
\begin{eqnarray}
 dw=df(z_0)=f'(z_0)dz
\end{eqnarray}-ийг нэрлэнэ. Дифференциал нь функцийн өөрчлөлтийн гол хэсэг.

Комплекс хувьсагчийн функцийн дифференциалчлагдах нөхцөл бодит хувьсагчийн функцийн дифференциалчлагдах нөхцөлөөс гадна нэмэлт нөхцөл шаардахыг дараах теорем харуулна.
\begin{thm}
Комплекс $w=f(z)=u(x,y)+iv (z,y)$ функц $z_0=x_0+iy_0$ цэгийн орчинд тодорхойлогдсон байг. Энэ функц $z_0$ цэг дээр дифференциалчдаглах зайлшгүй бөгөөд хүрэлцээтэй нөхцөл нь
\begin{enumerate}
    \item $u(x,y)$, $v(x,y)$ функцүүд $(x_0,y_0)$ цэг дээр дифференциалчлагдах,
    \item Эйлер-Даламберийн 
    \begin{equation}
\tag{E-D}
\dfrac{\partial u}{\partial x}=\dfrac{\partial v}{\partial y}, \qquad \dfrac{\partial u}{\partial y}=-\dfrac{\partial v}{\partial x}
\end{equation} нөхцөл $(x_0,y_0)$ цэг дээр биелэгдэх явдал мөн.
\end{enumerate}
\end{thm}
\begin{prf}
\textbf{Зайлшгүй:} (\ref{3.1}) биелэгдэж байг. 
$$\Delta f(z_0)=\Delta u+i\Delta v=f'(z_0)\Delta z+\circ (\Delta z)$$ бөгөөд $f'(z_0)=a+ib$, $\circ (\Delta z)= \alpha +i\beta$ гэж орлуулахдаа $a,b\in\mathbb{R}$, $\alpha$, $\beta$ нь $\Delta z\rightarrow 0$ үед дээд эрэмбийн төгсгөлгүй багасах бодит тоонууд гэж тооцъё.

Тэгтэл 
\begin{eqnarray}
 \Delta u+i \Delta v=(a+ib) (\Delta x+i\Delta y)+\alpha +i\beta
\end{eqnarray} болж 
\begin{eqnarray}
 \Delta u=a \Delta x-b\Delta y+\alpha, \qquad \Delta v=b\Delta x+a\Delta y +\beta
\end{eqnarray}

эндээс $u(x,y)$, $v(x,y)$ нь $(x_0,y_0)$ цэг дээр дифференциалчлагдах функцүүд болох нь батлагдаж түүнчлэн (3.5)-аас
$$a=\dfrac{\partial u}{\partial x}=\dfrac{\partial v}{\partial y}, \quad b=\dfrac{\partial u}{\partial y}=-\dfrac{\partial v}{\partial x}$$
гээд (E-D) гарч ирлээ. Ингэж теоремын зайлшгүй нөхцөл батлагдав. Теоремын нөхцөлүүдийн хүрэлцээтэйг баталъя.
1) ба 2) нөхцөл биелэгдсэн байг, өөрөөр хэлбэл
$$
\begin{cases}
\Delta  u=\dfrac{\partial u}{\partial x}\Delta x+\dfrac{\Delta u}{\Delta y}\Delta y+\alpha_1\Delta x+\alpha_2\Delta y\\
\Delta v=\dfrac{\partial v}{\partial x}\Delta x+\dfrac{\Delta v}{\Delta y}\Delta y+\beta_1\Delta x+\beta_2\Delta y
\end{cases}
$$ 
Энд $\alpha_1$, $\alpha_2$, $\beta_1$, $\beta_2$ нь $\Delta x$, $\Delta y$-ийн тэг рүү тэмүүлэх нөхцөлд тэг рүү тэмүүлэх хэмжигдэхүүнүүд байг. Түүнээс гадна
\begin{eqnarray}
 \dfrac{\partial u}{\partial x}=\dfrac{\partial v}{\partial y}=a, \qquad -\dfrac{\partial u}{\partial y}=\dfrac{\partial v}{\partial x}=b
\end{eqnarray}
гэе. Иймд
$$\Delta u=a\Delta x-b\Delta y+\alpha_1\Delta x+\alpha_2\Delta y$$
$$\Delta v=b\Delta x+a\Delta y+\beta_1\Delta x+\beta_2\Delta y$$ ба 
\begin{equation}
 \begin{split}
\Delta f(z)&=\Delta u+i\Delta с \\
&=a(\Delta x+i\Delta y)+ib(\Delta x+i\Delta y)+\\
&+(\alpha_1+i\beta_1)\Delta x+(\alpha_2+i\beta_2)\Delta y \\
&=(a+ib) \Delta z+\varepsilon \cdot \Delta z
 \end{split}  
\end{equation}
 болж эндээс 
$$|\varepsilon| =\left|(\alpha_1+i\beta_1)\frac{\Delta x}{\Delta z}+(\alpha_2+i\beta_2)\frac{\Delta y }{\Delta z}\right|\leq |\alpha_1|+|\beta_1|+|\alpha_2|+|\beta_2|$$ бөгөөд  $\Delta z=\Delta x+i\Delta y$-ийн тэг рүү тэмүүлэх үед $\alpha_1$, $\alpha_2$, $\beta_1$, $\beta_2$ тэг рүү тэмүүлэх тул улмаар $\varepsilon\rightarrow 0$ теорем бүрэн батлагдав.
\hfill $\blacksquare$\end{prf}
Энэ теорем ёсоор 
\begin{equation}
 \begin{split}
f'(z)&=a+ib=\dfrac{\partial u}{\partial x}+i\frac{\partial v}{\partial x}=\frac{\partial u}{\partial x}-i\frac{\partial u}{\partial y} \\
&=\dfrac{\partial v}{\partial y}+i\frac{\partial v}{\partial x}=\frac{\partial v}{\partial y}-i\frac{\partial u}{\partial y}
 \end{split}  
\end{equation}
байна. 

Хэрэв $z=re^{i\varphi}$ ба $f(z)=u(r,\varphi)+iv(r,\varphi)$ гэж өгөгдсөн байг. Тэгвэл Эйлер-Даламберийн (E-D) нөхцөл туйлын координатад
\begin{eqnarray} \label{3.9}
 \dfrac{\partial u}{\partial r}=\dfrac{1}{r}\cdot \dfrac{\partial v}{\partial \varphi}, \quad
 \dfrac{\partial v}{\partial r}=-\dfrac{1}{r}\cdot \dfrac{\partial u}{\partial \varphi}
\end{eqnarray}
дүрстэй байна. Үүнийг өөрсдөө батлаарай.(\ref{3.9}) ёсоор
\begin{eqnarray}
 f'(z)=\dfrac{r}{z}\left(\dfrac{\partial u}{\partial r}+i\dfrac{\partial v}{\partial r}\right)=\dfrac{1}{z}\left(\dfrac{\partial v}{\partial \varphi}-i\dfrac{\partial u}{\partial \varphi}\right)
 \end{eqnarray}
байна.

Формаль уламжлал ашиглаж (Е-D) нөхцөлийг өөр дүрстэй тавьж болно. 

$f(z)=u(x,y)+iv(x,y)$ функц $x,y$ хувьсагчаараа дифференциалчдаглах функц байг. $z=x+iy$ гэвэл
$$x=\frac{1}{2}(z+\overline{z}), \qquad y=-\frac{i}{2}(z-\overline{z})$$ тул 
$$\frac{\partial f}{\partial z}=\frac{1}{2}\left(\frac{\partial f}{\partial x}-i\frac{\partial f}{\partial y}\right)=\dfrac{1}{2} \left(\frac{\partial u}{\partial x}+\frac{\partial v}{\partial y}\right) +\frac{i}{2} \left(\frac{\partial v}{\partial x}-\frac{\partial u}{\partial y}\right),$$
$$\frac{\partial f}{\partial \overline{z}}=\frac{1}{2}\left(\frac{\partial f}{\partial x}+i\frac{\partial f}{\partial y}\right)=\dfrac{1}{2} \left(\frac{\partial u}{\partial x}-\frac{\partial v}{\partial y}\right) +\frac{i}{2} \left(\frac{\partial v}{\partial x}+\frac{\partial u}{\partial y}\right),$$
Эндээс харахад (E-D) биелэгдэж байвал
\begin{equation}
 \tag{$\overline{E-D}$}
 \dfrac{\partial f}{\partial \overline{z}}=0
\end{equation}
байна. Энэ нөхцөл биелсэн гэвэл $f(z)$-ийн уламжлал
$$\frac{\partial f}{\partial z}=\frac{\partial u}{\partial x}+i\frac{\partial v}{\partial x}=f'(z)$$
байна.
\begin{dfn}
$u(x,y)$, $v(x,y)$ функцүүд $z_0=(x_0, y_0)\in D$ цэг дээр тухайн уламжлалуудын хамт тасралтгүй байхын дээр мөнхүү цэг дээр дифференциалчлагдах функцүүд байвал $w=f(z)=u(x,y)+iv(x,y)$ функцийг $z_0\in D$ дээр бодит анализын утгаар буюу $\mathbb{R}$-дифференциалчлагдах функц гэнэ. \\ $w=f(z)=u(x,y)+iv(x,y)$ функц $z_0$ дээр $\mathbb{R}$-дифференциалчлаг-
дахаас гадна энэ цэг дээр $(E-D)$ нөхцөл биелэгдэж байвал уул функцийг $z_0\in D$-дээр комплекс анализын утгаар буюу $\mathbb{C}$-дифференциалчлагдах функц гэнэ.
\end{dfn}
\begin{exm}
 $w= \ch z$ нь $\mathbb{C}$-дээр $\mathbb{C}$ дифференциалчлагдах функц. Учир нь $u= \re z$ $\ch z=\ch x \ch y$; $v= \im \ch z= \sh x \sh y.$
\end{exm}
\begin{sol}
$$\frac{\partial u}{\partial x}=\sh x \cos y, \quad \frac{\partial v}{\partial y}=\sh x\cos y$$
$$\frac{\partial u}{\partial y}=-\ch x \sin  y, \quad \frac{\partial v}{\partial x}=\ch  x\sin y$$
Иймд $u$, $v$ функцүүд дифференциалчдагдахын хамт $(E-D)$ нөхцөл хангах $w=\ch z$ нь $\mathbb{C}$-дээр $\mathbb{C}$-утгаар дифференциалчлагдах функц. 
\end{sol}
\begin{dfn}
Нэг утгат $w=f(z)$ функц зөвхөн $z_0$ цэг дээр төдийгүй энэ цэгийн ямар нэг орчинд  $\mathbb{C}$- дифференциалчлагдаж байвал түүнийг $z_0$ цэг дээр голоморф буюу аналитик функц \index{Аналитик буюу голоморф функц}гэнэ. $D$ мужийн цэг бүр дээр голоморф функцийг $D$ муж дээр голоморф функц \index{Голоморф функц (цэг дээр, муж дээр)} гэнэ. $f(z)$ нь $D$ муж дээр голоморф гэхийг $f(z)\in H(D)$ гэж тэмдэглэдэг.
\end{dfn}
\begin{exm}
 $w=e^z=e^x(\cos y +i\sin y)$ функц авбал 
\end{exm}
\begin{sol}
$$\frac{\partial u}{\partial x}=e^x\cos y, \quad \frac{\partial u}{\partial y}=-e^x\sin y $$
$$\frac{\partial v}{\partial x}=e^x\sin y, \quad \frac{\partial v}{\partial y}=e^x\cos y $$ болж $(E-D)$ нөхцөл биелэгдэх тул 
$$(e^z)'=e^x \cos y+ie^x\sin y= e^z, \qquad \forall z\in \mathbb{C}$$ болов.

Иймд $w=e^z$ функц $\mathbb{C}$ хавтгай дээр голоморф, өөрөөр хэлбэл $e^z\in H(\mathbb{C}).$ 
\end{sol}
\begin{exm}
 $w= \re z$ функц авч үзье.
\end{exm}
\begin{sol}
$u(x,y)=x$, $v(x,y)=0$ байх тул
$$\frac{\partial u}{\partial x}=1, \quad \frac{\partial u}{\partial y}=0, \quad \frac{\partial v}{\partial x}=0, \quad \frac{\partial v}{\partial y}=0$$
болж $(E-D)$ нөхцөл биелэгдэхгүй $\left(\dfrac{\partial u}{\partial x}\neq \dfrac{\partial v}{\partial y}\right)$ тул $w= \re z$ функц хаа ч дифференциалчлагдахгүй улмаар хаа ч голоморф биш функц. 
\end{sol}
\section{Голоморф ба гармоник функцүүдийн холбоо} 
$x,y$  аргументийн хоёр хувьсагчийн $u(x,y)$ функц $D$ муж дээр нэгдүгээр ба хоёрдугаар эрэмбийн тухайн уламжлалынхаа хамтаар тасралтгүй функцүүд байхын гадна $D$-мужийн цэг бүр дээр Лапласын
\begin{eqnarray} \label{3.11}
 \Delta u=\frac{\partial ^2u}{\partial x^2}+\frac{\partial ^2u}{\partial y^2}=0
\end{eqnarray}
тэгшитгэлийг \index{Лапласын тэгшитгэл} хангаж байвал уул функцийг $D$ муж дээр гармоник функц \index{Гармоник функц} гэнэ.

Лапласын тэгшитгэл ба гармоник функцүүд математик физикт онцгой үүрэгтэй. Жишээлбэл цэнэгээс үл хамаарах муж дахь цахилгаан орны потенциал, цахилгаан гүйдлээс үл хамаарах стационар соронзон орны скаляр потенциал зэрэг нь гармоник функцүүд байдаг. Голоморф ба гармоник функцүүдийн хоорондын холбоо голоморф функцийн физик хэрэглээнд ашиглагдана.

$D$ мужид голоморф $w=f(z)=u+iv$ функц өгөгдсөн байг. Тэгэхдээ түүний бодит ба хуурмаг хэсгүүд болох  \\ $u=u(x,y)$,$v=v(x,y)$ функцүүд нэгдүгээр ба хоёрдугаар эрэмбийн тухайн уламжлалуудын хамтаар тасралтгүй функцүүд гэж үз. (E-D) нөхцөл $D$ мужид биелэгдэж байг. (E-D)- ийн эхнийхийг $x$-ээр, нөгөөг $y$-ээр уламжлал авч хооронд нь нэмбэл (\ref{3.11}) гарч $u(x,y)$ функц гармоник функц болох нь тогтоогдоно. Яг үүнчлэн $v(x,y)$ функц ч гармоник болох нь батлагдана. Иймд дараах теорем томьёологдоно: 
\begin{thm}
Голоморф функцийн бодит ба хуурмаг хэсгүүд гармоник функц байна. Харин урвуу өгүүлбэр үнэн биш. Өөрөөр хэлбэл гармоник функцүүдийн аливаа хос голоморф функц үүсэх албагүй.

Гармоник $v$ функц $(E-D)$ буюу (\ref{3.9}) нөхцөлөөр гармоник $u$ функцтэй холбогдож байвал $v$ функцийг $u$-тэй хосмог функц гэнэ. Иймд дараах теорем хүчинтэй байна.
\end{thm}
\begin{thm}
Хоёр гармоник $u=u(x,y)$ ба $v=v(x,y)$ функц голоморф $w=f(z)=u+iv$ функц бүтээх зайлшгүй бөгөөд хүрэлцээтэй нөхцөл нь $v$ нь $u$-ийн хувьд хосмог байх явдал мөн.
\end{thm}
Одоо мэдэгдэж байгаа бодит буюу хуурмаг хэсгээр голоморф функцийг байгуулах асуудал шийдвэрлэе. Нэг холбоост $D$ муж дээр гармоник $u=u(x,y)$ функц нь олох гэж буй $f(z)=u+iv$ функцийн бодит хэсэг байг. Хуурмаг хэсэг $v(x,y)$ функцийг хоёрдугаар төрлийн муруй шугаман интеграл дүрстэй эрье:
$$v(x,y)=\int\limits_{(x_0,y_0)}^{(x,y)} \frac{\partial v}{\partial x}dx+\frac{\partial v}{\partial y}dy+C,$$
энд $(x_0,y_0)$ нь $D$-д бэхлэн авсан цэг, $C$-тогтмол, харин интегралыг $(x_0,y_0)$, $(x,y)$ цэгүүдийг холбосон $D$-д орших ямар нэг муруй (Зураг 5)-аар авна.
\begin{center}
    \includegraphics[width=4cm]{Zurag/5.jpg}
    \captionof{figure}{}
        \label{ref hiih ner}
\end{center}
(E-D) нөхцөл ашиглавал
\begin{eqnarray} \label{3.12}
 v(x,y)=\int\limits_{(x_0,y_0)}^{(x,y)} -\frac{\partial u}{\partial y}dx+\frac{\partial u}{\partial x}dy+C.
\end{eqnarray}
(\ref{3.12})дэх интеграл замаас хамаарахгүй, учир нь $u(x,y)$ функцийн гармоник чанар ёсоор 
$$\frac{\partial}{\partial y}\left(-\frac{\partial u}{\partial y}\right) =\frac{\partial}{\partial x}\left(\frac{\partial u}{\partial x}\right)$$
нөхцөл биелэгдэнэ. 

Гармоник $v=v(x,y)$ функц $D$-д өгөгдсөн байвал түүнтэй хосмог гармоник $u(x,y)$ функц
\begin{eqnarray}
 u(x,y)=\int\limits_{(x_0,y_0)}^{(x,y)} \frac{\partial v}{\partial y}dx-\frac{\partial v}{\partial x}dy+C
\end{eqnarray}
гэж олдох нь мэдээж.
\begin{exm}
 $u=x^2-y^2+2x$ функцээр голоморф $f(z)=u+iv$ функцийг $f(i)=2i-1$ нөхцөлд байгуул.
\end{exm}
\begin{sol}
$$\frac{\partial u}{\partial x}=2x+2,\quad \frac{\partial^2 u}{\partial x^2}=2, \quad \frac{\partial u}{\partial y}=-2y, \quad \frac{\partial^2 y}{\partial y^2}=-2$$ болж $\forall x$, $\forall y$ ийн хувьд $\Delta u=0$ байх тул дурдсан функц нийт $Oxy$ дээр гармоник байх ба хялбарыг харгалзан $x_0=0$, $y_0=0$ гэе.
(\ref{3.12}) ёсоор
\begin{eqnarray} \label{3.14}
 v(x,y)=\int\limits_{(0,0)}^{(x,y)} -(-2y)dx+(2x+2)dy+C
\end{eqnarray} болно. $\dfrac{\partial }{\partial x}\left(\dfrac{\partial u}{\partial x}\right) =\dfrac{\partial }{\partial y}\left( -\dfrac{\partial u}{\partial y}\right)$ 
байх тул (\ref{3.14}) нь замаас хамаарахгүй, зураг 6-д үзүүлснээр $OA\cup AB$ замаар интегралчилж $v(x,y)$ функцийг бодож гаргана. \\

    \begin{figure}[h]
    \centering
\begin{tikzpicture}[line cap=round,line join=round,>=triangle 45,x=1cm,y=1cm]
\clip(0,0) rectangle (6,4.5);
\coordinate (x1) at (0,1);
\coordinate (y1) at (1,0);
\coordinate (x2) at (5.5,1);
\coordinate (y2) at (1,4);
\coordinate (o) at (1,1);
\coordinate (z) at (4,3);
\draw[-stealth,line width=1pt] (x1)  -- (x2) node[right] {$x$};
\draw[-stealth,line width=1pt] (y1)  -- (y2) node[right] {$y$};
\draw (o) node[anchor=north east] {$O$};
\draw[-stealth,line width=1pt] (4,1) node[below] {$A$} -- (z) node[above right ] {$B(x,y)$};
\draw[-stealth,line width=1pt] (o)  -- (4,1);
\end{tikzpicture}
    \captionof{figure}{}
        \label{ref hiih ner}
\end{figure}
Иймд 
$$z\in OA: 0\leq x, \quad y=0 \Rightarrow dy=0,$$
$$z\in AB: x=const \Rightarrow dx=0, \quad 0\leq y$$ болж 
$$v(x,y)=\int\limits_0^x 0\cdot dx+\int\limits_0^y (2x+2)dy+C=(2x+2)y+C$$
бөгөөд 
\begin{align*}
    f(z)&=x^2-y^2+2x+[(2x+2)y+C]i\\
    &=x^2-y^2+2xyi+2(x+iy)+Ci \\
    &= z^2+2z+Ci ,\quad f(i)=2i-1
\end{align*}
нөхцөлийг ашиглавал 
$C=0$ тул $f(z)=z^2+2z\in H(\mathbb{C}).$
\end{sol}
Голоморф функц байгуулахад муруй шугаман интеграл бодоход хүргэхгүй байж болно. Үүнийг саяын жишээн дээр харуулъя:
$$\frac{\partial v}{\partial x}=-\frac{\partial u}{\partial y}=2y\Rightarrow v(x,y)=\int 2ydx+\varphi (y)=2yx+\varphi (y)$$
тул $\dfrac{\partial v}{\partial y}=2x+\varphi ' (y)$ ба $\dfrac{\partial v}{\partial y}=\dfrac{\partial u}{\partial x}$ ёсоор 
$$2x+\varphi ' (y)=2x+2\Rightarrow \varphi ' (y)=2\Rightarrow \varphi (y)=2y+C$$
болж улмаар $v(x,y)=2xy+2y+C$ гарлаа. Цаашид нөхцөлөө ашиглан $f(z)=z^2+2z$ функцээ гаргана.
\section{Уламжлалын геометр утга. Конформ буулгалт }
Комплекс $\circled{z}$ хавтгай дээр $\Gamma$-муруй $t$ параметрт
\begin{eqnarray}
  \Gamma:z=z(t)=x(t)+iy(t), \qquad \alpha\leq t\leq \beta
\end{eqnarray}
тэгшитгэлээр, өөрөөр хэлбэл бодит аргументийн $x=x(t)$, $y=y(t)$ гэсэн параметрт тэгшитгэлээр тодорхойлогдоно.
\begin{eqnarray} \label{3.16}
  z'(t_0)=\lim\limits_{\Delta t\rightarrow 0}\frac{\Delta z(z_0)}{\Delta t}=x'(t_0)+iy'(t_0), \qquad t_0\in (\alpha, \beta)
\end{eqnarray}
бөгөөд (\ref{3.16})-аас харахад хэрэв $\exists z'(t_0)$, $z'(t_0)\neq 0$ байвал комплекс (\ref{3.16}) вектор нь $\Gamma$ муруйн $z_0=z(t_0)=(x(t_0), y(t_0))$ цэг дээр татсан шүргэгчийн дагуу чиглэнэ.

$w=f(z)$ нь $z_0$ цэг дээр голоморф функц байг. Бас \\ $w_0=f(z_0)$, $f'(z_0)\neq 0$ байг. Энэ үед $z_0$-ийн орчны $\forall z$ цэг бүр
$w_0$-ийн орчны ямар нэг $w$ цэгт бууна.

$z_0$ цэгийг дайруулан түүн дээр огтлолцсон толигор \\ $\Gamma_1:z_1=z_1(t)$, $\Gamma_2:z_2=z_2(t)$ шугам авъя. 
\begin{center}
    \includegraphics[width=9cm]{Zurag/7.PNG}
    \captionof{figure}{}
        \label{ref hiih ner}
\end{center}
$z_0=z_1(t_{01})$ ба $z_0=z_2t(_{02})$ байг. $w=f(z)$ буулгалтаар эдгээр шугам $\circled{w}$ дээр $w_0$ цэгт огтолцсон $L_1$, $L_2$ шугамд бууна (Зураг 7). $$L_1:w=w_1(t_1)=f(z_1(t_1));$$ $$L_2:w=w_2(t_2)=f(z_2(t_2))$$ бөгөөд харин $w_0=w_1(t_{01})=f(z_0)$; $w_0=w_2(t_{02})=f(z_0).$

Давхар $w=f[z(t)]$ функцийн уламжлалын дүрмээр:
\begin{equation} \label{3.17}
  \begin{split}
      &w'_1(t_{01})=f'(z_0)\cdot z_1'(t_{01}); \quad w'_2(t_{02})=f'(z_0)\cdot z_2'(t_{02})\\
      &dw'_1(t_{01})=f'(z_0)\cdot z_1'(t_{01})dt=f'(z_{0})\cdot dz_1(t_{01})\\
      & dw'_2(t_{02})=f'(z_0)\cdot z_2'(t_{02})dt_1=f'(z_0)\cdot dz_2(t_{02})
  \end{split}
\end{equation}
Бид одоо, уламжлалын модуль ба аргументийн геометр утгын талаар дараах хоёр теорем батлах боломжтой боллоо. 
\begin{thm}
Голоморф $w=f(z) \quad (f'(z_0)\neq 0)$ функцээр буулгалт гүйцэтгэхэд $\Gamma$ муруйн дүр болох $L$-ийн хувьд олдох $|dw|$ ба $|dz|$-ийн харьцаа нь $z_0$ цэг дээр тогтмол бөгөөд $|f'(z_0)|(>0)$-тай тэнцүү байна.

Өөрөөр хэлбэл $w=f(z)$, $f'(z_0)\neq 0$ функцээр буулгалт гүйцэтгэхэд $z_0$ цэг дээрх шугаман суналтын коэффициент нь $|f'(z_0)|$ байна.
\end{thm}
\begin{prf}
(\ref{3.17}) томьёог ашиглавал 
$$\frac{|dw_1(t_{01})|}{|dz_1(t_{01})|}=|f'(z_0)|; \qquad \frac{|dw_2(t_{02})|}{|dz_2(t_{02})|}=|f'(z_0)|$$
$\Gamma_1$ ба $\Gamma_2$ шугамыг дураараа сонгоно. 
\hfill $\blacksquare$\end{prf}
\begin{thm}
Голоморф $w=f(z)$, $f(z_0)\ne 0$ функцээр буулгалт гүйцэтгэхэд $z_0$ цэг дээрх шүргэгч векторын эргэлтийн өнцөг $\Arg  f'(z_0)$-тэй тэнцүү бөгөөд муруйнуудын $z_0$ дээрх өнцгүүд хадгалагдана.
\end{thm}
\begin{prf}
 $z_0$ цэг дээр $\Gamma_1$ ба $\Gamma_2$ шугамуудын хоорондох өнцөг (нар буруу тойролтоор тоологдох)
 $$\varphi =\Arg  dz_2(t_{02})-\Arg  dz_1(t_{01})$$
 байх бөгөөд (\ref{3.17}) ёсоор харгалзах $L_1$ ба $L_2$ муруйнуудын $w_0$ цэг дээрх $\theta$ өнцөг
\begin{align*}
    \theta&= \Arg  dw_2(t_{02})-\Arg  dw_1(t_{01}) \\
    &=\Arg  f'(z_0)+\Arg  dz_2(t_{02})-\Arg  f'(z_0)-\Arg  dz_1(t_{01})\\
    &=\Arg  dz_2(t_{02})-\Arg  dz_1(t_{01})=\varphi
\end{align*}
болж, вектор бүхний хувьд
$$\Arg  dw (t_0)=\Arg  f'(z_0)+\Arg  z(t_0)$$
байна. Теорем батлагдав. 
\hfill $\blacksquare$\end{prf}
\begin{hnt}
Энэ тохиолдолд муруйнуудын хоорондох өнцгүүд тэнцүү төдийгүй тэдгээрийн тоологдох чиглэл адилхан байна (цөм нар буруу тойрч тоологдох).
\end{hnt}
\begin{dfn}
$z_0$ цэгийг дайрсан муруйнуудын хоорондох өнцгийг хэвээр хадгалж, чиглэл бүрдээ шугаман суналт тогтмол байх $w=f(z)$ буулгалтыг $z_0$ цэг дээр конформ буулгалт \index{Конформ буулгалт} гэнэ. Өнцгийн тоологдох чиглэл таарч байвал I төрлийн, \index{I төрлийн конформ буулгалт} эсрэг байвал II төрлийн конформ буулгалт \index{II төрлийн конформ буулгалт} гэдэг. 
\end{dfn}
\begin{thm}
$w=f(z)$ функц $z_0$ цэг дээр голоморф бөгөөд $f'(z_0)\neq 0$ байвал тэр нь $z_0$ цэг дээр конформ буулгалт өгнө.

$D$ мужийн цэг бүр дээр конформ $w=f(z)$ буулгалтыг $D$ дээр конформ гэнэ. $D$ муж дээр голоморф $w=f(z)$ функц $D$ мужийг $D^*$ муж руу конформ буулгахын тулд $\forall z\in D$ дээр $f'(z)\neq 0$, түүнчлэн $D$ дээр уул функц нэг хуудаст $(\forall z_1,z_2\in D$ үед $f(z_1)\neq f(z_2)$ байх нь зөвхөн $z_1\neq z_2$ үед) байхыг шаардана.
\end{thm}
\begin{hnt}
$w=f(z)$ функц $z_0$ цэгийн ямар нэг орчинд нэг хуудаст байвал $z_0$ цэг дээр нэг хуудаст гэнэ. Голоморф $w=f(z)$ функц $f'(z_0)\neq 0$ байвал тэр нь $z_0$ дээр нэг хуудаст байна. Харин $D$ мужид нэлэнхийдээ $f'(z)\neq 0$ байхаас $w=f(z)$ функц $D$ дээр нэг хуудаст байх нь алба биш. Жишээ нь $w=z^2$ функц цагираг хэлбэрийн $D:1<|z|<2$ мужийн цэг бүр дээр нэг хуудаст боловч тэр нь $D$-дээр нэг хуудаст биш. Учир нь $\forall z\in D$-ийн хувьд $f(-z)=f(z)$ байна. 
\end{hnt}
\begin{exm}
$f(z)=z^2$ буулгалтын $1+\sqrt{3}i$ цэг дээрх шугаман суналтын коэффициент ба эргэлтийн өнцгийг ол.
\end{exm}
\begin{sol}
$w'=2z$, $|2(1+\sqrt{3}i)|=4$, $\arg  2(1+\sqrt{3}i)=\dfrac{\pi}{3}.$\\ $f(z)=z^2$ функц нийт $\mathbb{C}$ дээр голоморф бөгөөд $z_0$ дээр конформ буулгалт өгнө.
$$|f'(1+\sqrt{3}i)|=4, \quad \arg  f'(1+\sqrt{3}i)=\frac{\pi}{3}$$ тул өгөгдсөн цэг дээр шугаман суналтын коэффициент 4 ба $\dfrac{\pi}{3}$ өнцгөөр эргэлт явагдана.
\end{sol}
\begin{exm}
$w=\sin z$ функц $1+i$ цэг дээр ямар суналт, эргэлт өгөх вэ?
\end{exm}
\begin{sol}
$w'=\cos z$ тул $k=|\cos (1+i)|=\sqrt{ \ch^2 1-\sin^2 1}$; \qquad $\theta = \arccos (1+i)=-\arctan 11$ байна. 
\end{sol}
\begin{center}
    Дасгал
\end{center}
\begin{enumerate}
    \item [13.]
    Дараах функцийн бодит ба хуурмаг хэсгийг ол.
    \colgur
    {\begin{enumerate}
        \item [a)] $w=z^3,$
        \item [b)] $w=ze^{-z},$
        \item [c)] $w=\sin\overline{z}.$
    
    \end{enumerate}}
    \item[14.] Дараах функцийн бодит ба хуурмаг хэсгийг ол.
    \colgur{
    \begin{enumerate}
    \item [a)]$\sin z$
    \item [b)] $ch z$
    \item[c)] $\Ln i$
    \item [d)] $ln (1-i)$
    \item [e)] $2^i$
    \item [(f)] $(1+i)^i$
    \end{enumerate}
    }
    \item[15.] Мэдэгдэж байгаа бодит буюу хуурмаг хэсгээр голоморф функц $w= f(z)$ байгуул.
    \begin{enumerate}[a)]
        \item  $u=\dfrac{x}{x^2+y^2}, \qquad f(\pi)=\dfrac{1}{\pi}$
        \item $v=xy-2y, \qquad f(0)=0$
        \item $v=2 \ch x \sh y-xy, \qquad f(0)=0.$
    \end{enumerate}
\end{enumerate}
%\chapter{Интеграл} \index{Интеграл}
\section{Интегралын тодорхойлолт}
$\gamma :[\alpha,\beta]\rightarrow \mathbb{C}$ (Ийм буулгалтыг зам гэж нэрлэдэг) ба $f\in C(D)$ функц авч үзье. $\gamma ([\alpha, \beta])\subset D$ байг. $T$ нь $[\alpha, \beta]$ интервалийн хуваалт байг.
$$T=\{t_k:k\in\overline{0,n} \quad \text{ба} \quad \alpha =t_0<t_1<\ldots <t_n=\beta\}$$
$z_k=\gamma(t_k)$ ба $\textrm{\normalfont{diam} }T=\max\limits_{k\in \overline{1,n}} (t_k-t_{k-1})$ тэмдэглээнүүд оруулъя. Мөн $k\in\overline{1,n}$ бүрийн хувьд $\tau _k \in [t_{k-1}, t_k]$ сонгосноор $\zeta =\{ \zeta_k:\zeta_k=\gamma(\tau_k)\}$ олонлог байгуулъя (Үүнийг сонголт гэж нэрлэдэг).
\begin{eqnarray}
  S(f,\gamma ,T,\zeta)=\sum\limits_{k=1}^n f(\zeta_k)\cdot (z_k-z_{k-1})
\end{eqnarray}
хэмжигдэхүүнийг $f$ функцийн хувьд $\gamma$ замын $T$ хуваалт ба $\zeta$ сонголтоор зохиогдсон интеграл нийлбэр гэнэ.
\begin{dfn}
 \textbf{(Комплекс интеграл)}. Хэрэв $T$ ба $\zeta$ сонголтыг хэрхэн байгуулснаас үл хамаарч
 \begin{eqnarray}
   \exists \lim\limits_{\textrm{\normalfont{diam} }T\rightarrow 0}S(f,\gamma, T,\zeta)=I\in\mathbb{C}
 \end{eqnarray}
 бол I комплекс тоон утгыг $f$ функцээс $\gamma$ замаар авсан интеграл гэж нэрийдээд  $$I=\int\limits_\gamma f(z)dz$$ гэж тэмдэглэдэг.
\end{dfn}
\begin{thm}
 \textbf{(Интеграл орших зайлшгүй нөхцөл.)} $\gamma$ нь гөлгөр зам ба $f$ нь уг замын муруй дээр тасралтгүй бол
 \begin{eqnarray} \label{4.3}
   \exists \int\limits_{\gamma} f(z)dz=\int\limits_\alpha^\beta f(\gamma(t))\cdot \gamma '(t)dt 
 \end{eqnarray}
 байна. (\ref{4.3}) томьёог интеграл бодох параметрчилэлийн томьёо гэдэг.
\end{thm}
\begin{prf}
$\gamma (t)=a(t)+ib(t)$ ба $\gamma$ гөлгөр гэдгээс $a(t)$, $b(t)\in C^{(1)}[\alpha,\beta]$ болж Лагранжийн томьёог ашиглах боломжтой.
$$
\begin{cases}
a(t_k)-a(t_{k-1})=a'(\xi_k)(t_k-t_{k-1}) \\
b(t_k)-b(t_{k-1})=b'(\xi_k')(t_k-t_{k-1}), \quad (\xi_k,\xi_k'\in[t_{k-1},t_k]) 
\end{cases}
$$ 

болно.

$a'$, $b'\in C[\alpha,\beta]$ учир Канторын теоремоор жигд тасралтгүй байх ёстой. $\forall \varepsilon >0 \exists \delta >0 \quad \forall t_1,$ $t_2\in [\alpha, \beta]$ хувьд $|t_1-t_2|<\delta$ бол $|a'(t_1)-a'(t_2)|<\varepsilon$; $|b'(t_1)-b'(t_2)|<\varepsilon$ болно. Хуваалтын диаметрийг $\delta$-с багаар авбал 
\begin{align*}
    &\left|\sum\limits_{k=1}^n f(\zeta_k)(\gamma(t_k)-\gamma(t_{k-1}))-\sum\limits_{k=1}^nf(\zeta_k)\gamma'(\tau_k)(t_k-t_{k-1})\right|=\\
    &=\left| \sum\limits_{k=1}^n f(\zeta_k)(a'(\xi_k)-a'(\tau_{k}))+i\sum\limits_{k=1}^nf(\tau_k)(b'(\xi'_k)-b(\tau_{k}))\right|\leq \\
    &\leq 2\varepsilon \max\limits_{z\in \gamma [\alpha,\beta]}|f(z)|\cdot (\beta-\alpha)
\end{align*}
болно.

$\varepsilon$-г дураар сонгож болох учир
$$\left|\int\limits_\gamma f(z)dz-\int\limits_\alpha^\beta f(\gamma(t))\gamma '(t)dt\right|=0$$
болж теорем батлагдав. 
\end{prf} 
\begin{cor}
$$\int\limits_\gamma f(z)dz=\int\limits_\gamma udx-vdy+i\int\limits_\gamma vdx+udy.$$
\end{cor}
\begin{cor}
$\gamma$ нь хэсэгчилсэн гөлгөр зам үед дээрх теорем хүчинтэй.
\end{cor}
\begin{exm} \label{Example z^n int}
$\gamma$ нь $a\in\mathbb{C}$ төвтэй $R$ радиустай тойргийг цагийн зүүний эсрэг 1 эргэсэн $\gamma(t)=a+R\cdot e^{it}(t\in[0;2\pi])$ зам байг.
\begin{align*}
    \int\limits_\gamma (z-a)^n dz&=\int\limits_{0}^{2\pi} ((a+Re^{it})-a)^n\cdot(a+Re^{it})'_tdt \\
    &=R^{n+1}\cdot i\int\limits_0^{2\pi} [\cos (n+1)t+i\sin (n+1)t]dt \\
    &=
    \begin{cases}
    0, \quad n\neq -1, \\
    2\pi i, \quad n=-1
    \end{cases}
\end{align*}
\end{exm}
\begin{exm}
$\gamma(t)=t+it^2(t\in [0;1])$ параболын нумаар 
\begin{align*}
    \int\limits_\gamma \re (z^2)dz&=\int\limits_0^1 \re (t+it)^2\cdot (t+it^2)_t'dt \\
    &=\int\limits_0^1 (t^2-t^4)(1+2itdt) \\
    &=\dfrac{2}{15}+\dfrac{1}{6}\cdot i
\end{align*}
\end{exm}
\begin{exm}
$\gamma$ нь $z_0$ цэгээс $z_1$ хүрсэн хэрчим зам \\ $\gamma(t)=z_0+(z-1-z_0)t$ \quad $(t\in[0,1])$ байг.
\begin{align*}
    \int\limits_{\gamma^+}(az+b)dz&= \int\limits_0^1 [a(z_0+(z_1-z_0)t)+b]\cdot (z_0+(z_1-z_0)t)'_tdt \\
    &=a\dfrac{z_1^2-z_0^2}{2}+b(z_1-z_0).
\end{align*}
\end{exm}
\begin{exm} \label{Жишээ 4}
Гурвалжны хилээр тойрсон зам контур $\partial \Delta^+$ (битүү замыг контур гэж нэрлэх нь байдаг) бол $$\bint\limits_{\partial\Delta^+} (az+b)dz=0$$ байхыг жишээ 4.3-аас гаргана уу.
\end{exm}
\section{Комплекс интегралын чанарууд}
Энд $\gamma$ нь хэсэгчилсэн гөлгөр зам ба $f$, $g\in C(\gamma [\alpha, \beta])$ (замын муруй дээр тасралтгүй) функцүүд байг.
\begin{sen}
(Шугаман чанар) $\forall a,b\in\mathbb{C}$ хувьд
$$\int\limits_{\gamma} [af(z)+bg(z)]dz=a\int\limits_\gamma f(z)dz+b\int\limits _\gamma g(z)dz$$
байна.
\end{sen}
\begin{sen}
(Аддитив чанар) $\gamma =\gamma_1 \wedge \gamma_2$ бол 
$$\int\limits_{\gamma} f(z)dz=\int\limits_{\gamma_1} f(z)dz+\int\limits _{\gamma_2} f(z)dz$$
байна.
\end{sen}
\begin{sen} \label{Чанар 3}
$\gamma^{-}$ эсрэг замын хувьд
$$\int\limits_{\gamma^-}f(z)dz =-\int\limits_\gamma f(z)dz$$
\end{sen}
\begin{hnt}
Битүү замын (контур) хувьд зүүн гарын дүрмээр ориентацлаж чиглэл тогтоовол $\gamma^+$ гэж тэмдэглэдэг.
\end{hnt}
\begin{center}
 \includegraphics[width=6cm]{Zurag/8.jpg}  
 \captionof{figure}{}
        \label{ref hiih ner}
\end{center}
\begin{prf}
Комплекс интеграл нь муруй шугаман 2-р төрлийн интегралаар илэрхийлэгддэг нь теорем 1-ийн мөрдлөгөө болсон. Мөн дээрх чанарууд түүнийх нь хувьд биелдгийг Математик анализын хичээл дээр үзсэнээ санавал чанарууд батлагдаж байна.
\end{prf}
\begin{sen}
(Интегралыг үнэлэх).
$$\left|\int\limits_\gamma f(z)dz\right|\leq \int\limits_\gamma |f|dl\leq |\gamma|\cdot \max\limits_{z\in\gamma [\alpha,\beta]}|f|.$$
\end{sen}
\begin{prf}
$|dz|=\sqrt{(dx)^2+(dy)^2}=dl$ (Уртын дифференциал) байна.

Иймд гурвалжны тэнцэл биш ашиглавал
$$\left|\int\limits_\gamma f(z)dz\right|\leq \int\limits_\gamma \left|f(z)dz\right|=\int\limits_l |f(z)|dl$$ болж эхний тэнцэл биш
батлагдав.

Эерэг функцээс уртын дифференциалаар авч байгаа учир
$$\int\limits_\gamma \left(\max\limits_{z\in \gamma[\alpha, \beta]}|f|-|f(z)|\right)dl\leq 0\Rightarrow$$
$$\int\limits_\gamma|f(z)|dl\leq \int\limits_\gamma \max_{z\in \gamma [\alpha,\beta]}|f|dl=\max\limits_{z\in \gamma[\alpha,\beta]}|f|\int\limits_\gamma dl=\left(\max\limits_{z\in\gamma [\alpha,\beta]}|f|\right) \cdot |\gamma|$$
болж дараагийн тэнцэл биш батлагдав.
\hfill \end{prf}
\section{Кошийн интеграл теорем} \index{Кошийн интеграл теорем}
\begin{thm}
\textbf{ (Гурвалжин контурын Кошийн теорем).} $f\in H(D)$ ба дотоод муж нь хилийн хэрчмүүдтэйгээ $D$ мужид агуулагдсан гурвалжин $\Delta$ авъя $(\overline{\Delta}\subset D)$. Тэгвэл 
$$\int\limits_{\partial \Delta^+}f(z)dz=0$$
байна. 
\end{thm}
\begin{prf}
Эсрэгээс нь $\exists \Delta_0:\overline{\Delta_0}\subset D$ боловч \\ $\left|\bint\limits_{\partial \Delta_0^+}f(z)dz\right|=M>0$ гэж үзье.
\begin{center}
 \includegraphics[width=6cm]{Zurag/9.jpg}  
 \captionof{figure}{}
        \label{ref hiih ner}
\end{center}
Гурвалжныг зурагт үзүүлсэн аргаар дундаж шугамуудаар нь хуваавал Чанар 4.3 ёсоор $\bint\limits_{\partial \Delta_0^+} f(z)dz$ нь дөрвөн интегралын нийлбэр болж сална. Иймд эдгээрээс модуль нь $\dfrac{M}{4}$ багагүй байх дор хаяж нэг гурвалжин сонгогдож чадна. Үүнийг нь $\Delta_1$ гэе. 
$$\left|\int\limits_{\partial \Delta_1^+} f(z)dz\right|\geq \frac{M}{4}$$
Энэ алхмыг үргэлжлүүлэн $n$ удаа хийсэн гэвэл $\exists \Delta_n:\Delta_n$ нь $\Delta_0$ гурвалжингаас $2^n$ дахин төсөөтэйгөөр багассан ба 
\begin{equation}
    \tag{$*$} \label{*}
    \left|\int\limits_{\partial \Delta_n^+} f(z)dz\right|\geq \frac{M}{4^n}
\end{equation}
байна.

Багталцсан бөмбөрцгүүдийн зарчим ёсоор 
$$\exists z_0\in D:z_0=\bigcap\limits_{n=1}^\infty \overline{\Delta_n} $$
байна.
$f\in H(d)$ нөхцөл ёсоор 
$$\forall \varepsilon >0: \exists \delta>0 \forall z\in U=\{|z-_0|<\delta\}\cap D$$ хувьд
$$
\begin{cases}
 f(z)=f(z_0)+f'(z_0)(z-z_0)+\alpha (z)(z-z_0) \\
 |\alpha (z)|<\varepsilon
\end{cases}
$$ байна. 
$\Delta_0\subset U$  гэж үзэхэд явцуурахгүй нь ойлгомжтой. Иймд
$$\int\limits_{\partial \Delta_n}f(z)dz=\int\limits_{\partial \Delta_n} [f(z_0)+f'(z_0)(z-z_0)]dz+\int\limits_{\partial \Delta_n}\alpha (z)(z-z_0)dz$$
ба жишээ 4.4 ёсоор эхний нэмэгдэхүүн тэг болно. Харин дараагийн нэмэгдэхүүнийг үнэлбэл
\begin{equation}
\tag{$**$} \label{**}
    \begin{split}
    \left| \int\limits_{\partial \Delta_n}\alpha(z)(z-z_0)dz\right|&\leq \max |\alpha(z)|\cdot \max_{z\in \partial \Delta_n}|z-z_0|\cdot P_{{\Delta_n}}< \\
    &<\varepsilon \cdot P_{\Delta_n}\cdot P_{\Delta_n}=\varepsilon\cdot \frac{P^2_{\Delta_0}}{4^n}
\end{split}
\end{equation}
байна. (Энд $P_{\Delta_n}$ нь $\Delta_n$ гурвалжны периметрийг авав). ($*$) (\ref{**}) тэнцэл бишүүдээс
$$\frac{M}{4^n}<\varepsilon \frac{P^2_{\Delta_0}}{4^n}\Rightarrow \frac{M}{P^2_{\Delta_0}}<\varepsilon$$
болж $\varepsilon$ дураараа сонгож болно гэдэгтэй харшилдав. Иймээс Кошийн теорем гурвалжин контурын хувьд батлагдаж байна.
\hfill \end{prf}
\begin{thm} \label{Теорем 3}
\textbf{(Тахир шугаман контурын Кошийн теорем).} $f\in H(D)$, $\gamma$ нь битүү жорданы тахир шугам (дотоод мужийн хамт $D$ мужид оршдог) байг. Тэгвэл $\bint\limits_{\gamma^+}f(z)dz=0$ байна.
\end{thm}
\begin{prf}
Аливаа тахир шугамыг дотоод диагональ татсаар төгсгөлөг тооны гурвалжнуудад хувааж (триангуляци хийж) болдгийг дасгал болгон батална уу. Тэгвэл интегралын аддитив чанар ёсоор $\bint\limits_{\gamma^+} f(z)dz$ нь гурвалжны хилээр авсан интегралуудын нийлбэрт задарна. Одоо Теорем 4.2 ашиглавал энэ теорем батлагдана. 
\hfill  \end{prf}
\begin{thm} \label{Теорем 4}
\textbf{(Гөлгөр контурын Кошийн теорем).} \\ $f\in H(D)$, $\gamma$ нь битүү жорданы гөлгөр зам (дотоод мужтайгаа хамт $D$ мужид оршдог) байг. Тэгвэл $\bint\limits_{\gamma^+}f(z)dz=0$ байна.
\end{thm}
\begin{prf}
Интеграл нийлбэр зохиох дурын $T$ хуваалт нь битүү тахир шугам өгнө. Энд Теорем 4.3-г ашиглах хэрэгтэй. 
\hfill  \end{prf}
\begin{dfn}
 $D$ муж өгөгдсөн.$\gamma_1:[0,1]\rightarrow D$,\\ $\gamma_2:[0,1]\rightarrow D$ замуудын хувьд $\gamma_1(0)=\gamma_2(0)$, $\gamma_1(1)=\gamma_2(1)$  ба 
 $$\exists \gamma (t,s)\in C ([0,1]^2\rightarrow D):\gamma_1(t)=\gamma (t,0); \quad \gamma_2(t)=\gamma(t,1)$$
 бол  $\gamma_1$, $\gamma_2$ замуудыг $D$ мужид нэг эхлэл төгсгөлтэй гомотоп замууд гэж нэрлээд $\gamma_1 \undertilde{D} \gamma_2$ тэмдэглэдэг.
\end{dfn}
\begin{dfn}
 $D$ муж өгөгдсөн. $\gamma_1:[0,1]\rightarrow D$,\\ $\gamma_2:[0,1]\rightarrow D$ замуудын хувьд $\gamma_1(0)=\gamma_1(1)$ ба $\gamma_2(0)=\gamma_2(1)$  ба 
 $$\exists \gamma (t,s)\in C ([0,1]^2\rightarrow D):\gamma_1(t)=\gamma (t,0); \quad \gamma_2(t)=\gamma(t,1)$$
 бол  эдгээр замуудыг $D$ мужид битүү замын гомотоп гэж нэрлээд $\gamma_1\undertilde{D} \gamma_2$ тэмдэглэдэг.
\end{dfn}
\begin{center}
 \includegraphics[width=8cm]{Zurag/10.PNG}  
 \captionof{figure}{}
        \label{ref hiih ner}
\end{center}
\begin{cor} 
Теорем 4.4-г дараах байдлаар өөрчлөн томьёолж болно. $f\in H(D)$ нь битүү гөлгөр Жорданы зам ба $\gamma \undertilde{D} 0$ бол
$\bint\limits _\gamma f(z)dz=0$ байна.
\end{cor}
\begin{thm} \label{Теорем 5}
\textbf{(Кошийн теоремын гомотоп хэлбэр).} \\ $f\in H(D)$, $\gamma_1 \undertilde {D} \gamma_2$ гөлгөр замууд байг. Тэгвэл $$\int\limits_{\gamma_1} f(z)dz =\int\limits_{\gamma_2} f(z)dz$$ байна.
\end{thm}
\begin{prf}
Эхлээд $\gamma_1\undertilde{D}\gamma_2$ эхлэл төгсгөл давхацсан гомотоп байг. Теорем 4.4 ёсоор $\bint\limits_{\gamma_1\wedge \gamma_2^{-}} f(z)dz=0$. Улмаар 
$$\int\limits_{\gamma_1}f(z)dz-\int\limits _{\gamma_2} f(z)dz=0$$
болж теорем батлагдана.

Одоо $\gamma_1\undertilde{D}\gamma_2$ битүү замын гомотоп байг.
\begin{center}
 \includegraphics[width=8cm]{Zurag/11.PNG}  
 \captionof{figure}{}
        \label{ref hiih ner}
\end{center}
Зургийг ажиглан энэ тохиолдолд Теорем 4.4-г хэрхэн хэрэглэхийг олно уу.
\hfill  \end{prf}
\begin{thm}\label{removeablediscontinousfunctionint}
$f\in H(D)$, $\gamma$ битүү гөлгөр жорданы зам ба түүний дотоод мужийг $int \gamma$ гэж тэмдэглэе. $z_0\in int \gamma$ хувьд $z_0\notin D$ ба $(\overline{in t \gamma}\setminus \{z_0\})\subset D$ байг. Харин $f$ функцийн хувьд $z_0$ засагдах тасралтын цэг (өөрөөр хэлбэл $\exists \lim\limits_{z\rightarrow z_0}f(z)=a\in \mathbb{C})$ байг. Тэгвэл $\bint\limits_\gamma f(z)dz=0$ байна.
\end{thm}
\begin{prf}
$C_\varepsilon (z_0)$ нь $z_0$ төвтэй $\varepsilon$ радиустай тойрог байг. $\varepsilon$ хангалттай бага үед $\gamma \undertilde{D} C_{\varepsilon} (z_0)$ тул Теорем 4.5 ёсоор $$\int\limits_{\gamma}f(z)dz=\int\limits_{C_\varepsilon (z_0)} f(z)dz$$
ба $$\left|\int\limits_{C_\varepsilon(z_0)} f(z)dz\right|\leq 2\pi \cdot \varepsilon \cdot \max_{z\in C_{\varepsilon} (z_0)}|f(z)| =2\pi \cdot \varepsilon (|a|+\circ (1))$$
болно. $\varepsilon$-г хэчнээн ч багаар сонгож болох тул байна. 
$$\int\limits_{\gamma} f(z)dz=0$$ байна. 
\hfill  \end{prf}
\begin{thm}
\textbf{(Олон холбоост мужийн Кошийн теорем).} \index{Олон холбоост муж дахь Кошийн интеграл теорем} $D$ нь $n$ холбоост зааглагдсан муж ба $f\in H(D)\cap C(\overline{D})$ байг. Тэгвэл $$\bint\limits_{\partial D^+} f(z)dz=0$$ байна.
\end{thm}
\begin{prf}
Теорем 4.4-рүү хэрхэн шилжүүлж байгааг зургаас харна уу.
\begin{center}
 \includegraphics[width=8cm]{Zurag/12.jpg}  
 \captionof{figure}{}
        \label{ref hiih ner}
\end{center}
\hfill  \end{prf}
\begin{hnt}
Кошийн интеграл теоремууд нь гармоник функцийн хувьд хүчинтэй байдаг. 
\end{hnt}
\begin{exm}
$I=\bint\limits_0^{+\infty}\cos x^2dx$; $J=\bint\limits_0^{+\infty}\sin x^2dx$ интегралуудын утгыг тооцоол.
\end{exm}
\begin{sol}

Зурагт үзүүлсэн секторын хил контурыг $\gamma^+$ гэе.
$$\gamma^+=\gamma^+_1\vee \gamma_2^+\vee \gamma_3^+$$ 
$\gamma^+$ нь 2 хэрчим, 1 тойргийн нумаас тогтоно.
\begin{center}
    \includegraphics [width=6cm]{Zurag/13.jpg}
 \captionof{figure}{}
        \label{ref hiih ner}
\end{center}
Кошийн теорем ёсоор
\begin{equation}
\tag{$*$} \label{7*}
  0=\int\limits_{\gamma^+} e^{iz^2}dz=\int\limits_{\gamma_1^+} e^{iz^2}dz+\int\limits_{\gamma_2^+} e^{iz^2}dz+\int\limits_{\gamma_3^+} e^{iz^2}dz
\end{equation}
 $\gamma^+_1 (x)=x(0\leq x\leq R)$ учир 
$$\int\limits_{\gamma_1^+} e^{iz^2}dz=\int\limits_{0}^R(\cos x^2+i\sin x^2)dx$$
$\gamma^+_2 (t)=Re ^{it}(0\leq t\leq \dfrac{\pi}{4})$ учир 
$$\int\limits_{\gamma_2^+} e^{iz^2}dz=\int\limits_{0}^{\frac{\pi}{4}}e^{-R^2\sin 2t+iR^2\cos 2t}\cdot Re^{it}\cdot idt$$
ба 
$$\left|\int\limits_{\gamma_2^+}e^{iz^2}dz\right|  \leq R\int\limits_{0}^\frac{\pi}{4} r^{-R^2\sin 2t}dt\leq R\int\limits_0^{\frac{\pi}{4}}e^{-R^2\frac{2}{\pi}2t}dt=\dfrac{\pi}{4}\cdot \frac{1-e^{-R^2}}{R}$$
$\gamma_3^+(t)=t\left(\dfrac{1}{\sqrt{2}}+i\dfrac{1}{\sqrt{2}}\right) \quad (R\geq t\geq 0)$ учир 
\begin{align*}
    \int\limits_{\gamma_3^+}e^{iz^2}dz&=\int\limits_R^0 e^{-t^2}\cdot \left(\frac{1}{\sqrt{2}}+i\frac{1}{\sqrt{2}}\right) dt \\
    &=- \left(\frac{1}{\sqrt{2}}+i\frac{1}{\sqrt{2}}\right)\int\limits_0^Re^{-t^2}dt.
\end{align*}
 Одоо (\ref{7*}) тэнцэтгэлд $R\rightarrow +\infty$ шилжвэл
 $$0=(I+iJ)+0+\left(-\frac{1}{\sqrt{2}}-i\frac{1}{\sqrt{2}}\right) \int\limits_0^{+\infty} e^{-t^2} dt.$$
 Иймд
 $$J=I=\frac{1}{\sqrt{2}}\int\limits_0^{+\infty} e^{-t^2} dt=\frac{\pi}{2\sqrt{2}}.$$

\end{sol}
\begin{center}
    \textbf{Дасгал}
\end{center}
$f(z)=\dfrac{e^{iz}}{z}$ ба зурагд үзүүлсэн $\gamma^+$ контурыг ашиглан $\bint\limits_0^{+\infty} \dfrac{\sin x}{x}dx$	 тооцоол.
\begin{center}
    \includegraphics[width=6cm]{Zurag/14.jpg}
     \captionof{figure}{}
\end{center}
\section{Эх функц. Ньютон-Лейбницийн томьёо} \index{Ньютон-Лейбницийн томьёо}
\begin{thm}
$D$ гүдгэр муж ба $z_0\in D$ байг. $f\in H(D)$ үед $F(z)=\bint\limits_{z_0}^zf(z)dz$ томьёогоор функц тодорхойлогдох ба 
$F'(z)=f(z)$ байна. (Өөрөөр хэлбэл голоморф функц гүдгэр муж дээр эх функцтэй байна).
\end{thm}
\begin{prf}
Гүдгэр муж нэг холбоост байна. Иймд $z_0$ цэгээс гарч $z$ цэг дээр төгсгөлтэй бүх замууд гомотоп байна. Ийнхүү $\bint\limits_{z_0}^z f(z)dz$ ганц утгатай, функц тодорхойлж өгнө. Энэ замуудаас $[z_0,z]$ хэрчим дагуу шугаман зам авахад хангалттай (Энд гүдгэр ашиглагдана).
\begin{center}
    \includegraphics[width=4cm]{Zurag/15.jpg}
     \captionof{figure}{}
\end{center}
\begin{align*}
    \frac{f(z+\Delta z)-F(z)}{\Delta z}&=\frac{\int\limits_z^{z+\Delta z}f(z)dz}{\Delta z} \\
    &=\frac{\int\limits_0^1 f(z+\Delta z\cdot t)\Delta z dt}{\Delta z}=\int\limits_0^1 f(z+\Delta zt) dt;
\end{align*}
\begin{align*}
    \lim\limits_{\Delta z\rightarrow 0}\frac{F(z+\Delta z)-F(z)}{\Delta z}&=\int\limits_0^1 \left(\lim\limits_{\Delta z\rightarrow 0}f(z+\Delta z\cdot t)\right) dt \\
    &=\int\limits_0^1 f(z)dt= f(z)
\end{align*}
болж батлагдав. 
\hfill  \end{prf}
\begin{cor}
$D$ нэг холбоост муж $f\in H(D)$ бол $f$ нь $D$ муж дээр эх функцтэй байна.
\end{cor}
\begin{hnt}
Голоморф биш функцүүд муж дээр эх функцтэй байж чаддаггүй. $(w=|z| \quad w= \re z;\quad w=x+iy^3$ гэх мэт)
\end{hnt}
\begin{hnt}
$f\in H(D)$ боловч $D$ нэг холбоост бус үед эх функц олддог $\left( D=\mathbb{C}\setminus \{0\}, \quad f(z)=\dfrac{\sin z}{z}\right)$ ба эс оршдог \\ $\left(D=\mathbb{C}\setminus \{0\}, \quad f(z)=\dfrac{\cos z}{z}\right)$ жишээнүүд байдгийг анхаарах хэрэгтэй.
\end{hnt}
\begin{thm}
$\gamma:[\alpha, \beta]\rightarrow \mathbb{C}$ гөлгөр зам ба $f$ функц $\gamma$ замын дагуу $\Phi(z)$ эх функцтэй байг. Тэгвэл
\begin{eqnarray}
  \int\limits_\gamma fdz=\Phi (\gamma (\beta))-\Phi (\gamma (\alpha))
\end{eqnarray} байна.
\end{thm}
\begin{prf}
$$\int\limits_\gamma f(z)dz=\int\limits_\alpha^\beta f(\gamma (t))\gamma'(t)dt =\Phi (\gamma(t))\Big|_{\alpha}^\beta$$
\hfill  \end{prf} болно.
\section{Кошийн интеграл томьёо} \index{Кошийн интеграл томьёо}
\begin{thm}
\textbf{(Кошийн интеграл томьёо).} \\ $f\in H(D)$ ба $z_0\in D$ байг. 
$$V_{z_0} (r)=\{z:|z-z_0|<r\} \quad \text{ба} \quad \gamma^+\undertilde{D} \partial V_{z_0}^+ (r)$$ байг. Өөрөөр хэлбэл   $\gamma^+ \text{ ба } \partial V_{z_0}^+ (r)$ гомотопи замууд байг. Тэгвэл $$ \bint\limits_{\gamma^+}\dfrac{f(z)}{z-z_0}dz=2\pi i f(z_0).$$
\end{thm}
\begin{prf}
$$\frac{f(z)}{z-z_0}=\frac{f(z)-f(z_0)}{z-z_0}+\frac{f(z_0)}{z-z_0}$$ учир
$$\int\limits_{\gamma^+}\frac{f(z)}{z-z_0}dz=\int\limits_{\gamma^+}\frac{f(z)-f(z_0)}{z-z_0}dz+\int\limits_{\gamma+}\frac{f(z_0)}{z-z_0}dz$$ болно.

$f(z)$ функц $z_0$ цэг дээр голоморф учир $\varphi(z)=\dfrac{f(z)-f(z_0)}{z-z_0}$ функц нь $z_0$ цэг дээр засагдах тасралтттай. Иймд Теорем \ref{removeablediscontinousfunctionint} ёсоор $$\int\limits_{\gamma^+}\frac{f(z)-f(z_0)}{z-z_0}dz=0$$ болно.

Хэрэв Жишээ \ref{Example z^n int}-ийг $n=-1$ үед хэрэглэвэл 
$$\int\limits_{\gamma+}\frac{f(z_0)}{z-z_0}dz=\int\limits_{\partial V^+_{z_0}(r)}\frac{f(z_0)}{z-z_0}dz=f(z_0)\cdot\int\limits_{\partial V^+_{z_0}(r)}\frac{1}{z-z_0}dz=2\pi i f(z_0)$$ болно. Иймд 
$$\int\limits_{\gamma^+}\frac{f(z)}{z-z_0}dz =2\pi i f(z_0)$$ болж батлагдав.
\hfill  \end{prf}
\begin{cor}
\textbf{(Дунджийн теорем).} $f\in H(D)$ ба $z_0\in D$ байг. Тэгвэл $f(z_0)$ утга нь $z_0$ төвтэй хангалттай жижиг тойргийн цэгүүд дээрх функцийн утгуудын дундаж байна. Өөрөөр хэлбэл: $$f(z_0)=\frac{1}{2\pi}\int\limits_0^{2\pi}f(z_0+\varepsilon \cdot  e^{it} dt)$$ (Энд $\overline{V_{z_0}}\subset D$ байх ёстой).
\end{cor}
\begin{prf}
$\partial V^+_{z_0}(\varepsilon)$ хувьд $f(z_0)=\dfrac{1}{2\pi i}\int\limits_{\partial V^+_{z_0}(\varepsilon)}\dfrac{f(z)}{z-z_0}dz$ гэж Кошийн интеграл томьёо бичигдэнэ. Одоо уг тойргийг \\ $z=z_0+\varepsilon e^{it}\quad (0\leq t\leq 2\pi)$ гэж параметрчилбэл
\begin{align*}
    f(z_0)&=\frac{1}{2\pi i}\int\limits_{\partial V^+_z (\varepsilon)}\frac{f(z)}{z-z_0}dz=\frac{1}{2\pi i}\int\limits_0^{2\pi}\frac{f(z_0+\varepsilon e^{it})}{\varepsilon e^{it}}\cdot (\varepsilon e^{it})'_{t}dt\\
    &=\dfrac{1}{2\pi}\int\limits_{0}^{2\pi} f(z_0+\varepsilon e^{it})dt
\end{align*}
болж батлагдав. 
\hfill  \end{prf}
\begin{exm}
$\gamma$ нь $3i$ төвтэй $R\neq 6$ радиустай тойрог байг. $I(R)=\bint\limits_{\gamma^+}\dfrac{dz}{z^2+9}$ тооцоол.
\end{exm}
\begin{sol}
$\gamma_1=\partial V^+_{3i} (1)$ ба $\gamma_2=\partial V^+_{-3i} (1)$ байг. 
    \begin{enumerate}
    \item [a)] $0<R<6$ үед $\gamma^+\sim \gamma_1^+$ учир 
    \begin{align*}
        \int\limits_\gamma \frac{dz}{z^2+9}&=\int\limits_{\gamma_1^+}\frac{\left(\frac{1}{z+3i}\right)dz}{z-3i}\\
        &=2\pi i\cdot \frac{1}{3i+3i}=\frac{\pi}{3}
    \end{align*}
 гарна.
    \end{enumerate}
    \begin{center}
        \includegraphics[width=5cm]{Zurag/16.PNG}
 \captionof{figure}{}
        \label{ref hiih ner}
    \end{center}
\begin{enumerate}
    \item [б)] $R>6$ үед $\gamma^+\sim \gamma_1^+\wedge \gamma_2^+$ тул 
    $$\int\limits_{\gamma^+}\frac{dz}{z^2+9}=\int\limits_{\gamma_1^+}\frac{\left(\frac{1}{z+3i}\right)dz}{z-3i} +\int\limits_{\gamma_2^+}\frac{\left(\frac{1}{z-3i}\right)dz}{z+3i}=0$$ байна. 
\end{enumerate}
\end{sol}
\begin{thm}
$f$ нь $D$ мужид $n$ удаа дифференциалчлагддаг байг. Тэгвэл 
$$\int\limits_{\gamma^+}\frac{f(z)}{(z-z_0)^{n+1}}dz=2\pi i\frac{f^{(n)}(z_0)}{n!}$$ байна.
\end{thm}
\begin{prf}
$\bint\limits_{\gamma^+}\frac{f(z)}{z-z_0}dz=2\pi i f(z_0)$ тэнцэтгэлийг $z_0$ параметрээр $n$ удаа дифференциалчилвал батлах тэнцэтгэл гарч ирнэ. Харин хэдийд тодорхой интегралыг ингэж дифференциалж болохыг санаж уг интеграл дээр нөхцөлийг шалгаарай.
\hfill  \end{prf}
\begin{thm}
$D\subset \mathbb{C}$ компакт муж ба $\partial D$ нь хэсэгчилсэн гөлгөр муруйнууд байг. $f$ нь $\overline{D}$ дээр $\mathbb{R}$ дифференциалчиллагддаг байг. $z_0\in D$ бол
\begin{eqnarray} \label{4.5}
  f(z_0)=\frac{1}{2\pi i}\int\limits_{2\pi i}\int\limits_{\partial D^+}\frac{f(z)}{z-z_0}dz-\frac{1}{\pi}\int\int\limits_D \frac{\partial f}{\partial \overline{z}} \cdot \frac{dxdy}{z-z_0}
\end{eqnarray}
\end{thm}
\begin{prf}
Гриний томьёоноос \index{Коши-Гриний томьёо}
$$\int\limits_{\partial D^+}\frac{f(z)}{z-z_0}dz-\int\limits_{C^+_\varepsilon (z_0)}\frac{f(z)}{z-z_0}dz=2i\int\int\limits \frac{\partial f}{\partial \overline{ z}}\frac{dx dy}{z-z_0}$$
(энд $\tilde{D}=D\setminus \overline{V}_{\varepsilon} (z_0)).$
\begin{align*}
    \int\limits_{C^+_\varepsilon (z_0)}\frac{f(z)}{z-z_0}dz&=\int\limits_{C^+_\varepsilon (z_0)}\frac{f(z)}{z-z_0}dz+\\
    &+\int\limits_{C^+_\varepsilon (z_0)}\frac{f(z)-f(z_0)}{z-z_0}dz \xrightarrow[\varepsilon \rightarrow 0 ]{} 2\pi i f(z_0)+0
\end{align*} болж (\ref{4.5}) томьёог хялбар батална. 
\hfill  \end{prf}
\begin{center}
    \textbf{Дасгал}
\end{center}
Энэ бодлогуудад бүх замыг шулуусгагдах гэж үзнэ.
\begin{enumerate}
    \item $\bint\limits_\gamma \dfrac{dz}{1+z^2}=\dfrac{\pi}{4}+k\pi $ байхыг батал. 
    
    Энд $\gamma$ нь $z=0$ эхтэй $z=1$ төгсгөлтэй $z=\pm i$ цэгүүдийг дайраагүй зам ба $k$ бүхэл тоо.
    \item  $C$ битүү муруйн бүх боломжит байрлалд $\bint\limits_C \dfrac{dz}{z(z^2-1)}$ ямар утгуудад авч болох вэ?
\item $\bint\limits_C \dfrac{dz}{w_n(z)}; w_n(z)=(z-z_1)(z-z_2)\ldots (z-z_n)$ \quad $(z_i\neq z_j)$ ба $C$ нь $z_i$ цэгүүдийг дайраагүй битүү контур. Интеграл ямар утгууд авах боломжтой вэ? 
\item $C$ энгийн битүү муруй ба $z=0$, $z=1$ цэгүүдийг дайраагүй бол $\dfrac{1}{2\pi i}\bint\limits_C \dfrac{ze^z dz}{(z-a)^3}$ бод.

\textbf{Заавар.} Интеграл дор уламжлал авахыг ашигла.
\item $\bint\limits_{|z-a|=a}\dfrac{zdz}{z^4-1}$ \quad $(a>1)$ бод. 
\item $\dfrac{1}{2\pi i}\bint\limits_{C^+}\dfrac{e^z dz}{z^2+a^2}$ бод. $(C=\{z:|z|=2a\}).$
\end{enumerate}
%\chapter{Цуваа} \index{Цуваа}
\section{Тоон цуваа}
Гишүүд нь комплекс $z_k=x_k+iy_k$ тоо байх
\begin{eqnarray} \label{5.1}
   z_1+z_2+\cdots +z_k+\cdots =\sum\limits_{k=1}^\infty z_k
\end{eqnarray}
цуваа авч үзье.

Бодит гишүүдтэй цувааных шиг
\begin{eqnarray}
   S_n=\sum\limits_{k=1}^n z_k
\end{eqnarray}
тоог цувааны хэсгийн нийлбэр гэнэ. Хэрэв хэсгийн нийлбэрийн хязгаар $S=\lim\limits_{n\rightarrow \infty} S_n$ оршин байвал (\ref{5.1}) цувааг нийлэх гэх ба $S$ тоог уул цувааны нийлбэр гэнэ. Өөрөөр хэлбэл $S=\sum\limits_{k=1}^\infty z_k$ гэнэ. Хэсгийн нийлбэрт төгсгөлтэй хязгаар оршин байхгүй бол (\ref{5.1}) цувааг сарних гэнэ.

Нийлэх цувааны ерөнхий $z_n$ гишүүн тэг рүү тэмүүлнэ. Учир нь (\ref{5.1}) нийлэх цуваа гэвэл

$$\lim\limits_{n\rightarrow \infty}z_n=\lim\limits_{n\rightarrow \infty} (S_n-S_{n-1})=S-S=0$$
Үүнийг цувааны нийлэлтийн зайлшгүй нөхцөл гэнэ. (\ref{5.1}) цувааны нийлэхийн зайлшгүй бөгөөд хүрэлцээтэй нөхцөл нь $\sum\limits_{k=1}^\infty x_k$ ба $\sum\limits_{k=1}^\infty y_k$-цуваануудын нийлэх явдал мөн. Иймд $\sum\limits_{k=1}^\infty z_k=S$, $\sum\limits_{k=1}^\infty x_k=\alpha$, $\sum\limits_{k=1}^\infty y_k=\beta$ гэвэл $S=\alpha +i\beta$ 
байна.
\begin{thm}
\textbf{(Нийлэлтийн Кошийн шинжүүр) } \index{Кошийн шинжүүр} (\ref{5.1}) цуваа нийлэх зайлшгүй бөгөөд хүрэлцээтэй нөхцөл нь аливаа $\varepsilon >0$ тооны хувьд натурал $N=N(\varepsilon)$ олдож $n>N$ ба ямар ч натурал $p$ тооны хувьд
\begin{eqnarray}
   |S_{n+p}-S_n|=|z_{n+1}+z_{n+2}+\cdots +z_{n+p}|=\left|\sum_{k=n+1}^{n+p}z_k\right| <\varepsilon
\end{eqnarray}
биелэгдэх явдал мөн. (\ref{5.1}) цувааны гишүүдийн модулиудаас зохиосон
\begin{equation} \label{5.4}
    \sum\limits_{k=1}^\infty |z_k|
\end{equation}
цуваа нийлж байвал (\ref{5.1}) цуваа нийлнэ.

Энэ нь Кошийн шинжүүрээр нотлогдоно:
$$\left|\sum_{k=n+1}^{n+p}z_k\right|\leq \sum_{k=n+1}^{n+p}|z_k|.$$
Энэ үед (\ref{5.1}) цувааг абсолют нийлэлттэй цуваа гэнэ. Харин (\ref{5.4}) цуваа сарниж байхад (\ref{5.1}) цуваа нийлэх тохиолдол байдаг. Энэ үед (\ref{5.1}) цувааг нөхцөлт нийлэлттэй цуваа гэнэ.

Цувааны абсолют нийлэлт тогтоохын тулд эерэг гишүүдтэй бодит цувааны нийлэлтийн хүрэлцээт шинжүүдийг ашиглана. Дэлгэр хэрэглэгддэг зарим шинжүүдийг баталгаагүйгээр тэмдэглэе:
\end{thm} 
\textbf{Жишилтийн шинж.} $\forall k>N$-ийн хувьд $|z_k|\leq |w_k|$ байхад $\sum\limits_{k=1}^\infty |w_k|$ цуваа нийлэх бол $\sum\limits_{k=1}^\infty |z_k|$ цуваа нийлнэ.\\
\textbf{Даламберийн шинж.} $\forall k>N$-ийн хувьд $\dfrac{|z_{k+1}|}{|z_k|}\leq q<1$ байвал (\ref{5.4}) цуваа нийлнэ. Харин $\forall k>N$-ийн хувьд $\dfrac{|z_{k+1}|}{|z_k|}\geq 1$  байвал (\ref{5.4}) цуваа сарнина. \\
\textbf{Кошийн шинж.} $\forall k>N$-ийн хувьд $\sqrt[k]{|z_k|}\leq q<1$ байвал (\ref{5.4}) цуваа нийлнэ. Харин $\forall k>N$-ийн хувьд $\sqrt[k]{|z_k|}\geq 1$ байвал (\ref{5.4}) цуваа сарнина.
\section{Функцэн цуваа} \index{Функцэн цуваа}
$f_k(z)\in H(D)$ функцүүдээс зохиосон
\begin{eqnarray} \label{5.5}
   \sum\limits_{k=1}^\infty f_k(z)
\end{eqnarray}
дүрсийн цуваа авч үзье. Энэ цувааны нийлэх $z\in D$ цэгүүдийн олонлогийг түүний нийлэлтийн муж гэнэ. 

Нийлэлтийн энэ муж дээр (\ref{5.5}) цувааны нийлбэр
$$S(z)=\sum\limits_{k=1}^\infty f_k(z)$$
ба түүний үлдэгдэл
$$r_n(z)=S(z)-S_n(z)=\sum\limits_{k=n+1}^\infty f_k(z)$$
гэж тодорхойлогдсон байг. (\ref{5.5}) цуваа ямар нэг $G$ муж дээр $S(z)$ рүү нийлдэг байг. Функцэн цувааны хувьд жигд нийлэлтийн асуудлыг чухалчилдаг.
\begin{dfn}
 Функцэн (\ref{5.5}) цуваа $\forall \varepsilon>0$-ийн хувьд зөвхөн $\varepsilon$ тооноос хамааран $N=N(\varepsilon)$ дугаар олдож $N<n$ байх бүх $n$-ийн үед $G$ муж дээр
 $$|S(z)-S_n(z)|=|r_n (z)|<\varepsilon$$
нөхцөл $(S_n(z)=  \sum\limits_{k=1}^n f_k(z)$- хэсгийн нийлбэр) биелэгдэж байвал  функцэн (\ref{5.5}) цуваа $G$ муж дээр жигд нийлж байна гэдэг.
\end{dfn}
Нийлэх функцэн цуваа бүр жигд нийлэх нь албагүй.

Жишээ нь хэсгийн нийлбэр нь
$$S_n(z)=\sum\limits_{k=0}^n z^n=\frac{1-z^{n+1}}{1-z}$$
байх
$$\sum\limits_{k=0}^n z^k$$
геометр прогресс нэгж $|z|<1$ дугуй дээр $S(z)=\dfrac{1}{1-z}$ рүү абсолют нийлнэ. Учир нь
$$\lim\limits_{n\rightarrow \infty} S_n(z)=\lim\limits_{n\rightarrow \infty} \frac{1-z^{n+1}}{1-z}=\frac{1}{1-z}\qquad |z|<1.$$
Гэтэл энэ цуваа $|z|<1$ дээр жигд биш нийлнэ. Учир нь үлдэгдэл гишүүн нь $\dfrac{1}{2}<|z|<1$ үед доороосоо
$$|r_n (z)|=\dfrac{|z|^{n+1}}{|1-z|}\geq \frac{|z|^{n+1}}{1-|z|}\geq \frac{1}{2^{n+1}(1-|z|)}$$
гэж үнэлэгдэнэ.

Иймд дурын бэхлэгдсэн хүрэлцээт их $n$-авахад $|z|$-ийг нэгжид хүрэлцээтэй ойр сонгох замаар үлдэгдлийн модулийг урьдчилан өгөгдсөн ямар ч тооноос ихээр авах боломж олдоно.

\begin{thm}
\textbf{(Жигд нийлэлтийн Вейерштрассын шинжүүр). } \index{Жигд нийлэлтийн Вейерштрассын шинжүүр}(\ref{5.5}) цувааны $k$ тооны гишүүдийн хувьд $G$ муж дээр $|f_k(z)|\leq c_k(k=1,2,3,\ldots)$ тэнцэл биш биелэгдэхэд $\sum\limits_{k=1}^\infty c_k$ нь  нийлэх тоон цуваа (мажорант цуваа) байвал уул функцэн цуваа $G$ муж дээр абсолют бөгөөд жигд нийлнэ.
\end{thm}
\begin{prf}
Абсолют нийлэх нь жишилтийн шинжээр шууд тогтоогдоно. Цааш нь (\ref{5.5}) цувааны нийлбэр ба хэсгийн нийлбэрийг харгалзан $S(z)$, $S_n(z)$ гэж тэмдэглэвэл 
$$|S(z)-S_n(z)|\leq \sum\limits_{k=n+1}^\infty |f_{k}(z)|\leq \sum\limits_{k=n+1}^\infty c_k$$
бөгөөд аливаа $\varepsilon >0$ авахад $N=N(\varepsilon)$ дугаар олдож $n\geq N$ байх $n$-ийн хувьд $\sum\limits_{k=n+1}^\infty c_k<\varepsilon$ болж теорем батлагдав. 
\end{prf}
\begin{thm} \label{Теорем 2}
$f_n(z)=H(G)$ байг. (\ref{5.5}) цуваа $G$ дээр жигд нийлж байвал $S(z)$ нийлбэр нь $G$ дээр тасралтгүй байна.
\end{thm}
\begin{prf}
$\varepsilon >0$ ба $z_0\in G$ өгөгдсөн байг. Цувааны хэсгийн нийлбэрийн хувьд $|S(z)-S_n(z)|<\dfrac{\varepsilon}{3}$ тэнцэл биш биелэгдэж байхуйц $n$ тоо сонгон авъя. $S_n(z)$ нийлбэр $G$ дээр тасралтгүй тул $|z-z_0|<\delta$ байвал $|S_n(z)-S_n(z_0)|<\dfrac{\varepsilon}{3}$ байхаар $\delta > 0$ олдоно.

Одоо $S(z)-S(z_0)$ ялгаврын хувьд
\begin{align*}
    |S(z)-S(z_0)|&=|S(z)-S_n(z)+S_n(z)- \\
    &-S_n(z_0)+S_n(z_0)-S(z_0)|<\frac{\varepsilon}{3}+\frac{\varepsilon}{3}+\frac{\varepsilon}{3}=\varepsilon
\end{align*}
болж теорем батлагдав.
\end{prf}
Энэ теорем $f_n(z)$ нь $G$ дээр зөвхөн тасралтгүй байх тохиолдолд ч үнэн байна.

Функцүүдийн төгсгөлтэй нийлбэрийг гишүүнчлэн интегралчилж болдгийг бид интеграл тооллоос мэднэ. Харин цувааг гишүүнчлэн интегралчлах нь өвөрмөц нөхцөлөөр шийдэгдэхийг дараах теоремоос танилцъя.
\begin{thm}
$\Gamma \subset G$ шугам дээр тасралтгүй функцүүдийн жигд нийлэх цувааг гишүүнчлэн интегралчлана: 
\begin{eqnarray}
   \int\limits_{\Gamma} S(z)dz=\sum\limits _{k=1}^\infty \int\limits_{\Gamma} f_k(z)dz
\end{eqnarray}
\end{thm}
\begin{prf}
$S(z)=\sum\limits_{k=1}^\infty f_k(z)$ байг. $\displaystyle\int\limits_{\Gamma} S(z)dz$ оршин байна. Учир нь $S(z)$ функц $\Gamma$ дээр тасралтгүй (өмнөх Теорем 5.3) интегралчлагдсан цувааны хэсгийн нийлбэр $\sigma_n$ гэвэл
$$\sigma_n= \sum\limits_{k=1}^n\int\limits_{\Gamma} f_k(z)dz=\int\limits_\Gamma S_n(z)dz$$
бөгөөд одоо
\begin{eqnarray} \label{5.7}
   \lim\limits_{n\rightarrow \infty} \sigma_n =\int\limits_\Gamma S(z)dz
\end{eqnarray}
гэж батлах хэрэгтэй.

$\varepsilon > 0$ байг. Энэ тохиолдолд теоремын нөхцөл ёсоор $N$ тоо авахад $n>N$ үед $|S(z)-S_n(z)|<\varepsilon$ байна. Тийм $n$-ийн хувьд 
\begin{align*}
    \left|\int\limits_{\Gamma} S(z)dz-\sigma _n \right|&= \left|\int\limits_{\Gamma} [S(z)-S_n(z)]dz\right|\leq \\
    &\leq \int\limits_{\Gamma}|S(z)-S_n(z)|\cdot |dz|<\varepsilon l.
\end{align*}
Үүнд $I$ нь $\Gamma$ шугамын урт. Эндээс (\ref{5.7}) шүтэлцэл батлагдав.
 \end{prf}
Комплекс фундаментал цуваануудын онцлог шинжийг тусгасан теорем авч үзье.
\begin{thm}
\textbf{(Вейерштрасс).}  Хэрэв $f_k(z)\in H(G)$ функцүүдээс зохиосон (\ref{5.5}) цуваа $G$ мужид нийлбэр $S(z)$ рүүгээ жигд нийлж байвал $S(z)\in H(G)$ байна. (\ref{5.5})-ийг гишүүнчлэн дифференциалчилбал 
\begin{eqnarray} \label{5.8}
   \sum\limits_{k=1}^\infty f_k^{(n)} (z)
\end{eqnarray}
цуваа үүсэх ба энэ нь харгалзан $S^{(n)}(z)$ уламжлал руугаа $G$ дээр жигд нийлнэ.
\end{thm}
\begin{prf}
Дураараа $z_0\in G$ цэг авъя. $\Delta >0$ тоог хүрэлцэхүйц бага өөрөөр хэлбэл $z_0$ цэгт төвтэй $2\delta$ радиустай $|z-z_0|=2\delta$ тойрог бүхлээрээ $G$ мужид оршин байхаар авъя. Цувааны нийлбэр $S(z)$ функц $z_0$ цэгийн $U=\{z:|z-z_0|<\delta\}$ орчинд $S(z)\in H(U)$ гэж баталъя (\ref{z.17}-г үз).
\begin{center}
    \includegraphics[width=6cm]{Zurag/17.PNG}
    \captionof{figure}{}
        \label{z.17}
\end{center}

$\zeta \in \Gamma=|z-z_0|=\delta$ байг. (\ref{5.5}) цувааны $z$-ийг $\zeta$-аар сольж $\dfrac{1}{2\pi i}\cdot \dfrac{1}{(\zeta - z)}$-аар үржүүлж авсан
\begin{eqnarray} \label{5.9}
   \dfrac{1}{2\pi i}\cdot \dfrac{S(\zeta)}{(\zeta -z)}=\dfrac{1}{2\pi i} \sum\limits_{k=1} ^\infty \dfrac{f_k(\zeta)}{\zeta -z}
\end{eqnarray}
цуваа сонирхъё.

(\ref{5.5}) цуваа $|z-z_0|=\delta$ тойрог дээр жигд нийлэх тул түүнийг гишүүнчлэн интегралчилж болно. Кошийн интеграл теорем ёсоор (\ref{5.9}) цувааны гишүүн бүрээс авсан интеграл нь $f_k(z)$-тэй тэнцэх тул
\begin{eqnarray} \label{5.10}
   \dfrac{1}{2\pi i} \int\limits_{\Gamma} \dfrac{S(\zeta)}{(\zeta -z)}d\zeta= \sum\limits_{k=1} ^\infty f_k(z)=\rho (z)
\end{eqnarray}
болно. 

(\ref{5.5})	цувааны нийлбэр $z_0$ цэгийн $\delta$-орчинд (\ref{5.10}) байдлаар Кошийн интегралаар тавигдах болж эндээс голоморф байх нь тогтоогдов. $S(z)\in H(G)$ болов. Теоремын хоёрдугаар хэсэг болох (\ref{5.8}) цуваа нь $G$ дээр $S^{(n)}(z)$ функц рүү жигд нийлэхийн Коши маягийн интегралаас авах уламжлалын томьёо ёсоор батлагдана. 
 \end{prf}
\section{Зэрэгт цуваа}
\textbf{1. Тейлорын цуваа.} Комплекс хувьсагчийн функцийн онолд зэрэгт цуваа онцгой үүрэг гүйцэтгэнэ.
\begin{eqnarray} \label{5.11}
   \sum\limits_{n=0}^\infty c_n(z-z_0)^n
\end{eqnarray}
дүрсийн цувааг зэрэгт цуваа \index{Зэрэгт цуваа} гэнэ.

$c_0, c_1, \ldots ,c_n,\ldots$ тогтмол комплекс тоонууд, $z_0$ бэхэлсэн тоо, $z$-комплекс хувьсах хэмжигдэхүүн. (\ref{5.11}) цуваа нь функцэн (\ref{5.5}) цувааны тухайн тохиолдол дээр нь (\ref{5.11}) цуваа нийлж байх бүх $z$-цэгүүдийн олонлогийг уул цувааны нийлэлтийн муж гэнэ. (\ref{5.11}) дүрсийн аливаа цуваа $z=z_0$ цэг дээр нийлэх нь илэрхий.

Зэрэгт цувааны үндсэн ухагдахуунуудын нэг болох дараах теорем авч үзье.
\begin{thm}
(Тейлор) \index{Тейлорийн теорем}. $f(z)\in H(D)$ бол энэ функц дурын \\ $U=\{z:|z-z_0|<R\}\subset D$ дугуй дээр нийлэх зэрэгт цувааны нийлбэр 
$$f(z)=\sum\limits_{n=0}^\infty c_n(z-z_0)$$
дүрстэй тавигдана.
\end{thm}
\begin{prf} 
$z\in U$ нь $z_0$-оос ялгаатай дурын цэг, $r$ тоог \\ $|z-z_0|<r<R$ байхаар сонгон авъя. $\{\zeta:|\zeta -z_0|=r\}$ тойргийг $\gamma$-ээр тэмдэглэе (зураг 18). Кошийн интеграл томьёо ёсоор
\begin{center}
    \includegraphics[width=5cm]{Zurag/18.PNG}
    \captionof{figure}{}
        \label{}
\end{center}
\begin{eqnarray}
    f(z)=\frac{1}{2\pi i}\int\limits_{\gamma_r}\frac{f(\zeta)d\zeta}{\zeta -z}
\end{eqnarray}
бөгөөд энэ функцийг цуваанд задлахын тулд $\dfrac{1}{\zeta -z}$ үржигдэхүүнийг $z-z_0$-ийн зэргээр геометр прогресс байдалтай тавих хэрэгтэй:
\begin{eqnarray} \label{5.13}
    \frac{1}{\zeta -z}=\frac{1}{(\zeta -z_0)\left(1-\frac{z-z_0}{\zeta -z_0}\right)}=\sum\limits_{n=0}^\infty \frac{(z-z_0)^n}{(\zeta -z_0)^{n+1}}
\end{eqnarray}
$\zeta \in \gamma_r$ байх бүхий л $\zeta$-ийн хувьд
$$\frac{|z-z_0|}{|\zeta -z_0|}=\frac{|z-z_0|}{r}=q<1$$
тул (\ref{5.13}) прогресс $\gamma_r$ дээр абсолют бөгөөд жигд нийлнэ. Одоо (\ref{5.13})-ийн хоёр талыг $\dfrac{1}{2\pi i}\cdot f(\zeta)$-аар үржүүлж $\gamma_r$-ээр гишүүнчлэн интегралчилъя. Ингэхэд уул цувааны жигд нийлэлтэд нөлөөлөхгүй. Учир нь $f(\zeta)$ функц $\gamma_r$ дээр тасралтгүй, улмаар зааглагдсан функц. Иймд
$$f(z)=\frac{1}{2\pi i}\int \gamma_r \sum\limits_{n=0}^\infty \frac{f(\zeta) d\zeta}{(\zeta -z_0)^{n+1}}(z-z_0)^n=\sum\limits_{n=0}^\infty c_n (z-z_0)^n$$
энд 
\begin{eqnarray} \label{5.14}
    c_n=\dfrac{1}{2\pi i} \int\limits_{\gamma_r} \frac{f(\zeta) d\zeta}{(\zeta -z_0)^{n+1}}, \quad n=0,1, \ldots
\end{eqnarray}
 \end{prf}
\begin{dfn}
 Коэффициент нь (\ref{5.14}) томьёонуудаар илэрхийлэгдэх (\ref{5.11}) цувааг $z_0$ төвт дугуй дээрх $f(z)$ функцийн Тейлорын цуваа \index{Тейлорын цуваа} гэнэ.
\end{dfn}
\begin{cor}
\textbf{(Кошийн тэнцэтгэл биш). } \index{Кошийн тэнцэтгэл биш}
$f(z)$ функц битүү $\overline{U}=\{|z-z_0|\leq r\}$ дугуй дээр голоморф. Мөн $\gamma_r =\partial U$ тойрог дээр түүний модуль $M$-ээс үл хэтрэх байг. Тэгвэл $z_0$ төвт дугуй дээрх Тейлорын цувааны коэффициентууд 
\begin{eqnarray}
    |c_n|\leq \frac{M}{r^{n}}, \quad (n=0,1,\cdots)
\end{eqnarray}
тэнцэл бишийг хангана.
\end{cor}
\begin{prf}
$\forall \zeta \in \gamma$ дээр $|f(\zeta)| \leq M$ байхыг тооцон (\ref{5.14}) томьёоноос олбол
$$|c_n|\leq \frac{1}{2\pi}\cdot \frac{M}{r^{n+1}}\cdot 2\pi r =\frac{M}{r^n}$$
 \end{prf}
Кошийн тэнцэл биш ашиглаж дараах сонирхолтой теорем батлагдана.
\begin{thm}
\textbf{(Лиувилль).} \index{Лиувиллийн теорем} $|f(z)|\leq M$ ба $f(z)\in H(\mathbb{C})$, $(\forall z\in \mathbb{C})$ байвал $f(z)=const $ байна. 
\end{thm}
\begin{prf}
Тейлорын теорем ёсоор $\overline{ U}=\{|z|\leq R\}$, $R<\infty$ дугуй дээр $f(z)$ функц Тейлорын цуваанд задрах тул \\ $f(z)=\sum\limits_{n=0}^\infty c_n z^n$ бөгөөд үүний коэффициентууд нь $R$-ээс үл хамаарах, нөгөө талаар $f(z$) нь $\mathbb{C}$ дээр зааглагдсан учир Кошийн тэнцэл биш ёсоор аливаа $n=0,1,\ldots$-ийн хувьд $|c_n|\leq \dfrac{M}{R^n}$.
Сүүлчийн тэнцэл бишийн баруун гар тал $n=1,2,\ldots$ үед \\ $R\rightarrow \infty$ байхад тэг рүү тэмүүлнэ. Харин зүүн гар тал $R$-ээс үл хамаарах, иймд $n=1,2,\ldots$ үед $c_n=0$ болж улмаар $f(z)=c_0.$
 \end{prf}

Лиувиллийн энэ теоремыг өргөтгөн томьёолж болно:
\begin{thm}
$f(z)\in H(\overline{\mathbb{C}})$ бол $f=const$ байна.
\end{thm}
\begin{prf}
$f(z)\in H(\infty)$ тул $\lim\limits_{z\rightarrow \infty} f(z)$ оршин байх ба төгсгөлтэй. Энд $f(z)$ нь $\infty$-ийн зааглагдсан олонлог дээр тасралтгүй байх функц хавтгай үлдсэн $\{|z|>R\}$ хэсэг дээр зааглагдсан нь илэрхий тул өмнөх теорем ёсоор $f(z)=const$
 \end{prf}
\begin{thm}
Абелийн теорем.\index{Абелийн теорем} Лемм.Хэрэв зэрэгт (\ref{5.11}) 
$$\sum\limits_{n=0}^\infty c_n(z-a)^n$$
цувааны гишүүд ямар нэг $z_0$ цэг дээр зааглагдсан, өөрөөр хэлбэл
\begin{equation} \label{5.16}
    |c_n (z_0-a)|\leq M, \quad n=0,1, \ldots
\end{equation}
байвал энэ цуваа $U=\{z:|z-a|<|z_0-a|\}$ дугуй дээр нийлэх ба компакт дэд олонлог $K\Subset U$ бүр дээр абсолют бөгөөд жигд нийлнэ.
\begin{center}
    \includegraphics[width=4cm]{Zurag/19.jpg}
    \captionof{figure}{}
        \label{}
\end{center}
\end{thm}
\begin{prf}
$a\neq z_0$ байхаар $z_0$ авбал $|z_0-a|=\rho >0$. Ийм байхгүй бол $U$ хоосон олонлог.

$K\Subset U$ байг. Тэгвэл $\forall z\in K$ авахад $\dfrac{|z-a|}{\rho}\leq < 1$ (зураг 9) бөгөөд дурын $n=0,1, \ldots$ тооны хувьд
$$|c_n(z-a)^n|\leq |c_n|\cdot \rho^nq^n$$
байна. (\ref{5.16}) нөхцөлөөр $|c_n|\rho ^n\leq M$ болж $\forall z\in K$-ийн хувьд (\ref{5.11}) цуваа нийлэх геометр $M\sum\limits_{n=0}^\infty q^n$ прогрессоор мажорантлагдах тул тэр нь $K$ дээр абсолют бөгөөд жигд нийлнэ. Ингээд леммийн сүүлчийн өгүүлбэр батлагдав. Эхнийх нь эндээс мөрдөгдөн гарна. Учир нь $\forall z'\in U$ авахад тэр нь ямар нэг $|z'-a|<\rho '$ дугуйд оршино. $\rho '<\rho$ тул энэхүү шинэ дугуй нь $U$-д компактаар оршин байна. 
 \end{prf}
\begin{thm}
\textbf{(Абель). } Хэрэв зэрэгт (\ref{5.11}) цуваа $z_0\in \mathbb{C}$ цэг дээр нийлж байвал тэр нь $U=\{z:|z-a|< |z_0-a|\}$ дугуй дээр нийлэх ба $U$-д компактаар орших ямар нэг дэд олонлог дээр абсолют бөгөөд жигд нийлнэ.
\end{thm}
\begin{prf}
(\ref{5.11}) цуваа $z_0$ цэг дээр нийлж байвал харгалзах тоон цувааны ерөнхий гишүүн $c_n(z_0-a)^n$ тэг рүү тэмүүлнэ. Аливаа нийлэх дараалал зааглагдсан байх тул леммийн нөхцөл биелэгдэж теорем батлагдав. 
 \end{prf}
\begin{thm}
\textbf{Коши-Адамарын томьёо. } \index{Коши-Адамарын томьёо} Зэрэгт (\ref{5.11}) цуваа өгөгдсөн байг. $\overline{\lim\limits_{n\rightarrow \infty}} \sqrt{ |c_n|}= \wedge$ байвал $\wedge =0$ үед (\ref{5.11}) цуваа бүх хавтгай дээр абсолют нийлнэ.

$\wedge =\infty$ үед тэр нь зөвхөн $z=z_0$ цэг дээр нийлнэ. $z\neq z_0$ дээр сарнина.

$0<\wedge <\infty$ үед уул цуваа $|z-z_0|<\dfrac{1}{\wedge}$ дугуй дээр абсолют нийлэх ба энэ дугуйн гадна $\left(|z-z_0|>\dfrac{1}{\wedge}\right)$ сарнина.
\end{thm}
\begin{prf}
Зэрэгт (\ref{5.11}) цуваа $z_0$ дээр нийлэх тул энэ цувааны абсолют нийлэлтийг $z\neq z_0$ цэг дээр Кошийн шинж ашиглан
шинжилье.
\begin{eqnarray} \label{5.12}
    \sum\limits_{n=0}^\infty |c_n|\cdot |z-z_0|^n
\end{eqnarray} цуваа авъя. 
\begin{enumerate}
    \item [a)] $\wedge = 0$ гэвэл
    $$\overline{\lim} \sqrt[n]{|c_n||z-z_0|^n}=|z-z_0|\overline{\lim}\sqrt[n]{|c_n|}=0$$
    болж Кошийн шинжээр бүх хавтгай дээр (\ref{5.11}) цувааны абсолют нийлэлт тогтоогдлоо.
    \item [б)] $\wedge =\infty$ байг. $z\neq z_0$ байх аливаа $z$-ийн хувьд
    $$|z-z_0|\overline{\lim}\sqrt[n]{|c_n|}=\infty$$
    болж, өөрөөр хэлбэл (\ref{5.12}) цувааны ерөнхий гишүүн тэг рүү үл тэмүүлэх нь илэрхий тул (\ref{5.11}) цуваа $ z\neq z_0$ байх бүх цэг дээр сарнина.
    \item [в)] Эцэст нь $0<\wedge <\infty$ гэе. $z$ бүхний хувьд:
    $$\overline{\lim}\sqrt[n]{|c_n||z-z_0|^n}=|z-z_0|\cdot \wedge$$
Хэрэв $|z-z_0|<\dfrac{1}{\wedge}$ байвал $\overline{\lim}\sqrt[n]{|c_n||z-z_0|^n}<1$ болж (\ref{5.12}) цуваа Кошийн шинж ёсоор нийлнэ. \\ Иймд $|z-z_0|<\dfrac{1}{\wedge}$ дугуй дээр (\ref{5.11}) цуваа абсолют нийлнэ. Харин $|z-z_0|>\dfrac{1}{\wedge}$ үед $\overline{\lim\limits_{n\rightarrow\infty}}\sqrt[n]{|c_n||z-z_0|^n}>1$ болж (\ref{5.12}) цувааны ерөнхий гишүүн
тэг рүү үл тэмүүлэх учир улмаар (\ref{5.11}) цуваа $|z-z_0|>\dfrac{1}{\wedge}$ үед сарних нь тогтоогдов. 
\end{enumerate}
 \end{prf}
\begin{hnt}
\begin{eqnarray} \label{5.18}
    R=\frac{1}{\wedge}, \quad \wedge =\overline{\lim\limits_{n\rightarrow \infty}} \sqrt[n]{|c_n|}
\end{eqnarray}
шүтэлцлийг Коши-Адамарын томьёо гэнэ.
\end{hnt}
\begin{dfn}
 Зэрэгт (\ref{5.11}) цуваа нийлэх $z\in \mathbb{C}$ цэгүүдийн $E$ олонлог дахь задгай $\mathring{E}$ цөм (дотоод цэгүүдийн бөөгнөрөл)-ийг энэ цувааны нийлэлтийн муж гэнэ.
\end{dfn}
\begin{thm}
Зэрэгт (\ref{5.11}) цувааны нийлэлтийн муж нь \\ $|z-z_0|<R$ дугуй байна. $R$ тоо нь (\ref{5.18}) томьёогоор тодорхойлогдоно.
\end{thm}
\begin{prf}
Өмнөх тодорхойлолтоос харахад зэрэгт (\ref{5.11}) цувааны нийлэх цэгүүдийн $E$ олонлог $\{|z-z_0|=R\}$ тойргийн зарим цэгүүдээр сэлбэгдсэн $\{|z-z_0|<R\}$ дугуй байх учиртай. Иймд задгай $\mathring{E}$ цөм нь $|z-z_0|<R$ болой. 
 \end{prf}
Оршин байх нь саяны теоремоор нотлогдсон (задгай) дугуйг (\ref{5.11}) цувааны нийлэлтийн дугуй гэх ба $R$ тоог нийлэлтийн радиус гэнэ.
\begin{exm}
Дараах зэрэгт цуваануудын нийлэлтийн радиус ол.
\colhoy{
\begin{enumerate}
    \item [a)] $\displaystyle\sum\limits_{n=1}^\infty \dfrac{n!}{n^n}z^n$;
    \item [б)] $\displaystyle\sum\limits_{n=0}^\infty (3+(-1)^n)^nz^n;$
    \item [в)] $\displaystyle\sum\limits_{n=0}^\infty z^{2^n}$
\end{enumerate}
} 
\end{exm}
\begin{sol}
\begin{enumerate}
    \item  [a)] $$R=\lim\limits_{n\rightarrow \infty} \frac{n!}{n^n}\cdot \frac{(n+1)^{n+1}}{(n+1)!}=\lim\limits_{n\rightarrow \infty} \left(1+\frac{1}{n}\right)^n =e.$$
    \item [б)] $R=\overline{\lim\limits_{n\rightarrow \infty}}\sqrt[n]{\frac{1}{(3+(-1)^n)^n}}=\lim\limits_{n\rightarrow \infty} \sqrt[2k]{[3+(-1)^{2k}]^{2k}}=\dfrac{1}{4}.$
    \item [в)] Цувааг $\sum\limits_{k=0}^\infty a_kz^k$ дүрстэй бичиж болох ба үүнд 
    
    $a_k= \begin{cases}
     1, \quad \text{хэрэв} \quad k=2^n, \\
      0, \quad \text{хэрэв} \quad k\neq 2^n
    \end{cases}$ Иймд $R=\dfrac{1}{\overline{\lim\limits_{k\rightarrow \infty}}\sqrt[k]{|a_k|}}=1$
\end{enumerate}
\end{sol}
\begin{exm}

Дараах цуваануудыг нийлэлтийн дугуйн хил дээр шинжил.
\colgur{
\begin{enumerate}
    \item [a)] $\sum\limits_{n=1}^\infty \dfrac{(-1)^n}{n}z^n$;
    \item [б)] $\sum\limits_{n=0}^\infty \dfrac{n^2+1}{(z-i)^n}z^n;$
    \item [в)] $\sum\limits_{n=1}^\infty \dfrac{z^n}{n}.$
\end{enumerate}}
    
\end{exm}
\begin{sol}
\begin{enumerate}
    \item  [a)] $\wedge =1$, $R=1$. $|z|=1$ тойргийн $z=-1$-ээс ялгаатай цэг бүр дээр абсолют бишээр нийлнэ. 
    \item [б)] $\wedge =1$, $R=1$. $|z-i|=1$ тойрог дээр абсолют хэмжигдэхүүний цуваа $\sum\limits_{n=0}^\infty \dfrac{1}{n^2+1}$-нийлэх тул уул цуваа $|z-i|\leq 1$ гэсэн битүү дугуй дээр абсолют нийлнэ.
    \item [в)] $\wedge =1$, $R=1$. $|z|=1$  тойрог дээр цувааны нийлэлт шинжилье. Абсолют хэмжигдэхүүний  $\sum\limits_{n=1}^\infty \dfrac{1}{n}$ цуваа $|z|=1$ байх бүх цэг дээр сарнина. 
    
    $z\neq 1$, $|z|=1$ байхад $\sum\limits_{n=1}^\infty \dfrac{z^n}{n}$ цуваа нийлэхийг харуулъя. Дирихле-н шинж ашиглая. $b_n=\dfrac{1}{n}$, $a_n=z^n$, $|z|=1$ гэе. $b_n=\dfrac{1}{n}$ дараалал тэг рүү монотон тэмүүлнэ. Харин 
$$\left|\sum\limits_{k=1}^n z^k \right| = \left|\frac{z-z^{n+1}}{1-z}\right| <\frac{2}{|1-z|}, \quad z\neq 1$$ болж цувааны хэсгийн нийлбэрүүдийн дараалал зааглагдсан тул цуваа нийлнэ. 
\end{enumerate}
\end{sol}
\begin{thm}
\textbf{(Нийлбэрийн голоморф байх тухай). } Зэрэгт
\begin{eqnarray} \label{5.19}
    f(z)=\sum\limits_{n=0}^\infty c_n(z-a)^n 
\end{eqnarray} 
цувааны нийлбэр түүний нийлэлтийн дугуй дээр голоморф функц байна.
\end{thm}
\begin{prf}
$R>0$ гэе. Тэгэхгүй бол батлах зүйлгүй. (\ref{5.19}) цувааны хийсвэр уламжлал болох
\begin{eqnarray} \label{5.15}
    \sum\limits_{n=1}^\infty nc_n (z-a)^{n-1}=\varphi (z)
\end{eqnarray}
цувааны нийлэх буюу сарних нь $\sum\limits_{n=1}^\infty nc_n(z-a)^n$ цуваатай  адил. Учир нь
$$\overline{\lim\limits_{n\rightarrow \infty}} \sqrt[n-1]{n|c_n|} = \overline{\lim\limits_{n\rightarrow \infty}} (n|c_n|)^{\frac{1}{n}\cdot \frac{n}{n-1}} =\overline{\lim\limits_{n\rightarrow \infty}} \sqrt[n]{n|c_n|}.$$
Харин 
$$\overline{\lim\limits_{n\rightarrow \infty}} \sqrt[n]{n|c_n|} =\overline{\lim\limits_{n\rightarrow \infty}} \sqrt[n]{|c_n|}$$ байх тул (\ref{5.15}) цувааны нийлэлтийн радиус бас л $R$ байна. $U=\{|z-a|<R\}$ дугуйн компакт дэд олонлогууд дээр (\ref{5.15}) цуваа жигд нийлнэ. Иймд $\varphi (z)$ функц энэ дугуй дээр тасралтгүй. Мөн тийм шалтгаанаар (\ref{5.15}) цувааг $\Delta \Subset U$ байх дурын гурвалжны хилээр гишүүнчлэн интегралчилж болно:
$$\int\limits_{\partial \Delta} \varphi(z)dz=\sum\limits_{n=1}^\infty nc_n\int\limits_{\partial \Delta}(z-a)^{n-1}dz=0.$$
Баруун гар дахь бүх интеграл Кошийн теорем ёсоор тэгтэй тэнцүү, улмаар бас зүүн гар тал. Иймд, эх функц оршин байх тухай Теорем 5.13 ёсоор ашиглаж гаргах
$$\int\limits_{[a,z]} \varphi (\zeta)d\zeta=\sum\limits_{n=1}^\infty nc_n\int\limits_{[a,z]}(\zeta-a)^{n-1}d\zeta=\sum\limits_{n=1}^\infty c_n$$
функц (бид жигд нийлэлтийг дахин хэрэглэв) $z\in U$ цэг бүр дээр $\varphi(z)$-тэй тэнцэх уламжлалтай байна. Харин
$$f(z)=c_0+\int\limits_{[a,z]}\varphi (\zeta)d\zeta$$ 
$z\in U$ цэг бүр дээр уламжлалтай: $f'(z)=\varphi (z).$ 
 \end{prf}
\textbf{Голоморф функцийн чанарууд.} Зэрэгт цувааны нийлбэрийн голоморф байх тухай теоремын мөрдлөг байдлаар зарим теорем тэмдэглэе.
\begin{thm}
Аливаа $f(z)\in H(D)$ функцийн уламжлал $D$ дээр голоморф байна.
\end{thm}
\begin{prf}
$\forall z_0\in D$ дээр $U=\{|z-z_0|<R\}\subset D$ орчин байгуулъя. Тэгвэл Теорем 5.13 ёсоор $f'(z)=\varphi (z)$ уламжлал нь мөнхүү дугуй дээр нийлэх цуваа байдлаар тавигдана. Иймд $\varphi(z)$ дээр Теорем 5.13-ийг дахин хэрэглэвэл $\varphi(z)$ функц $U$ дээр комплекс анализын утгаар дифференциалчилагдах болов.
 \end{prf}
\begin{cor}
Тасралтгүй $f(z)$ функц $D$ дээр $F(z)$ хэмээх эх функцтэй байвал $f(z)\in H(D)$ байна. Теорем 5.14-ийг дахин хэрэглэж батлагдлаа.
\end{cor}
\begin{thm}
Аливаа $f(z)\in H(D)$ функц $D$ дээр бүхий л эрэмбийн уламжлалтай, тэдгээр нь $H(D)$ ангид харьяалагдана: $f^{(n)} (z)\in H(D)$.
\end{thm}
\begin{thm}
Хэрэв $f(z)$ функц $\{|z-z_0|<R\}$ дугуй дээр зэрэгт цувааны нийлбэр
\begin{eqnarray} \label{5.21}
    f(z)=\sum\limits_{n=0}^\infty c_n(z-z_0)^n
\end{eqnarray}
дүрсээр тавигдсан байвал энэ цувааны коэффициентууд
\begin{eqnarray} \label{5.22}
    c_n=\frac{f^{(n)}z_0}{n!}, \quad n=0,1,\ldots
\end{eqnarray}
томьёогоор нэгэн утгатай тодорхойлогдоно.
\end{thm}
\begin{prf}
(\ref{5.21}) тэнцэлд $z=z_0$ гэвэл $f(z_0)=c_0$ бөгөөд (\ref{5.21}) цувааг гишүүнчлэн дифференциалчилж $z=z_0$ гэвэл $f'(z_0)=c_1$ гэх мэтээр $f^{(n)}(z)=n!c_n+b_1(z-z_0)+\cdots $ болж $z=z_0$ гэвэл $f^{(n)}(z_0)=n!cn.$ 
 \end{prf}

Одоо (\ref{5.22}) томьёо ашиглан хялбар трансцендент $e^z$, $\sin z$, $\cos z$, $\sinh z$, $\cosh z$ функцүүдийн $z=0$ цэг дээрх зэрэгт цуваанд бичье:
\begin{eqnarray} \label{5.23}
    e^z=1+z+\frac{z^2}{2}+\cdots =\sum\limits_{n=0}^\infty \frac{z^n}{n!}, \quad R=\infty 
\end{eqnarray}
\begin{eqnarray} \label{5.24}
    \sin z=z-\frac{z^3}{3!}+\frac{z^5}{5!}-\cdots =\sum\limits_{n=0}^\infty (-1)^n\frac{z^{2n+1}}{(2n+1)!}, \quad R=\infty
\end{eqnarray}
\begin{eqnarray} \label{5.25}
    \cos z=1-\frac{z^2}{2!}+\frac{z^4}{4!}-\cdots =\sum\limits_{n=0}^\infty (-1)^n\frac{z^{2n}}{(2n)!}, \quad R=\infty
\end{eqnarray}
\begin{eqnarray} \label{5.26}
    \sinh z=z+\frac{z^3}{3!}+\frac{z^5}{5!}+\cdots =\sum\limits_{n=0}^\infty \frac{z^{2n+1}}{(2n+1)!}, \quad R=\infty
\end{eqnarray}
\begin{eqnarray} \label{5.27}
    \cosh z=1+\frac{z^2}{2!}+\frac{z^4}{4!}+\cdots =\sum\limits_{n=0}^\infty \frac{z^{2n}}{(2n)!}, \quad R=\infty
\end{eqnarray}
Энд (\ref{5.24})-ийг $i$-ээр үржүүлж (\ref{5.25}) дээр нэмэхэд Эйлерийн $e^{iz}=\cos z +i \sin z$ батлагдана. Хэрэв $f(z)$ функц $D$ муж дээр голоморф ба $G\Subset D$, $\forall z\in G$ байвал уламжлалын Кошийн томьёо 
\begin{eqnarray}
    f^{(n)}(z)=\frac{n!}{2\pi i}\int\limits_{\partial G}\frac{f(\zeta)d\zeta}{(\zeta -z)^{n+1}}, \quad n=0,1,2,\ldots
\end{eqnarray}
дүрстэй бичигдэнэ. Энэ томьёонууд Кошийн интеграл 
$$f(z)=\frac{1}{2\pi i}\int\limits_{\partial G} \frac{f(\zeta)d\zeta}{\zeta -z} $$
томьёоноос интеграл дор параметр $z$-ээр нь дифференциалчлах замаар гарна.
\begin{thm}
\textbf{(Морер).} \index{Морерийн теорем} Хэрэв $f(z)$ функц $D$ мужид тасралтгүй ба $$ \int\limits_{\partial \Delta}f(z)dz=0, \quad \forall \Delta \Subset D$$ байвал $f(z)\in H(D)$.
\end{thm}
\begin{prf}
$\forall a\in D$  цэгийн хувьд $U=\{|z-a|<r\}\subset D$ дугуй байгуулъя. $F(z)=\bint\limits_{[a,z]}f(\zeta) d\zeta$ функц $U$ дээр голоморф ба $z\in U$ цэг бүр дээр $F'(z)=f(z)$ тул өмнөх Теорем 5.14 ёсоор $f(z)\in H(D)$
 \end{prf}

Өгөгдсөн цэг дээр функцийн голоморф байхын өөр өөр тодорхойлолт хоорондоо чацуу чанартай.
\begin{thm}
Дараах гурван өгүүлбэр эквивалент:\\
$(R):f(z)$ функц $a$ цэгийн ямар нэг $U$ орчинд $\mathbb{C}$- дифференциалчлагдаж байвал $(C):f(z)$ функц $a$ цэгийн ямар нэг $U$ орчинд тасралтгүй ба $\bint\limits _{\partial \Delta}f(z)dz=0$, $\forall \Delta \Subset D$ байвал $(W):f(z)$ функц $a$ цэгийн ямар нэг $U$ орчинд нийлэх зэрэгт цуваанд задарч байвал энэ гурван өгүүлбэрийн аль нь ч биелэгдсэн $f(z)$ функц $a$ цэг дээр голоморф байна.\\
Эдгээр нотолгоо нь голоморф функцийн онолыг бүтээсэн гурван урсгалыг тусгана.

$(R)$ нөхцөлийг хангах функцийг Риманы утгаар голоморф \index{Риманы утгаар голоморф функц}, $(C)$ нөхцөлийг Кошийн утгаар голоморф \index{Кошийн утгаар голоморф функц}, $(W)$ нөхцөлийг Вейерштрассын утгаар голоморф \index{Вейерштрассын утгаар голоморф функц} байх функцүүд гэнэ.

$(R)\Rightarrow (C)$ нь Кошийн интеграл теоремоор, $(C)\Rightarrow (W)$ нь Тейлорын теоремоор. $(W)\Rightarrow (R)$ нь зэрэгт цувааны нийлбэрийн голоморф байх теоремоор тус тус батлагдана. 
\end{thm}
\begin{dfn}
 Дээр нь $f(z)$ функц тэг утга авах $\forall a\in\overline{\mathbb{C}}$ цэгийг, өөрөөр хэлбэл $f(z)=0$ тэгшитгэлийн язгуурыг уул \textbf{функцийн тэг} гэнэ.
\end{dfn}
\begin{thm}
Хэрэв $a$ цэг дээр голоморф бөгөөд тэр цэгийн ямар ч орчинд адилтган тэгтэй тэнцэхгүй $f(z)$ функцийн тэг нь $a$ байвал
\begin{eqnarray} \label{5.29}
     f(z)=(z-a)^n\varphi (z),
\end{eqnarray}
(энд $\varphi (z)$ функц нь $a$ цэг дээр голоморф, мөн түүнчлэн энэ цэгийн ямар нэг орчинд тэгээс ялгаатай утгатай байх функц) нөхцөлд тохирох натурал $n$ тоо оршин байна.
\end{thm}
\begin{prf}
$a$ цэгийн ямар нэг орчинд $f(z)$ функц зэрэгт цуваанд задарна. Сул гишүүн нь тэгтэй тэнцүү, учир нь $f(a)=0.$ Харин түүний бүх коэффициент тэгтэй тэнцэхгүй, хэрэв тэгвэл $f(z)=0$ байхад хүрнэ. Иймд тэгээс ялгаатай коэффициент олдох бөгөөд түүний хамгийн бага дугаарыг $n$-ээр тэмдэглэвэл задаргаа
\begin{eqnarray}
     f(z)=c_n(z-a)^n+c_{n+1}(z-a)^{n+1}+\cdots +c_n\neq 0
\end{eqnarray}
дүрстэй бичигдэнэ.


$a$ цэгийн ямар нэг орчинд нийлэх
$$c_n+c_{n+1}(z-a)+\cdots =\varphi (z)$$ цувааны нийлбэрийг $\varphi(z)$ гэж тэмдэглэвэл тэр нь энэ орчинд голоморф байна. $\varphi (a)=c_n\neq 0$ байх тул
$$f(z)=(z-a)^n\varphi (z)$$
нөхцөлд тохирох п дугаар оршин байна.
 \end{prf} 
\begin{thm}
\textbf{(Цор ганц байх тухай).} \index{Ганц ч чанарын теорем}\\
Хэрэв $f_1(z)$, $f_2(z)\in H(D)$ функцүүд $D$ мужид орших ядаж нэг хязгаарын $a$ цэгтэй $D$ олонлогт утгууд нь таарч байвал $D$ мужид газар сайгүй $f_1(z)=f_2(z)$ байна.
\end{thm}
\begin{prf}
$f(z)=f_1(z)-f_2(z)\in H(D)$ функц авъя. $D$ мужид $f(z)=0$ гэж батлах хэрэгтэй. Өөрөөр хэлбэл\\
$F=\{z\in D:f(z)=0\} \supset E$ олонлог $D$-тэй таарна. $f(z)$ функц тасралтгүй учраас хязгаарын $a$ цэг нь $f(z)$-ийн тэг байна. Теорем 5.14 ёсоор а-ийн ямар нэг орчинд $f(z)\equiv 0$. Тэгэхгүй бол энэ цэг $f(z)$-ийн тэгүүдийн олонлогийн хязгаарын цэг болж чадахгүй байна. Тийм учраас $F$ олонлогийн задгай $\mathring{F}$ цөм (өөрөөр хэлбэл түүний дотоод цэгүүдийн бөөгнөрөл) хоосон биш, тэр нь $a$ цэгийг агуулна. Байгуулалт
ёсоор  $\mathring{F}$  нь задгай, бас тэр нь битүү ($D$ мужийн харьцангуй топологит). Учир нь хэрэв $b\in D$ нь $\mathring{F}$ олонлогийн орчинд $b$ байх өөрөөр хэлбэл $b\in \mathring{F}$. $D$ нь мужийн тодорхойлолтоор холбоост байх тул I бүлэг 4. Теорем 2 ёсоор $\mathring{F}=D$. 
 \end{prf}
\begin{dfn}
$f(z)$ функцийн хувьд 
\begin{eqnarray}
     f(a)=\cdots =f^{(n-1)}(a)=0, \quad f^{(n)}(a)\neq 0(n  \geq 1)
\end{eqnarray}
байвал $a\in\mathbb{C}$ цэгийг түүний $n$-эрэмбийн тэг \index{Тэгийн эрэмбэ} гэнэ. Тейлорын цувааны коэффициент олох $c_n=\dfrac{f^{(n)}}{n!}$ томьёоноос харахад тэгийн эрэмбэ нь уул функцийн энэ цэг дээрх Тейлорын задаргааны коэффициентын тэгээс ялгаатай бага дугаартай таарна, өөрөөр хэлбэл теорем 5.14-ээр томьёолсон $n$ тоотой тэнцэнэ. Тэгтэй адилтган тэнцүү биш голоморф функц төгсгөлгүй эрэмбийн тэгтэй байж чадахгүйг цор ганц байх тухай теорем зааж байна.

Олон гишүүнтийн үлгэрээр хуваагдахын шинжийн тусламжтайгаар тэгийн эрэмбийг тодорхойлж болно. 
\end{dfn}
\begin{thm}
Голоморф $f(z)$ функцийн $a\in\overline{\mathbb{C}}$ тэгийн эрэмбэ нь $\dfrac{f(z)}{(z-a)^k}$ ноогдвор $a$ дээр голоморф байхуйц $(z-a)^k$ зэргийн  хамгийн их эрэмбэтэй таарна.
\end{thm}
\begin{prf}
$a$ тэгийн эрэмбийг $n$-ээр, $f(z)$-ийн хуваагдах $z-a$ биномын хамгийн их эрэмбийг $N$-ээр тус тус тэмдэглэе. (\ref{5.29}) томьёо ёсоор $f(z)$ аливаа $k\leq n$ зэрэгт хуваагдах тул
$$\frac{f(z)}{(z-a)^k}=(z-a)^{n-k}\varphi (z).$$
 \end{prf}
Иймд $N\geq n$. $f(z)$ функц $(z-a)^n$ зэрэгт хуваагдана гэе.
Өөрөөр хэлбэл
$$\frac{f(z)}{(z-a)^N}=\psi(z)$$
ноогдвор $a$ цэг дээр голоморф функц дүрсэлнэ. $\psi (z)$ функцийг $(z-a)$-ийн зэргээр цуваанд задлаад улмаар $f(z)$ функцийн $a$ цэг дээрх Тейлорын задаргаа $N$-ээс багагүй зэрэгтээр бичигдэх нь илэрхий. Иймд $N\leq n$ болж өмнөхтэй нэгтгэхэд $N=n$ байх нь тогтоогдов.
\begin{exm}
$f(z)=\dfrac{z^7}{1-\cos z}$ функцийн $z=0$ байх тэгийн эрэмбийг ол.
\end{exm}
\begin{sol}
$$f(z)=\frac{z^7}{1-\left(1-\frac{z^2}{2!}+\frac{z^4}{4!}-\cdots \right)}=\frac{z^5}{\frac{1}{2!}-\frac{z^2}{4!}+\cdots}=z^5\varphi(z),$$ энд $\varphi (z)=\dfrac{1}{\frac{1}{2}-\frac{z^2}{4}+\cdots}$ тул $\varphi (0)\neq 0$ болж (\ref{5.29}) томьёо ёсоор $z=0$ нь 5-р эрэмбийн тэг юм. Энэ бодлогыг 
$$\left(\frac{z^7}{1-\cos z}\right)^{(v)}\neq 0$$
байхаар ч шийдэж болно. 
\end{sol}
\begin{hnt}
$f(z)$ функц $\infty$ дээр голоморф бөгөөд $\infty$ нь түүний тэг байг. Энэ $\infty$ хэмээх тэгийн эрэмбэ нь $\varphi (z)=f\left(\dfrac{1}{z}\right)$ функцийн $a=\infty$ гэсэн тэг цэгийн эрэмбэ гэдэг нь ойлгомжтой. $a=\infty$ цэгийн хувьд өмнөх теорем хүчинтэй, харин $(x-a)^k$-д хуваахын оронд $z^k(k=1,2,\ldots)$-аар үржүүлж авна.

Функцийн тэг цэгийн эрэмбийн ухагдахууныг голоморф функцийн $A$-цэгийнх болгон өргөтгөж болно.

$a\in\overline{\mathbb{C}}$ цэгийн хувьд $f(a)=A$ байвал түүнийг $f(z)$ функцийн $A$ цэг \index{A-цэг} гэнэ. $a$ дээр голоморф $f(z)$ функцийн $A$-цэгийн эрэмбэ \index{A-цэгийн эрэмбэ} нь $f(z)-A$ функцийн тэг цэгийн эрэмбэ байна.
\end{hnt}
\begin{thm}
\textbf{(Вейерштрасс).}  $f(z)\in H(D)$ функцүүдийн
\begin{eqnarray} \label{5.32}
     f(z)=\sum\limits_{n=0}^\infty f_n(z)
\end{eqnarray} 
цуваа $\forall G\Subset D$ дээр жигд нийлж байвал
\begin{enumerate}
    \item $f(z)\in H(D)$
    \item Цувааг $\forall z\in D$ дээр гишүүнчлэн дифференциалчилагдана. 
\end{enumerate}
\end{thm}
\begin{prf}
$\forall a\in D$ байг. $U=\{z:|z-a|<r\}\Subset D$ дугуй авъя. (\ref{5.32}) цуваа өгөгдсөн нөхцөл ёсоор $U$ дээр жигд нийлэх ба гишүүд нь $U$ дээр тасралтгүй, улмаар $f(z)$ нь $U$ дээр тасралтгүй байна.

$\Delta \subset U$ авахад $\partial \Delta$ дээр (\ref{5.32}) цуваа жигд нийлэх тул түүнийг гишүүнчлэн интегралчилж болох тул
$$\int\limits_{\partial \Delta}f(z)dz=\sum\limits_{n=0}^\infty \int\limits_{\partial \Delta}f(z)dz=0$$
болж Морера-ийн теорем ёсоор $f(z)\in H(U)$ болж 1) нь батлагдав.

Одоо 2)-ийг батлахын тулд $a\in D$ цэгийн хувьд дахиад л $U=\{|z-a|<r\}\Subset D$ дугуй авъя. $|z-a|=r$ тойргийг $\partial U$ гэвэл уламжлалын Кошийн теорем ёсоор
\begin{eqnarray} \label{5.33}
     f^{(n)}(a)=\frac{n!}{2\pi i}\int\limits_{\partial U} \frac{f(z)dz}{(z-a)^{n+1}}, \quad n=0,1,\ldots
\end{eqnarray}
байна. 

$U$ дээр 
\begin{eqnarray} \label{5.34}
     \frac{f(z)dz}{(z-a)^{k+1}}=\sum\limits_{k=0}^\infty \frac{f_n(z)}{(z-a)^{k+1}}
\end{eqnarray}
(энэ нь (\ref{5.32})-аас $\forall \in\partial U$-ийн хувьд модуль нь $\dfrac{1}{r^{k+1}}$-тай тэнцэх үржигдэхүүнээр ялгагдана) цуваа жигд нийлэх учраас (\ref{5.34})-ийг (\ref{5.33}) интегралд тавьж болно:
$$f^{(k)}(a)= \sum\limits_{k=0}^\infty \frac{k!}{2\pi i} \int\limits_{\partial U} \frac{f_{n}(z)dz}{(z-a)^{k+1}}=\sum\limits_{k=0}^\infty f_n^{(k)}(a)$$ болж 2) батлагдав. 
 \end{prf}

Анализад олон гишүүнтээс зохиосон цуваа онцгой үүрэг гүйцэтгэдэг. Вейерштрассын теорем ёсоор хэрэв олон гишүүнтүүдийн 
\begin{eqnarray}
     f(z)=\sum\limits_{n=0}^\infty P_n(z)
\end{eqnarray}
цуваа ямар нэгэн $D$ мужид компактаар орших дэд олонлог бүр дээр жигд нийлж байвал түүний нийлбэр $f(z)\in H(D)$ гэж батлагдана.

Практикт, $f(z)\in H(D)$ байх аливаа функцэд олон гишүүнтүүдийн дарааллаар дөхөх урвуу бодлого чухал байдаг. Олон гишүүнтээр жигд дөхөх тухай бодлогын тавилттай танилцъя:\\
$D\subset \overline{\mathbb{C}}$ муж болон $f(z)\in H(D)$ гэж өгөгдсөн байг.

Дурын $K\Subset D$ олонлог ба дурын $\varepsilon >0$ тооны хувьд
\begin{eqnarray}
     ||f(z)-P(z)||_k=\sup_{z\in K}|f(z)-P(z)|<\varepsilon
\end{eqnarray}
байхаар $P(z)$ олон гишүүнт байгуулах шаардлагатай. Энэ бодлого нь $D$ мужийн компакт дэд олонлог бүр дээр $f(z)$-рүү жигд нийлэх олон гишүүнтүүдийн $\sum\limits_{n=0}^\infty P_n(z)$ цуваа байгуулахтай чацуу чанартай бодлого юм.
\begin{exm}
$f(z)=(z^2-1)\sinh z$ функцийг $z_0=0$ цэг дээр цуваанд задал.
\end{exm} 
\begin{sol}
\ref{5.26} ёсоор 
\begin{align*}
    (z^2-1)\sinh z&=(z^2-1)\sum\limits_{n=0}^\infty \frac{z^{2n+1}}{(2n+1)!} \\
    &=\sum\limits_{n=0}^\infty \frac{z^{2n+3}}{(2n+1)!}-\sum\limits_{n=0}^\infty\dfrac{z^{2n+1}}{(2n+1)!}
\end{align*}
\end{sol}
\begin{exm}
$e^{-z^2}$ функцийг $z_0=0$ цэг дээр цуваанд задал.
\end{exm} 
\begin{sol}
(\ref{5.23}) задаргаанд $-z^2$-ийг $z$-ийн оронд тавьж
$$e^{-z^2}=1-z^2+\frac{z^4}{2!}+\cdots =\sum\limits_{n=0}^\infty (-1)^n\frac{z^{2n}}{n!}$$
\end{sol}
\begin{exm}
$\sin^2z$ функцийг $z_0=0$ цэг дээр цуваанд задал.
\end{exm}
\begin{sol}
$\sin^2 z=\dfrac{1}{2}-\dfrac{1}{2}\cos 2z$ тэнцэлд (\ref{5.25}) томьёо ашиглавал 
$$\sin^2 z=\sum\limits_{n=0}^\infty (-1)^{n+1}\frac{2^{2n-1}z^{2n}}{(2n)!}.$$
\end{sol}
\begin{exm}
$f(z)=\ln (1+z)$, $f(0)=1$ функцийг $z$-ийн зэргээр цуваанд задал.
\end{exm}
\begin{sol}
Геометрийн прогресс $\dfrac{1}{1+z}=\sum\limits_{n=0}^\infty (-1)^n z^n; \ \ |z|<1$ дугуй дээр нийлэх тул гишүүнчлэн интегралчилбал
\begin{align*}
    f(z)&=\ln (1+z)=\bint\limits_0^z \frac{dz}{1+z}=z-\frac{z^2}{2}+\frac{z^3}{3}-\cdots \\
    &=\sum\limits^{n=0}_\infty (-1)^n\frac{z^{n+1}}{n+1}, \quad R=1.
\end{align*}
\end{sol}
\begin{exm}
$\dfrac{z}{z+2}$ функцийг $z_0=1$ цэгийн орчинд өөрөөр хэлбэл $z-1$-ийн зэргээр цуваанд задал.
\end{exm}
\begin{sol}
$\dfrac{z}{z+2}=\dfrac{z-1+1}{3\left(1+\frac{z-1}{3}\right)}$ тул 
$$\frac{z}{z+2}\sum\limits _{n=0}^\infty \frac{(-1)^n}{3^{n+1}}(z-1)^{n+1} +\sum\limits_{n=0}^\infty \frac{(-1)^n}{3^{n+1}}(z-1)^n, \quad R=3.$$
\end{sol}
\begin{exm}
$\dfrac{1}{3-2z}$ функцийг $z_0=3$ цэг дээр цуваанд задалж, 
нийлэлтийн радиусыг ол.
\end{exm}
\begin{sol}
$$\frac{1}{3-2z}=\frac{1}{3-2(z-3+3)}=\frac{1}{-3-2(z-3)}=-\frac{1}{3}\frac{1}{1+\frac{2(z-3)}{3}}$$ бөгөөд 
$\left|\dfrac{2(z-3)}{3}\right| <1$, өөрөөр хэлбэл $|z-3|<\dfrac{3}{2}$ үед 
$$\frac{1}{3-2z}=-\frac{1}{3}\sum\limits_{n=0}^\infty (-1)^n\left(\frac{2}{3}\right)^n (z-3)^n=\sum\limits_{n=0}^\infty \frac{(-1)^{n+1}2^n}{3^{n+1}}(z-3)^n, $$
$$R=\frac{3}{2}.$$
\end{sol}
\begin{exm}
$\dfrac{1}{1-z+z^2}$ функцийг $z$-ийн зэргээр задал.
\end{exm}
\begin{sol}
\begin{align*}
  \dfrac{1}{1-z+z^2}&=\dfrac{1+z}{1+z^3} \\
  &=\sum\limits_{n=0}^\infty (-1)^n z^{3n}+\sum\limits_{n=0}^\infty (-1)^nz^{3n+1}, \quad R=1.  
\end{align*}
\end{sol}
\begin{center}
    \textbf{Дасгал}
\end{center}
\begin{enumerate}
    \item 
Дараах функцүүдийг $(z-z_0)$-ийн зэргээр Тейлорын цуваанд задал. Нийлэлтийн радиусыг ол.
\colhoy{
\begin{enumerate}
    \item [a)] $\cos 2z, z_0=1$
     \item [б)] $ze^{-z}, z_0=1$
      \item [в)] $\dfrac{1}{z-1}, z_0=2$
      \item [г)] $\dfrac{z}{z^2+4}, z_0=2$
       \item [д)] $\cos^2z, z_0=\pi$
        \item [е)] $\dfrac{z}{z^2-2z-3}, z_0=0$
         \item [ж)] $e^z, z_0=\dfrac{1}{2}$
          \item [з)] $\cos z, z_0=-\dfrac{\pi}{4}$
\end{enumerate}
}
\item Дараах функцүүдийн тэг цэг олж, тэдгээрийн эрэмбийг тогтоо.
\colgur{
\begin{enumerate}
    \item [a)] $1+\cos z$
    \item [б)] $z^4+4z^2$
    \item [в)] $z\sin z -z^2$
\end{enumerate}
}
\end{enumerate}

\chapter{Лораны цуваа
}
\section{Лораны цуваа} \index{Лораны цуваа}
$|z-a|<R$ дугуй дээр голоморф функцийг $z-a$-ийн сөрөг биш зэрэгтээр бичигдэх Тейлорын цуваатай өмнөх бүлэгт танилцсан. $z-a$-ийн эерэг ба сөрөг зэрэгтээр бичигдсэн
$$\sum\limits_{n=-\infty}^\infty c_n(z-a)^n$$
(энд $a$-бэхэлсэн комплекс тоо, $c_n$-тогтмол комплекс коэффициентууд) цувааг Лораны цуваа гэдэг.
\begin{thm}
\textbf{(Лоран).} \index{Лораны теорем}Цагириг хэлбэрийн 
$$V=\{z:r<|z-a|<R\}, \quad r\geq 0, \quad R\leq \infty$$
муж дээр голоморф аливаа $f(z)$ функц нийлэх
\begin{eqnarray} \label{6.1}
     f(z)=\sum\limits_{n=-\infty}^\infty c_n(z-a)^n
\end{eqnarray}
цувааны нийлбэр дүрстэй тавигдана. Коэффициентууд нь дараах томьёогоор олдоно:
\begin{eqnarray} \label{6.2}
     c_n=\frac{1}{2\pi i}\int\limits_{|z-a|=\rho} \frac{f(\zeta)d\zeta}{(\zeta -a)^{n+1}}, \quad n=0,\pm 1, \pm 2, \ldots
\end{eqnarray}
үүнд $r<\rho < R$ байна.
\begin{center}
    \includegraphics[width=4cm]{Zurag/20.jpg}
    \captionof{figure}{}
        \label{}
\end{center}
\end{thm}
\begin{prf}
Бэхлэн авсан $z$ цэг дотор нь байхаар $V'\Subset V$ гэсэн шинэ цагираг муж байгуулъя: $V'=\{\zeta :r'<|\zeta-a|<R'\}$ энд $r<r'$, $R'<R$. Кошийн интеграл томьёо ёсоор
\begin{eqnarray} \label{6.3}
     f(z)=\frac{1}{2\pi i} \int\limits_{\partial V'} \frac{f(\zeta)d\zeta}{(\zeta -z)}= \frac{1}{2\pi i} \int\limits_{\Gamma'} \frac{f(\zeta)d\zeta}{(\zeta -a)}-\frac{1}{2\pi i} \int\limits_{\gamma'} \frac{f(\zeta)d\zeta}{(\zeta -a)}
\end{eqnarray}
энд $\Gamma '=\{|\zeta -a|=R'\}$ ба $\gamma '=\{|\zeta -a|=r'\}, \quad \zeta\in \Gamma '$ байх үед $\left|\dfrac{z-a}{\zeta -z}\right|=q<1$ тул 
$$\frac{1}{\zeta-z}=\frac{1}{(\zeta -z)\left(1-\frac{z-a}{\zeta -a}\right)}=\sum\limits_{n=0}^\infty \frac{(z-a)^n}{(\zeta -a)^{n+1}}$$
$\zeta \in \Gamma'$-ээрээ абсолют бөгөөд жигд нийлнэ. Үүнийг зааглагдсан $\dfrac{1}{2\pi i}f(\zeta)$ функцээр үржүүлж (жигд нийлэлт алдагдахгүй тул) дараа нь $\Gamma '$ тойргоор гишүүнчлэн интегралчлавал
\begin{eqnarray} \label{6.4}
    \frac{1}{2\pi i}\int\limits_{\Gamma'} \frac{f(\zeta)d\zeta}{(\zeta -z)}= \sum\limits_{n=-\infty}^\infty c_n(z-a)^n
\end{eqnarray}
\begin{eqnarray} \label{6.5}
    c_n=\frac{1}{2\pi i}\int\limits_{\Gamma'} \frac{f(\zeta)d\zeta}{(\zeta -a)^{n+1}}, \quad (n=0,1,\ldots)
\end{eqnarray}
(\ref{6.3}) томьёоны хоёрдох интеграл арай өөрөөр задлагдана. $\zeta \in \gamma'$ байх үед $\left|\dfrac{z-a}{\zeta -z}\right| =q<1$ тул бид $\gamma'$ дээр абсолют бөгөөд жигд нийлэх геометр прогрессыг 
$$-\frac{1}{\zeta -z}=\frac{1}{(\zeta -a)\left(1-\frac{\zeta -a}{z-a}\right)}=\sum\limits_{n=1}^\infty \frac{(\zeta-a)^{n-1}}{(z -a)^{n}}$$
байдалтай гарган авна. Үүнийг мөн л зааглагдсан $\dfrac{1}{2\pi i}f(\zeta)$
функцээр үржүүлж $\gamma '$-ээр гишүүнчлэн интегралчилбал
\begin{eqnarray} \label{6.6}
    -\frac{1}{2\pi i}\int\limits_{\gamma'} \frac{f(\zeta) d\zeta}{\zeta -z}=\sum\limits_{n=1}^\infty \frac{d_n}{(z-a)^n},
\end{eqnarray}
энд 
\begin{eqnarray} \label{6.7}
    d_n= \frac{1}{2\pi i}\int\limits_{\gamma'} f(\zeta)(\zeta -z)^{n-1}d\zeta, \quad (n=1,2,\ldots)
\end{eqnarray}
(\ref{6.7}) ба (\ref{6.5}) томьёонуудаас
\begin{eqnarray}
    c_n=d_{-n}= \frac{1}{2\pi i}\int\limits_{\gamma'} f(\zeta)(\zeta -z)^{-n-1}d\zeta, \quad (n=-1,-2, \ldots)
\end{eqnarray}
болж (\ref{6.6})
\begin{equation}
    \tag{6.6'} \label{6.6'}
    -\frac{1}{2\pi i}\int\limits_{\gamma'} \frac{f(\zeta)d\zeta}{(\zeta -a)}=\sum\limits_{n=-1}^{-\infty} c_n(z-a)^n
\end{equation}
(\ref{6.4}) ба (\ref{6.6'})-ийг (\ref{6.3}) тавибал (\ref{6.1}) задаргаа гарна. Кошийн интеграл томьёо ёсоор $\Gamma'$ ба $\gamma '$ тойргуудыг тэдэнтэй гомотоп дурын тойргоор (тухайлбал $|\zeta -a|=\rho, \quad r<\rho <R$.) сольж болох тул (\ref{6.2}) томьёо гарна.
 \end{prf}
\begin{dfn}
Коэффициентууд нь (\ref{6.2}) томьёонуудаар тодорхойлогдох (\ref{6.1}) цувааг цагираг $V$ муж дахь $f(z)$ функцийн Лораны цуваа гэнэ. Энэ цувааны сөрөг биш зэрэгтээр тавигдсан хэсгийг зөв хэсэг, сөрөг зэрэгтээр тавигдсан хэсгийг гол хэсэг \index{Лораны цувааны гол ба зөв хэсэг} гэж нэрлэнэ. Иймд
\begin{eqnarray}
    \sum\limits_{n=-\infty}^\infty c_n(z-a)^ n
\end{eqnarray}
цуваа
\begin{eqnarray}
    S_1=\sum\limits_{0}^\infty c_n(z-a)^n \quad \text{ба} \quad S_2=\sum\limits_{n=-1}^{-\infty} c_n(z-a)^ n
\end{eqnarray}
цуваануудын нийлбэр байдлаар тодорхойлогдоно. \\
$S_1$ цуваа нь ердийн зэрэгт цуваа тул түүний нийлэлтийн радиус $R$
\begin{eqnarray}
    \frac{1}{R}=\overline{\lim\limits_{n\rightarrow \infty}} \sqrt[n]{|c_n|}
\end{eqnarray}
томьёогоор $S_2$ цувааны нийлэлтийн радиус $r$
\begin{eqnarray}
    r=\overline{\lim\limits_{n\rightarrow \infty}} \sqrt[n]{|c_{-n}|}
\end{eqnarray}
томьёогоор тус тус тогтоогдохыг төвөггүй харуулж болно. $R>r$ байж болохгүй. Учир нь тийм байвал цувааны нийлэх муж хоосон олонлог байхад хүрнэ.

Абелийн теорем ёсоор $V$ цагираг мужид компактаар орших аливаа дэд олонлог дээр жигд нийлэх тул Вейерштрассын теорем ёсоор түүний нийлбэр $V$ дээр голоморф байна. Эдгээр тэмдэглэл дээр тулгуурлан Лораны цувааны задаргааны цор ганц байхыг тогтоож болно.
\end{dfn}
\begin{thm} \label{Теорем 6.2}
Функцийн Лораны задаргаа цор ганц \index{Лораны задаргааны цор ганц}, өөрөөр хэлбэл $f(z)$ функц $0\leq r<|z-a|<R$ цагираг дээр 
\begin{eqnarray} \label{6.13}
    f(z)=\sum\limits_{k=-\infty}^\infty a_n(z-a)^ k
\end{eqnarray}
дүрстэй цуваанд задарсан бол $a_k$ коэффициентууд (\ref{6.5}) томьёогоор илэрхийлэгдэнэ.
\end{thm}
\begin{prf}
$\gamma =|z-a|=\rho$ тойрог авъя. (\ref{6.13})-ын хоёр талыг 
$\dfrac{1}{2\pi i}(z-a)^{-n-1}, \quad (n=-1,-2,\ldots)$-ээр үржүүлж $\gamma$ тойргоор гишүүнчлэн интегралчилбал 
$$\sum\limits_{n=-\infty}^\infty a_n \int\limits_\gamma (z-a)^{k-n-1}dz= \int\limits_\gamma \frac{f(z)dz}{(z-a)^{n+1}}$$
бөгөөд зэргийн ортогональ чанараар $k\neq n$ байх бүх интеграл тэг болж $k=n$ тохиолдолд $2\pi i$ болох тул
$$2\pi i a_n=\int\limits_\gamma  \frac{f(z)dz}{(z-a)^{n+1}}, \quad (n=0, \pm 1, \ldots)$$ эндээс $c_n\equiv a_n$ болов. 
 \end{prf}
Теорем 6.2-ыг өөрөөр томьёолбол $z-a$-ийн эерэг, сөрөг зэрэгтээр задарсан нийлэх аливаа цуваа нь өөрийн нийлбэрийнхээ Лораны цуваа байна.

Лораны цувааны коэффициентуудыг (\ref{6.2}) томьёогоор бодох нь нэн ховор. Учир нь интегралчлах ажиллагаа шаардагдана. Задаргаа цор ганц байх теоремын үндсэн дээр хуулийн хүчинтэй аливаа арга хэрэглэнэ. Ийм арга бүхэн хүсэн хүлээсэн ганц л үр дүнд хүргэнэ.
\begin{exm}
$f(z)=\dfrac{1}{(z-1)(z-2)}$ функцийг $z$-ийн зэрэгтээр Лораны цуваанд задал.
\end{exm}
\begin{sol}
Энэ функц $z=1$, $z=2$ цэгүүд дээр голоморф чанараа алдана. Иймд 
$$D_1:|z|<1, \quad D_2:1<|z|<2, \quad D_3:2<|z|<\infty$$ цагирагууд дээр голоморф функц юм. $f(z)=\dfrac{1}{z-2}-\dfrac{1}{z-1}$ дүрстэй тавигдах тул 
$$\frac{1}{z-2}=-\frac{1}{2}\cdot \frac{1}{1-\frac{z}{2}}=-\frac{1}{2}\sum\limits_{n=0}^\infty \left(\frac{z}{2}\right)^n \quad (|z|<2)$$
$-\dfrac{1}{z-1}=\sum\limits_{n=0}^\infty z^n \qquad (|z|<1).$ Иймд $|z|<1$ байхад 
$$f(z)=\sum\limits_{n=0}^\infty \left(1-\frac{1}{2^{n+1}}\right) z^n \qquad \text{(Энэ нь Тейлорын цуваа)}$$
Одоо $D_2$ цагирагт
$$-\frac{1}{z-1}=-\frac{1}{z}\frac{1}{1-\frac{1}{z}}=-\frac{1}{z}\sum\limits_{n=0}^\infty \frac{1}{z-n}=-\sum\limits_{n=0}^\infty \frac{1}{z^n+1}\quad (1<|z|)$$
Иймд $1<|z|<2$ байхад 
$$f(z)=-\sum\limits_{n=-1}^{-\infty} z_n-\sum\limits_{n=0}^\infty \frac{z^n}{2^{n+1}}.$$
Одоо $D_3$ цагирагт
$$\frac{1}{z-2}=\frac{1}{z}\frac{1}{1-\frac{2}{z}}=\frac{1}{z}\sum\limits_{n=0}^\infty \left(\frac{2}{z}\right)^n=\sum\limits_{n=1}^\infty \frac{2^{n-1}}{z^n}\quad (2<|z|).$$
Иймд $2<|z|<\infty$ байхад 
$$f(z)=\sum\limits_{n=2}^\infty (2^{n-1}-1)z^{-n}.$$
\end{sol}
\begin{exm}
$f(z)=\dfrac{1}{z(1-z)}$ функцийг $z=0$, $z=1$, $z=\infty$ цэгийн орчинд Лораны цуваанд задал.  
\begin{center}
    \includegraphics[width=4cm]{Zurag/21.jpg}
    \captionof{figure}{}
        \label{}
\end{center}
\end{exm}
\begin{sol}
$z=0$ ба $z=1$ нь $\dfrac{1}{z(1-z)}=\dfrac{1}{z}+\dfrac{1}{1-z}$ функцийн голоморф алдагдах цэгүүд. \\
$D_1:z=0$ цэгийн цоорхойт орчин тул $0<|z|<1$ бөгөөд 
$$\frac{1}{z(1-z)}=\frac{1}{z}+\frac{1}{1-z}=\frac{1}{z}+\sum\limits_{n=0}^\infty z^n$$ болж задаргаа гол хэсэгтээ ганц гишүүнтэй байна. \\
$D_3:z=\infty$ цэгийн орчин $1<|z|$ 
\begin{center}
    \includegraphics[width=3cm]{Zurag/22.jpg}
    \captionof{figure}{}
        \label{}
\end{center}
$$\frac{1}{z(1-z)}=-\frac{1}{z^2}\frac{1}{1-\frac{1}{z}}=-\sum\limits_{n=2}^\infty \frac{1}{z^n} \quad \left(\left|\frac{1}{z}\right|<1 \ \ \text{тул} \right)$$ 
$D_2:z=1$ цэгийн орчин $|z-1|<1$ 
$$\frac{1}{z(1-z)}=-\frac{1}{z-1}+\frac{1}{z-1+1}=-\frac{1}{z-1}+\sum\limits_{n=0}^\infty (-1)^n (z-1)^n.$$ \end{sol}
Лораны цувааны коэффициентын (\ref{6.2}) томьёо $n\geq 0$ үед Тейлорын цувааны коэффициент олох интеграл томьёотой тохирно.

$f(z)$ функц $V=\{r<|z-a|<R\}$ цагираг дээр голоморф, $\gamma =\{|z-a|=\rho\}, \quad r<\rho <R$ тойрог дээр $|f(z)|<M$ байвал задаргаа коэффициент
\begin{eqnarray} \label{6.14}
    |c_n|<\frac{M}{\rho^n},\quad n=0, \pm 1,\ldots
\end{eqnarray}
(\ref{6.14})-ийг Лораны цувааны Кошийн тэнцэл биш \index{Лораны цувааны Кошийн тэнцэл биш} гэнэ.

Одоо Лораны цуваа, Фурьегийн цуваатай хэрхэн холбогдохыг \index{Лораны цуваа, Фурьегийн цуваатай холбогдох нь} авч үзье. $[0,2\pi]\subset \mathbb{R}$ дээр интегралчилагдах $\varphi (z)$ функц Фурьегийн цуваагаар 
\begin{eqnarray}
    \varphi(t)=\frac{a_0}{2}+\sum\limits_{n=1}^\infty (a_n\cos nt +b\sin nt),
\end{eqnarray}
гэж тодорхойлогдоно. Энд
\begin{eqnarray}
\begin{split}
     &a_n=\frac{1}{\pi} \int\limits_0^{2\pi} \varphi(t)\cos nt dt \quad n=0,1,\ldots\\
    &b_n= \frac{1}{\pi} \int\limits_0^{2\pi} \varphi(t)\sin nt dt, \quad n=0,1,\ldots
\end{split}
\end{eqnarray}
Энэ цувааг Эйлерийн $\cos nt =\dfrac{e^{int}+e^{-int}}{2}$, $\sin nt =\dfrac{e^{int}-e^{-int}}{2}$ томьёо ашиглах замаар комплекс хэлбэртэй бичиж болно:
$$\frac{a_0}{2}+\sum\limits_{n=1}^\infty \left(\frac{a_n-ib_n}{2} e^{int}+\frac{a_n+ib_n}{2}e^{-int}\right) =\sum\limits_{n=-\infty}^\infty c_ne^{int}$$
энд $$c_n=\frac{a_n-ib_n}{2}=\dfrac{1}{2\pi }\int\limits_{0}^{2 \pi}\varphi (t)e^{-int}dt, \quad n=0,1,\ldots$$
$$c_n=\frac{a_{-n}+ib_n}{2}=\dfrac{1}{2\pi }\int\limits_{0}^{2 \pi}\varphi (t)e^{-int}dt, \quad n=-1,-2,\ldots$$
байна.

Коэффициентууд нь
\begin{eqnarray}
c_n=\dfrac{1}{2\pi }\int\limits_{0}^{2 \pi}\varphi (t)e^{int}dt
\end{eqnarray}
томьёонуудаар олдох
\begin{eqnarray} \label{6.18}
\sum\limits_{n=-\infty}^\infty c_ne^{int}
\end{eqnarray}
цуваа $\varphi$ функцийн комплекс хэлбэртэй тавигдсан Фурьегийн цуваа байна.

$e^{int}=z$ ба $\varphi (t)=f(e^{it})=f(z)$ гэж авбал \ref{6.18} цуваа 
\begin{eqnarray} 
\sum\limits_{n=-\infty}^\infty c_n z^n
\end{eqnarray}
болж коэффициентууд нь 
\begin{eqnarray}
c_n=\dfrac{1}{2\pi }\int\limits_{0}^{2 \pi} f (e^{it})e^{-int}dt= \dfrac{1}{2\pi i}\int\limits_{|z|=1} f(z)\frac{dz}{z^{n+1}}
\end{eqnarray}
Иймд $\varphi (t)$, $t\in [0, 2\pi]$ функцийн комплекс хэлбэртэй бичигдсэн Фурьегийн цуваа нь $f(z)=\varphi (t)$, $z=e^{it}$ функцийн нэгж тойрог $(|z|=1)$ дээр тавигдсан Лораны цуваа юм.

Фурьегийн цуваа $[0,2\pi]$ хэрчмийн цэг бүр дээр $\varphi (t)$-рүү нийлэх тохиолдолд ч харгалзах Лораны цувааны хувьд $R=r=1$ болж нийлэлтийн муж нь хоосон олонлог болно.

$\varphi (t)$ функц нэлээд хатуу шаардлага хангах нөхцөлд харгалзах Лораны цуваа нийлэлтийн хоосон биш мужтай байдаг.
\section{Тусгаарлагдсан онцгой цэгүүд,\\
тэдгээрийг ангилах нь
}
\begin{dfn}
$a\in\overline{\mathbb{C}}$ цэг $f(z)$ функцийн онцгой цэг, харин түүний цоорхойт ($a$-төгсгөлтэй бол $0<|z-a|<r$, $a=\infty$ бол $R<|z|<\infty)$ орчинд уул функц голоморф байвал $a$-цэгийг тусгаарлагдсан онцгой цэг \index{Тусгаарлагдсан онцгой цэг} гэнэ.

Тийм цэг рүү дөхөхөд $f(z)$ функцийн төрх ямар байхаас хамаарч тусгаарлагдсан онцгой цэгийг гурав ангилна. 
\end{dfn}
\begin{dfn}
Тусгаарлагдсан онцгой $a$ цэгийн хувьд
\begin{enumerate}
    \item [I.] $\lim\limits_{z\rightarrow a} f(z)=A$-төгсгөлтэй байвал $a$-засагдах онцгой цэг. \index{Засагдах онцгой цэг}
    \item [II.] $\lim\limits_{z\rightarrow a} f(z)=\infty$ байвал $a$-туйл.
     \item [III.] $\exists f(z)$ нь $z\rightarrow a$ үед төгсгөлтэй ч, төгсгөлгүй ч хязгаар байхгүй байвал $a$-гоц онцгой цэг гэж нэрлэнэ.
\end{enumerate}

Тусгаарлагдсан онцгой а цэгийн төрх нь түүний цоорхойт орчинд задалсан Лораны цуваатай нягт холбоотой байдаг. Төгсгөлтэй а цэгийн хувьд энэ холбоог харуулсан дараах гурван теоремтой танилцъя:
\end{dfn}
\begin{thm}
Тусгаарлагдсан онцгой $a\in\mathbb{C}$ цэг засагдах онцгой цэг байх зайлшгүй бөгөөд хүрэлцээтэй нөхцөл нь түүний орчин дахь Лораны задаргаа гол хэсэг агуулахгүй байх явдал мөн. Өөрөөр хэлбэл
\begin{eqnarray} \label{6.21}
f(z)=\sum\limits_{n=0}^\infty c_n(z-a)^n
\end{eqnarray}
\end{thm}
\begin{prf}
\textbf{Зайлшгүй нь- } $a$ засагдах онцгой цэг байг. Тэгвэл төгсгөлтэй $\lim\limits_{z\rightarrow a}=A$ хязгаар оршин байж эндээс $a\in\mathbb{C}$-ийн цоорхойт $(0<|z-a|<R)$ орчинд зааглагдсан $(|f(z)|<M)$ байна. $0<\rho <R$ байх аливаа $\rho$ авч Кошийн тэнцэл биш бичвэл
$$|c_n|\leq \frac{M}{\rho^n}, n=0, \pm 1, \ldots$$
бөгөөд $n<0$ тохиолдолд $\rho\rightarrow 0$ үед баруун гар тал нь тэг рүү тэмүүлэх тул $c_{-n}=0$, $n=1,2,\ldots$ болж гол хэсэггүй Лораны цуваа үүснэ.\\
\textbf{Хүрэлцээтэй нь}. $a$ цэгийн цоорхойт орчинд Лораны (\ref{6.21}) цуваа гол хэсэггүй задарсан байг. Иймд $\lim\limits_{z\rightarrow a}f(z)=c_0$ болж $a$ нь засагдах онцгой цэг байх болов.
 \end{prf}
\begin{thm} \label{Теорем 6.2.2}
Тусгаарлагдсан онцгой $a\in\mathbb{C}$ цэг туйл байх зайлшгүй бөгөөд хүрэлцээтэй нөхцөл нь түүний орчин дахь Лораны задаргааны гол хэсэг нь $0$-ээс ялгаатай төгсгөлтэй бөгөөд эерэг тооны гишүүн агуулах явдал мөн:
\begin{eqnarray}\label{6.22}
f(z)=\sum\limits_{n=-N}^\infty c_n(z-a)^n, \quad N>0
\end{eqnarray}
\end{thm}
\begin{prf}
\textbf{Зайлшгүй нь}. $a\in\mathbb{C}$ цэг туйл байг. $\lim\limits_{z\rightarrow a}f(z)=\infty$ тул $\varphi(z)=\dfrac{1}{f(z)}$ функц $a$ цэгийн орчинд голоморф бөгөөд $\lim\limits_{z\rightarrow a}\varphi(z)=0$ байна. Иймд $z=a$ нь $\varphi(z)$-ийн засагдах онцгой цэг учир
$$\varphi(z)=b_N (z-a)^N+b_{N+1}(z-a)^{N+1}+\cdots \quad (b_N\neq 0)$$
болно. Эндээс
\begin{eqnarray} \label{6.23}
f(z)=\frac{1}{\varphi(z)}=\frac{1}{(z-a)^N}\cdot \frac{1}{b_N+b_{N+1}(z-a)+\cdots}
\end{eqnarray}
болж хоёрдугаар үржигдэхүүн нь $z=a$ дээр голоморф байх тул тэр нь $z=a$ дээр тейлороор
$$\frac{1}{b_N+b_{N+1}(z-a)+\cdots}=c_{-N}+c_{-N+1}(z-a)+\cdots$$
гэж задрах бөгөөд $C_{-N}=\dfrac{1}{b_N}\neq 0$ болох ба үүнийг (\ref{6.23})-т тавибал 
$$f(z)=\frac{c_{-N}}{(z-a)^N}+\frac{v_{-N+1}}{(z-a)^{N-1}}+\cdots \sum\limits_{n=0}^\infty c_n(z-a)^n.$$
\textbf{Хүрэлцээтэй нь.} $a\in\mathbb{C}$ цэгийн цоорхойт орчинд  функц (\ref{6.22}) байдлаар гол хэсэгтээ төгсгөлтэй гишүүн агуулж задарсан байг. Мөн $c_{-N}\neq 0$ байг. Тэгвэл $\varphi(z)=(z-a)^N\cdot f(z)$ функц $a$-ийн орчинд голоморф
$$\lim\limits_{z\rightarrow a}\varphi(z) =c_{-N}\neq 0$$
байна. Нөгөөтээгүүр $\lim\limits_{z\rightarrow a}f(z) =\lim\limits_{z\rightarrow a}\dfrac{\varphi (z)}{(z-a)^N}=\infty$  болж $a$ нь $f(z)$-ийн туйл ажээ. 
 \end{prf}

Функцийн туйл ба тэгийн хоорондох холбооны тухай авч үзье.
\begin{thm}
$a$ цэг $f(z)$-ийн туйл байх зайлшгүй бөгөөд хүрэлцээтэй нөхцөл нь
 $\varphi(z)=\dfrac{1}{f(z)}$, $\varphi(z) \not \equiv  0$ функц $a$ цэгийн орчинд голоморф бөгөөд $\varphi(a)=0$ байх явдал мөн.
 \end{thm}
 \begin{prf}
 Энэ өгүүлбэрийн зайлшгүйг өмнөх теоремд баталсан. Харин хүрэлцээтэй нь: $\varphi (z)\not\equiv 0$ функц $a$ цэг дээр голоморф бөгөөд $\varphi(a)=0$ байвал цор ганц байх чанарын теорем ёсоор $\varphi(z)\neq 0$ байхуйц $a$ цэгийн цоорхойт орчин олдоно. Энэ орчинд $f(z)=\frac{1}{\varphi(z)}$ функц голоморф байх тул $a$ нь $f(z)$ функцийн тусгаарлагдсан онцгой цэг юм. Гэтэл $\lim\limits_{z\rightarrow a}f(z)=\infty$ учир $a$ нь $f(z)$-ийн туйл. 
  \end{prf}
 \begin{dfn}
 $f(z)$ функцийн туйлын эрэмбэ \index{Туйлын эрэмбэ} нь $\varphi(z)=\dfrac{1}{f(z)}$ функцийн тэгийн эрэмбэ байна.
 
 Өмнөх Теорем 6.4-ын баталгаанаас харахад туйлын эрэмбэ нь энэхүү туйлын цоорхойг орчин дахь Лораны задаргааны гол хэсгийн ахлах гишүүний дугаар болох $N$ тоотой тэнцэх нь ил байна.
 \end{dfn}
 \begin{thm}
 Тусгаарлагдсан онцгой $a$ цэг $f(z)$ функцийн гоц онцгой цэг байх зайлшгүй бөгөөд хүрэлцээтэй нөхцөл нь түүний орчин дахь Лораны задаргаа төсгөлгүй тооны гишүүн бүхий гол хэсэг агуулна.
 \end{thm}
 \begin{prf}
 Энэ теорем өмнөх 6.3 ба 6.4 теоремд агуулагдана. Гоц онцгой цэгийн орчин дахь функцийн төрх байдлыг харуулах дараах теоремтой танилцъя. 
  \end{prf}
 \begin{thm}
 \textbf{(Вейерштрасс-Сохоцкий).} \index{Вейерштрасс-Сохоцкийн теорем} $a$ нь $f(z)$-ийн функцийн гоц онцгой цэг байвал аливаа $A\in\overline{\mathbb{C}}$ тооны хувьд $a$-руу нийлэх $\{z_n\}$ дараалал 
 $$\lim\limits_{n\rightarrow \infty} f(z_n)=A$$ байхаар олдоно.
 \end{thm}
 \begin{prf}
 $A=\infty$ байхад теорем үнэн. Учир нь 
  $$\overline{\lim\limits_{n\rightarrow a}} |f(z_n)|=\infty$$
  $A\neq \infty$ байг. Баталгааг эсрэгээс нь хийе. $a$-цэгийн аливаа орчинд өөр дээр нь $f(z)$ функцийн авч чадах утга $A$ тоонд аль болох ойр байх цэгүүд олдохгүй байвал $U(a,\delta)$ орчин болон $a\neq z\in U(a,\delta)$ үед $|f(z)-A|>\alpha$ байхуйц $\alpha>0$ тоо тус тус оршин байна. Тэхэд цоорхойт $0<|z-a|<\delta$ орчинд $\varphi(z)=\dfrac{1}{f(z)-A}$ функц голоморф ба зааглагдсан: 
  $$|\varphi(z)|<\frac{1}{\alpha}$$
  Иймд $\lim\limits_{z\rightarrow a}\varphi(z)$ төгсгөлтэй бөгөөд харин $f(z)$ функцийн $a$ цэгийн аливаа орчинд үл зааглагдах явдлаас үүдэж энэ хязгаар зөвхөн тэг:
  $$\lim\limits_{z\rightarrow a}\varphi(z)=\lim\limits_{z\rightarrow a}\frac{1}{f(z)-A}=0$$
  $f(z)-a$ функц Теорем 6.4 ёсоор $a$ гэсэн туйлтай болов. Гэтэл энэ нь теоремын нөхцөлд харшлах тул теорем батлагдав.
  \end{prf}
 
 $a$ руу нийлэх янз бүрийн $\{z_n\}$ дарааллуудаар $f(z)$ функцийн хязгаарын утгуудын бөөгнөрлийг энэ функцийн $a$ цэг дээрх тодорхойгүй байдлын олонлог гэж нэрлэнэ. Хэрэв $a$ цэг нь $f(z)$ функцийн засагдах онцгой цэг буюу туйл байвал энэ цэг дээрх тодорхойгүй байдлын олонлог зөвхөн ганц (төгсгөлтэй буюу төгсгөлгүй) цэгээс тогтоно. Вейерштрасс-Сохоцкийн теорем гоц онцгой цэгийн хувьд нөгөө захын тохиолдол: энэ цэг дээрх тодорхойгүй байдлын олонлог нь бүхий л битүү комплекс хавтгай $\overline{\mathbb{C}}$-ийг дүүргэнэ гэж баталж байна.
 
 Одоо төгсгөлгүй алсын цэг тусгаарлагдсан онцгой байх тухай цөөн зүйл тэмдэглэе. Төгсгөлгүй алсын $(a=\infty)$ цэг дээрх Лораны задаргааны эерэг биш зэрэгт нь зөв, эерэг зэрэгт нь гол хэсгүүдийг бүрдүүлнэ.Энэ нь $z=\dfrac{1}{w}$ оруулгаар $f(z)=f\left(\dfrac{1}{w}\right)=\varphi(w)$ болж $\lim\limits_{z\rightarrow \infty}f(z)=\lim\limits_{w\rightarrow 0}\varphi (w)$ байдгаас шууд илэрхий.\\
 Жишээ нь $\varphi(w)$ функцийн туйлын хувьд $0<|w|<\rho$ байхад 
 $$\varphi (w)=\frac{b_{-N}}{w^N}+\cdots +\frac{b_{-1}}{w}+\sum\limits_{n=0}^\infty b_nw^n, \quad b_{-N}\neq 0$$
 бөгөөд одоо $w=\dfrac{1}{z}$ гэвэл $f(z)$ функц $\left(R<|z|<\infty, \quad R=\dfrac{1}{2}\right)$ үед
 $$f(z)=\sum\limits_{n=-1}^{-\infty} c_nz^n+c_0+c_1z+\cdots +c_Nz^N,$$
 үүнд, $c_n=b_{-N}$, $c_N\neq 0$ болох гол хэсэг нь төгсгөлтэй тооны гишүүн агуулна. Энэчлэн засагдах ба гоц онцгой цэгийн хувьд задална.
 
 Одоо онцгой цэгийн төрхөөр нь голоморф функцийг ангилахтай танилцъя.
 \begin{dfn}
 $f(z)\in H(\mathbb{C})$ байх, өөрөөр хэлбэл \textbf{төгсгөлтэй} онцгой цэггүй функцийг бүхэл функц \index{Бүхэл функц} гэнэ. Иймд $a=\infty$ нь бүхэл функцийн тусгаарлагдсан онцгой цэг байна. Хэрэв энэ нь засагдах онцгой цэг байвал $f\equiv \infty$ (Лиувиллийн теорем ёсоор) хэрэв энэ нь туйл байвал $f(z)$-ийн орчин дахь Лораны задаргааны гол хэсэг $g(z)=c_1 z+\cdots c_N z^N$ гэсэн олон
гишүүнт байна. Энэ гол хэсгийг $f(z)$ функцээс хасвал бас л бүхэл $f(z)-g(z)$ функц гарах ба харин $\infty$ нь засагдах онцгой байна. Энэ нь тогтмол тоо болов. Иймд бүхэл функц $\infty$ цэг дээр туйлтай байвал тэр нь ямагт олон гишүүнт байна.

Төгсгөлгүй алсын цэг дээр гоц онцгой цэгтэй бүхэл функцүүдийг бүхэл трансцендент функцүүд (жишээ нь $e^z$, $\sin z$, $\cos z$ ) гэнэ. 
 \end{dfn}
 \begin{dfn}
 Задгай $\mathbb{C}$ хавтгайд туйлаас өөр онцгой цэг үгүй функцийг мероморф функц гэнэ.
 \end{dfn}
 \section{Үлдэцийн онол}
 
 $a$ нь $f(z)$-ийн тусгаарлагдсан онцгой цэг, $L$ нь $a$ цэгийг дотроо агуулсан шулуутгагдах энгийн битүү муруй, $f(z)$ нь $L$ дээр ба түүний доторх $a$-аас бусад цэг дээр голоморф байх функц гэвэл
 $$\frac{1}{2\pi i}\oint\limits_{L} f(z)dz$$
 илэрхийллийг $f(z)$ функцийн $a$-цэгт хамаарах үлдэц\index{Үлдэц} гэж нэрлэх ба $Res [f(z),a]$ (Residu-үлдэгдэл гэсэн франц үг) тэмдэглэл оруулна. Иймд
 \begin{eqnarray}
Res [f(z),a]=\frac{1}{2\pi i}\oint\limits_{L} f(z)dz
 \end{eqnarray}
 \begin{thm}
\textbf{ (Үлдэцийн үндсэн теорем).\index{Үлдэцийн үндсэн теорем} } $f(z)$ функц төгсгөлөг тооны $a_k(k=\overline{1,n})$ онцгой цэгээс бусад $D$ муж ба $\partial D$ хил дээр голоморф (Зураг 23)
\begin{center}
    \includegraphics[width=4cm]{Zurag/23.jpg}
    \captionof{figure}{}
        \label{}
\end{center}
байвал 
\begin{eqnarray} \label{6.25}
\int\limits_{\partial D}f(z)dz=2\pi i\sum_{k=1}^n Res [f(z), a_k]
\end{eqnarray} байна.
 \end{thm}
 \begin{prf}
 $a_k(k=\overline{1,n})$ цэгүүд дээр төвтэй хүрэлцээтэй бага $\rho$ радиустай сөрөг чиг барьсан $\gamma_k:|z-a_k|=\rho$ тойргууд татъя. Мөн $|z-a_k|\leq \rho$ дугуйнууд $D$-д агуулагдаж, $\gamma_k$ тойрог бүхэн бусад тойргийнхоо гадна оршин байвал олон холбоост мужийн Кошийн теорем ёсоор
 $$\oint\limits_{\partial D} f(z)dz+\sum_{k=1}^n \oint\limits_{\gamma_k}f(z)dz=0$$
 болж эндээс 
 $$\oint\limits_{\partial D} f(z)dz=\sum_{k=1}^n \oint\limits_{\gamma_k}f(z)dz=2\pi i\sum_{k=1}^n Res [f(z), a_k] $$ 
  \end{prf}
 Онцгой цэг дээрх үлдэцийг олж чаддаг бол (\ref{6.25}) томьёо ашиглан комплекс интеграл бодох боломж нээгдэнэ.
 \begin{thm}
 $a\in\mathbb{C}$ цэг $f(z)$-ийн тусгаарлагдсан онцгой цэг гэвэл
\begin{eqnarray}
Res [f(z), a]=c_{-1}
\end{eqnarray}
 \end{thm}
 \begin{prf}
 (\ref{6.25}) томьёонд $n=-1$ гэвэл
 $$c_{-1}=\frac{1}{2\pi i}\int\limits_{\Gamma'}f(\zeta) d\zeta=Res [f(z),a] \quad \Gamma':|z-a|=R$$ 
  \end{prf}
 \begin{cor}
 $a$ засагдах онцгой цэг бол
 $$Res[f(z),a]=0$$
 \end{cor}
 \begin{thm}
 $f(z)$-функцийн хувьд $a$ нь $n\geq 1$ эрэмбийн туйл байг. Тэгвэл:
 \begin{eqnarray} \label{6.27}
 Res [f(z),a]=\dfrac{1}{(n-1)!}\lim\limits_{z\rightarrow a}\frac{d^{(n-1)}}{dz^{n-1}}[(z-a)^n\cdot f(z)].
 \end{eqnarray}
 \end{thm}
 \begin{prf}
$a$ цэгийн цоорхойт орчинд Лораны цуваанд $f(z)$ функцийг задалбал
$$f(z)=\frac{c_{-n}}{(z-a)^n}+\cdots +\frac{c_{-1}}{z-a}+\sum\limits_{n=0}^\infty c_n(z-a)^n$$
бөгөөд хоёр талыг $(z-a)^n$-ээр үржүүлж $(n-1)$ удаа дифференциалчилбал
$$\frac{d^{(n-1)}}{dz^{n-1}}[(z-a)^n\cdot f(z)]=(n-1)!c_{-1}+n(n-1)!+2c_0(z-a)+\cdots$$
бөгөөд $z\rightarrow a$ хязгаарт шилжвэл (\ref{6.27}) гарна.

$z=a$ нь $f(z)$ функцийн энгийн туйл байхад (\ref{6.27}) томьёо ёсоор ($0!=1$  гэж үзэх тул)
$$Res [f(z),a]=\lim\limits_{z\rightarrow a}(z-a)f(z)$$
юм. Хэрэв $f(z)=\dfrac{\varphi (z)}{\psi(z)}$, $(\varphi, \psi\in H(a)$, $\varphi(a)\neq 0$, $\psi(a)=0)$ бол 
$$Res[f(z),a]=\lim\limits_{z\rightarrow a}f(z)(z-a)=\lim\limits_{z\rightarrow a}\frac{\varphi (z)}{\frac{\psi(z)-\psi (a)}{z-a}}=\frac{\varphi(a)}{\psi '(a)}$$ 
  \end{prf}
 \begin{exm}
 $\tan z$ функцийн $\dfrac{\pi}{2}$ дээрх үлдэц ол.
 \end{exm}
 \begin{sol}
 $\tan z=\dfrac{\sin z}{\cos z}$, $\sin\dfrac{\pi}{2}=1$, $\cos\dfrac{\pi}{2}=0$ тул $$Res \left[\tan z, \dfrac{\pi}{2}\right] =-\sin \frac{\pi}{2}=-1$$
 \end{sol}
 \begin{exm}
 $(z-1)\cdot \sin\dfrac{1}{z-1}$ функцийн $z=1$ дээрх үлдэц ол.
  \end{exm}
\begin{sol}
	(\ref{5.19}) томьёонд $z$-ийн оронд $\dfrac{1}{z-1}$тавибал
	$$\sin\frac{1}{z-1}=\sum\limits_{n=0}^\infty (-1)^n \frac{1}{(2n+1)!(z-1)^{2n+1}}$$ тул 
	$$Res\left[(z-1)\cdot \sin \frac{1}{z-1}, 1\right] =-\frac{1}{6}$$
\end{sol}
\begin{exm}
$\dfrac{\sin 2z}{(z+1)^3}$ функцийн $z=-1$ дээрх үлдэц ол.
\end{exm}
\begin{sol}
    $z=-1$ нь III эрэмбийн туйл тул 
    $$Res \left[\frac{\sin 2z}{(z+1)^3}, -1\right]=\lim\limits_{z\rightarrow -1}\frac{1}{2!}(\sin 2z)''=2\sin 2$$
\end{sol}
\textbf{Төгсгөлгүй алсын цэг дээр үлдэц ямар байх вэ? } 
\begin{dfn}
$$\frac{1}{2\pi i}\int\limits_{\gamma_R}-f(z)dz$$
хэмжигдэхүүнийг $f(z)$ функцийн $\infty$ дээрх үлдэц гэнэ. Үүнд $\gamma^-_R$ нь сөрөг чиг барьсан (нар зөв тойрох) хүрэлцээтэй их радиустай $|z|=R$ тойрог. $f(z)$ функцийг $\infty$-ийн цоорхойт \\ $R<|z|<\infty$ орчинд задалбал
$$f(z)=\sum\limits_{n=-\infty}^\infty c_n z^n$$
бөгөөд $\gamma^-_R$тойргоор гишүүнчлэн интегралчилж, зэрэгтийн ортогональ чанар ашиглавал
$$Res [f(z),\infty]=-c_{-1}$$
\end{dfn}
\begin{thm}
\textbf{(Үлдэцүүдийн бүтэн нийлбэр)} \index{Үлдэцүүдийн бүтэн нийлбэр} $f(z)$ нь $\mathbb{C}$ дээр $a_k(k=\overline{1,n})$ онцгой цэгтэй байвал
\begin{eqnarray}
\sum\limits_{k=1}^nRes [f(z),a_k]+Res [f(z),\infty]=0
\end{eqnarray} байна. 

Энэ хялбар теоремыг өөрөө батална уу.
\end{thm}
\begin{exm}
$\dfrac{e^z}{z^2(z^2+9)}$ функцийн бүх онцгой цэг болон төгсгөлгүй алсын цэг\index{Төгсгөлгүй алсын цэг дээрх үлдэц} дээрх үлдэц ол.
\end{exm}
\begin{prf}
$z=0$ II эрэмбийн туйл,$z=\pm 3i$ энгийн туйлууд.
\begin{align*}
    Res \left[\frac{e^z}{z^2(z^2\pm 9)}, 0\right] &=\frac{1}{2!}\lim\limits_{z\rightarrow 0}\left(\frac{e^z}{z^2+9}\right)'\\
    &=\lim\limits_{z\rightarrow 0} \frac{e^z(z^2+9)-2ze^z}{(z^2+9)^2}=\frac{1}{18};
\end{align*}
\begin{align*}
    &Res \left[\frac{e^z}{z^2(z^2+ 9)}, 3i\right]=\lim\limits_{z\rightarrow 3i}\frac{e^z}{z^2(z+3i)}=-\frac{1}{54}(\sin 3-i\cos 3)\\
   &Res \left[\frac{e^z}{z^2(z^2+ 9)}, -3i\right]=\lim\limits_{z\rightarrow -3i}\frac{e^z}{z^2(z-3i)}=-\frac{1}{54}(\sin 3+i\cos 3)\\
   &Res \left[\frac{e^z}{z^2(z^2+ 9)}, \infty\right]=\frac{1}{27}(\sin 3-3)
\end{align*}
 \end{prf}
\begin{exm}
$\bint\limits_{|z|=1} \dfrac{e^zdz}{z^2(z^2-9)}$ бод. 
\end{exm}
\begin{sol}
    $|z|=1$ тойрог дотор $z=0$ гэсэн II эрэмбийн туйл байх тул
    \begin{align*}
        \bint\limits_{|z|=1} \dfrac{e^zdz}{z^2(z^2-9)}&=2 \pi i Res\left[ \dfrac{e^zdz}{z^2(z^2-9)}, 0\right]\\
        &=2 \pi i \lim\limits_{z\rightarrow 0}\left(\frac{e^z}{z^2-9}\right)'=2\pi i \lim\limits_{z\rightarrow 0}
        \frac{e^z(z^2-9)-2ze^z}{(z^2-9)^2} \\
        &=-\frac{2\pi i}{9}
    \end{align*}
    (үлдэцийн үндсэн теорем ашиглав). 
\end{sol}
\section{Интеграл бодоход үлдэц ашиглах нь}
\begin{enumerate}
    \item [\textbf{I.}] $\bint\limits_0^{2\pi} R(\cos t,\sin t)dt$, $R(u,v)$ нь бодит $u,v$ гэсэн хоёр хувьсагчийн рациональ функц.
    
    
$R(\cos t, \sin t)$ нь $0\leq t\leq 2\pi$ завсар дээр тасралтгүй функц.
Эйлерийн томьёо ёсоор
$$\cos t=\frac{e^{it}+e^{-it}}{2}, \sin t=\frac{e^{it}-e^{-it}}{2}$$
бөгөөд $z=e^{it}$ гэвэл $|z|=1$, $t\in[0,2\pi]$ үед $z=e^{it}$ цэг $z=1$ цэгээс нар буруу (цагийн зүүний эсрэг) чиглэлээр нэгж тойргийг нэг удаа тойрно. $\cos t=\dfrac{z+z^{-1}}{2}$, $\sin t=\dfrac{z-z^{-1}}{2i}$, $dt=\dfrac{dz}{iz}$ байх тул 
\begin{eqnarray}
\begin{split}
    \int\limits_0^{2\pi}R(\cos t,\sin t)dt&=\oint\limits_{|z|=1}R\left(\frac{1}{2}\left[z+\frac{1}{z}\right], \frac{1}{2}\left[z-\frac{1}{z}\right]\right)\frac{dz}{iz} \\
    &=\oint\limits_{|z|=1} \tilde{R}(z)dz,
\end{split}
\end{eqnarray}
$\tilde{R}(z)$ нь рациональ функцүүдийн композиц тул рациональ функц байна. Энэ л функц нь $|z|=1$ дээр голоморф, түүний дотор төгсгөлтэй тооны онцгой цэгтэй учир үлдэцийн онолоор бодогдоно.
\begin{exm}
$\bint\limits_{0}^{2\pi} \dfrac{dt}{\cos t-2}$ бодож гарга.
\end{exm}
\begin{sol}
$$\int\limits_{0}^{2\pi} \dfrac{dt}{\cos t-2} =\int\limits_{|z|=1} \dfrac{dz}{iz\left(\frac{z+z^{-1}}{2}-2\right)}=\frac{2}{i}\oint\limits_{|z|=1}\frac{dz}{z^2-4z+1}$$ болж $\dfrac{2}{i}\dfrac{1}{z^2-4z+1}$ функцийн зөвхөн $z_1=2-\sqrt{3}$ цэг нь $|z|<1$-д орших тул 
\begin{align*}
    \int\limits_{0}^{2\pi} \dfrac{dt}{\cos t-2}& =2\pi i Res\left[ \dfrac{1}{z^2-4z+1}, 2-\sqrt{3}\right] \\
    &=4 \pi \frac{1}{2(2-\sqrt{3})-4} \\
    &=-\frac{2\pi}{\sqrt{3}}
\end{align*}

\end{sol}
\item[\textbf{II.}]
$\bint \limits_{-\infty}^\infty R(x)dx$, $R(x)$ нь бодит тэнхлэг дээр туйл үгүй рациональ функц, $R(x)=\dfrac{P_m(x)}{Q_n(x)}$ гэвэл $P_m(x)$, $Q_n(x)$ олон гишүүнтүүд ба $m+2\leq n$. Өөрөөр хэлбэл $z=\infty $ нь уул $R(x)$ функцийн 2-оос доошгүй эрэмбийн тэг байг.

Бодит тэнхлэгийн $[-R;R]$ хэрчим $|z|=R$, $\im z>0$ хагас тойргоор битүүрэхдээ $R(z)$ функц комплекс хавтгайд орших бүх туйлыг хамран авахаар радиус $R$-ыг сонгон авна.
\begin{center}
    \includegraphics[width=4cm]{Zurag/24.jpg}
    \captionof{figure}{}
        \label{}
\end{center}
$\infty$-ийн орчинд $R(x)$ функцийн Лораны цуваанд задалбал 
$R(z)=\dfrac{c_{-2}}{z^2}+\dfrac{c_{-3}}{z^3}+\cdots$. Иймд $\lim\limits_{z\rightarrow \infty} R(z)\cdot z^2=c_{-2} \ \text{ба} \ \lim\limits_{z\rightarrow \infty}|R(z)|=0,$
$|z|\rightarrow \infty$ үед $|R(z)|\cdot |z|=0(|z|).$ Эндээс аливаа $\varepsilon >0$ тооны хувьд хүрэлцээт их $R$-утганд $\left|\bint\limits_{|z|=R}R(z)dz\right|\leq \varepsilon$. Иймд $R\rightarrow \infty$ үед $\lim\limits_{R\rightarrow \infty}\bint\limits_{|z|=R}R(z)dz=0$ болж улмаар
\begin{eqnarray} \label{6.30}
\begin{split}
    \int\limits_{-\infty}^\infty R(x)dx&= \lim\limits_{R\rightarrow \infty} \int\limits_{-R}^R R(x)dx \\
    &=2\pi i\sum\limits_{k=1}^n Res [R(z), a_k]
\end{split}
\end{eqnarray}
\begin{exm}
$f(z)=\dfrac{1}{(z^2+1)^2}$ функц дээд хагас хавтгайд цор ганц $z=i$ гэсэн II эрэмбийн туйлтай, бодит тэнхлэг дээр туйлгүй бөгөөд $\infty$ нь IV эрэмбийн тэг байх тул (\ref{6.30}) ёсоор 
$$\int\limits_{-\infty}^\infty \frac{dz}{(z^2+1)^2}=2\pi i Res \left[\frac{1}{(z^2+1)^2}, i\right]$$
гэж бодогдоно.

Харин
$$Res \left[\frac{1}{(z^2+1)^2}, i\right]=\lim\limits_{z\rightarrow i} \left[\frac{(z-i)^2}{(z-i)^2(z+i)^2}\right]'=\frac{1}{4i}$$
тул 
$$\int\limits_{-\infty}^\infty \frac{dx}{(x^2+1)^2}=\frac{\pi}{2}.$$
\end{exm}
\item[III.]\textbf{Жорданы лемм.} \index{Жорданы лемм} \\
$F(z)=e^{\lambda iz}\cdot f(z)$, $\lambda>0$, $f(z)$ функц дараах 3 нөхцөлийг хангадаг байг.
\begin{enumerate}
    \item Дээд хагас хавтгайд төгсгөлтэй $n$-тооны онцгой \\ $a_k(k=\overline{1,n})$ цэгтэй,
    \item $Ox$ тэнхлэг дээр төгсгөлтэй $m$-тооны зөвхөн энгийн $x_j(j=\overline{ 1,m})$ туйлтай, 
    \item $\im z \geq 0$ байхаар $z\rightarrow \infty$ үед $f(z)\rightarrow 0$ байвал
\begin{align*}
    \int\limits_{-\infty}^\infty e^{\lambda ix}f(x)dx&= 2\pi i\left\{\sum\limits_{k=1}^n Res[F(z),a_k]+\right.\\
    &\left.+\frac{1}{2}\sum\limits_{j=1}^m Res[F(z), x_j]\right\}
\end{align*}
\end{enumerate}
\begin{center}
    \includegraphics[width=6cm]{Zurag/25.PNG}
    \captionof{figure}{}
        \label{}
\end{center}
\end{enumerate}
\begin{prf}
Бүх $a_k(k=\overline{1,n})$ ба $x_j(j=\overline{1,m})$ онцгой цэгүүдийг хүрээндээ хамрахуйц $C_R:|z|=R$, $\im z>0$ гэсэн хагас тойрог татъя.

Мөн $x_j(j=\overline{1,m})$ цэгүүдэд төвтэй дээд хагас хавтгайд орших, зэргэлдээ хоёр нь ерөнхий цэггүй байхаар $\gamma^-_\rho: |z-x_j|=\rho$, $\im z>0(j=\overline{1,m})$ гэсэн сөрөг тойролттой хагас тойрог авбал үлдэцийн үндсэн теорем ёсоор 
\begin{equation}
    \tag{$*$} \label{*}
    \begin{split}
        &\int\limits_{-R}^{x_1-\rho}F(z)dz+\int\limits_{x_1+\rho}^{x_2-\rho}F(z)dz+\cdots +\\
        &+\int\limits_{x_m+\rho}^RF(z)dz+\sum\limits_{j=1}^m \int\limits_{\gamma_\rho^-}e^{\lambda i z}f(z)dz+\int\limits_{C_R}e^{\lambda iz}f(z)dz = \\
        &=2\pi i\sum\limits_{k=1}^n Res[e^{\lambda iz} f(z), a_k]
    \end{split} 
\end{equation} бөгөөд 
\begin{align*}
    &\lim\limits_{\tiny{\begin{array}{c}
R\to\infty  \\
 \rho\to 0 
\end{array}}}  \left(\int\limits_{-R}^{x_1-\rho} F(z)dz+\cdots +\int\limits_{x_m+\rho}^R F(z)dz\right)=\int\limits_{-\infty}^\infty F(x)dx, \\
    &\lim\limits_{\tiny{\begin{array}{c}
R\to\infty  \\
 \rho\to 0 
\end{array}}}\sum\limits_{j=1}\int\limits_{\gamma^-_\rho}F(z)dz=-\pi i\sum\limits_{j=1}^m Res[F(z), x_j]
\end{align*}
байхыг тооцож хязгаарт шилжсэний дараа (\ref{*}) нь 
\begin{align*}
    &\int\limits_{-\infty}^\infty e^{\lambda ix}f(x)dx+\lim\limits_{\tiny{\begin{array}{c}
R\to\infty  \\
 \rho\to 0 
\end{array}}} \int\limits_{C_R} e^{\lambda iz}f(z)dz= \\
    &=2\pi i\left\{\sum\limits_{k=1}^n Res[F(z), a_k]+\frac{1}{2}\sum\limits_{j=1}^m Res [F(z), x_j]\right\}
\end{align*}
болно. $C_R:z=R(\cos \varphi +i\sin\varphi)$, $0\leq \varphi \leq \pi$ байх тул 
\begin{align*}
    |I_R|&=\left|\int\limits_{C_R}e^{\lambda iz}f(z)dz\right| \\ &=\left|\int\limits_0^\pi e^{\lambda i R\cos \varphi}\cdot e^{-\lambda R\sin \varphi}\cdot f(Re^{i\varphi})\cdot Rid\varphi\right|.
\end{align*}
Гуравдугаар нөхцөл ёсоор $\max\limits_{0\leq \varphi \leq \pi}|f(\re^{i\varphi})|=\sigma(R)\rightarrow 0 (R\rightarrow \infty)$ ба  $|e^{\lambda iR\cos \varphi}|=1$ учир 
$$|I_R|\leq \sigma (R)\int\limits_0 e^{-\lambda R\sin \varphi}d\varphi $$
улмаар хялбар $\sin \varphi >\dfrac{2\varphi}{\pi} \left(\forall \varphi \in\left[0;\dfrac{\pi}{2}\right]\right)$ тэнцэл биш ашиглавал 
$$|I_R|\leq \sigma (R)\int\limits_0^{\frac{\pi}{2}}Re^{-\lambda R\cdot\frac{2\varphi}{\pi}}d\varphi=\frac{\pi \sigma(R)}{\lambda}(1-e^{-\lambda R})\rightarrow 0 (R\rightarrow \infty)$$ 
болж (\ref{*}) тэнцэл:
\begin{eqnarray}\label{6.31}
\begin{split}
    \int\limits_{-\infty}^\infty e^{\lambda ix}f(x)dx& =2\pi i \left\{\sum\limits_{k=1}^nRes[e^{\lambda iz}f(z),a_k]+\right.\\
    &\left.+\frac{1}{2}\sum\limits_{j=1}^m Res[e^{\lambda iz}f(z), x_j]\right\}
\end{split}
\end{eqnarray} 
\begin{exm}
$\int\limits_{0}^\infty \dfrac{\sin ax}{x(x^2+b^2)}dx, \quad (a,b$-эерэг тоонууд) интегралыг бод. 
\end{exm}
 \end{prf}
\begin{sol}
Интеграл доорх функц тэгш, нөгөө талаар \\ $\sin ax=\im e^{iax}$ тул (\ref{6.31}) ёсоор 
\begin{align*}
    \int\limits_0^\infty \frac{\sin ax}{x(x^2+b^2)}dx&=\im \left\{\frac{1}{2}\int\limits_{-\infty}^\infty \frac{e^{iax}}{x(x^2+b^2)}dx\right\} \\
    &=\im \left\{\pi i\left(Res\left[\frac{e^{iaz}}{z(z^2+b^2)}, bi\right]\right.\right.+\\
    &\left.\left.+\frac{1}{2}Res \left[\frac{e^{iaz}}{z(z^2+b^2)},0\right]\right)\right\}
\end{align*} бөгөөд үлдэцүүдийг жич олбол:
$$Res\left[\frac{e^{iaz}}{z(z^2+b^2)}, bi\right] =\lim\limits_{z\rightarrow bi}\frac{e^{iaz}}{z(z+bi)}=-\frac{e^{-ab}}{2b^2}$$
$$Res\left[\frac{e^{iaz}}{z(z^2+b^2)}, 0\right] =\lim\limits_{z\rightarrow 0}\frac{e^{iaz}}{z^2+b^2}=\frac{1}{b^2}$$ болж 
$$\int\limits_0^\infty \frac{\sin ax dx}{x(x^2+b^2)}=\im\left\{\left(\frac{-e^{-ab}}{2b^2}+\frac{1}{2}\frac{1}{b^2}\right)\cdot \pi i\right\}=\frac{\pi}{2b^2}(1-e^{-ab})$$
\end{sol}
\begin{center}
   \textbf{ Дасгал}
\end{center}
Үлдэц ашиглан дараах интегралуудыг бод.

\begin{enumerate}
\colhoy{
    \item $\bint\limits_{0}^{2\pi} \dfrac{d\varphi}{\cos \varphi +3}$;
    \item $\bint\limits_{0}^{2\pi} \dfrac{dx}{\sin x +2}$;
    \item $\bint\limits_{-\infty}^{\infty} \dfrac{dx}{(x^2+1)(x^2+2)}$;
       \item $\bint\limits_{-\infty}^{\infty} \dfrac{(x-1)dx}{(x^2+1)^2};$ }
       \item $\bint\limits_{-\infty}^{\infty} \dfrac{\cos xdx}{x^4+10x^2+9}$;
       \item $\bint\limits_{0}^{\infty} \dfrac{x\sin  ax dx}{x^2+b^2}, \quad (a>0, b>0)$; \colhoy{
          \item $\bint\limits_{-\infty}^{\infty} \dfrac{x\cos x dx}{x^2-5x+6}$;
            \item $\bint\limits_{-\infty}^{\infty} \dfrac{\sin x dx}{x^3-x^2+4x-4}$}
\end{enumerate}

%\chapter{Шугаман ба бусад хялбар хувиргалтууд}
\section{Ерөнхий шугаман функц}
 \subsection{Бүхэл шугаман функц}
 
\begin{eqnarray} \label{7.1}
W=az+b \ \ a\neq 0 (\text{үүнд:} a,b \ - \text{тогтмол комплекс тоонууд,)}
\end{eqnarray}
дүрстэй тавигдана. $\dfrac{1}{a}\neq 0$ тул (\ref{7.1})-ийн урвуу нь бас шугаман функц байна. $z=\infty$ байхад $w=\infty$ гэж урьдчилан тодорхойлолт хийж болох учир (\ref{7.1}) нь $\overline{\mathbb{C}}\rightarrow \overline{\mathbb{C}}$ харилцан нэгэн утгатай буулгана.

$a=|a|e^{i\alpha}$, $\alpha =\Arg  a$ байх тул $w=az+b=|a|e^{i\alpha}z+b$ болж бүхэл шугаман (\ref{7.1}) функц
\begin{enumerate}
    \item [$1^\circ$] $w_1=|a|z,$
     \item [$2^\circ$] $w_2=e^{i\alpha }w_1,$
      \item [$3^\circ$] $w=w_2+b$ гэсэн 3 хялбар буулгалтаар задлагдана.
\end{enumerate}
\begin{enumerate}
    \item [$1^\circ$]. $z$ цэгийг төвт төсөө (гомотет)-өөр $w_1$ цэгт буулгана. Энд төсөөгийн төв нь $z=0$, коэффициент нь $|a|$ байна (Зураг 26а).
    \begin{center}
    \includegraphics[width=8cm]{Zurag/26.jpg}
    \captionof{figure}{}
        \label{}
\end{center}
\item [$2^\circ$] $w_1=u_1+iv_1$, $w_2=u_2+iv_2$, $e^{i\alpha}=\cos \alpha +i\sin \alpha$ тус тус авбал 
$$
\begin{cases}
 u_2=u_1\cos \alpha-v_1\sin \alpha \\
 v_2=u_1 \sin \alpha +v_1\cos\alpha
\end{cases}
$$
болж $w_1$ нь $\alpha$ өнцгөөр эргэж $w_2$ цэгт бууна (Зураг 26б).
\item [$3^\circ$] $w=u+iv$, $w_2=u_2+iv_2$, $\overline{b}=b_1+ib_2$ гэж авбал
$$
\begin{cases}
 u=u_2+b_1 \\
 v=v_2+b_2
\end{cases}
$$ болж $w_2$ цэг өгөгдсөн тогтмол $\overline{b}$ вектортой параллель зөөлтөөр $w$ цэгт бууна (Зураг 26в).

Иймд бүхэл шугаман (\ref{7.1}) буулгалт нь дурдсан 3 хялбар буулгалтыг цувруулан гүйцэтгэсний үр дүн, өөрөөр хэлбэл, эдгээр буулгалтын үржвэрээр тодорхойлогдох болой.

Шугаман (\ref{7.1}) буулгалтыг 
\begin{eqnarray}
w-A=a(z-A)
\end{eqnarray}
байдалтай тавигдана гэвэл $A=\dfrac{b}{1-a}$ байна. $a=1$ бол ердийн параллель зөөлт.\\
$a\neq 0$ бол $A$-цэгийн орчинд эргүүлэхийг төсөөтэй буулгалттай хослуулан гүйцэтгэнэ.
\end{enumerate}
\begin{center}
    \includegraphics[width=6cm]{Zurag/27.PNG}
    \captionof{figure}{}
        \label{}
\end{center}
\section{Бутархай шугаман функц} \index{Бутархай шугаман функц}
\begin{eqnarray} \label{7.3}
w(z)=\frac{az+b}{cz+d}, \qquad ad-bc\neq 0
\end{eqnarray}
дүрстэй бичигдэх буулгалтыг бутархай шугаман функц \index{Шугаман функц} гэнэ. $a$, $b$, $c$, $d$-бэхлэгдсэн тогтмол комплекс тоонууд, $ad-bc\neq 0$ нь $w(z)$-ийн тогтмол тоо бишийг илтгэнэ. $(ad-bc=0$ байвал (\ref{7.3})-ын хүртвэр, хуваарь пропорциональ болж $w(z)$ нь $z$-ийн функц биш болоход хүргэнэ).

$c=0$, $d\neq 0$ тохиолдолд (\ref{7.3}) нь өмнө танилцсан бүхэл шугаман функц үүсгэнэ.

$c\neq 0$ үед (\ref{7.3}) нь $z\neq -\dfrac{d}{c}$, $\infty$ цэгүүд дээр тодорхойлогдоно. Харин онцгой $-\dfrac{d}{c}$ ба $\infty$ цэгүүд дээр $w(z)$-ийн тасралтгүй чанараар урьдчилан тодорхойлолт хийж $w\left(-\dfrac{d}{c}\right)=\infty$, $w(\infty)=\dfrac{a}{c}$ гэж авдаг.

$c=0$ үед шууд л $w(\infty)=\infty$ гэж авна.
\begin{thm}
(\ref{7.3}) нь $\overline{\mathbb{C}}$-г $\overline{\mathbb{C}}$-рүү гомеоморф (харилцан нэгэн утгатай, тасралтгүй) буулгана.
\end{thm}
\begin{prf}
$c=0$ үед илэрхий. $c\neq 0$ гэе. (\ref{7.3}) функц нийт $\overline{\mathbb{C}}$ дээр нэгэн утгатай тодорхойлогдоно. Нөгөөтээгүүр 
\begin{eqnarray} \label{7.4}
z=\frac{dw-b}{-cw+a}
\end{eqnarray}
өмнө тэмдэглэсэн урьдчилсан тодорхойлолт ёсоор $w=\dfrac{a}{c}$ цэгт $z=\infty$, харин $w=\infty$ цэгт $z=-\dfrac{d}{c}$ харгалзаж (\ref{7.3})-ийн $\overline{\mathbb{C}}$-г $\overline{\mathbb{C}}$- рүү харилцан нэгэн утгатай буулгах нь тогтоогдлоо. Одоо тасралтгүйг батлах үлдэв.

$z\neq -\dfrac{d}{c}$, $z\neq \infty$ үед (\ref{7.3}) тасралтгүй байх нь илэрхий. Харин эдгээр цэг дээр 
$$\lim\limits_{z\rightarrow -\frac{d}{c}}w(z)=\infty, \qquad \lim\limits_{z\rightarrow \infty} w(z)=\frac{a}{c}$$
хязгаарууд оршин байх тул мөн тасралтгүй. 
\hfill $\blacksquare$\end{prf}

Одоо бид, (\ref{7.3}) буулгалт $\overline{\mathbb{C}}$-ийн бүх цэг дээр өнцөг хадгалахыг батлан харуулъя. $z\neq -\dfrac{a}{c}$, $\infty$ цэгүүд дээр 
$$\frac{dw}{dz}=\frac{ad-cb}{(cz+d)^2}\neq 0$$
байх тул өнцөг хадгалагдана. Харин онцгой цэгүүдийн (хоёулаа төгсгөлгүйтэй холбоотой, нэг нь өөрөө төгсгөлгүй, нөгөөгийн дүр нь төгсгөлгүй) хувьд мөнхүү чанартайг тогтоохын тулд төгсгөлгүй цэг дээрх өнцгийн тухай ойлголт оруулъя.
\begin{dfn}
$z=\infty$ цэгийг дайрах замуудын (тэдгээр нь Риманы бөмбөрцгийн хойт туйл дээр огтлолцох шүргэгчүүдтэй) хооронд үүсэх өнцөг гэж
\begin{eqnarray}
Z=\dfrac{1}{z}
\end{eqnarray}
буулгалтын $Z=0$ дээрх $\Gamma_1$ ба $\Gamma_2$  дүрүүдийн хоорондох өнцгийг ойлгоно.
\end{dfn}
\begin{thm}
Бутархай шугаман (\ref{7.3}) буулгалт $\overline{\mathbb{C}}$-ийн цэг бүр дээр конформ буулгалт ($\infty$ дээр өнцөг хадгалагдах) өгнө.
\end{thm}
\begin{prf}
Онцгой биш цэгүүдийн хувьд теорем өмнө батлагдсан. $\gamma_1$ ба $\gamma_2$ нь $z=-\dfrac{d}{c}$ цэгийг дайрсан замууд бөгөөд энэ цэг дээр $\alpha$ өнцгөөр огтлолцсон байг. Эдгээр замын (\ref{7.3}) буулгалт дахь $w=\infty(z=-\dfrac{d}{c}$ цэгт харгалзах) цэг дээр $\gamma_1^*$ ба $\gamma_2^*$ дүрүүдийн хоорондох өнцөг нь тодорхойлолт ёсоор $\gamma_1^*$ ба $\gamma_2^*$-ийн $W=\dfrac{1}{w}$ буулгалт дахь $W=0$ цэг дээрх $\Gamma^*_1$ ба $\Gamma_1^*$ дүрүүдийн хоорондох өнцөгтэй тэнцүү. Харин
$$W=\frac{cz+d}{az+b}$$
бөгөөд иймд $\Gamma_1^*$ ба $\Gamma_2^*$ нь энэ буулгалт дахь $\gamma_1$ ба $\gamma_2$-ийн дүр.
$$\frac{dW}{dz}=\frac{bc-ad}{(az+b)^n}$$
уламжлал нь $z=-\dfrac{d}{c}$ цэг дээр $0$-ээс ялгаатай тул  $\Gamma_1^*$ ба $\Gamma_2^*$-ийн $W=0$ цэг дээр үүсгэх $\alpha$ өнцөгтэй тэнцүү. $z=-\dfrac{d}{c}$ цэгийн хувьд теорем батлагдлаа. $z=\infty$ цэгийн хувьд теоремыг батлахын тулд (\ref{7.3})-ийн урвуу (\ref{7.4}) функцэд өмнөх адил хандах явдал хүрэлцээтэй.
\hfill $\blacksquare$\end{prf}
\begin{thm}
Бутархай шугаман буулгалтуудын $\wedge$ цогц бүлэг үүсгэнэ.
\end{thm}
\begin{prf}
$$\zeta=\frac{a_1z+b_1}{c_1z+d_1} \ \ \text{ба }\ \ w=\frac{a_2\zeta+b_2}{c_2\zeta+d_2}$$
гэсэн бутархай шугаман хоёр буулгалт авч $\zeta$-ийн илэрхийллийг $w$-ийн томьёонд тавибал
$$w=\frac{az+b}{cz+d}$$ болж, энд
$$a=a_1a_2+d_2c_1 \qquad b=a_2b_1+b_2d_1$$
$$c=a_1c_2+c_1d_2 \qquad d=b_1c_2+d_1d_2$$
байх бөгөөд хэрэв $a_1d_1-b_1c_1\neq 0$ ба $a_2d_2-b_2c_2\neq 0$ байвал $ad-bc\neq 0$ байна.

Бид $L_1$,$L_2\in\wedge$ бол $L_1\circ L_2\in \wedge$ гэж баталлаа. \\
Одоо бүлгийн аксиомуудын биелэгдэхийг шалгая.
\begin{enumerate}
    \item [\textbf{a)}] 
    Дамжих чанар: $L_1$,$L_2$,$L_3\in \wedge$ бол
    \begin{eqnarray} \label{7.6}
    L_1\circ (L_2\circ L_3)=(L_1\circ L_2)\circ L_3
    \end{eqnarray}
    учир нь (\ref{7.6})-ийн хоёр тал $L_1(L_2[L_3(z)])$ гэж тавигдана.
    \item [\textbf{б)}] Нэгж буулгалт орших: нэгж нь адилтгал буулгалт болох нь мэдээж.
    \begin{eqnarray}
    \wedge \ni E:z\rightarrow z
    \end{eqnarray}
     \item [\textbf{в)}] Урвуу элемент орших: $\forall L\in\wedge$ байвал 
     \begin{eqnarray}
     L^{-1}\circ L=L\circ L^{-1}=E
     \end{eqnarray}
     байх $L^{-1}\in \wedge$ олдоно. Учир нь (\ref{7.3}) буулгалтын хувьд түүний урвуу элемент болох (\ref{7.4}) олдоно. 
\end{enumerate}
\hfill $\blacksquare$\end{prf}
\begin{hnt}
$\wedge$ бүлэг байр солих хуульд үл захирагдана.\\
Жишээ нь $L_1:z\rightarrow \dfrac{1}{z}$, $L_2:z\rightarrow z+1$ гэвэл $L_2\circ L_1:z\rightarrow \dfrac{1}{z}+1$, $L_1\circ L_2:\dfrac{1}{z+1}$
\end{hnt}
\begin{thm}
Аливаа бутархай шугаман буулгалт $\overline{\mathbb{C}}$ дээрх дурын тойргийг $\overline{\mathbb{C}}$ дээрх мөн л тойрогт хувиргана. (Үүнийг бутархай шугаман буулгалтын тойргийн чанар гэнэ).
\end{thm}
\begin{prf}
$c=0$ үед (\ref{7.3}) буулгалт $w=\dfrac{a}{d}z+\dfrac{b}{d}$ болж тойргийн төсөөтэй хувиргах, эргүүлэх ба параллель зөөлт хийхэд тойрог хэвээрээ байх билээ. Харин $c\neq 0$ 
\begin{eqnarray}
L:z\rightarrow \frac{az+b}{cz+d}=\dfrac{a}{c}-\frac{ad-bc}{c(cz+d)}=A+\frac{B}{z+C}
\end{eqnarray}
болж дараах гурван хялбар буулгалт
$$L_1:z\rightarrow A+Bz, \quad L_2:z\rightarrow \frac{1}{z}, \quad L_3:z\rightarrow z+C$$
гүйцэтгэсний дүнд $L$ хийгдэнэ $(L=L_1\circ L_2\circ L_3)$. Эхний ба эцсийн буулгалт $\overline{\mathbb{C}}$ дээр тойргийг хэвээр үлдээх нь мэдээж. $L_2:z\rightarrow \dfrac{1}{z}$ буулгалтын хувьд энэ чанар хэвээр хадгалагдахыг батлах л үлдлээ.

$\overline{\mathbb{C}}$ дээр аливаа тойргийн буюу шулууны тэгшитгэл, тэгш өнцөгт декартын координатын систем дээр 
\begin{equation} \label{7.10}
 A(x^2+y^2)+2D_1x+2E_1y+F=0   
\end{equation} 

дүрстэй тавигдана. Энд $A$, $D_1$, $E_1$, $F$ бодит тоонууд, харин $A$, $D_1$, $E_1$ нь нэгэн зэрэг тэгтэй тэнцэхгүй. Тухайлбал $A=0$ бол шулуун дүрсэлнэ. Харин $A_1\neq 0$ ба  $D_1^2+E_1^2>AF$ байвал тойрог. \z  гэж авбал \hz, $x^2+y^2=z\overline{z}$ ба \\ $2x=z+\overline{z}$, $2yi=z-\overline{z}$ болж эдгээрийг (\ref{7.10})-т тавихад шулуун буюу тойргийн комплекс хэлбэртэй
тэгшитгэл
\begin{eqnarray} \label{7.11}
Az\overline{z}+Dz+\overline{D}\overline{z}+F=0
\end{eqnarray}
бөгөөд энд $\overline{D}=D_1-E_1i$, $D=D_1+E_1 i.$

(\ref{7.11}) тойрог нь $L_2$ буулгалтаар ямар дүртэй болохыг мэдэхэд $z=\dfrac{1}{w}$, улмаар $\overline{z}=\dfrac{1}{\overline{w}}$ гэж орлуулах явдал хангалттай, иймд
\begin{eqnarray}
Fw\overline{w}+D\overline{w}+\overline{D}w+A=0
\end{eqnarray}
болж (\ref{7.10})-ийн адил тэгшитгэл гарлаа. Энэ нь $\mathbb{C}$ дээр мөн л тойрог дүрсэлнэ. 
\hfill $\blacksquare$\end{prf}

Одоо бутархай шугаман буулгалтын өөр нэг геометр чанартай танилцъя.
\begin{dfn}
\textbf{а). } Хэрэв $\Gamma =\{z:|z-z_0|=R\}$ тойргийн хувьд $z$ ба $z^*$ цэгүүд $z_0$-оос эхлэх цацраг дээр оршиж $z_0$ хүртэлх зайнуудын үржвэр $R^2$-тай тэнцэх: $(\arg  (z-z_0))=arg (z^*-z_0)$ ба $|z-z_0||z^*-z_0|=R^2$ байвал тэдгээрийг $\Gamma$ тойрогтой тэгш хэмтэй цэгүүд гэнэ.  \\
\textbf{б).} Эдгээр $z$, $z^*$ цэгүүдийг дайрсан $\overline{\mathbb{C}}$ дээрх аливаа тойрог $\Gamma$ тойрогтой ортогональ байвал \\
\begin{minipage}[b]{0.5\linewidth}
 \includegraphics[width=5cm]{Zurag/28.jpg}
 \captionof{figure}{}
\end{minipage}
\hfill
\begin{minipage}[b]{0.5\linewidth}
    \includegraphics[width=5cm]{Zurag/29.jpg}
 \captionof{figure}{}
        \label{ref hiih ner}
\end{minipage}
\end{dfn}
Эдгээр тодорхойлолтын чацуу чанартай нь зураг 28-аас илэрхий. Хэрэв $z$, $z^*$ цэгүүд а) тодорхойлолт ёсоор $\Gamma$-тай тэгш хэмтэй бөгөөд $\gamma$ нь тэдгээрийг дайрсан аливаа тойрог байвал $z_0$ цэгээс $\gamma$-тойрогт татсан шүргэгчийн уртын квадрат нь элементар геометрийн теорем ёсоор огтлогч $|z_0-z^*|$ ба түүний гаднах хэсэг $|z_0-z|$-ийн үржвэртэй тэнцүү.

Иймд $z_0$ цэгээс $\gamma$-д татсан шүргэгч нь $\Gamma$-ийн радиус байх тул энэ тойргууд ортогональ. (Хэрэв $\gamma$ нь шулуун ч байсан тэр нь $z_0$-ийг дайрах учир $\Gamma$-тай ортогональ). Урвуугаар, $\overline{\mathbb{C}}$ дээр орших $z$ ба $z^*$ цэгүүдийг дайрсан аливаа $\gamma$-тойрог (тухайн тохиолдлын $zz^*$-шулуун) $\Gamma$-д ортогональ байвал нэгдүгээрт, $z$ ба $z^*$ цэгүүд $z_0$-оос эхлэх цацраг дээр орших бөгөөд хоёрдугаарт, тэдгээрээс $z_0$ хүртэлх зайнуудын үржвэр $R^2$-тай тэнцүү (элементар геометрийн мөнхүү теорем ёсоор) иймд $z$ ба $z^*$ цэгүүд  $\Gamma$-тай тэгш хэмтэй байна.

Тодорхойлолт б)-ийн нэг ашиг нь $\Gamma$ тойргийг $\overline{\mathbb{C}}$ дээр дэлгэрүүлэн авч болоход оршино: $\Gamma$ нь шулуун гэвэл тэр нь тэгш хэмийн ердийн ойлголтод хүргэнэ. Тодорхойлолт а)- ийн ашиг нь тэгш хэмт цэгүүдийг холбосон хялбар томьёо бичихэд тусална:
$$arg(z^*-z_0)=arg(z-z_0) \ \ \text{ба} \ \ |z-z_0|\cdot |z^*-z_0|=R^2$$
нөхцөлүүдээс
\begin{eqnarray} \label{7.13}
z^*=z_0+\frac{R^2}{\overline{ z-z_0}}
\end{eqnarray}
томьёо гарна. 

Зураг 29, тэгш хэмт цэгүүдийг хэрхэн байгуулахыг заана: $z$ нь $\Gamma$ дотор байвал $z$-ээс $z_0z$ цацрагт перпендикуляр татаж тойрогтой огтлолцох $T$ цэг олж $T$-ээс $\Gamma$-д татсан шүргэгчийн цацрагтай огтлолцоход $z^*$ цэг олдоно.

Хэрэв $z$ нь $\Gamma$-ийн гадаа байвал байгуулалт урвуу эрэмбээр хийгдэнэ (Баталгаа нь $z_0zT$ ба $z_0Tz^*$ гурвалжнуудын төсөөгөөс үүдэн гарна).

$z\in \overline{\mathbb{C}}$ цэг бүрийг $\Gamma$ тойргийн түүнтэй тэгш хэмтэй $z^*$ цэгт буулгах $z\rightarrow z^*$ буулгалтыг энэ тойрогтой харьцуулсан тэгш хэм буюу инверси \index{Инверси} гэж нэрлэнэ.

(\ref{7.13})-аас харахад $\overline{\mathbb{C}}$ дээрх тойрогтой харьцуулсан инверси нь бутархай шугаман функцийн хосмог функцээр илэрхийлэгдэнэ.

Одоо тэгш хэмт цэгүүдийн хадгалагдах чанарыг харуулсан дараах теоремтой танилцъя.

\begin{thm}
Аливаа бутархай шугаман $L$ буулгалт $\overline{\mathbb{C}}$ дээрх ямар нэгэн $\Gamma$ тойрогтой харьцуулсан тэгш хэмт дурын $z$ ба $z^*$ цэгүүдийг энэ тойргийн $L(\Gamma)$ дүртэй харьцуулсан тэгш хэмт $w$ ба $w^*$ цэгүүдэд хувиргана (тэгш хэмт цэгүүдийн хадгалагдах чанар).
\end{thm}
\begin{prf}
$z$ ба $z^*$ цэгүүдийг дайрах $\overline{\mathbb{C}}$ дээрх $\Gamma$ тойрогт ортогональ бүхий л тойргуудын бүлийг $\{\gamma\}$ гэж тэмдэглэе.

Теорем 7.1 ёсоор $\gamma$ тойргууд $\overline{\mathbb{C}}$ дээрх $L(\gamma)$ тойргуудад буух бөгөөд харин $L$-ийн конформлог чанараар бүх $L(\gamma)$ тойргууд $L(\Gamma)$-д ортогональ байна. Эндээс бүх $L(\gamma)$ тойргуудын дайрах $w$ ба $w^*$ цэгүүд $L(\Gamma)$-тай харьцуулсан тэгш хэмт байх нь тогтоогдов. 
\hfill $\blacksquare$\end{prf}
\section{Бутархай шугаман изоморфизм ба автоморфизм}
Бутархай шугаман буулгалтын
\begin{eqnarray}
L:w=\frac{az+b}{cz+d}
\end{eqnarray}
томьёонд $a$, $b$, $c$, $d$ гэсэн дөрвөн комплекс коэффициент орох боловч түүнийг гурван комплекс параметрээс хамаарах мэт үндэс байна. Учир нь бутархайн хүртвэр ба хуваарийг $0$-ээс ялгаатай аль $1$ коэффициентод хувааж болно. Үүнийг ашигласны дүнд дараах теорем батлагдана.
\begin{thm}
Аливаа ялгаатай $3$-н $z_1$, $z_2$, $z_3\in\mathbb{C}$ цэг ба мөнхүү ялгаатай $3$-н $w_1, w_2, w_3\in\mathbb{C}$ цэг авахад $L(z_k)=w_k$, $k=1,2,3$ байх бутархай шугаман $L$ буулгалт олдох ба тэр нь цор ганц.
\end{thm}
\begin{prf}
Тийм $L$ буулгалт оршин байх нь хялбархан батлагдана: $z_k$ ба $w_k(k=1,2,3)$ цэгүүдийг харгалзан $\zeta$ хавтгайн $0,\infty,1$ цэгүүдэд буулгах бутархай шугаман $L_1$ ба $L_2$ буулгалт байгуулъя:
\begin{eqnarray} \label{7.15}
L_1:\zeta=\frac{z-z_1}{z-z_2}\cdot \frac{z_3-z_2}{z_3-z_1}, \quad
L_2:\zeta=\frac{w-w_1}{w-w_2}\cdot \frac{w_3-w_2}{w_3-w_1}
\end{eqnarray}
гэвэл 
\begin{eqnarray} \label{7.16}
L=L^{-1}_2 \cdot L_1
\end{eqnarray} буулгалт
\begin{eqnarray} \label{7.17}
\frac{z-z_1}{z-z_2}\cdot \frac{z_3-z_2}{z_3-z_1}=\frac{w-w_1}{w-w_2}\cdot \frac{w_3-w_2}{w_3-w_1}
\end{eqnarray}
шүтэлцэл $w=w(z)$ функц тодорхойлох ба тэр нь бидний эрж буй функц. Учир нь (\ref{7.17})-бутархай шугаман функц байхын хамт $z_k$-ийг $w_k$-д $(k=1,2,3)$ буулгах нь илэрхий.

Одоо буулгалтын цор ганц байхыг баталъя.

$\lambda$ нь $\lambda(z_k)=w_k \quad (k=1,2,3)$ байх ямар нэг бутархай шугаман буулгалт байг. (\ref{7.15}) томьёогоор тодорхойлогдох $L_1$ ба $L_2$-ийн хувьд $\mu =L_1\circ \lambda\circ L_1^{-1}$ буулгалт авбал тэр нь $0,\infty, 1$ цэгүүдийг үл хөдөлгөх нь ойлгомжтой. $\mu (\infty)=\infty$ нөхцөлөөс $\mu$ нь бүхэл шугаман функц: $\mu(\zeta)=\alpha\zeta +\beta$, $\mu(0)=0$ нөхцөлөөс $\beta =0$ ба $\mu(1)=1$ нөхцөлөөс $\alpha =1$ . Иймд $\mu\zeta =\zeta$ бөгөөд өөрөөр хэлбэл $L_2\cdot \lambda \cdot L_1^{-1}=E$ болох тул $\lambda =L_2^{-1}\cdot L_1$ буюу (\ref{7.16}) ёсоор $\lambda =L.$
\hfill $\blacksquare$\end{prf}
\begin{hnt}
$z_k$ ба $w_k$ цэг бүр (\ref{7.17}) шүтэлцэлд хоёр удаа орохдоо нэг удаа хүртвэрт, нөгөө удаа хуваарьт байх тул $z_k$ буюу $z_k$-ийн аль нэг нь $\infty$ байхад (\ref{7.17}) нь хүчинтэй байна. Бутархайн хүртвэр ба хуваарийн хаана нь $\infty$ тавигдана түүнийг 1-ээр солих хэрэгтэй. Жишээ нь $z_1=w_2=\infty$ байвал 
$$\frac{1}{z-z_2}\cdot \frac{z_3-z_2}{1}=\frac{w-w_1}{w-w_2}\cdot \frac{1}{1}$$

(\ref{7.17}) тэгшитгэлээс бутархай шугаман хувиргалтын бас нэг ерөнхий чухал чанар батлагдана. $a$, $b$, $c$, $d$ нь хоорондоо ялгаатай (төгсгөлгүй) дурын комплекс тоонууд байг.
$$\frac{c-a}{c-b}:\frac{d-a}{d-b}$$
гэсэн харьцаануудын харьцааг $a$, $b$, $c$, $d$ дөрвөн тооны буюу цэгийн (ангармоник буюу) давхар харьцаа \index{Ангармоник буюу) давхар харьцаа} гэж нэрлэнэ. Энэ харьцааг ($a$, $b$, $c$, $d$) гэж тэмдэглэх бөгөөд
$$(a,b,c,d)=\frac{c-a}{c-b}:\frac{d-a}{d-b}$$
юм.

Давхар харьцааны дээрх тодорхойлолтыг эдгээр $4$ цэгийн аль нэг нь $\infty$ байх тохиолдолд өргөтгөж болно. Дотор нь төгсгөлгүй алсын цэг багтах $4$ цэгийн давхар харьцаа гэж $3$ нь өгөгдсөн цэгүүдтэй таарах, $4$ дэх нь $\infty$ рүү тэмүүлэх хязгаарын утгыг нэрлэнэ. Энэ тодорхойлолт ёсоор
$$(\infty, b,c,d)=\frac{1}{c-b}:\frac{1}{d-b}$$
$$(a,\infty, c, d)=(c-a):(d-a)$$
$$(a,b,\infty, d)=1:\frac{d-a}{d-b}$$
$$(a,b,c,\infty) =\frac{c-a}{c-b}$$
байна.

Бутархай шугаман $w=w(z)$ функц хоорондоо ялгаатай $a$, $b$, $c$, $d$ $4$ цэгийг $A$, $B$, $C$, $D$ $4$ цэгт буулгадаг. 
\end{hnt}
\begin{thm}
Аливаа бутархай шугаман буулгалт $4$ цэгийн давхар харьцааг үл өөрчилнө:
$$(a,b,c,d)=(A,B,C,D)$$
Өөрөөр хэлбэл, $4$ цэгийн давхар харьцаа нь бутархай шугаман буулгалтын инвариант чанар юм. Теорем 7.1 болон тойргийн чанарыг үндэслэн $\overline{\mathbb{C}}$ дээрх аливаа $\Gamma$ тойргийг бутархай шугаман буулгалтаар аливаа өөр  $\Gamma^*$ тойрогт хувиргаж болно (Ингэхдээ $\Gamma$ тойргийн $3$ цэгийг $\Gamma^*$ тойргийн $3$ цэгт буулгаад тойргийн чанар ашиглах явдал л хангалттай). Топологи ухааралтын үүднээс $\Gamma$-аар зааглагдсан $B$ дугуй $\Gamma^*$-аар зааглагдсан хоёр дугуйн нэгэнд (чухамхүү алинд нь буухыг ямар нэгэн $z_0\in B$ цэг хаана бууж байгаагаар мэднэ) буух нь мэдээж.

Эндээс ямар нэгэн $B\subset \overline{\mathbb{C}}$ дугуй бутархай шугаман хувиргалтаар ямар нэгэн өөр  $B^*\subset \overline{\mathbb{C}}$ дугуйд буух нь ч ойлгомжтой болов.

$D$ мужийг $D^*$ мужид буулгах бутархай шугаман буулгалтыг бутархай шугаман изоморфизм гэх ба $D$ ба $D^*$ мужуудын хувьд тийм изоморфизм оршин байвал тэдгээрийг бутархай шугаман изоморф \index{Изоморфизм} гэнэ.

Дөнгөж саяхан нотломжлон хэлсэн өгүүлбэрийг теорем байдлаар томьёолж болно.
\end{thm}
\begin{thm}
$\overline{\mathbb{C}}$ дахь аливаа $2$ дугуй бутархай шугаман изоморф.

Одоо жишээ болгож дээд хагас хавтгай $(\im z>0)$-ийг нэгж $(|w|<1)$ дугуйд буулгах бүхий л тийм изоморфизм олъё. Үүний тулд $a$, $\im a>0$ цэг бэхлэн авч түүнийг $|w|<1$ дугуйн төвд буулгасан гэж бодъё. Тэгвэл $a$-тай бодит тэнхлэгийн хувьд тэгш хэмтэй $\overline{a}$ цэг өмнөх зүйлийн теорем 7.5 ёсоор $|w|=1$ тойргийн хувьд $w=0$ цэгтэй тэгш хэмтэй $w=\infty$ цэгт бууна. Иймд эдгээр цэгийг $0$ ба $\infty$-д харгалзан буулгах бутархай шугаман (тогтмол $k$ үржигдэхүүний нарийвчлалтайгаар) буулгалт
$$w=k\cdot \frac{z-a}{z-\overline{a}}$$
дүрстэй тавигдана.

Бодит $z=x$ үед $|z-a|=|z-\overline{a}|$ тэнхлэг нэгж тойрогт буухын тулд $|k|=1$ гэж авах хэрэгтэй. Өөрөөр хэлбэл $k=e^{i\theta}$. Иймд дээд хагас хавтгай $(\im z>0)$ нэгж дугуй $(|w|<1)$-д буулгах бутархай шугаман изоморфизмүүд 
\begin{eqnarray} \label{7.18}
w=e^{i\theta}\cdot \frac{z-a}{z-\overline{a}}
\end{eqnarray}
томьёогоор тодорхойлогдоно. ($a$ дээд хагас хавтгайн аливаа цэг, $\theta$-бодит тоо)

Мужийг өөрт нь буулгадаг бутархай шугаман изоморфизмыг бутархай шугаман автоморфизм\index{Автоморфизм} гэнэ. Ямар нэгэн бутархай шугаман бүх автоморфизмүүдийн бөөгнөрөл бутархай шугаман буулгалтуудын $\wedge$-бүлгийн дэд бүлэг үүсгэх нь ойлгомжтой.

$\overline{\mathbb{C}}$-г $\overline{\mathbb{C}}$-рүү буулгах бүх бутархай шугаман автоморфизмүүдийн бөөгнөрөл $\wedge$ бүлэгтэй давхацна. $\mathbb{C}$-г $\mathbb{C}$-рүү буулгах бутархай шугаман автоморфизмүүдийн бөөгнөрөл $z\rightarrow az+b$, $a\neq 0$ гэж буулгах бүхэл шугаман хувиргалтуудын дэд $\wedge_0$ бүлэгтэй давхцах нь ч бас мэдээж. Одоо нэгж дугуйн автоморфизмүүдийн бүлгийг бодож гаргая. Тойргийн төв $w=0$ цэгт буух $a$, $|a|<1$ цэгийг бэхлэн авъя. $a$-цэгтэй нэгж тойргийн $\{|w|<1\}$ хувьд тэгш хэмтэй $a^*=\dfrac{1}{\overline{a}}$ цэг $w=\infty$ цэгт буух учиртай. Иймд бидний эрж буй буулгалт
$$w=k\cdot \frac{z-a}{z-\frac{1}{\overline{a}}}=k_1\cdot \frac{z-a}{1-\overline{a}\cdot z}$$
дүрстэй ($k$, $k_1$ нь тогтмол) бичигдэнэ. $z=1$ цэг нэгж тойргийн цэгт буух тул $|k_1|\cdot \left|\frac{1-a}{1-\overline{a}}\right|=|k_1|=1$ болно. Өөрөөр хэлбэл, $k_1=e^{i\theta}$, $\theta$-бодит тоо. Иймд нэгж дугуйн автоморфизм
\begin{eqnarray}
w=e^{i\theta}\frac{z-a}{1-\overline{a}z}
\end{eqnarray} болой.
\end{thm}
\section{Зарим рациональ функц}
\subsection{Зэрэгт функц\index{Зэрэгт функц}}

Натурал $n$ тооны хувьд $\mathbb{C}$ дээр голоморф
\begin{eqnarray} \label{7.20}
w=z^n
\end{eqnarray}
функцийг зэрэгт функц гэнэ. Түүний уламжлал
$\dfrac{dw}{dz}=nz^{n-1}$, $(n>1)$ нь  $z\neq 0$ үед $0$-ээс ялгаатай тул
$z\in\mathbb{C}\setminus \{0\}$ цэг бүр дээр конформ буулгалт өгнө.

Зэрэгт функц $0$ цэгийн орчинд ямар төрхтэй байхыг авч үзье. Үүний тулд туйлын координат:
$$z=re^{i\varphi}, \qquad w=\rho e^{i\theta}$$
ашиглавал (\ref{7.20}) тэнцлээр
\begin{eqnarray} \label{7.21}
\begin{cases}
 \rho=r^n \\
 \theta =n\varphi
\end{cases}
\end{eqnarray}
болно. (\ref{7.21})-ийн $2$дугаархаас өнцөг $0$ цэг дээр хэвээр хадгалагдахгүй, харин $n$ дахин өсөх нь тодорхой байна. Өнцгийн хадгалагдах чанаргүй байх нь $\infty$ дээр ч мөн адил. Учир нь $\dfrac{1}{w\left(\frac{1}{z}\right)}$ функц  $z=0$ цэгийн орчинд өгөгдсөн функцтэй давхацна.

$0$ ба $\infty$ цэг нь $z=\sqrt[n]{w}$ функцийн салаалалтын цэгүүд болдог. (\ref{7.21})-ээс харахад $\circled{z}$ хавтгайн $0$ цэгт төвтэй тойрог \\ $(r=const)$ $\circled{w}$ хавтгай $(\rho =consrt)$-т, $0$ цэгээс гарсан цацраг ($\varphi=const$) мөн цацраг ($\theta =const$)-т буух нь тодорхой байна.

$\circled{z}$ хавтгай дээр бодит эерэг тэнхлэг ба $0$ цэгээс гарсан цацрагаар үүсгэгдсэн $\dfrac{2\pi}{n}$ өнцөгт харгалзах сектор (зураг 30) авъя. Энэ сектор (\ref{7.20}) функцээр $\circled{w}$ хавтгайн эерэг хагас тэнхлэгийн дагуу хагалбар бүхий бүх хавтгай (зураг 31)-д бууна. Учир нь $\varphi =0$ бол $\theta =0$, $\varphi =\dfrac{2\pi}{n}$ бол $\theta =2\pi $ өнцөг харгалзана. \\
\begin{minipage}[b]{0.5\linewidth}
 \includegraphics[width=5cm]{Zurag/30.PNG}
 \captionof{figure}{}
\end{minipage}
\hfill
\begin{minipage}[b]{0.5\linewidth}
    \includegraphics[width=5cm]{Zurag/31.PNG}
 \captionof{figure}{}
        \label{ref hiih ner}
\end{minipage}
Өөрөөр хэлбэл (\ref{7.20}) функц $D=\left\{0<\arg  z<\dfrac{2 \pi}{n}\right\}$ секторийг
$D^*=\{0< \arg  w<2 \pi \}$ мужид гомеоморф буулгана.
\subsection{Жуковскийн функц.} \index{Жуковскийн функц}
$\mathbb{C}\setminus \{0 \}$ мужид голоморф
\begin{eqnarray} \label{7.22}
w=\frac{1}{2}\left(z+\frac{1}{z}\right)
\end{eqnarray}
рационал функцийг Жуковскийн функц гэнэ. Түүний хувьд
$$\frac{dw}{dz}=\frac{1}{2}\left(1-\frac{1}{z^2}\right)$$ 
бөгөөд $z\neq \pm 1$ цэг дээр $0$ ялгаатай тул (\ref{7.22}) функц $z\neq 0$, $\pm 1$ цэг бүр дээр конформ буулгалт өгнө. 

Ямар нэг $D$ муж дээр манай функцийн нэг хуудаст байх нөхцөлийг тогтооё. Энэ функц $z_1$, $z_2$ цэгийг нэгэн цэгт буулгалаа гэвэл
$$z_1+\frac{1}{z_1}-\left(z_2+\frac{1}{z_2}\right) =(z_1-z_2)\left(1-\frac{1}{z_1z_2}\right)=0$$
байх учиртай бөгөөд $z_1\neq z_2$ бол $z_1z_2=1$ байх аж. Иймд
Жуковскийн функц ямар нэг $D$ муж дээр нэг хуудаст байх зайлшгүй бөгөөд хүрэлцээтэй нөхцөл нь
\begin{eqnarray}
z_1\cdot z_2=1
\end{eqnarray}
гэсэн ямар ч хос $D$-мужид агуулагдахгүй байх явдал мөн.

Нэг хуудаст мужийн жишээ нь нэгж дугуйн гаднах хэсэг $D=\{z\in\overline{\mathbb{C}}\}:|z|>1$ юм. (\ref{7.22}) буулгалтыг ил харуулахын үүднээс $z=re^{i\varphi}$, $w=u+iv$ гэж авбал (\ref{7.22}) 
\begin{eqnarray}
u=\frac{1}{2}\left(r+\frac{1}{r}\right) \cos \varphi, \qquad v=\frac{1}{2}\left(r-\frac{1}{r}\right) \sin \varphi
\end{eqnarray}
дүрстэй бичигдэнэ.

Эдгээр шүтэлцлээс Жуковскийн функц $\{|z_1|<r_0\}, r_0>1$ байх аливаа тойрог $a_{r_0}=\dfrac{1}{2}\left(r_0+\dfrac{1}{r_0}\right)$ ба $b_{r_0}=\dfrac{1}{2}\left(r_0-\dfrac{1}{r_0}\right)$ гэсэн хагас тэнхлэгүүд бүхий $\pm 1$ цэгүүдэд фокустай ($\forall r_0$- ийн хувьд $a^2_{r0}-b_{r0}^2=1$ тул) эллипст буух нь харагдаж байна (зураг 23 үз). \\
\begin{center}
    \includegraphics[width=7cm]{Zurag/32.jpg}
    \captionof{figure}{}
        \label{}
\end{center}
$\lim\limits_{r_0\rightarrow 1} b_{r_0}=0$ байх тул эллипс $[-1;1]\subset \mathbb{R}$ хэрчим рүү унтрана. $a_{r_0}-b_{r_0}$ ялгавар $\dfrac{1}{r_0}$ болж их хэмжээний $r_0$-ийн хувьд эллипс нь тойргоос багаар ялгагдана.

$\{\varphi =\varphi_0$, $1<r<\infty\}$ цацраг $\dfrac{u^2}{\cos\varphi}-\dfrac{v^2}{\sin^2\varphi}=1$ гиперболын хэсэгт буух ба мөн  $\pm 1$ цэгүүдэд фокустай ($\forall \varphi_0$-ийн хувьд $\cos^2_0+\sin^2\varphi_0=1$ тул) зураг 32-г үз.

Конформ чанараар эдгээр гиперболын бүл уул эллипсүүдийн бүлтэй ортогональ.

Дээр дурдсанаас харахад Жуковскийн функц нэгж тойргийн гаднах (төгсгөлгүй алсын цэгийг агуулах) хэсгийг бодит тэнхлэгийн $[-1;1]$ хэрчмийн гаднах хэсэгт харилцан нэгэн утгатай, конформ буулгах нь илэрхий болов.

$\pm 1$ цэгүүдэд (\ref{7.22}) буулгалт конформ биш. Үүнд үнэмшихийн тулд Жуковскийн функцийг
\begin{eqnarray} \label{7.25}
\frac{w-1}{w+1}=\left(\frac{z-1}{z+1}\right)^2
\end{eqnarray}
дүрстэй бичье. (\ref{7.25})-ийг
\begin{eqnarray} \label{7.26}
\zeta =\frac{z-1}{z+1}, \quad w=\zeta^2,\quad w=\frac{1+w}{1-w}
\end{eqnarray}
байдалтай бичиж (сүүлчийнх нь $w=\dfrac{w-1}{w+1}$-ийн урвуу) болно. (\ref{7.26})-ийн эхнийх ба гуравдах нь бутархай шугаман тул нийт $\overline{\mathbb{C}}$ дээр конформ буулгалт өгнө. Харин $w=\zeta ^ 2$ буулгалт $\zeta =0$, $\infty$ цэгүүдэд ($z=\pm 1$) дээрх өнцгийг хоёр дахин өсгөнө.

(\ref{7.25}) задаргаанаас Жуковскийн функц зураг 24-т үзүүлсэнчлэн ($\pm 1$ цэгүүдийг дайрахдаа бодит тэнхлэгтэй $\alpha$ өнцөг үүсгэх) $\gamma$-тойргийн гаднах хэсгий тойргийн нумын ($\pm 1$ цэгүүдэд төгсгөл бүхий бодит тэнхлэгтэй $w=1$ цэг дээр $2\alpha$ өнцөг үүсгэх) гаднах хэсэгт нэг хуудаст бөгөөд конформ буулгах нь илт байна.

Мөн энэ буулгалтаар $\gamma$ тойрогтой $\pm 1$ цэгүүдийн аль нэг дээр гадаад шүргэх $\delta$-тойргийг нисэх онгоцны далавчны хөндлөн огтлол санагдуулах өвөрмөц битүү $\delta '$ муруйд буух (зураг 33-г үз) нь харагдаж байна. \\
\begin{center}
    \includegraphics[width=8cm]{Zurag/33.PNG}
    \captionof{figure}{}
        \label{}
\end{center}
Энэ санамж далавчны аэродинамик тооцооны анхны аргыг боловсруулахад нь Жуковскийд тусалжээ.
\section{Зарим трансцендент функц, түүний буулгалт}
\subsection{Илтгэгч функц \index{Илтгэгч функц, (экспонент)} тодорхойлох нь.}
$\dfrac{df}{dz}=f(z)$ ба $f(0)=u(0,0)+iv(0,0)=1$ нөхцөлүүд хангах голоморф функц олъё. 
$$\frac{df}{dz}=\frac{\partial u}{\partial x}+i\frac{\partial v}{\partial x}=\frac{\partial v}{\partial y}-i\frac{\partial u}{\partial y}$$ 
байх тул манай дифференциал тэгшитгэл
$$\frac{\partial u}{\partial x}+i\frac{\partial v}{\partial x}=\frac{\partial u}{\partial y}-i\frac{\partial u}{\partial y}= u+iv$$
дүрстэй болж эндээс
$$\frac{\partial u}{\partial x}=\frac{\partial v}{\partial y}=u,\qquad \frac{\partial v}{\partial x}=-\frac{\partial u}{\partial y}= v$$
бөгөөд $\dfrac{\partial u}{\partial x}=u$ ба $u(0,0)=1$ нөхцөлөөс $u(x,y)=e^x\lambda (y)$ болно. (Энд $\lambda (y)$ нь $\lambda(0)=1$ нөхцөл хангах дифференциалчлагдах функц.) Яг энэчлэн $\dfrac{\partial v}{\partial x}=v$ ба $u(0,0)=0$ нөхцөлөөс $v(x,y)=e^x\cdot \mu (y)$ болно. (энд $\mu(y)$ нь $\mu(0)=0$ нөхцөл хангах дифференциалчлагдах функц.) Иймд
$$\frac{\partial v}{\partial y}=e^x \mu'(y)=u=e^x\lambda(y), \qquad -\frac{\partial u}{\partial y}=-e^x \lambda'(y)=v=e^x\mu(y). $$
Өөрөөр хэлбэл
$$\mu'(y)=\lambda (y), \qquad -\lambda '(y)=\mu (y).$$
Эндээс $\lambda(y)$ ба $\mu(y)$ функцүүд ерөнхий шийд нь
$$\varphi (y)=c_1\cos y+c_2\sin y$$
байх хоёрдугаар эрэмбийн шугаман дифференциал
$$\frac{d^2\varphi(y)}{dy^2}+\varphi (y)=0$$
тэгшитгэлийг хангах нь мэдээж юм. $\mu (y)$-ийн анхны нөхцөлийг тооцвол $\mu(y)=C\sin y$ болж улмаар $\lambda (y)=\mu '(y)=C\cos y$ бөгөөд эндээс $C=1$-ийн анхны нөхцөл ёсоор $\lambda(y)=\cos y$, $\mu (y)=\sin y$ болно.

Иймд  $\lambda (y)=\cos y$, $\mu(y)=\sin y$ эндээс
$$f(z)=u(x,y)+iv(x,y)=e^x(\cos y +i\sin y)$$ болов. 

Бид дурдсан нөхцөлүүд хангах цор ганц голоморф функц оллоо.

Энэ функцийг комплекс $z$ хувьсагчийн экспонент буюу илтгэгч функц гэх ба $\exp z$ гэж тэмдэглэнэ:
\begin{eqnarray} \label{7.27}
\exp z=e^z=e^x(\cos y +i\sin y)
\end{eqnarray}
Эндээс 
\begin{eqnarray}
|e^z|=e^{\re z}, \qquad \arg  e^z=\im z
\end{eqnarray}
бөгөөд (\ref{7.27})-д $z=0$ байвал Эйлерийн\index{Эйлерийн томьёо }
\begin{eqnarray}
e^{iy}=\cos y +i\sin y
\end{eqnarray}
томьёо батлагдлаа.

I бүлэгт $e^{iy}$ гэдгийг баруун $pq$ илэрхийллийн товчилсон тэмдэглэгээ төдий хэрэглэсэн бол одоо түүнийг $e$-ийн хуурмаг зэрэг гэж ойлгох боломж бидэнд олдлоо.

Илтгэгч функцийн үндсэн чанаруудыг тоочъё:
\begin{enumerate}
    \item [$1^\circ.$] $e^z$ нь нийт $\mathbb{C}$ дээр голоморф функц.
    \item [$2^\circ.$] $e^z$ функц $\forall z\in\mathbb{C}$ дээр конформ буулгалт өгнө. Учир нь $|e^z|=e^x>0$ тул
    $$(e^z)'=e^z\neq 0$$
     \item [$3^\circ.$] $e^z$ функцийн хувьд нийлбэрийн
     \begin{eqnarray} \label{7.30}
     e^{z_1+z_2}=e^{z_1}\cdot e^{z_2}
     \end{eqnarray}
     теорем хүчинтэй.
     
Учир нь $z^k=x_k+iy_k \quad (k=1,2)$ гэвэл
\begin{align*}
    e^{x_1}(\cos y_1+i\sin y_1)&\cdot e^{x_2}(\cos y_2+i\sin y_2)\\
    &=e^{x_1+x_2}[\cos(y_1+y_2)+i\sin (y_1+y_2)]
\end{align*}
Иймд комплекс $z_1$ ба $z_2$ тоонуудын нийлбэрт тэдгээрийн дүрүүдийн $e^{z_1}\cdot e^{z_2}$ үржвэр харгалзана. Өөрөөр хэлбэл илтгэгч $e^z$ комплекс тоон талбарын аддитив бүлгийг энэ талбарын мультипликатив бүлэгт
\begin{eqnarray}
z\rightarrow e^z:z_1+z_2\rightarrow e^{z_1}\cdot e^{z_2}
\end{eqnarray}
гэж буулгана.
 \item [$4^\circ.$]  $e^z$ нь $2\pi i$ гэсэн хуурмаг үет функц. Учир нь Эйлерийн томъёо ёсоор $e^{2\pi i}=\cos 2\pi +i\sin 2\pi =1$ болж (\ref{7.30}) ёсоор
 $$e^{z+2\pi i}=e^z\cdot e^{2\pi i}=e^z$$
 юм. Нөгөө талаар $e^{z+T}=e^z$ гэвэл $e^{-z}$-ээр үржүүлсний дараа $e^T=1$ болох ба эндээс $T=T_1+iT_2$ гэвэл $e^{T_1}(\cos T_2+i\sin T_2)=1$
\end{enumerate}
 Харин $e^{T_1}=1$ тул $T_1=0$, $\cos T_2=1$ тул $1\sin T_2=0$ болж улмаар $T_2=2n \pi$, $n$-бүхэл тоо. Иймд $T=2n \pi i$. $2\pi i$ нь үндсэн үе.
    
Эндээс тунгаан үзвэл $w=e^z$ буулгалтын ямар нэг $D$ мужид нэг хуудаст байх зайлшгүй бөгөөд хүрэлцээтэй нөхцөл нь
 \begin{eqnarray} \label{7.32}
    z_1-z_2=2n \pi i \quad (n=\pm 1, \pm 2, \ldots)
    \end{eqnarray}
    шүтэлцлээр холбогдох нэг ч хосыг энэ муж агуулахгүй байх явдал мөн.
    
Энэ нөхцөлийг хангах мужийн жишээ нь $\{0<\im z<2\pi \}$ зурвас юм.
\begin{center}
    \includegraphics[width=8cm]{Zurag/34.jpg}
    \captionof{figure}{}
        \label{}
\end{center}
$z=x+iy$ ба $w=\rho e^{i\theta}$ гэвэл (\ref{7.32}) ёсоор $w=e^z$ буулгалтыг туйлын координатад 
\begin{eqnarray}
\rho =e^x, \quad \theta =y
\end{eqnarray}
дүрстэй бичье. Эндээс энэ буулгалтыг $\{y=y_0\}$ шулууныг \\ $\{\theta =y_0\}$ цацрагт, харин $x=x_0$, $0<y<2\pi $ хэрчмийг цэгээр цоорхойт тойрогт тус тус буулгана (зураг 25). $\{0<y<2\pi \}$ зурвас бодит хагас тэнхлэгийн дагуу хагалбар бүхий $W$ хавтгайд буух нь тогтоогдлоо. 2 дахин нарийн $\{0<y<\pi\}$ зурвас $\{\im z>0\}$ дээд хагас хавтгайд бууна.

\subsection{Тригонометр функцүүд.} \index{Тригонометр функцүүд.} Эйлерийн томьёоноос дурын $x$-бүхний хувьд $e^{ix}=\cos x+i\sin x$, $e^{-ix}=\cos x-i\sin x$ ба эндээс
$$\cos x=  \frac{e^{ix}+e^{-ix}}{2}, \qquad \sin x  =\frac{e^{ix}-e^{-ix}}{2i}.$$

Энэ томьёонуудыг комплекс хавтгай дээр $\forall z\in\mathbb{C}$-ийн хувьд тодорхойлогдох косинус ба синусийн голоморф үргэлжлэл байдлаар 
\begin{eqnarray} \label{7.34}
\cos z =\frac{e^{iz}+e^{-iz}}{2}, \qquad \sin z=\frac{e^{iz}-e^{-iz}}{2i}
\end{eqnarray}
ашиглаж болно. 

(Баруун гар дахь илэрхийллүүдээс $\mathbb{C}$ дээр голоморф байх нь илэрхий).

Эдгээр функц илтгэгч функцээр тодорхойлогдох тул түүний бүхий л чанар хадгалагдана. Хоёулаа үет функц, үндсэн үе нь $2\pi $ (илтгэгч функц $2\pi i$ үетэй (\ref{7.34}) томьёонд $z$-ийн үржигдэхүүн нь $i$ байх учраас). Косинус-тэгш, синус-сондгой функц. Эдгээр функцүүдийн хувьд уламжлалын ердийн томьёонууд хүчинтэй:
$$(\cos z)'=i\frac{e^{iz}-e^{-iz}}{2}=-\sin z$$
энэчлэн $(\sin z)'=\cos z$. Мөн тригонометр
$$\cos^2 z+\sin^2 z=1, \qquad \cos z=\sin\left(z+\frac{\pi}{2}\right)$$
шүтэлцлүүд, нийлбэрийн теорем гэх мэтчилэн хүчинтэй хэвээр байхыг (\ref{7.34}) томьёоноос төвөггүй тогтоогдоно.

$\forall z\in\mathbb{C}$-ийн хувьд
\begin{eqnarray} \label{7.35}
\ch z=\frac{e^z+e^{-z}}{2}, \qquad \sh z=\frac{e^z-e^{-z}}{2i}
\end{eqnarray}
гэж тодорхойлогдох гипербол функцүүд тригонометр функцүүдтэй нягт холбоотой. 

Энэ холбоо (\ref{7.34}) ба (\ref{7.35}) томьёонуудыг жиших замаар гарах 
\begin{eqnarray} \label{7.36}
\begin{split}
    &\ch z=\cos iz, \qquad \sh z=-i\sin iz \\
    &\cos z=\ch iz, \qquad \sin z=-\sh iz
\end{split}
\end{eqnarray}
шүтэлцлүүдээр илэрхийлэгдэнэ.

Нийлбэрийн теорем ба (\ref{7.36}) томьёонууд ашиглан бичвэл 
\begin{align*}
    &\cos z =\cos(x+iy)=\cos x\ch y-i\sin x\sh y\\
    &\sin z= \sin(x+iy)=\sin x\ch y-i\sh y \cos x
\end{align*}
бөгөөд эндээс
\begin{eqnarray}
\begin{split}
    &|\cos z|=\sqrt{\cos^2x+\sh ^2y} \\
    &|\sin z|=\sqrt{\sin^2x+\sh^2 y}
\end{split}
\end{eqnarray}
байна. (Энд $\cos^2x+\sin^2x=1$, $\ch^2y-\sh^2y=1$ адилтгалууд ашиглав.) Энэ томьёонууд косинус ба синус функцүүдийн модулийн гадаргуу (рельеф) байгуулах боломж олгоно.\\
Жишээ болгож косинусын рельефийг\index{Рельеф косинусын} зураг 35 дээр дүрслэв. \\
\begin{center}
    \includegraphics[width=8cm]{Zurag/35.PNG}
    \captionof{figure}{}
        \label{}
\end{center}

$w=\sin z$ функц дээд хагас хавтгайд байх \\ $D=\left\{-\dfrac{\pi}{2}<x<\dfrac{\pi}{2}, \quad y=0\right\}$-хагас зурвас мужийг хэрхэн буулгахтай танилцъя. Энэ буулгалт бидний урьдын танил хялбар
$$z_1=iz,\quad z_2=e^{z_1}, \quad z_3=\frac{z_2}{i}, \quad w=\dfrac{1}{2}\left(z_3+\frac{1}{z_3}\right)$$
буулгалтууд болон задарна. 

Иймд $w=\sin z$ функц $D$ мужийг дээд хагас хавтгайд нэг хуудаст (бөгөөд конформ) буулгана. 

Ингэхдээ $\{x=x_0, \quad 0<y<\infty\}$ цацрагийг $\pm 1$ цэгт фокус бүхий гиперболын хэсэгт,  $х\left\{-\dfrac{\pi}{2}<x<\dfrac{\pi}{2}, \quad y=y_0\right\}$ хэрчмийг мөн тийм фокусууд бүхий эллипсийн хэсэгт бууна. Зураг $36$-аас харахад \\
\begin{center}
    \includegraphics[width=9cm]{Zurag/36.PNG}
    \captionof{figure}{}
        \label{}
\end{center}
хагас зурвас мужийн босоо хилүүд $w=\sin z$-функцээр бодит тэнхлэгийн модуль 1-ээс их хэсэгт буух нь тодорхой байна.

Тангенс ба котангенс функцүүд комплекс аргументийн хувьд 
\begin{eqnarray}
\tan z=\frac{\sin z}{\cos z}, \qquad \cot z=\frac{\cos z}{\sin z}
\end{eqnarray}
томьёонуудаар тодорхойлогдоно. Эдгээр нь итгэгч функцээр рационал байдалтай
\begin{eqnarray} \label{7.38}
\tan z=-i\frac{e^{iz}-e^{-iz}}{e^{iz}+e^{-iz}}, \quad \cot z=i \frac{e^{iz}+e^{-iz}}{e^{iz}-e^{-iz}}
\end{eqnarray}
илэрхийлэгдэнэ. Эдгээр функцүүд (\ref{7.38}) томьёонуудын хуваарийн $0$-тэй тэнцэх цэгүүд (тэдгээр дээр хүртвэр нь $0$-ээс ялгаатай)-ээс бусад нийт $\mathbb{C}$ дээр голоморф байна. 

Жишээ нь $\tan z$ функцийн тийм цэгүүдийг олъё. Тэдгээр дээр $\cos z=0$, өөрөөр хэлбэл $e^{iz}+e^{-iz}=0$ бөгөөд эндээс $z=(2z-1)\dfrac{\pi}{2}$, $(n=0,\pm 1, \ldots)$ болно.

Тангенс ба котангенс нь комплекс хавтгайд бодит $\pi$ үетэй функцүүд, тэдгээрийн хувьд дифференциалчлалын болон тригонометр шүтэлцлийн ердийн томьёонууд хүчинтэй хэвээрээ бөгөөд энэ бүхэн нь (\ref{7.38}) томьёонуудаас шууд гарна.

Тангенс функцийг бидний танил хялбар
\begin{eqnarray} \label{7.39}
z_1=2iz, \quad z_2=e^{z_1}, \quad w=-i\frac{z_2-1}{z_2+1}
\end{eqnarray}
функцүүдийн тусламжтайгаар түүний буулгалтыг харуулъя. $w=\tan z$ функц $\left\{-\dfrac{\pi}{4}<x<\dfrac{\pi}{4}\right\}$ зурвасыг (\ref{7.39}) ёсоор нэгж дугуйд $1$ хуудаст бөгөөд конформ буулгана.

Энэ буулгалтаар $\{x=x_0\}$ шулуунууд $\pm 1$ цэгүүдийн дайрах тойргийн нумд харин  $\left\{-\dfrac{\pi}{4}<x<\dfrac{\pi}{4}, \quad y=y_0\right\}$хэрчмүүд түүний төгсгөлийн цэгүүдийг тэгш хэмтэйгээр агуулах тойргийн нумд хувиргана.

Зарим жишээ авч үзье.
\begin{exm}
Дээд хагас $P=\{z:\im z>0\}$ хавтгайг өөрт нь $w(0)=1$,$w(i)=2i$ нөхцөлтэйгөөр буулга.
\end{exm}
\begin{sol}
2 дахь нөхцөлөөс харахад $-i$ цэг $2i$-тэй тэгш хэмтэй $-2i$ цэгт буух тул (\ref{7.17}) ёсоор 
$$\frac{w-1}{w-2i}\cdot \frac{-2i-2i}{-2i-1}=\frac{z-0}{z-i}\cdot \frac{-i-i}{-i-0}$$ буюу $w=-2\cdot \dfrac{2z+1}{z-12}.$
\end{sol}
\begin{exm}
Дээд хагас $P=\{z:\im z>0\}$хавтгайг нэгж $K=\{w:|w|<1\}$ дугуйд $w(i)=0$, $\arg  w'(i)=-\dfrac{\pi}{2}$ нөхцөлтэй буулга.
\end{exm}
\begin{sol}
(\ref{7.18}) томьёогоор 
$$w=e^{i\theta}\cdot \frac{z-i}{z+i}$$
бөгөөд $$w'(z)=e^{i\theta}\cdot \dfrac{2i}{(z+i)^2}, \quad w'(i)=e^{i\theta} \left(-\dfrac{i}{2}\right)$$ болж $\arg  w'(i)=-\dfrac{\pi}{2}$ нөхцөлөөс $\theta -\dfrac{\pi}{2}=-\dfrac{\pi}{2}$ буюу $\theta= 0$ тул $w=\dfrac{z-i}{z+i}.$
\end{sol}
\begin{exm}
$K=\{z:|z|<1\}$ тойргийг $\overline{K}=\{w:|w-i|<1\}$ тойрогт шилжүүлэх $w$ буулгалтыг $w(0)=\dfrac{i}{2}$ ба $w(1)=0$ байхаар ол.
\end{exm}
\begin{sol}
(\ref{7.13}) томьёо ёсоор $w=\dfrac{i}{2}$ цэгт тэгш хэмт $w^*$ цэг олбол
$$w^*=i+\frac{1}{\overline{\frac{i}{2}-i}}=-i$$
бөгөөд эндээс $w(\infty)=-i$ улмаар (\ref{7.17}) ёсоор
$$\frac{w-\frac{i}{2}}{w-0}\cdot \frac{-i-0}{-i-\frac{i}{2}}=\frac{z-0}{z-1}\cdot \frac{\infty -1}{\infty -0}.$$ болж 
$w=\dfrac{i-iz}{z+2}.$
\end{sol}
\begin{exm}
$D=\{z:|z|<1\}\cap \{z:|z-1|>1\}$ мужийг дээд хагас $P=\{w:\im w>0\}$ хавтгайд буулга.
\end{exm}
\begin{sol}
Тойргийн нумуудаар зааглагдсан саран $D$ мужийг дээд хагас $P$ хавтгайд буулгахын тулд эхлээд түүнийг бутархай шугаман функцээр төгсгөлгүй сектор мужид буулгана.  

Үүний тулд
$\left.\begin{array}{c}
     |z|=1  \\
     |z-1|=1
\end{array}\right\}$ системээс
$$z_1=\frac{1}{2}+\frac{\sqrt{3}}{2}i, \quad z_2=\frac{1}{2}-\frac{\sqrt{3}}{2}i$$
цэгүүдийг олж $z_2\rightarrow 0$, $z_1\rightarrow \infty$ байх бутархай шугаман $w_1$ буулгалт байгуулбал:
$$w_1=\frac{z-z_2}{z-z_1}=\frac{2z-1+\sqrt{3i}}{2z-1-\sqrt{3i}}$$
\begin{align*}
    &w_1(0)=\frac{-1+\sqrt{3i}}{-1-\sqrt{3i}}=-\frac{1}{2}-\frac{\sqrt{3}}{2}i\\
    &w_1(1)=\frac{2-1+\sqrt{3i}}{2-1-\sqrt{3i}}=-\frac{1}{2}+\frac{\sqrt{3}}{2}i
\end{align*}
цэгүүдийн дайрах цацрагууд $D$-ийн хилд буух тул эдгээрийн хоорондох өнцөг $(w_1(0)^{\wedge}w_1(1))=\dfrac{2\pi}{3}$ байх учраас $w_2=e^{-i\frac{2\pi}{3}}w_1$ болно.

Эндээс $w=w_2^{\frac{3}{2}}$ болж бидний эрж байгаа буулгалт:
$$w=-\left(\frac{2z-1+\sqrt{3}i}{2z-1-\sqrt{3}i}\right)^{\frac{3}{2}}$$
\end{sol}
\begin{center}
    \textbf{Дасгал}
\end{center}
\begin{enumerate}
    \item $K=\{z:|z|<1\}$ дугуйг дээд хагас $P=\{w:\im w>0\}$ хавтгайд буулгах бутархай шугаман функцийг $w(-1)=\infty$, $w(1)=0$, $w(i)=1$ нөхцөлөөр ол.
    \item Нэгж $K=\{z:|z|<1\}$ дугуйг өөрт нь өөрөөр хэлбэл
$\overline{K}=\{w:|w|<1\}$ буулгах функцийг $w\left(\frac{1}{2}\right)$, $\arg  w'\left(\frac{1}{2}\right)=\dfrac{\pi}{2}$ нөхцөлөөр ол.
\item $\gamma=\{z:|z-i|=3\}$ тойргийн хувьд $2+i$ цэгтэй тэгш
хэмтэй цэгийг ол.
\item Дээд хагас $P=\{z:\im z>0\}$ хавтгайг нэгж $K=\{w:|w|<1\}$ дугуйд $w(2i)=0$, $\arg  w'(2i)=0$ нөхцөлөөр буулга.
\item 	Дараах хоёр өнцөгт (саран) мужийг дээд хагас хавтгайд буулга:
\begin{enumerate}
    \item [a)] $D_1=\{z:|z|<1\}\cap \{z:|z-i|<1\}$
     \item [b)] $D_2=\{z:|z|>1\}\cap \{z:|z-i|<1\}$
\end{enumerate}
\end{enumerate}
\chapter{Голоморф функц ба төгсгөлгүй үржвэр
}
\section{Үндсэн ойлголтууд}
\begin{dfn}
$u_k\in\mathbb{C}(k=1,2,3,\ldots)$ байг. 
$$P_n=(1+u_1)(1+u_2)(1+u_3)\cdot \ldots \cdot (1+u_n)$$ гэж тодорхойлъё.
\begin{enumerate} \label{8.1}
    \item [a)] Хэрэв $\exists \lim\limits_{n\rightarrow \infty} P_n=P\in \mathbb{C}$ ба $P\neq 0$ бол 
    \begin{eqnarray}
     P=\prod\limits_{k=1}^\infty (1+u_k)
    \end{eqnarray}  төгсгөлгүй үржвэр \index{Төгсгөлгүй үржвэр} нийллээ гэж үздэг.
 \item [б)] Хэрэв дээрх хязгаар оршдоггүй эсвэл $0$ байвал сарниж байна гэж үздэг.
  \item [в)] Хэрэв $\exists \tilde{P}=\prod\limits_{k=1}^\infty (1+|u_k|)\in\mathbb{R}\setminus \{0\}$ байвал абсолют нийлдэг төгсгөлгүй үржвэр\index{Абсолют нийлдэг төгсгөлгүй үржвэр} гэж үздэг.
  \item [г)] Төгсгөлгүй үржвэр нийлдэг боловч абсолют нийлэхгүй
  $$\left( \text{жишээ нь} \left(1+\frac{1}{2}\right)\left(1-\frac{1}{2}\right) \left(1+\frac{1}{3}\right) \left(1-\frac{1}{4}\right)\cdots \right)$$
  бол нөхцөлт нийлдэг үржвэр гэдэг.
\end{enumerate}
\end{dfn}
\textbf{Дасгал 1}. Дараах үржвэр ба цуваануудын нийлэлт сарнилт нь ижил байна. 
\begin{eqnarray}
 \prod\limits_{k=1}^\infty (1+|u_k|)
\end{eqnarray}
\begin{eqnarray}
 \sum\limits_{k=1}^\infty \ln (1+|u_k|)
\end{eqnarray}
ба 
\begin{eqnarray} \label{8.4}
 \sum\limits_{k=1}^\infty |u_k|
\end{eqnarray}
\begin{cor}
(\ref{8.1})-ийн абсолют нийлэлтийг шинжихдээ (\ref{8.4}) цувааг судалж болно.

Тоон цуваанаас функцэн цуваанд шилжих алхмыг үржвэрт бас хэрэглэж болно. Тухайлбал:
\end{cor}
\begin{thm}
$u_k(z)\in H(D)(k=1,2,3,\ldots)$ байг. $\forall k\in\mathbb{N}:|u_k(z)|\leq a_k$ байдаг ба $a_1+a_2+\cdots +a_k+\cdots$ цуваа нийлдэг байг. Тэгвэл  
$$\exists P(z)=\prod\limits_{k=1}^\infty (1+u_k(z))\in H(D)$$ байна.
\end{thm}
\begin{prf}
$a_1+a_2+\cdots +a_n+\cdots$ эерэг тоон цуваа нийлэх тул $\forall z\in D$ хувьд жиших шинжүүр ёсоор $$|u_1(z)|+|u_2(z)|+\cdots +|u_k(z)|+\cdots$$
цуваа нийлнэ. Иймд $P(z)$ утга оршино.

Одоо $$P(z,n)=(1+u_1(z))(1+u_2(z)\cdot \ldots \cdot(1+u_n(z))$$
$\tilde{P}=\prod\limits_{k=1}^\infty (1+a_k)$ ба $\tilde{P}(n)=\prod\limits_{k=1}^n (1+a_k)$ гэж тэмдэглэе.
\begin{equation}
    \tag{$*$} \label{8*}
    \begin{split}
        &|P(z,n)|=|1+u_1(z)|\cdot |1+u_2(z)|\cdot \ldots \cdot |1+u_n(z)|\leq \\
        &\leq (1+|u_1(z)|)(1+|u_2(z)|)\cdot \ldots \cdot (1+|u_n(z)|)\leq \tilde{P}(n)
    \end{split}
\end{equation}
\begin{equation}
    \tag{$**$} \label{8**}
    \begin{split}
    &\left|\frac{P(z)}{P(z,n)}-1\right|=\lim\limits_{k\rightarrow\infty}\left|\frac{P(z,n+k)}{P(z,n)}-1\right|=\\
        &\leq\lim\limits_{k\rightarrow\infty}(|u_{n+1}(z)|+ |u_{n+2}(z)|+ \cdots + |u_{n+k}(z)|+\cdots \\
        &\cdots +|u_{n+1}(z)\cdot u_{n+2}(z)\cdot \ldots \cdot u_{n+k}(z)|)\leq\\
        &\leq \lim\limits_{k\rightarrow \infty} (a_{n+1}+a_{n+2}+\cdots +a_{n+k}+a_{n+1}\cdot a_{n+2}+\cdots\\
        &\cdots +a_{n+1}\cdot a_{n+2}\cdot \ldots \cdot a_{n+k})= \\
        &=\lim\limits_{k\rightarrow \infty}\left(\frac{\tilde{P}(n+k)}{\tilde{P}(n)}-1\right) =\frac{\tilde{P}}{\tilde{P}(n)}-1\Rightarrow \\
        &\Rightarrow \left|\frac{P(z)}{P(z,n)}-1\right|\leq \frac{\tilde{P}}{\tilde{P}(n)}-1
    \end{split}
\end{equation}
болно. (\ref{8*}), (\ref{8**}) ашиглавал
\begin{align*}
    |P(z)-P(z,n)|&=|P(z,n)|\cdot \left|\frac{P(z)}{P(z,n)}-1\right| \leq \\
    &\leq \tilde{P}(n)\cdot \left(\frac{\tilde{P}}{\tilde{P}(n)}-1\right) =\tilde{P}-\tilde{P}(n)\xrightarrow[n\rightarrow\infty]{} 0
\end{align*}
учир $P(z,n)$ функцэн дараалал $P(z)$-рүү жигд нийлж байна. Иймд Вейерштрассын теорем ёсоор $P(z)\in H(D)$ болж батлагдав.
\hfill $\blacksquare$\end{prf}
\section{Төгсгөлгүй үржвэр ба бүхэл функц}
$a_1,a_2,\ldots ,a_n,\ldots$ тэгээс ялгаатай гишүүдтэй комплекс тоон дарааллын хувьд 
\begin{eqnarray} \label{8.5}
 |a_k|\leq |a_{k+1}| \quad \text{ба} \quad |a_k|\rightarrow +\infty
\end{eqnarray} байг. 
\begin{thm}
\textbf{(Вейерштрасс).}  Тэгвэл $p\in\mathbb{N}$ үед
\begin{eqnarray} \label{8.6}
 G(z)=z^{p-1}\cdot \prod\limits_{k=1}^\infty \left(1-\frac{z}{a_k}\right)\cdot e^{\frac{z}{a_k}+\frac{1}{2}\left(\frac{z}{a_k}\right)^2+\cdots +\frac{1}{k-1}\cdot \left(\frac{z}{a_k}\right)^{k-1}}
\end{eqnarray}
функц
\begin{enumerate}
    \item [a)] Бүхэл функц байна (өөрөөр хэлбэл $G(z)\in H(G))$
      \item [б)] $G(a_k)=0 \quad (\forall k).$
     \item [в)] $p-1$ нь $z=0$ цэг $G(z)=0$ тэгшитгэлийн шийд болох эрэмбэ нь байна.
     \item [г)] Өөр тэглэгч цэг $G(z)$ функцэд байхгүй.
\end{enumerate}
\end{thm}
Энэ теоремын баталгааг орхив.\\
Харин $exp$ үржигдэхүүн нь төгсгөлгүй үржвэрийн нийлэлтийг хангахад шаарддагийг тэмдэглэе. Теоремын урвуу нь дараах хэлбэрээр батлагдсан байдаг.
\begin{thm}
$\forall f(z)\in H(\mathbb{C})$ бүхэл функц авч үзье. Түүний тэглэгч цэгийн тоо төгсгөлгүй $a_1,a_2, \ldots, a_k,\ldots$  дараалал үүсгэдэг бол
$$f(z)=e^{H(z)}\cdot z^{p-1}\prod\limits_{k=1}^\infty \left(1-\frac{z}{a_k}\right) e^{\frac{z}{a_k}+\frac{1}{2}\left(\frac{z}{a_k}\right)^2+\cdots +\frac{1}{n-1}\cdot \left(\frac{z}{a_kk}\right)^{n-1}} $$ байна. Үүнийг Вейерштрассын томьёо \index{Вейерштрассын томьёо} гэнэ. Энд $p\in\mathbb{N}$ ба $H(z)$ ямар нэг бүхэл функц байна. 
\end{thm}
\begin{exm}
$f(z)=\sin z$ гэж авъя. Тэглэгч цэгүүд нь $z=k\pi $ учир
\begin{align*}
    \sin z&=e^{H(z)}\cdot z\cdot \prod\limits_{k=-\infty}^{+\infty} \left(1-\frac{z}{k\pi}\right) e^{\frac{z}{k\pi}+\frac{1}{2}\left(\frac{z}{k\pi}\right)^2+\cdots +\frac{1}{n-1}\cdot \left(\frac{z}{k\pi}\right)^{n-1}}=\\
    &=e^{\tilde{H}(z)}\cdot z\cdot \prod\limits_{k=1}^{\infty} \left(1-\frac{z}{k\pi}\right)\left(1-\frac{z}{-k\pi}\right)=\\
    &=e^{\tilde{H}(z)}\cdot z\cdot \prod\limits_{k=1}^{\infty} \left(1-\frac{z^2}{k^2\pi^2}\right)
\end{align*}
(энд $k$ ба $-k$ гишүүдээр нь бүлэглэснээр абсолют нийлдэг үржвэр үүссэн тул ехр үржигдэхүүнүүдийн илтгэгчийг $H(z)$ дээр нэмж $\tilde{H}(z)$-р тэмдэглэв).

$\tilde{H}(z)$  функцийг олж тогтооход дээрх теорем тус болохгүй. Тухайн тохиолдолд
$$\cot (z)=\frac{1}{z}+\sum\limits_{k=1}^\infty \frac{2z}{z^2-k^2\pi^2}$$
гэж байдгийг шууд ашиглаж $(\ln\sin z)'$ уламжлалыг судалснаар $H(z)=1$ болохыг тогтоож болно. Үүнийг дасгал болгон ажиллаарай. Ийнхүү
$$\sin z=z\cdot \prod\limits_{k=1}^\infty\left(1-\frac{z^2}{k^2\pi^2}\right)$$
\end{exm}
\begin{center}
  \textbf{  Дасгал}
\end{center}
\begin{enumerate}
    \item $e^z-1=e^{\frac{z}{2}}\cdot z\cdot \prod\limits_{k=1}^\infty \left(1+\dfrac{z^2}{4\pi^2k^2}\right)$
    \item $\cos z=\prod\limits_{k=1}^\infty \left(1-\dfrac{4z^2}{\pi^2(2k-1)^2}\right);$
    \item $\dfrac{2\sin iz}{e^{2z}-1}=e^{-z+\frac{\pi}{2}i}$ 
\end{enumerate} байдгийг тус тус харуул.
\section{Мерорморф функц\index{Мерорморф функц}}
$\mathbb{C}$ хавтгайд туйлын цэгүүдээс өөр онцгой цэггүй функцийг мероморф функц гэж тодорхойлдог. \\
\textbf{Дасгал.} Мероморф функцийн туйлын цэгүүд нь төгсгөлтэй хязгаарын цэгтэй байж болохгүй.\\
\textbf{Заавар.} Тийм хязгаарын цэг оршдог гэвэл тэр нь онцгойрлын цэг болно. Гэтэл мероморф функцийн тодорхойлолт ёсоор уг цэг туйлын онцгой цэг байна. Энэ цэг дээрх Лораны цувааг ерөнхий хуваарьт оруулж харвал ямар зөрчил үүсэхийг ол.\\
\begin{cor}
Мероморф функц нь төгсгөлтэй эсвэл (\ref{8.5}) нөхцөлийг хангасан төгсгөлгүй олон туйлын цэгүүдтэй байна.
\end{cor}
$f(z)$ нь мероморф функц бол $G(z)$ бүхэл функцийг эсвэл олон гишүүнт (\ref{8.6}) томьёогоор байгуулахдаа $f(z)$-н туйлын цэг бүр нь яг тэр эрэмбээрээр $G(z)$-н тэглэгч цэг байхаар байгуулж болно. Энэ үед
\begin{eqnarray}
 G_1(z)=f(z)\cdot G(z)
\end{eqnarray}
бүхэл функц болно. Ийнхүү:

\begin{thm}
Аливаа мероморф функц нь хоёр бүхэл функцийн ноогдвор хэлбэрт тавигдана.
\end{thm}
\section{Аналитик үргэлжлэлийн тодорхойлолт} \index{Аналитик үргэлжлэл}
Морерын теорем ёсоор хоёр голоморф функц тодорхойлогдох мужуудынхаа огтлолцол $D=D_1\cap D_2$; $int D\neq 0$ дээр ижил байвал $D_1\cup D_2$ муж дээр голоморф функц тодорхойлдог.

 Мөн $D_1$ муж дээрх голоморф функцийг $D_2\setminus D$ муж руу өргөтгөх өөр функц байхгүй нь нэг утгын теоремоор батлагдсан. Ингэж өгөгдсөн функцийн голоморф чанарыг алдуулалгүй тодорхойлох мужийг нь өргөтгөж байгааг аналитик үргэлжлэл хийгдэж байна гэдэг.
 \begin{center}
     \includegraphics[width=5cm]{Zurag/37.jpg}
 \captionof{figure}{}
        \label{ref hiih ner}
 \end{center}
\begin{exm}
$y=\sin x$ функц $\mathbb{R}$-дээр тодорхойлогдсон байг. Үүний хувьд функц байх нь $\mathbb{C}$ дээрх аналитик үргэлжлэл байх нь 
$$\sin x=\sum\limits_{k=0}^\infty (-1)^k\cdot \frac{x^{2k+1}}{(2k+1)!} \quad (x\in\mathbb{R})$$
тэнцэтгэлээс харагдаж байна.
\end{exm}
\begin{exm}
$$f(z)=1+z+z^2+\cdots \in H(|z|<1)$$ бол $g(z)=\dfrac{1}{1-z}\in H(\mathbb{C}\setminus \{1\})$ нь аналитик үргэлжлэл гэдгийг батал.
\end{exm}
Эндээс харвал ямар нэг $D$ мужид голоморф функц өгөгдсөнөөр түүнийг аналитик үргэлжлүүлсээр цааш өргөжүүлэх боломжгүй болсон үед үүсэж байгаа объектыг бүрэн аналитик функц гэнэ. Бүрэн аналитик функц өөрийн өчүүхэн хэсгээрээ ганц утгатай тодорхойлогддог.

\section{Бүрэн аналитик функц тодорхойлох Вейерштрассын арга }\index{Вейерштрассын арга}
$f\in H(D)$ байх функц өгөгдөв. Дурын $z_1\in D$ цэгийг авснаар энэ цэгийн орчинд өгөгдсөн функц
\begin{eqnarray} \label{8.8}
 f(z)=\sum\limits_{n=0}^\infty a_n(1)(z-z_1)^n
\end{eqnarray}
Тейлорын цуваанд задарна. Энэ цувааны нийлэлтийн радиус $R_1$ хязгааргүй бол функц бүхэл функц болон өргөтгөгдөнө. Харин (\ref{8.8}) цувааны нийлэлтийн радиус төгсгөлөг бол уг дугуйгаас $z_2(z_2\neq z_1)$ цэг сонгон авч
\begin{eqnarray}
 f(z)=\sum\limits_{n=0}^\infty a_n(2)(z-z_2)^n
\end{eqnarray}
цуваанд задалъя.

Хэрэв $R_2>R_1-|z_2-z_1|$ бол функц анхны дугуйн $z_2-z_1$ радиус векторын дагуу аналитик үргэлжлэл зохиогдоно. Эсрэг тохиолдолд уг чиглэлийн дагуу үргэлжлэхгүй байна. Одоо (\ref{8.8}) элементийг бүх чиглэлийн дагуу үргэлжлүүлсэн гэвэл цааш үргэлжлэхгүй $$\left (z_k; \quad \sum\limits_{n=0}^\infty a_n(k)(z-z_k)^n\right)$$ элементүүдээс тогтсон бүтэц үүснэ. Энэ бүтцийн бүрэн аналитик функц Вейерштрассын аргаар өгөгдлөө гэдэг. Дугуйгаар дамжин үргэлжилсээр нэг эгшинд дугуй өмнө таарч байсан дугуйтай огтлолцох нь бий. Энэ үед элементүүдийн зэрэгт цувааны утга ерөнхий хэсэгтээ таарах эсвэл эс таарах тохиолдол үүснэ.
\begin{center}
    \includegraphics[width=6cm]{Zurag/38.PNG}
    \captionof{figure}{}
        \label{}
\end{center}
\section{Бүрэн аналитик функц тодорхойлох \\
Риманы арга \index{Риманы арга}
}
Риманы арга нь $\mathbb{C}$ хавтгайн оронд Риманы гэж нэрлэгдэх тусгай гадаргуунууд авч үзсэнээр олон утгат бүрэн аналитик функцийг нэг утгатай функцийн тодорхойлолтод нийцүүлдэг.
\begin{exm}
$w=\sqrt{z}$ буулгалт нь $w_1=\sqrt{\rho}\cdot e^{i\frac{\varphi}{2}}$; $w_2=\sqrt{p}e^{i\left(\frac{\varphi}{2}+\pi \right)}$ хоёр мөчир функцэд задардгийг өмнө үзсэн. $\mathbb{C}$ хавтгайг сөрөг хагас тэнхлэгийн дагуу хайчилсан мужийг $D$ гэе. Дээрх мөчир функцүүд нь $D$ мужийг харгалзан $\mathbb{C}$ хавтгайн баруун ба зүүн хагас хавтгайд буулгана. ($E_1$ ба $E_2$ гэж эдгээрийг тэмдэглэе). Одоо $D$ мужийг 2 ширхэг авч зурагт үзүүлсэн наалтыг хийе:\\
\begin{center}
    \includegraphics[width=8cm]{Zurag/39.PNG}
\end{center}
\begin{center}
    \includegraphics[width=6cm]{Zurag/39.2.PNG}
    \captionof{figure}{}
        \label{}
\end{center}
Хар цэгээр тэмдэглэсэн ирмэгүүдийг хооронд нь, цагаанаар тэмдэглэснийг нь хооронд нь нааснаар $\tilde{D}$ өвөрмөц гадаргуу үүснэ (Жич: Ердийн цаасыг ингэж наах нь боломжгүй, 2-р наалтын үед $\tilde{D}$) өөрөө өөрийгөө огтолж байгаа мэт боловч ингэж огтлохгүйгээр топологи утгаар наасан гэж сэтгэдэг).

Ийнхүү $w=\sqrt{z}$ буулгалт нь  $\tilde{D}\rightarrow E=E_1\cup E_2$ харилцан нэг утгатай функц болон хувирав. $D$ мужийг хуудас гэж нэрлэх ба $\tilde{D}$ мужийг 2 хуудаст гэдэг. $0;\infty \in\tilde{D}$ бусад цэгийн орчин дугуй учир $\tilde{D}$ муж нь 2 хэмжээст цогцтой дөхөж байна. Энэ чанар алдагдаж байгаа $0,\infty$ цэгүүдийг мөчирлөлтийн цэг гэдэг.
\end{exm}
\begin{exm}
$w=\sqrt[n]{z}$ буулгалтын риманы гадаргуу \index{Риманы гадаргуу} нь яг ижлээр наагдах ба ялгаа нь $n$ хуудаст риманы гадаргуу болгоё. Энэ гадаргууг зурж үзэхийг дасгал болгоё. Харин $w=z^n$ буулгалтын утгын муж нь уг риманы гадаргуу.
\end{exm}
\begin{exm}
$w=\ln z$ буулгалт нь төгсгөлгүй хуудаст Риманы гадаргуутай байдаг.
$w_k=\ln \rho +i(\varphi +2\pi k)(k\in\mathbb{Z})$ мөчир нь $D$ мужийг
$$E_k=\{w:2\pi k -\pi \leq \im w<2\pi k +\pi\}$$ зурваст буулгана гэдгийг өмнө үзсэн.
\begin{center}
    \includegraphics[width=5cm]{Zurag/40.PNG}
    \captionof{figure}{}
        \label{}
\end{center}
Иймээс хязгааргүй олон $D$ мужийг өмнөх жишээний зарчмаар төгсгөлгүй наасан гэж сэтгэвэл $\tilde{D}$ риманы гадаргуу нь үүснэ.

Эцэст нь дурдахад бүрэн аналитик функцийн энэхүү хоёр байгуулалт нь яг ижил, нэг нь нөгөөдөө шилждэг гэдэг нь нэг утгын теоремоос илт билээ.
\end{exm}
\begin{center}
   \textbf{ Дасгал}
\end{center}
\begin{enumerate}
    \item [4.] $f(z)=\bint\limits_{n=1}^\infty (-1)^{n+1}\dfrac{z^n}{n}$ функцийн хувьд 
    $$g(z)=\ln 2-\frac{1-z}{1\cdot 2}-\frac{(1-z)^2}{2\cdot 2^2}-\frac{(1-z)^3}{3\cdot 2^3}\cdot \ldots$$
    аналитик үргэлжлэл болохыг батал.
    \item [5.] $$1+az+a^2z^2+\cdots \quad \text{ба} \quad \frac{1}{1-z}-\frac{(1-a)^z}{(1-z)^2}+\frac{(1-a)^2z^2}{(1-z)^3}-\cdots$$
    функцүүд нь харилцан аналитик үргэлжлэл болдгийг батал.
\item [6.] $f(z)=\bint\limits_{n=0}^\infty z^{n!}$ функц нэгж дугуйн гадна аналитик үргэлжлэлгүй болохыг батал.
\end{enumerate}
\chapter{Геометр зарчмууд}
Энэ бүлэгт голоморф функцийн ерөнхий топологи шинжүүдийг үзүүлдэг геометр онолын зарчмуудыг оруулав.
%\section{Аргументийн зарчим} \index{Аргументийн зарчим}
\begin{thm}
$f$ нь $D\subset \mathbb{C}$ муж дээр мероморф байг. \\
$\overline{G}\subset D$ байх $G$ мужийн хил $\partial G$ нь битүү муруйнуудаас тогтох ба $f$ функцийн туйлын ба тэглэгч цэгүүдийг дайраагүй байг.
$N$ ба $P$ нь $G$ муж дахь $f$ функцийн нийт тэглэгч ба туйлын тоог (эрэмбэ харгалзана) зааж байг.\\
Тэгвэл 
$$N-P=\frac{1}{2\pi i}\int\limits_{\partial G^+}\frac{f'(z)}{f(z)}dz$$
байна.
\end{thm}
\begin{prf}
$f$ функцийн туйлын цэг $a$ (эрэмбэ $p$) тэглэгч цэг $b$ (эрэмбэ $n)$ гэе. 

$Res \left(\dfrac{f'}{f}; a\right)=-p$, $Res \left(\dfrac{f'}{f}; b\right)=n$
болохыг үлдэц бодох томьёогоор гаргахыг дасгал болгоё.

$\partial G$ нь $\dfrac{f'}{f}$ функцийн онцгой цэгийг дайраагүй тул үлдэцийн тухай теорем хэрэглэвэл 
\begin{align*}
    \frac{1}{2\pi i}\int\limits_{\partial G^+}\frac{f'(z)}{f(z)}dZ&=\sum\limits_{a_k} Res\left(\frac{f'}{f}; a_k\right) +\sum\limits_{b_k}Res \left(\frac{f'}{f};b_k\right) \\
    &=N-P
\end{align*}

болж теорем батлагдав. 
\hfill $\blacksquare$\end{prf}
\begin{cor}
\textbf{(Аргументын зарчим). } Дээрх теоремын нөхцөлд
$$N-P=\frac{1}{2\pi}\Delta_{\partial G+}\arg f(z)$$
байна. (Энд $\Delta_{\partial G+}\arg f(z)$ гэдэг нь $\partial G^+$ муруйн дагуу бүтэн тойроход $\arg f(z)$ хэмжигдэхүүний өөрчлөгдөх хэмжээ билээ. Үүнийг $\partial G^+$ муруйн дагуу аргументийн өөрчлөлт гэдэг).
\end{cor}
\begin{prf}
Ньютон-Лейбницийн томьёо ёсоор
\begin{align*}
    \frac{1}{2\pi i}\int\limits_{\partial G^+}\frac{f'(z)}{f(z)}dz&=\frac{1}{2\pi i}\Delta_{\partial G^+}\ln f(z)\\
    &=\frac{1}{2\pi i}[\ln |f(z)|-\ln |f(z)|+i\Delta_{\partial G^+}\arg f(z)]\\
    &=\frac{1}{2\pi i}\Delta_{\partial G^+}\arg f(z)
\end{align*} болж батлагдав. 
\hfill $\blacksquare$\end{prf}
\begin{thm}
\textbf{ (Руше).} $f,g\in H(D)$ ба $G$ муж нь Теорем 9.1-ийн нөхцөлд тохирч байна. $|f(z)|>|g(z)| (\forall z\in\partial G)$ бол $f$ ба $f+g$ функцүүд нь $G$ мужид ижил тооны тэглэгчтэй байна. 
\end{thm}
\begin{prf}
$f,g\in H(D)$ учир $P_f=P_{f+g}=0$ байна. \\ $|f(z)|>|g(z)| (\forall z\in \partial G)$ нөхцөлөөс $\Delta_{\partial G^+}\arg \left(1+\dfrac{g}{f}\right) =0$  болно. (өөрөөр хэлбэл $1+\dfrac{g}{f}$ функцийн вектор $w=0$ цэгийг тойрч чадахгүй).

Иймд 
\begin{align*}
    \Delta_{\partial G^+}\arg (f+g)&=\Delta_{\partial G^+}\left[\arg f+\arg f\left(1+\frac{g}{f}\right)\right]=\\
    &=\Delta_{\partial G^+}\arg f
\end{align*}
\hfill $\blacksquare$\end{prf}
болж теорем батлагдлав.
\begin{cor}
\textbf{(Алгебрийн үндсэн теорем).} \index{Алгебрийн үндсэн теорем}
$$P_n(z)=z^n+a_1z^{n-1}+a_2z^{n-2}+\cdots +a_n=z^n+P_{n-1}(z)$$
олон гишүүнт өгөгдсөн байг. Тэгвэл $P_n(z)$ нь $\mathbb{C}$ огторгуйд яг $n$ ширхэг шийдтэй (давталт тооцоод).
\end{cor}
\begin{prf}
$\dfrac{P_{n-1}(z)}{z^n}\xrightarrow[z\rightarrow \infty]{}0$ учир $\{|z|<R\}$ дугуйн дотор бүх шийд агуулагдсан ба $$|z^n|>|P_{n-1}(z)| \qquad (z\in\{|z|=R\})$$ байхаар хангалттай их $R$ зааж болно. Иймд Рушегийн теорем ёсоор $N_{P_n(z)}=N_{z^n}=n$ болж батлагдав.
\hfill $\blacksquare$\end{prf}
\section{Муж хадгалагдах зарчим} \index{Муж хадгалагдах зарчим}
\begin{thm}
\textbf{(Муж хадгалагдах зарчим).} Хэрэв $f\in H(D)$, $D$ муж ба $f$ нь тогтмол биш бол $E=f(D)$ муж байна.\\
(Өөрөөр хэлбэл тогтмол биш голоморф функц холбоост задгай олонлогийг холбоост задгай олонлогт буулгана).
\end{thm}
\begin{prf}
Холбоост чанар хадгалагдахыг харуулах дасгал болгон үлдээе. Одоо $E$ задгай болохыг харуулъя. $w_0\in E$ ба $z_0$ нь түүний аль нэг эх дүр байг $(f(z_0)=w_0)$.\\
$D$ задгай тул $\{|z-z_0|<r\}\subset D$ ба уг дугуйд $f$ функцийн $w_0$-цэг $z_0$-оос гадна байхгүй байхаар $r>0$ радиус зааж болно. 

Иймд $\mu =\min\limits_{|z|=r}|f(z)-w_0|>0$ байх болно. Хэрэв $\{|w-w_0|<\mu \}\subset E$  гэж баталбал теорем батлагдана. \\
$\forall w_1\in \{|w-w_0|\}<\mu$ гэвэл $f(z)-w_1=(f(z)-w_0)+(w_0-w_1)$ ба $|f(z)-w_0|\geq \mu>|w_0-w_1|$ учир Рушегийн теорем ёсоор $f(z)-w_1$ ба $f(z)-w_0$ функцүүдийн шийдийн тоо адилхан байх ёстой. Иймд $\exists z_1\in \{|z-z_0|\}<r$ (цорын ганц!) ба $f(z_1)=w_1$ болж улмаар $w_1\in E$ болов. Ийнхүү $\{|w-w_0|<\mu \}\subset E$ болж муж хадгалагдах зарчим батлагдлаа.
\hfill $\blacksquare$\end{prf}
\section{Хамгийн их модулийн зарчим,
Шварцын лемм } \index{Шварцын лемм}
\begin{thm}
\textbf{(Хамгийн их модулийн зарчим).} \index{Хамгийн их модулийн зарчим} $f\in H(D)$ ба $f$ тогтмол биш функц бол $|f(z)|$ функц $D$ мужийн дотоод цэг дээр локаль максимум утгатай байж чадахгүй.
\end{thm}
\begin{prf}
Эсрэгээс нь $\exists z_0\in V_{z_0}\subset D:\forall z\in V_{z_0}$ хувьд $|f(z)|\leq |f(z_0)|$ (локаль эрс биш максимум утга) байдаг гэе. $V_{z_0}$ нь мужийн тодорхойлолтод тохирох учир муж хадгалагдах зарчим ёсоор $f(V_{z_0})$ нь задгай байх ёстой. $f(z_0)$ нь уг мужийн дотоод цэг байх ёстой болов. Иймээс
$\exists w\in f(V_{z_0}):|w|>|f(z_0)|$ болж $|f(z_0)|$ нь максимум утга гэсэнтэй зөрчилдөв. Иймээс локаль максимум утга олдохгүй болж теорем батлагдлаа. 
\hfill $\blacksquare$\end{prf}
\begin{cor}
$f\in H(D)$, $K$-компакт $\subset D$ бол $\forall z\in K$ хувьд  $|f(z)\leq \max\limits_{z\in \partial K}|f(z)|$ байна.
\end{cor}
\begin{lem}
\textbf{(Шварц).}  $f\in H(V=\{|z|<1\})$; $|f(z)|\leq 1(\forall z\in V)$; $f(0)=0$ байг. Тэгвэл $\forall z\in V$ хувьд $|f(z)|\leq |z|$ байна. 
\end{lem}
\begin{prf}
$\varphi (z)=\dfrac{f(z)}{z}$ функц авч үзье. $z=0$ нь засагдах онцгой цэг 
тул  $\varphi(z)\in H(V)$ байна. $\forall z\in V$ сонгон авч бэхэлье. $\forall r\in ]|z|;1[$ хувьд дээрх мөрдөлгөөг хэрэглэвэл
$$|\varphi (z)|\leq \max\limits_{|z|=r}\frac{|f(z)|}{|z|}=\frac{1}{r}\max\limits_{|z|=r}|f(z)|\leq \frac{1}{r}$$ байна.

Энд $r\rightarrow 1-0$ гэж хязгаар авбал $|\varphi(z)|\leq 1$ болж улмаар $|f(z)|\leq |z|$ байна. $|\varphi (z)|=1$ тохиолдол биелсэн гэвэл  $\varphi (z)=const$ болох ёстой болж $f(z)=e^{i\alpha}z (\alpha \in\mathbb{R})$ байхад хүрнэ.
\hfill $\blacksquare$\end{prf}
\section{Конформ изоморфизм ба автоморфизм,
Риманы теорем 
} \index{Риманы теорем}
$D_1$ ба $D_2$ мужууд өгөгдөөд, $f:D_1\rightarrow D_2$ харилцан нэг
утгатай, конформ $(f\in H(D_1))$ ба $f'(z)\neq 0(\forall z\in D_1)$ бол $f$
функцийг конформ изоморфизм гэдэг. Ийм бүх функцүүдийн олонлогийг $Iso (D_1, D_2)$ гэж тэмдэглэе. Харин $D_1=D_2$ бол изоморфизмийг автоморфизм \index{Автоморфизм} гэж нэрлээд 
$Aut (D_1)=Iso (D_1, D_2)$ гэж тэмдэглэдэг.\\
\textbf{Дасгал.} $Aut (D)$ нь функцийн композиц үйлдлээр бүлэг болно.
\begin{thm}
$V=\{|z|<1\}$ мужийн хувьд
$$Aut (V)=\left\{e^{i\alpha}\cdot \frac{z-a}{1-\overline{a}z}:|a|<1, \alpha\in\mathbb{R}\right\}$$
байна.
\end{thm}
\begin{prf}
$\forall a\in\mathbb{C}:|a|<1$ ба $\forall\alpha\in\mathbb{R}$ хувьд $f(z)=e^{i\alpha}\dfrac{z-a}{1-\overline{a}z}$ функцийг авч үзье. $a$ ба $\dfrac{1}{\overline{a}}$ нь $|z|=1$ тойргийн хувьд тэгш  хэмтэй.

$f(a)=0$; $f\left(\dfrac{1}{\overline{a}}\right)=\infty$ тул бутархай шугаман буулгалтын тэгш хэм хадгалах чанар ёсоор $f(|z|=1)$ нь  $z=0$ төвтэй тойрог байх ба $V$ нь түүний дотоод мужид бууна.\\
$|f(1)|=1$ болохыг дасгал болгон гүйцэтгэвэл тэр тойрог нь нэгж радиустай болох нь тогтоогдож $\{f(z)\}\subset Aut (V)$ болох нь батлагдлаа. Одоо $\forall f(z)\in Aut (V)$ авъя. $f(0)=a$ байг.Тэгвэл дасгал ёсоор $g(z)=f\left(\dfrac{z-a}{1-\overline{a}z}\right) \in Aut (V)$ ба $g(0)=0$ байна. Иймээс Шварцын лемма ёсоор $|g(z)|\leq |z|$ байх ба түүний урвуу нь автоморфизм тул $g^{-1}(z)$ функц яг адилаар $|g^{-1}(z)|\leq |z|$ тэнцэл биш өгнө. Иймээс $|g(z)|=|z|$ болж Шварцын леммийн тэнцэтгэл тохиолдлыг өгнө.
$$g(z)=e^{i\alpha}z\Rightarrow f\left(\frac{z-a}{1-\overline{a}z}\right) =e^{i\alpha}z\Rightarrow f(z)=e^{i\alpha} \frac{z+a}{1+\overline{a}z}$$ болж нөгөө тийш агуулагдал батлагдана. Иймд функцийн олонлогууд тэнцүү байна.
\hfill $\blacksquare$\end{prf}
\begin{center}
    \textbf{Дасгал}
\end{center}
\begin{enumerate}
    \item $Aut (\overline{\mathbb{C}})=\left\{\dfrac{az+b}{cz+d}; \quad ad-bc\neq 0\right\}$
    \item $Aut (\mathbb{C})=\{az+b: a\neq 0\}$
    \item $Aut (\{\im z>0\})=\left\{\dfrac{az+b}{cz+d};\quad a,b,c,d\in \mathbb{R},\quad ad-bc>0\right\}$
\end{enumerate}
\begin{thm}
\textbf{(Риман).} $D$ нэг холбоост муж $(D\neq \mathbb{C}$ ба $D\neq \overline{\mathbb{C}}$ бол $\exists f\in Iso (D,V)$ байна. (энд $V$ нь нэгж дугуй болно). Энэ теоремын баталгаа нь нүсэр тул орхиж байна.
\end{thm}
\section{Тэгш хэмийн зарчим} \index{Тэгш хэмийн зарчим}
\begin{lem}
$D_1$ ба $D_2$ мужуудын хувьд $D_1\cap D_2= \emptyset$; $f\in H(D_1)\cap H(D_2)$ байг. $\gamma$ нь хэрчим ба $D_1\cup \gamma \cap D_2$ муж үүсгэдэг байг.
\begin{center}
    \includegraphics[width=5cm]{Zurag/41.PNG}
    \captionof{figure}{}
        \label{}
\end{center}
$f\in C(D_1\cup \gamma \cup D_2)$ бол $f\in H(D_1\cup \gamma \cup D_2)$ байна. 
\end{lem}
\begin{prf}
Мореры теоремын нөхцөлийг хангасан дурын гурвалжин $\Delta$ авахад $\bint\limits_{\partial \Delta}f(z)dz=0$ гэж харуулбал хангалттай. Уг $\Delta$ нь $\gamma$ хэрчимтэй огтолцоогүй тохиолдолд Кошийн теорем ёсоор $\bint\limits_{\partial \Delta} f=0$ болно.
\begin{center}
    \includegraphics[width=5cm]{Zurag/42.PNG}
    \captionof{figure}{}
        \label{}
\end{center}
Одоо $\Delta$ нь $\gamma$-г зурагт үзүүлсэн байдлаар огтолж байг. $\gamma$-тай параллель түүнээс $h$ зайтай хоёр шулуунаар $T_h$ трапец огтолж авъя. Кошийн теорем ёсоор $\bint\limits_{\partial \Delta} f=\int\limits_{\partial T_h}f$  байна. $h\rightarrow 0$ үед трапецын хажуу ирмэгийн урт тэг рүү тэмүүлэх тул түүгээр авсан интегралын утга багасаж барагдашгүй хэмжигдэхүүн үүсгэнэ. $\overline{\Delta}$ компакт тул $f$ функцийг жигд тасралтгүй болохыг ашиглаж $h\rightarrow 0$ үед трапецын дээд доод ирмэгээр авсан интегралын утгууд $\bint\limits_{(\Delta \cup \gamma)^+}f$ утгыг эсрэг тэмдэгтэйгээр дөхнө. Иймд $\bint\limits_{\partial T_h}f \xrightarrow[h\rightarrow 0]{} 0$ болж  $\bint\limits_{\partial \Delta}f=0$ батлагдав. \\
Харин $\overline{\Delta}$ нь $\gamma$-тай ирмэг ба оройгоор огтолцох тохиолдолд леммын баталгааг дасгал болгон үлдээв. 
\hfill $\blacksquare$\end{prf} 
\begin{thm}
$D_1$ ба $E_1$ нь дээд хагас хавтгайд орших мужууд ба $\gamma\subset \partial D_1$, $\Gamma\subset \partial E_1$ нь бодит тэнхлэгийн хэрчмүүд байг.$D_1$
ба $E_2$ нь $D_1$ ба $E_1$-ийн бодит тэнхлэгийн хувьд тэгш хэмтэй мужууд гэе. $f\in H(D_1)\cap C(D_1\cup \gamma)$ функцийн хувьд $D_1$ мужийг $E_1$ мужид конформ буулгадаг ба $\Gamma =f(\gamma)$ байг. Тэгвэл
$$
\tilde{f}(z)=\begin{cases}
 f(z), \quad z\in D_1\cup \gamma \\
 \overline{f(\overline{z})}, \quad z\in D_2
\end{cases}
$$
функц $D=D_1\cup \gamma \cup D_2$ мужийг $E=E_1\cup \Gamma \cup E_2$ мужид конформ буулгана.
\end{thm}
\begin{prf}
Уг теоремын баталгааг Морерийн теорем ашиглан хийе. $\Delta$ сонгогдоход дараах тохиолдлууд байна.
\begin{enumerate}
    \item $\Delta \subset D_1$ энэ үед Кошийн теорем ёсоор $\bint\limits_{\partial \Delta}\tilde{f}=0$ байна.
    \item $\Delta \subset D_2$ энэ үед уг гурвалжны дээд хагас хавтгай дахь тэгш хэмийг $\Delta_1$ гэвэл
    $$\int\limits_{\partial \Delta^+} \tilde{f}(z)dz=\int\limits_{\partial \Delta^+} \overline{f(\overline{z})}dz=\int\limits_{\partial \Delta^-} \overline{f(z)}d\overline{z}=\overline{\int\limits_{\partial \Delta_1^-} f(z) dz}=\overline{0}=0 $$
    болж батлагдана.
    \item $\Delta \cap \gamma \neq \emptyset$ тохиолдлыг дасгал болгон үлдээв.
\end{enumerate}
\hfill $\blacksquare$\end{prf}
\begin{thm}
 (Риман-Шварцын тэгш хэмийн зарчим). Дээрх теоремын нөхцөлийг $\gamma$ ба $\Gamma$ нь дурын байрлалтай хэрчим эсвэл тойргийн нум болгон өөрчлөн томьёолж болно. Үүнийг Риман- Шварцийн тэгш хэмийн зарчим \index{Риман- Шварцийн тэгш хэмийн зарчим} гэнэ. Энэ ерөнхий тохиолдол нь өмнөх тухайн тохиолдол руу бутархай шугаман буулгалтаар шилждэг.
\end{thm}
\section{Эллипслэг функц}
\begin{dfn}
Коллинеар биш 2 үетэй мероморф функцийг эллипслэг функц гэж тодорхойлдог. Өөрөөр хэлбэл $w=f(z)$ мероморф функц ба $T_1, T_2\in\mathbb{C}\setminus \{0\}:\dfrac{T_1}{T_2}\notin \mathbb{R}$ үед $\forall z:f(z+T_1)=f(z+T_2)=f(z)$  бол уг функцийг эллипслэг функц гэж нэрлэдэг. Дараах жишээг авч үзье. 
$$F(z,k)=\int\limits_{0}^z\frac{dz}{\sqrt{(1-z^2)(1-k^2z^2)}}$$ (энд $0<k<1$ ба $x>0$ үед $\sqrt{x}>0$ мөчир авна). 

Уг функц $\{\im z>0\}$ мужийг
 \begin{align*} \left\{ \right.
    \pm K; \pm K+iK':&K=\int\limits_0^1\frac{dx}{\sqrt{(1-x^2)(1-k^2x^2)}} , \\
    &\left.K'=\int\limits_1^{\frac{1}{k}}\frac{dx}{\sqrt{(x^2-1)(1-k^2x^2)}} \right\}
\end{align*} 
\end{dfn} оройнууд бүхий тэгш өнцөгтөд буулгана. 

Үүний урвуу функцийг авч үзвэл Риман-Шварцийн тэгш хэмийн зарчмаар аналитик үргэлжлэл зохиогдсоор $\mathbb{C}$-д мероморф, эллипслэг функц үүсэх ба үүнийг $w=sn (z,k)$ гээд эллипслэг синус гэж тэмдэглэдэг.
\begin{center}
   \textbf{ Дасгал}
\end{center}
\begin{enumerate}
    \item [4.] Рушегийн теорем ашиглан тэгшитгэлүүд $|z|<1$ дугуйд хэчнээн шийдтэй болохыг тогтоо.
    \begin{enumerate}
        \item [a)] $z^9-2z^6+z^2-8z-2=0;$
        \item [б)] $2z^5-z^3+3z^2-z+8=0;$
        \item [в)] $z^7-5z^4+z^2-2=0$
    \end{enumerate}
    \item [5.] Аргументийн зарчим ашиглан доорх тэгшитгэлүүд баруун хагас хавтгайд $(\re z>0)$ хэдэн шийдтэйг тогтоо.
    \begin{enumerate}
        \item [a)] $z^6+z^5+6z^4+5z^3+8z^2+4z+1=0;$
        \item [б)] $z^4+2z^3+3z^2+z+2=0$
    \end{enumerate}
    \item [6.] $z=0$ цэгийн орчинд голоморф ба $f\left(\dfrac{1}{n}\right)$ дараах дараалал үүсгэдэг функц орших уу?
    \begin{enumerate}
        \item [a)] $0,1,0,1,0,1,\ldots ;$
        \item [б)] $0,\dfrac{1}{2},0, \dfrac{1}{4},0,\dfrac{1}{6}, \ldots ,$
        \item [в)] $\dfrac{1}{2},\dfrac{1}{2},\dfrac{1}{4}, \dfrac{1}{4},\dfrac{1}{6},\dfrac{1}{6}, \ldots ,$
    \end{enumerate}
    \item [7.] $f\in H(D)$ ба $f\neq const$, $f(z)\neq 0 (\forall z\in D)$ бол $|f(z)|$ хамгийн бага утгаа авч чадахгүйг харуул.
    \item [8.] $F(z,k)$ функц $\{\im z>0\}$ мужийг заагдсан тэгш өнцөгтөд буулгадгийг батал.
\end{enumerate}
%
\chapter{Үйлдлийн тооллын үндэс} \index{Үйлдлийн тоолол}
\section{Лапласын хувиргалт. Уг эх ба дүрслэл} \index{Уг эх}
\begin{thm}
 $f:\mathbb{R}\rightarrow \mathbb{C}$ функцийн хувьд
 \begin{enumerate}
     \item $t<0$ үед $f(t)\equiv 0,$
     \item $t\geq 0$ үед $f(t)$ функц нь хэсэгчилсэн тасралтгүй ба дурын битүү интервал дээр төгсгөлөг ширхэг 1-р төрлийн тасралттай.
     \item $\exists M>0$, $\tilde{a}>0$ ба $\forall t>0$ хувьд $|f(t)|\leq Me^{\tilde{a}t} (inf \tilde{a}=a)$ тоог $f$-н өсөлтийн зэрэг гэдэг) нөхцөлүүд биелж байг.
     
     Тэгвэл
     \begin{eqnarray} \label{10.1}
      F(z)=\int\limits_0^{+\infty} e^{-zt}f(t)
     \end{eqnarray}
     функцийн хувьд
     \begin{enumerate}
         \item [a)] $\re z<0$ мужид утгагүй (интервал сарнина)
         \item [б)] $F(z)\in H (\{\re z>a)\}$ байна. 
         \item[В)] $\forall a_0>a$ хувьд $\re z\geq a_0$ мужид уг интервал жигд нийлнэ.
     \end{enumerate}
 \end{enumerate}
\end{thm}
\begin{dfn}
$f(t)$ функцийг уг эх, $F(z)$ функцийг дүрслэл \index{Дүрслэл}гээд энэ харгалзааг Лапласын хувиргалт \index{Лапласын хувиргалт} гэж нэрлэгдэх ба $\ml{f(t)}=F(z)$ тэмдэглэл хэрэглэдэг.
\end{dfn}
\section{Лапласын хувиргалтын чанарууд} 
Хэрэв $\ml{f(t)}=F(z)$ байвал 
\begin{enumerate}
    \item $\ml{\sum\limits_{k=1}^n a_k f_k(t)}=\sum\limits_{k=1}^n a_kF_k (z)$ (шугаман чанар). 
    \item $\ml{f(at)}=\dfrac{1}{a}F\left(\dfrac{z}{a}\right).$
    \item $f_\tau (t) = \begin{cases}
     0, \qquad \tau >t \\
     f(t-\tau), \qquad \tau \leq t
    \end{cases}$ бол $\ml{f_\tau}=e^{-\tau z}F(z.)$
    \item $\ml{f'(t)}=zF(z)-f(0)$
    \item $\ml{f^{(n)}(t)}=z^n\left(F(z)-\dfrac{f(0)}{z}-\dfrac{f'(0)}{z^2}-\cdots -\dfrac{f^{(n-1)}(z)}{z^n}\right).$
    \item $\ml{\int\limits_0^t f(\tau)d\tau}=\dfrac{1}{z}F(z)$
    \item $\ml{\int\limits_{0}^t f_1(\tau) \cdot f_2(t-\tau)d\tau}= F_1(z)\cdot F_2(z)$
    \item $\ml{(-1)^nt^n\cdot f(t)} = F^{(n)}(z)$
    \item $\ml{\int\limits_{z}^\infty F(z)dz}=\dfrac{f(t)}{t}$
    \item $\ml{e^{-\lambda t}f(t)}=F(z+\lambda)$
\end{enumerate}
чанарууд биелдэг.

Эдгээр чанарууд нь (\ref{10.1}) томьёоноос параметр агуулсан интегралыг хувиргах замаар гаргаж ирдэг. Жишээ болгож 7-р чанарыг баталъя.
\begin{align*}
    F'(z)&=\left(\int\limits_0^{+\infty} e^{-zt}f(t)dt\right)'=\int\limits_0^{+\infty} (e^{-zr})'_zf(t)dt=\\
    &=\int\limits_{0}^{+\infty} \ml{e^{-zt}(-tf(t))dt}= -tf(t)
\end{align*}
Лапласын хувиргалт нь ердийн дифференциал ба интеграл тэгшитгэлүүдийг ердийн тэгшитгэлд шилжүүлэх тухайн уламжлалт дифференциал тэгшитгэлийг ердийн дифференциал тэгшитгэлд шилжүүлэх зэргээр ашигладаг.
\begin{exm}
$
\begin{cases}
 y'+y=1 \\
 y(0)=0
\end{cases}$ байг. $\ml{y(x)}=F(z)$ гэвэл 4-р чанар ёсоор  
\end{exm}
\begin{prf}
$\ml{y'}=zF(z)-0$ ба томьёо 10.1 ёсоор $\ml{1}=\dfrac{1}{z}$ тул $zF(z)+F(z)=\dfrac{1}{z}$ болно. 
\begin{equation} \tag{$*$} \label{10*}
 F(z)=\ml{\frac{1}{z(z+1)}}=1-e^{-z}
\end{equation} энэ томьёог (\ref{10.1})-ээр шалгана уу.

Иймд $y(x)=1-e^{-x}$. Үүнээс харвал дүрслэлээр нь уг эхийг олох аргачлал шаардагдаж байна. Энд дээрх чанарууд ба нэгэнт бодогдсон Лапласын хувиргалтын хүснэгтийг ашиглаж болно. 

Лапласын хувиргалтын хүснэгт 
\begin{enumerate}
    \item $\ml{1}=\dfrac{1}{z}$
    \item $\ml{t^n}=\dfrac{\Gamma (n+1)}{z^{n+1}}, \qquad (n>-1)$
    \item $\ml{e^{at}}=\dfrac{1}{z-a}$
    \item $\ml{\sin at}=\dfrac{a}{z^2+a^2}$
    \item $\ml{\cos at}=\dfrac{z}{z^2+a^2}$
    \item $\ml{\sh at}=\dfrac{a}{z^2-a^2}$
    \item $\ml{\ch at}=\dfrac{z}{z^2+a^2}$
    \item $\ml{t^ne^{at}}=\dfrac{n!}{(z-a)^{n+1}}, \quad (n\in\mathbb{N})$
\end{enumerate}
Одоо энэ хүснэгтийг хэрхэн ашиглаж (\ref{10*}) томьёог гарган авч байгааг харъя.
$$\frac{1}{z(z+1)}=\ml{\frac{1}{z}-\frac{1}{z+1}}=1-e^{-t}$$
(энд 1 ба 3 томьёо ба чанар 1 ашиглагдсан болно). Дүрслэлээр уг эхийг олох ерөнхий томьёо нь голоморф функцийн интеграл ашигладаг ба энэ талаар 10.3-д үзнэ.
\hfill $\blacksquare$\end{prf}
\section{Меллиний томьёо}\index{Меллиний томьёо}
\begin{thm}
 $f(t)$ нь хэсэгчилсэн гөлгөр, $a$ өсөлтийн зэрэгтэй уг эх байг. $F(z)$ нь түүний дүрслэл гэвэл $\forall t>a$ бол 
 $$f(t)\doteq \frac{1}{2\pi i}\int\limits_{x-i\infty}^{x+i\infty} e^{zt}F(z)dz$$
 (энд $x:x>a$ байх дурын бодит тоо) (энд $\doteq$ Фурьеийн тэнцэтгэл, тасралтгүй цэгүүд дээр тэнцэнэ).
\end{thm}
\begin{prf}
$x>a$ байг.
\begin{equation} \tag{$*$} \label{11*}
    \varphi (t)=e^{-xt} f(t)
\end{equation}
функц хэсэгчилсэн гөлгөр ба $x\rightarrow +\infty$ үед экспонент хурдаар тэг рүү тэмүүлэх ёстой. Иймд Фурьегийн интеграл ашиглавал
$$\varphi (t)\doteq \frac{1}{2\pi} \int\limits_{\infty}^{+\infty} d\xi \int\limits_{-\infty}^{+\infty} \varphi (\eta) e^{i\xi (t-\eta)} d\eta$$
байна. Энд (\ref{11*}) томьёог тавивал 
$$e^{-xt}f(t)\doteq \frac{1}{2\pi} \int\limits_{-\infty}^{+\infty} e^{i\xi}d\xi \int\limits_{0}^{+\infty}  e^{-(x+i\xi)\eta } f(\eta)f(\eta)d\eta$$
байна. 
$$x+i\xi =z \quad \text{гэвэл} \quad \int\limits_0^{+\infty} e^{-(x+i\xi)\eta}f(\eta) d\eta =F(z) \quad \text{учир}$$
$$f(t)\doteq \frac{1}{2\pi}\int\limits_{-\infty}^{+\infty} e^{(x+i\xi)t} F(z)d\xi =\frac{1}{2\pi i}\int\limits_{x-i\infty}^{x+i\infty} e^{zt}F(t)F(z)dz$$
Иймд Меллиний хувиргалт нь Фурьегийн утгаар Лапласын хувиргалтын урвуу болж болно.
\hfill $\blacksquare$\end{prf}
\begin{thm}
 $F(z)\in H(\{\re z>a\})$ хувьд 
 \begin{enumerate}
     \item [a)] $|z|\rightarrow \infty$ үед $F(z)$  утга жигд тэг рүү тэмүүлж байг.
     \item [б)] $\forall x>a$ хувьд $\bint\limits_{x-i\infty}^{x+i\infty} |F(z)|dy$
жигд нийлдэг байг.

Тэгвэл $F(z)$ нь  $f(t)=\dfrac{1}{2\pi i}\bint\limits_{x-i\infty}^{x+i\infty} e^{zt}F(z)dz(z>a)$ функцийн дүр байна.
 \end{enumerate}
\end{thm}
\section{Меллиний томьёог хэрэглэх нь}
\begin{exm}
$\ml{f_1(t)}=F_1(z)$; $\ml{f_2(t)}=F_2(z)$ бол 
$$\ml{f_1(t)\cdot f_2(t)}=\frac{1}{2\pi i}\int\limits_{x-i\infty}^{x+i\infty} F_1(\zeta)\cdot F_2(z-\zeta)d\zeta $$
\end{exm}
\begin{sol}
\begin{align*}
    \ml{f_1(t) \cdot f_2(t)}&=\int\limits_0^\infty e^{-zt}f_1(t)f_2(t)dt \\
    &=\int\limits_0^\infty e^{-zt}\left(\frac{1}{2\pi i} \int\limits_{x-i\infty}^{x+i\infty} e^{\zeta t}F_1(\zeta)d\zeta \right) \cdot f_2(t)dt \\
    &=\frac{1}{2\pi i}\int\limits_0^\infty \int\limits_{x-i\infty}^{x+i\infty} e^{(\zeta-z)t}F_1(\zeta)f_2(t)d\zeta dt \\
    &=\frac{1}{2\pi i} \int\limits_{x-i\infty}^{x+i\infty} F_1(\zeta)\left(\int\limits_0^\infty e^{(\zeta -z)t}f_2(t)dt\right) d\zeta \\
    &=\frac{1}{2\pi i} \int\limits_{x-i\infty}^{x+i\infty} F_1(\zeta)F_2(\zeta-z)d\zeta.
\end{align*}
\end{sol}
\begin{exm}
$f_1(t)=\cos at$, $f_2(t)=t$ бол $f(t)=t\cos at$ дүрийг ол.
\end{exm}
\begin{sol}
$\ml{\cos at}=\dfrac{z}{z^2+a^2}$ $\ml{t}=\dfrac{1}{z^2}$ тул Жишээ 10.2 ёсоор 
$$F(z)=\frac{1}{2\pi i}\int\limits_{x-i\infty}^{x+i\infty} \frac{\zeta}{\zeta^2+a^2}\cdot \frac{1}{(\zeta -z)^2}d\zeta$$
 Зурагт үзүүлсэн контураар интегралчилбал хагас тойргийн дагуух интегралын утга радиус өсөхөд тэг болохыг харуулж болно. 
 \begin{center}
     \includegraphics[width=5cm]{Zurag/43.jpg}
 \captionof{figure}{}
        \label{ref hiih ner}
 \end{center}

Иймд
\begin{align*}
    (z)&=\frac{1}{2\pi i}\left(-2\pi i 
    {\tiny{\begin{array}{c}
Res  \\
 \zeta =z
\end{array}}}
    \frac{\zeta}{(\zeta^2+a^2)(\zeta -z)^2}\right) \\
    &=-\frac{d}{d\zeta}\left(\frac{\zeta}{\zeta^2+a^2}\right)\Big |_ {\zeta =z}=\frac{z^2-a^2}{(z^2+a^2)^2}
\end{align*} Иймд $\ml{t\cdot \cos at}=\dfrac{z^2-a^2}{(z^2+a^2)^2}$
\end{sol}
\begin{exm}
$F(z)=\dfrac{e^{a\sqrt{z}}}{z}$-функцийн эхийг ол.
\end{exm}
\begin{sol}
$\mathbb{C}$ хавтгайг $(-\infty, 0]$ дагуу огтлол хийсэн гээд энэ муж дээр $z>0$ үед $\sqrt{z}>0$ байх мөчрийн аналатик үргэлжлэлийг авч үзье.
\begin{center}
     \includegraphics[width=5cm]{Zurag/44.jpg}
 \captionof{figure}{}
        \label{ref hiih ner}
\end{center}
 Зурагт үзүүлсэн хэрчим ба тойргийн нумуудаар (радиус нь $R$, $r$ байх) хүрээлэгдсэн контур авч үзье. Кошийн теорем ёсоор
$$\left(\int\limits_{\Gamma}+\int\limits_{\Gamma'_R\wedge \Gamma''_R}+\int\limits_{\Gamma_1\wedge \Gamma_2}+\int\limits_{\Gamma_\tau}\right) \left(e^{zt}\frac{e^{-z\sqrt{z}}}{z}\right)dz=0$$ байна. 

Одоо $R\rightarrow \infty$, $r\rightarrow 0$ гэвэл 
\begin{align*}
    &\int\limits_{\Gamma} e^{zt}\dfrac{e^{-a\sqrt{z}}}{z}dz\rightarrow 2\pi i f(t) \qquad \text{(Меллиний хувиргалт)} \\
     &\int\limits_{\Gamma'_R\wedge \Gamma''_{R}} e^{zt}\dfrac{e^{-a\sqrt{z}}}{z}dz\rightarrow 0 \qquad \text{(Жорданы лемм)}
\end{align*}
$$\int\limits_{\Gamma_\tau} e^{zt}\frac{e^{-a\sqrt{z}}}{z}dz=\lim\limits_{\varepsilon \rightarrow 0}\int\limits_{-\pi +\varepsilon}^{\pi -\varepsilon} e^{rte^{i\varphi}}\frac{e^{-a\sqrt{r}e^{i\frac{\varphi}{2}}}}{re^{i\varphi}}\cdot ure^{i\varphi} d\varphi \xrightarrow[r\rightarrow 0]{}\int\limits_{-\pi}^\pi id\varphi=2\pi i$$
байна. 
\begin{align*}
    \int\limits_{\Gamma_1\wedge \Gamma_2}e^{zt}\frac{e^{-a\sqrt{z}}}{z}dz &\xrightarrow[]{} -\int\limits_0^\infty \frac{e^{-yt}e^{-ia\sqrt{y}}}{y}dy+\int\limits_0^\infty \frac{e^{-yt}e^{ia\sqrt{y}}}{y}dy\\
    &=2\pi i \int\limits_0^\infty \frac{e^{-yt}}{y}\sin a\sqrt{y}dy
\end{align*}
 байна. 

Иймд $f(t)=1-\dfrac{1}{\pi}\int\limits_0^\infty e^{-yt}\frac{\sin a\sqrt{y}}{y}dy$ гэсэн дүрсээр эхийг олно. Френелийн интеграл ашиглан
$$f(t)=1-\Phi \left(\frac{a}{2\sqrt{t}}\right) \qquad \left(\text{энд} \quad \Phi(t)=\frac{2}{\sqrt{\pi}}\int\limits_0^t e^{-y^2}dy\right)$$
болохыг харуулж болно.
$$\ml{\frac{e^{-a\sqrt{z}}}{z}}=1-\Phi \left(\frac{a}{2\sqrt{t}}\right)$$
\end{sol}
\begin{exm}
$F(z)=\sum\limits_{n=1}^\infty \dfrac{a_n}{z^n}$ бол 
$$
\ml{F(z)}=f(t)=
\begin{cases}
 0, \qquad t<0 \\
 \sum\limits_{n=0}^\infty C_{n+1}\frac{t^n}{n!}, \qquad t>0
\end{cases}
$$ 
\end{exm}
\begin{sol}
$\ml{\dfrac{t^n}{n!}}=\dfrac{1}{z^{n+1}}$ тул шугаман чанар ёсоор 
$$\ml{\sum\limits_{n=0}^\infty C_{n+1}\frac{t^n}{n!}}=\sum\limits_{n=0}^\infty C_{n+1}\frac{1}{z^{n+1}}=F(z)$$ батлагдав. 
\end{sol}
\begin{exm}
$$\frac{e^{-\frac{1}{z}}}{z}=\mathcal{L}\left\{\sum\limits_{n=1}^\infty \frac{(-1)^{n-1}}{(n-1)!\cdot z^n}\right\}=\sum\limits_{k=0}^\infty (-1)^k \frac{\left(\frac{2\sqrt{t}}{2}\right)^{2k}}{(k!)^2}=J_0(2\sqrt{t})$$
(Энд $J_0$ Бесселийн тэг эрэмбийн функц).
\end{exm}
\begin{exm}
$$y^{IV}+2y''+y=0$$
$$y(0)=y'(0)=y''(0)=0$$
$$y'''(0)=1 \quad \text{бод}.$$
\end{exm}
\begin{sol}
Хэрэв $\ml{y}=F(z)$ гэвэл  \\
$\ml{y^{IV}}=z^4\left(F(z)\dfrac{1}{z^4}\right) =z^4F(z)-1$\\
$\ml{y''}=z^2F(z)$ тул $z^4F(z)-1+2z^2F(z)+F(z)=0\Rightarrow $\\
$F(z)=\dfrac{1}{z^4+2z^2+1}$ байна. Иймд Меллиний томьёо ёсоор 
$$y(x)=\frac{1}{2\pi i}\int\limits_{t-i\infty}^{t+i\infty} e^{zx}\frac{dz}{z^4+2z^2+1}$$
Зурагт үзүүлсэн контурыг ашиглавал $\Gamma_R$ хувьд жорданий лемма хүчинтэй тул
$$y(z)= {\tiny{\begin{array}{c}
Res  \\
 \zeta =z
\end{array}}} \frac{e^{zx}}{(z+i)^2}+ {\tiny{\begin{array}{c}
 Res  \\
 \zeta =z
\end{array}}}\frac{e^{zx}}{(z-i)^2}=\frac{1}{2}(\sin x-x\cos x)$$
\end{sol}
\begin{exm}
$
\begin{cases}
 y''\neq y=\sin t \\
 y(0)=0 \\
 y'(0)=0
\end{cases} 
$ бод. 
\end{exm}
\begin{sol}
$z^2F(z)+F(z)=\dfrac{1}{z^2+1} \Rightarrow F(z)=\dfrac{1}{(z^2+1)^2}.$\\ Иймд
$y(x)=\dfrac{1}{2\pi i}\bint\limits_{t-i\infty}^{t+i\infty} e^{zx}\cdot \dfrac{1}{(z^2+1)^2}dz=\dfrac{1}{2}(\sin x-x\cos x).$
\end{sol}
\begin{exm}
$\begin{cases}
 u_t=a^2u_{xx}\\
 u(x,0)=0\\
 u(0,y)=q(t)
\end{cases} 
$ \\
тухайн уламжлалт дифференциал тэгшитгэлийг бод. $x$ бэхлээд $t$ хувьд Лапласын хувиргалт хийвэл
\begin{align*}
    \ml{u(x,t)}=U(x,z)\\
    \ml{u_t(x,t)}=zU(x,z)\\
    \ml{u_{xx}(x,t)}=U_{xx}(x,z)
\end{align*} тул 
$$\begin{cases}
 U_{xx}(x,z)+\dfrac{z}{a^2}U(x,z)=0\\
 U(0,z)=Q(z)
\end{cases}
$$ болж энгийн дифференциал тэгшитгэл үүсэх рүү шилжинэ.
$$U(x,z)=Q(z)\cdot e^{-\frac{\sqrt{z}}{a}x}$$ тул 
$$u(x,t)=\frac{1}{2\pi i}\int\limits_{t-i\infty}^{ti\infty} e^{zx}Q(z)e^{-\frac{\sqrt{z}}{a}x}dz$$ байна.
\end{exm}
\begin{center}
    \textbf{Дасгал}
\end{center}
$F(z)=\dfrac{w}{z^2+w^2}, \quad \re z>0$ бол дүрийг ол. 
\\
\newpage
\begin{center}
    \textbf{Ашигласан ном зохиол}
\end{center}
\begin{enumerate}
    \item Гурвиц А., Курант Р. Теория функций (Перев.с Нем.с Евграпова М.А.) “Наука”, М., 1968.
    \item Привалов И.И. Введение в теорию функций комп- лексного переменного, изд.12-ое “Наука”, М., 1977.
    \item Маркушевич А.И. Краткий курс теории аналити- ческих функций, “Наука”, М., 1978.
    \item Шабат Б.В. Введение в комплексный анализ. ч.I, Функции одного переменного ч.II. Функции нескольких переменных, изд.З-е, “Наука”, М., 1985.
    \item Стоилов С. Теория функций комплексного переменного, т.I, т.II, перев.с Рум. Берштейна И., ИЛ, М., 1962.
    \item Евграфов М.А. Аналитические функции, “Наука”, М., 1965
    \item Бицадзе А. В Основы теории аналитических функций комплексного переменного, “Наука”, М., 1969.
    \item Свешников А,Г., Тихонов А.Н. Теория функций комплексной переменной, “Наука”, М., 1967
    \item Волковыский Л.И., Лунц Г.Л., Араманович И.Г. Сборник задач по теории функций комплексного переменного, изд.2-е, “Наука”, М.,1970
\end{enumerate}
%\printindex
%% =========================
% Гэрийн даалгавар: Комплекс тоо (12 нүүрээс)
% =========================
\chapter*{Гэрийн даалгавар: Комплекс тоо}
\addcontentsline{toc}{chapter}{Гэрийн даалгавар: Комплекс тоо}
Дараах 20 бодлогыг бодож шалгуулна. Хүн бүрээс аль ч хамаагүй 5 бодлогыг бодуулж шалгана.
\section*{1.1 Суурь алгебр ба геометр чанарууд}
\begin{enumerate}[label=\textbf{1.\arabic*}, itemsep=8pt]

\item Дараах тэнцэтгэлүүдийг шалга.
\begin{enumerate}[label=\textbf{(\alph*)}]
\item $\big(\sqrt{2}-i\big)-i\big(1-\sqrt{2}i\big)=-2i$.
\item $(2-3i)(-2+i)=-1+8i$.
\end{enumerate}

\item Дараах илэрхийллийг хялбарчил:
\[
\frac{5i}{(1-i)(2-i)(3-i)}.
\]

\item Дараахыг батал:
\begin{enumerate}[label=\textbf{(\alph*)}]
\item $\Re(iz)=-\Im(z)$;
\item $\Im(iz)=\Re(z)$.
\end{enumerate}

\item Комплекс тоонуудын үржвэрийн \textbf{хаалт тавих} (associative) хууль хүчинтэйг батал:
\[
(z_1z_2)z_3=z_1(z_2z_3)\qquad(\forall z_1,z_2,z_3\in\mathbb C).
\]

\item Тооцоол:
\begin{enumerate}[label=\textbf{(\alph*)}]
\item $\displaystyle \frac{2+i}{2-i}$;
\item $(1-2i)^4$.
\end{enumerate}

\item Дараах дүрслэл өгөгдсөн байг:
\[
f:\mathbb{C}\to \mathbf{M}_{2,2} \,
, \; z=x+yi\ \mapsto\
\begin{pmatrix}
x & y\\
-y & x
\end{pmatrix}.
\]
Тэгвэл бүх $z_1,z_2\in\mathbb C$-ийн хувьд $f(z_1z_2)=f(z_1)f(z_2)$ болохыг батал.

\item Биномын томьёо ашиглан задла:
\begin{enumerate}[label=\textbf{(\alph*)}]
\item $(1+\sqrt{3}\,i)^{2011}$;
\item $(1+\sqrt{3}\,i)^{-2011}$.
\end{enumerate}

\item $z_1,z_2$ комплекс тоонууд ба $z_1\overline{z_2}$ нь тэг биш бодит бол $\exists r\in\mathbb{R},  z_1=r z_2$ болохыг батал.

\end{enumerate}

\section*{1.2 Модуль}
\begin{enumerate}[label=\textbf{2.\arabic*}, itemsep=8pt]

\item Бүх $z_1,z_2\in\mathbb C$-ийн хувьд
\[
|z_1-z_2|^2+|z_1+z_2|^2=2\big(|z_1|^2+|z_2|^2\big)
\]
болохыг батал.

\item Дараах тэнцэтгэл бишийг шалга:
\[
|\Re z|+|\Im z|\le \sqrt{2}\,|z|.
\]
\textit{Санамж:} Үүнийг $(|x|-|y|)^2\ge 0$.

\item Комплекс хавтгай дээрх муруйнуудыг \textbf{зурагла}:
\begin{enumerate}[label=\textbf{(\alph*)}]
\item $\Im(z)=-1$;
\item $|z-1|=|z+i|$;
\item $2|z|=|z-2|$.
\end{enumerate}

\item $|z|=R>1$ нөхцөлтэй бүх $z$-ийн хувьд дараах үнэлгээнүүдийг батал:
\[
\frac{R^4-R}{R^2+R+1}\le
\left|\frac{z^4+iz}{z^2+z+1}\right|
\le
\frac{R^4+R}{(R-1)^2}.
\]

\item $z\ne 0$ үед
\[
|\log(z)|\le |\ln|z||+\pi
\]
болохыг батал. (Энд $\log$ нь үндсэн логарифм.)

\end{enumerate}

\section*{1.3 Экспоненциал ба туйлын хэлбэр, комплекс язгуурууд}
\begin{enumerate}[label=\textbf{3.\arabic*}, itemsep=8pt]

\item Дараахыг $x+iy$ хэлбэрт (энд $x,y\in\mathbb R$) хувирга:
\begin{enumerate}[label=\textbf{(\alph*)}]
\item $\displaystyle \frac{i}{1-i}+\frac{1-i}{i}$;
\item $-8i$ тооны бүх \textbf{гуравдугаар зэрийн} язгуурууд;
\item $\left(\frac{i+1}{\sqrt2}\right)^{2026}$.
\end{enumerate}

\item Дараах тоонуудын \textbf{ аргумент} (principal argument) болон \textbf{илтгэгч хэлбэр}-ийг ол:
\begin{enumerate}[label=\textbf{(\alph*)}]
\item $\displaystyle z=\frac{i}{1+i}$;
\item $z=\sqrt3+i$;
\item $z=2-i$.
\end{enumerate}

\item Дараах тэгшитгэлүүдийн бүх комплекс шийдийг ол:
\begin{enumerate}[label=\textbf{(\alph*)}]
\item $z^6=-9$;
\item $z^2+2z+(1-i)=0$.
\end{enumerate}

\item $z^4+16=0$ тэгшитгэлийн 4 язгуурыг ол.
Дараа нь эдгээрийг ашиглан $z^4+16$-г \textbf{бодит коэффициенттэй} хоёр квадрат үржвэрт задла:
\[
z^4+16=(z^2+az+b)(z^2+cz+d)\quad (a,b,c,d\in\mathbb R).
\]

\item Дараахыг хий:
\begin{enumerate}[label=\textbf{(\alph*)}]
\item Илтгэгч хэлбэр ашиглан
\[
(1+\sqrt3\,i)^{2011},\qquad (1+\sqrt3\,i)^{-2011}
\]
-ийг ол.
\item Дараах хоёр тэнцэтгэлүүдийг батал:
\[
\sum_{m=0}^{1005}\binom{2011}{2m}(-3)^m=2^{2010},
\qquad
\sum_{m=0}^{1005}\binom{2011}{2m+1}(-3)^m=2^{2010}.
\]
\end{enumerate}

\item Эхлээд дараах нийлбэрийн томьёог батал:
\[
1+z+z^2+\cdots+z^n=\frac{1-z^{n+1}}{1-z}\qquad (z\ne 1),
\]
дараа нь үүнийг ашиглан \textbf{Лагранжийн тригонометрийн адилтгал}-ыг гарга:
\[
1+\cos q+\cos 2q+\cdots+\cos nq
=
\frac12+\frac{\sin\left(\frac{(2n+1)q}{2}\right)}{2\sin\left(\frac{q}{2}\right)}
\qquad (0<q<2\pi).
\]

\item Комплекс тоо ашиглан \textbf{косинусын теорем}-ыг батал:
Гурвалжин $\triangle ABC$-д $|BC|=a$, $|CA|=b$, $|AB|=c$, $\angle BCA=q$ бол
\[
a^2+b^2-2ab\cos q=c^2.
\]
\textit{Санамж:} $C$-г тэг дээр байрлуулаад $B=z_1$, $A=z_2$ гэж авч,
$z_1\overline{z_2}+z_2\overline{z_1}=2|z_1||z_2|\cos q$ гэдгийг харуул.

\end{enumerate}
%\input{Midtermexam2026spring}
\end{document}

     %\includegraphics[width=8cm]{image/94h1.PNG}
     
